\documentclass[11pt,a4paper]{article}
\usepackage[margin=27mm]{geometry}
\usepackage{fontspec}
\setmainfont{Latin Modern Roman}
\usepackage{amsmath,amssymb,amsthm,mathtools}
\usepackage{microtype,longtable,booktabs,array,enumitem}
\usepackage[colorlinks=true,linkcolor=blue,citecolor=blue,urlcolor=blue]{hyperref}
\usepackage{xurl}
\newtheorem{theorem}{Statement}[section]
\newtheorem{lemma}[theorem]{Lemma}
\newtheorem{proposition}[theorem]{Proposition}
\newtheorem{corollary}[theorem]{Corollary}
\theoremstyle{definition}
\newtheorem{definition}[theorem]{Definition}
\newtheorem{remark}[theorem]{Remark}
\DeclareMathOperator{\Pic}{Pic}
\DeclareMathOperator{\Spec}{Spec}
\DeclareMathOperator{\Proj}{Proj}
\DeclareMathOperator{\NS}{NS}
\DeclareMathOperator{\Ext}{Ext}
\DeclareMathOperator{\Hom}{Hom}
\DeclareMathOperator{\End}{End}
\DeclareMathOperator{\Lie}{Lie}
\DeclareMathOperator{\Ker}{Ker}
\DeclareMathOperator{\Coker}{Coker}
\providecommand{\mathscr}{\mathcal}
\newcommand{\OO}{\mathcal O}
\newcommand{\GG}{\mathbf G}
\newcommand{\ZZ}{\mathbf Z}
\newcommand{\QQ}{\mathbf Q}
\newcommand{\PP}{\mathbf P}
\renewcommand{\AA}{\mathbf A}
\newcommand{\status}[1]{\textbf{#1.}}
\title{Grothendieck's Picard and Jacobian dossiers\\\large RESEARCH\_SA14\_SA19}
\author{AI-produced by GPT-6 Astra\\OpenAI, ChatGPT Pro mode}
\date{20 November 2026}
\begin{document}
\maketitle
\begin{center}\textbf{INCOMPLETE: research in progress}\end{center}
\begin{abstract}
We investigate the questions in cotes 43, 44, and 57--60 using the supplied diplomatic transcriptions and checked primary sources. We distinguish proved results, counterexamples, conditional results, and undecided questions. This working version claims neither completion nor priority.
\end{abstract}
\tableofcontents
\section*{Conventions and evidence}
We write $\Pic(X)$ for the abelian group of invertible sheaves up to isomorphism, and $\mathbf{Pic}_{X/S}$ for the fppf Picard functor or its representing group algebraic space when existence is specified. $\Pic^0$ denotes the identity component, $\Pic^\tau$ the union of components whose image in the N\'eron--Severi group is torsion. Quotations are from the supplied transcriptions, with references in the form cote:scan. We have not inspected manuscript scans in this run.
\section{Lefschetz for Picard groups in positive characteristic}\label{sec:A}

\subsection{The question and its quantifiers}
At 57:97 the transcription asks:
\begin{quote}
``Surjectivité (donc bijectivité) de $\Pic(X)\to\Pic(Y)$ pour $n\geq4$
sans conditions cohomologiques'';
\end{quote}
the adjacent box reads ``noyau nul si $n\geq4$''.
We take $X$ smooth, connected and projective over an algebraically closed
field $k$, $L$ a specified very ample invertible sheaf, and $Y=V(h)$ a
smooth member of $|L|$.  We distinguish an arbitrary smooth member from
the geometric generic member.  A tangent-space assertion for a Picard
scheme is not an assertion about the kernel on its $k$-points.

\paragraph{Current state of A1--A3.}
A3 is proved affirmatively by Theorem~\ref{a:very-ample-example}: an
explicit characteristic-two surface has a very ample sheaf with
one-dimensional negative first cohomology. The least requested dimension
is two. The unrestricted isomorphism A1 and injectivity A2 remain
undecided. The finite obstruction description A4 is proved.

\subsection{A finite, exact obstruction description}
Write $I=\OO_X(-Y)$ and $Y_m=V(I^{m+1})$, so $Y_0=Y$.
The cohomology of sheaves of units below is Zariski cohomology.

\begin{theorem}[A4: exact finite description; Proved]\label{a:finite}
Suppose $\dim X=n\geq4$.  Choose $N\geq0$ such that
\[
 H^i(Y,L^{-m}|_Y)=0\qquad(i=1,2;\ m>N).
\]
Such an $N$ exists.  On the topological space $Y$ put
\[
 U_N=\ker(\OO_{Y_N}^{\times}\longrightarrow\OO_Y^{\times})
     =1+I/I^{N+1}.
\]
For $N=0$ this is the zero group sheaf.  Restriction gives a canonical
identification $\Pic(X)\simeq\Pic(Y_N)$ and an exact sequence
\begin{equation}\label{a:exact}
 0\longrightarrow H^1(Y,U_N)\longrightarrow\Pic(X)
 \longrightarrow\Pic(Y)\xrightarrow{\delta_N}H^2(Y,U_N).
\end{equation}
Consequently the complete answers for this particular pair are
\[
 \ker(\Pic(X)\to\Pic(Y))=H^1(Y,U_N),\qquad
 \operatorname{coker}(\Pic(X)\to\Pic(Y))=\operatorname{im}(\delta_N).
\]
The sheaf $U_N$, with its actual multiplication, has the filtration
\[
 U_N^{(m)}=1+I^m/I^{N+1},\qquad
 U_N^{(m)}/U_N^{(m+1)}\simeq L^{-m}|_Y\quad(1\leq m\leq N).
\]
Thus a finite tower with the additive layers of R4 determines both groups
and all obstruction maps.  Replacing this tower by the direct sum of the
cohomology of its layers is not justified.
\end{theorem}
\begin{proof}
We first check algebraization.  The effective Lefschetz theorem
\cite[Expos\'e XII, Corollary 3.4, pp.~119--120]{SGA2} applies: closed
points of $Y$ have depth $n-1\geq2$, those of $X$ have depth $n\geq3$,
and $Y$ is an effective ample Cartier divisor.  Every line bundle on
$\widehat X_Y$ algebraizes on an open neighborhood $V$ of $Y$;
completion is fully faithful on vector bundles on $V$.
Every positive-dimensional projective closed subset meets $Y$: otherwise
$h$ would trivialize the ample sheaf $L$ on that subset.  Thus
$X\setminus V$ has codimension at least two.  Since $X$ is regular,
closing a divisor representing a line bundle on $V$ extends it uniquely
to $X$, by local factoriality.  We obtain
\begin{equation}\label{a:formal}
 \Pic(X)\simeq\Pic(\widehat X_Y).
\end{equation}

By \cite[Expos\'e XII, Corollary 3.5, p.~120]{SGA2}, $Y$ is connected.
Hence $H^0(Y,\OO_Y)=k$.  The anti-ample sheaves $L^{-m}|_Y$ have no
nonzero section for $m>0$.  Induction using
$0\to L^{-m}|_Y\to\OO_{Y_m}\to\OO_{Y_{m-1}}\to0$ gives
$H^0(Y_m,\OO_{Y_m})=k$: its injection into $k$ contains the constants.
All transition maps on global units are consequently the identity on
$k^\times$.

A formal line bundle is a compatible system of line bundles on the $Y_m$.
Compatible isomorphism classes can be equipped successively with
compatible isomorphisms.  The ambiguity is the derived inverse limit of
their global automorphism groups; it vanishes because the transition maps
on $k^\times$ are surjective.  Therefore
$\Pic(\widehat X_Y)=\varprojlim_m\Pic(Y_m)$.
For $m\geq1$, multiplication on $1+I^m/I^{m+1}$ is addition, and the
calculation of global units gives
\begin{equation}\label{a:layer}
 0\to H^1(Y,L^{-m}|_Y)\to\Pic(Y_m)\to\Pic(Y_{m-1})
       \xrightarrow{o_m}H^2(Y,L^{-m}|_Y).
\end{equation}
For $m>N$ its middle transition is an isomorphism, so the inverse limit
is $\Pic(Y_N)$.  The bound exists by Serre duality and Serre vanishing
on the $(n-1)$-fold $Y$, since $(n-1)-i>0$ for $i=1,2$.
Finally the cohomology sequence of
$1\to U_N\to\OO_{Y_N}^{\times}\to\OO_Y^{\times}\to1$ and
surjectivity on global units give \eqref{a:exact}.  Multiplication of units
gives the stated filtration.  These are the actual restriction maps of
the given thickening, so no extension or obstruction has been discarded.
\end{proof}

\paragraph{An explicit finite obstruction procedure.}
Choose a finite affine cover of $Y_N$ trivializing a line bundle $M$ on
$Y$.  Lift its transition units $g_{ij}$ to units $\widetilde g_{ij}$ on
the pairwise intersections in $Y_N$.  The products
\[
 c_{ijk}=\widetilde g_{ij}\widetilde g_{jk}\widetilde g_{ki}\in U_N
\]
represent $\delta_N(M)$.  Different lifts change it by a coboundary;
$M$ extends if and only if this class vanishes.  Alternatively solve the
equations first modulo $I^2$, then modulo $I^3$, through $I^{N+1}$.
At stage $m$ the remaining additive cocycle lies in $L^{-m}|_Y$ and
represents $o_m$; all choices of a lift form a torsor under
$H^1(Y,L^{-m}|_Y)$.  Those choices must be retained because the next
obstruction can depend on them.  Quasi-coherent acyclicity on affine
intersections, applied to the finite filtration, justifies these
\v Cech computations.  This is a finite description using the layers and
their multiplication; the vector spaces alone do not include its maps.

\begin{corollary}[A4b: bounded finite $p$-groups; Proved]\label{a:pgroups}
In characteristic $p>0$, for $n\geq4$, both groups in
Statement~\ref{a:finite} are finite and killed by $p^e$ whenever
$p^e\geq N+1$.  For $n=3$ the kernel is still finite and killed by a
power of $p$.
\end{corollary}
\begin{proof}
For $a\in I/I^{N+1}$, $(1+a)^{p^e}=1+a^{p^e}=1$.  Thus $p^e$ kills
$U_N$, its cohomology, and the two groups in \eqref{a:exact}.
The $p^e$-torsion of $\Pic(X)$ is finite: $\NS(X)$ is finitely generated,
and the $k$-points of $\mathbf{Pic}^0_{X/k}$ are those of an abelian
variety, with finite multiplication-by-$p^e$ kernel.  This proves
finiteness of the kernel.  Multiplication by $p^e$ is surjective on the
$k$-points of an abelian variety.  Since the cokernel is killed by $p^e$,
the image contains $\Pic^0(Y)(k)$.  It is therefore a bounded torsion
quotient of the finitely generated $\NS(Y)$ and is finite.

For $n=3$, choose $N$ so that only $H^1(Y,L^{-m}|_Y)=0$ for $m>N$.
The transitions in \eqref{a:layer} are then injective.
The identification \eqref{a:formal} still holds and embeds $\Pic(X)$
into $\Pic(Y_N)$.  Its restriction kernel embeds in $H^1(Y,U_N)$;
the exponent and finiteness arguments apply unchanged.  No surjectivity
of the transitions has been asserted in dimension three.
\end{proof}

\begin{corollary}[A4c: vanishing criterion; Known, hypotheses checked]
\label{a:vanishing}
If $n\geq3$ and $H^1(Y,L^{-m}|_Y)=0$ for all $m>0$, restriction
$\Pic(X)\to\Pic(Y)$ is injective.  If $n\geq4$ and also all the
groups $H^2(Y,L^{-m}|_Y)$ vanish, restriction is an isomorphism.
\end{corollary}
\begin{proof}
Each transition in \eqref{a:layer} is injective in the first case and an
isomorphism in the second; we use \eqref{a:formal}.  This also checks
the smooth case of \cite[Expos\'e XII, Corollary 3.6, p.~121]{SGA2}:
its local depth and parafactoriality conditions follow from regularity.
The vanishing assumptions are absent from the unrestricted question A1.
\end{proof}

\subsection{Reduction of the main gates}
\begin{proposition}[A3r: reduction to surfaces; Proved]\label{a:surface}
If there is a smooth projective variety of dimension at least two with
a very ample $L$ and $H^1(X,L^{-1})\ne0$, there is such a surface.
If A3 has a positive existence answer, its least dimension is two.
\end{proposition}
\begin{proof}
When $\dim X>2$, Serre duality and Serre vanishing give
$H^1(X,L^{-a})=0$ for large $a$.  Choose the largest positive $a$ for
which it is nonzero, and a smooth integral $D\in|L|$ by Bertini for
the specified embedding.  The sequence
\[
 0\to L^{-a-1}\to L^{-a}\to L^{-a}|_D\to0
\]
gives $0\ne H^1(X,L^{-a})\hookrightarrow H^1(D,L^{-a}|_D)$.
The sheaf $L^a|_D$ is very ample.  Repeat until the dimension is two.
This proves an implication between existence questions; it does not
provide the initial example.
\end{proof}

\begin{proposition}[A1r: reduction to fourfolds; Proved]\label{a:fourfold}
The universal assertion A1 is equivalent to its dimension-four case.
A failure of injectivity in any dimension $n\geq4$ also gives one in
dimension four.
\end{proposition}
\begin{proof}
For $n>4$, choose large $d$ and a general smooth $D\in|L^d|$ transverse
to $Y$.  Serre duality and vanishing on $X$ give
\[
 H^j(X,L^{-a})=0\qquad(j=1,2,3;\ a\geq d)
\]
after increasing $d$.  The restriction sequences with exponents $md$
and $(m+1)d$ give, for every $m\geq1$,
\[
 H^1(D,L^{-md}|_D)=H^2(D,L^{-md}|_D)=0.
\]  The same argument on $Y$, whose dimension is at least four,
gives both vanishings on $Y\cap D$ for its normal powers in $Y$.
Corollary~\ref{a:vanishing} gives vertical isomorphisms in
\[
 \begin{matrix}
 \Pic(X)&\longrightarrow&\Pic(Y)\\
 \big\downarrow\wr&&\big\downarrow\wr\\
 \Pic(D)&\longrightarrow&\Pic(Y\cap D).
 \end{matrix}
\]
Here $Y\cap D$ is a smooth hyperplane for $L|_D$, the restriction of the
\emph{original} embedding.  Only the auxiliary cut used $L^d$.
Both rows have the same kernel and cokernel.  Iteration proves the claim.
\end{proof}

\begin{proposition}[A2b: dimension two fails; Proved]\label{a:segre}
In every characteristic, a smooth hyperplane section of a smooth
projective surface can have a nonzero restriction kernel.
\end{proposition}
\begin{proof}
Let $X=\PP^1\times\PP^1$ with the Segre embedding and
$L=\OO(1,1)$.  The diagonal is the smooth hyperplane section
$x_0y_1-x_1y_0=0$, including in characteristic two.
The map $\Pic(X)=\ZZ^2\to\Pic(Y)=\ZZ$ is $(a,b)\mapsto a+b$;
its kernel is $\ZZ(1,-1)$.  A universal injectivity threshold, if it
exists, is at least three.  No dimension-three or dimension-four
injectivity assertion follows from this example.
\end{proof}

\subsection{The infinitesimal kernel on 57:100}
The supplied transcription calls this kernel a ``$p$-groupe fini
unipotent, with the marginal qualification ``purement infinitésimal
(57:100). Its phrase ``sans partie fixe must also be retained: the
map on classes defined over $k$ is distinguished from the group-scheme
map after geometric-generic base change.
\begin{proposition}[R3/A3t: tangent calculation; Proved]\label{a:tangent}
For $\dim X\geq2$ and any smooth hyperplane $Y$, the scheme-theoretic
restriction kernel $K$ satisfies
\[
 \Lie K=\ker\bigl(H^1(X,\OO_X)\to H^1(Y,\OO_Y)\bigr)
        =h\,H^1(X,L^{-1}).
\]
Multiplication by $h$ on this last group is injective; therefore
$\dim_k\Lie K=h^1(X,L^{-1})$.  If the kernel under consideration is
finite and connected, it is trivial exactly when this group vanishes.
\end{proposition}
\begin{proof}
First-order line-bundle deformations have transition units
$1+\epsilon a_{ij}$, giving tangent space $H^1(\OO)$; tangent spaces
preserve kernels.  Apply the long exact sequence of
$0\to L^{-1}\xrightarrow{h}\OO_X\to\OO_Y\to0$, observing that
$H^0(X,\OO_X)\to H^0(Y,\OO_Y)$ is the identity of $k$.
For a finite connected group scheme over $k$, vanishing of the tangent
space gives $\mathfrak m/\mathfrak m^2=0$ in its local Artinian ring.
Nakayama's lemma gives $\mathfrak m=0$, so it is trivial.
\end{proof}

\begin{proposition}[R3g: geometric generic kernel; Proved]\label{a:generic}
For the geometric generic hyperplane, the kernel on $\mathbf{Pic}^\tau$
is finite and connected.  Thus the equivalence in R3 is valid for smooth
$X$: this kernel is nontrivial if and only if $H^1(X,L^{-1})\ne0$.
\end{proposition}
\begin{proof}
We first check finiteness without presupposing the manuscript's assertion.
For normal projective $X$, $\mathbf{Pic}^\tau$ is proper and of finite
type: its identity component is proper
\cite[Theorem 5.4 and Remark 5.6]{Kleiman2005}, and its component group
is the finite torsion subgroup of $\NS(X)$.  The restriction kernel is
therefore a proper group scheme of finite type.  If its reduced identity
component had positive dimension, it would be a positive-dimensional
abelian variety and would have nonzero $\ell$-torsion points for a prime
$\ell\ne\operatorname{char}k$.  Such a point is impossible: Kummer
theory identifies prime-to-characteristic torsion line bundles with
$H^1(\mu_\ell)$, and the surjectivity of the fundamental-group map for
an ample divisor gives injectivity of restriction on that group
\cite[Expos\'e XII, Corollary 3.5]{SGA2}.  Thus the kernel is finite.

We prove that it has no nonzero geometric point.  Every point of a finite
group has finite order.  For each integer $a>0$, the $a$-torsion points
of $\Pic^\tau(X)$ are finite and defined over the algebraically closed
field $k$.  Hence every point in the geometric generic kernel comes from
a torsion line bundle on $X$.  Prime-to-$p$ torsion has already been
excluded.  To exclude $p$-torsion, we use the injective map
\[
 d\log:\Pic(X)[p]\hookrightarrow H^0(X,\Omega_X^1).
\]
Indeed, trivialize the $p$th power of a line bundle, obtaining local
units $u_i$ whose ratios are $p$th powers.  The forms $du_i/u_i$ glue.
If their common form is zero, then each $u_i$ has a $p$th root in the
function field, since $k$ is perfect.  Normality makes these roots local
units, and they trivialize the bundle.  The construction commutes with
restriction to a smooth divisor.

No nonzero form defined over $k$ restricts to zero on the geometric generic
hyperplane.  To see this, use the incidence variety of pairs $(x,H)$ with
$x\in X\cap H$.  On the locus where the hyperplane is smooth at $x$,
the relative cotangent sheaf is locally free.  Vanishing on its generic
fibre would make the pulled-back form zero throughout this integral
smooth locus.  But, at a point $x$ where the form is nonzero, the tangent
spaces of hyperplanes through $x$ run through all codimension-one
subspaces of $T_xX$.  Since $\dim X\geq2$, some such subspace is not
annihilated by the form.  This is a contradiction.  Thus $p$-torsion is
excluded; successively taking $p$th powers excludes $p^a$-torsion for
all $a$.  The finite kernel has just its identity as a geometric point
and is connected.  Proposition~\ref{a:tangent} finishes the proof.
\end{proof}

This argument checks the finite-infinitesimal property actually needed
for R3; it does not need the stronger unipotence assertion written at
57:100.  Its incidence argument is only injectivity for forms defined
over $k$.  We used it for torsion points that have first been proved to
be defined over $k$; we did not incorrectly infer injectivity for all
linear combinations with geometric generic coefficients.

\paragraph{A transcription correction, kept explicit.}
The two displayed identifications on 57:100 both print
$\mathrm H^1(X,\underline{\mathcal O}_X)^{1-F}$, including in
the $\mu_p$ line. In the additive line the correct expression is
$\ker(F:H^1(\OO_X)\to H^1(\OO_X))$ for proper connected
$X$ over the perfect field; $1-F$ instead belongs to the
Artin--Schreier sequence for $\mathbf Z/p$. The second line also
writes the quotient $\mathcal O_X/\mathcal O_X^*$, whereas the
multiplicative quotient needed there is
$\mathcal O_X^\times/\mathcal O_X^{\times p}$.  As printed in the supplied
transcription it is not correct.  For smooth connected proper $X/k$ the
Kummer sequence gives $H^1_{\mathrm{fppf}}(X,\mu_p)=\Pic(X)[p]$,
since $k^\times$ is $p$-divisible.  The corresponding Zariski quotient
is $\OO_X^\times/\OO_X^{\times p}$, and its map is $d\log$.
We use this corrected mathematical statement, not a silent change of
the diplomatic text.

\subsection{A sharper surface gate for A3}
\begin{proposition}[A3n: a necessary Frobenius-nilpotence bound; Known deduction]
\label{a:nilpotence}
For a smooth projective connected variety $X/k$ of dimension at least two,
let $q_{\mathrm{nil}}(X)$ be the dimension of the subspace of
$H^1(X,\OO_X)$ killed by some iterate of absolute Frobenius.
If $L$ is very ample and $H^1(X,L^{-1})\ne0$, then
\[
 q_{\mathrm{nil}}(X)\geq h^0(X,L).
\]
In the surface case we furthermore have $h^0(X,L)\geq5$.
Consequently A3 reduces to smooth surfaces with
$q_{\mathrm{nil}}(X)\geq h^0(X,L)\geq5$.
\end{proposition}
\begin{proof}
Apply \cite[Theorem 4.20, pp.~15--16]{Enokizono2021} with the
integral ample divisor $D$ of $L$, $M=D$, and $E=0$.
The variety is normal and geometrically connected over an infinite
perfect field. General hyperplanes cut a smooth surface by Bertini;
$M$ remains ample, hence nef and big, on that surface. The theorem
therefore gives $H^1(X,L^{-1})=0$ whenever
$h^0(X,L)-1\geq q_{\mathrm{nil}}(X)$. Its contrapositive is the
first assertion. In dimension two, a complete very ample system
with $h^0\leq4$ embeds the surface as $\PP^2$ or a hypersurface
in $\PP^3$. The usual hypersurface exact sequence and the
cohomology of line bundles on projective space give
$H^1(X,\OO_X)=0$ in both cases, contrary to the first assertion.
Combining with Proposition~\ref{a:surface} proves the last claim.
\end{proof}
The theorem applies in characteristics two and three as well.
It rules out ordinary varieties, in the explicit sense that Frobenius
is injective on $H^1(\OO_X)$, but supplies no existence theorem for
the remaining surfaces. Its Frobenius-nilpotent dimension must not
be confused with the nonreduced dimension of the Picard scheme.

\begin{proposition}[A3s: general type and numerical bounds; Known deduction]
\label{a:generaltype}
Let $X/k$ be a smooth projective surface and let $L$ be very ample.
If $H^1(X,L^{-1})\ne0$, then $X$ is of general type. For its
smooth minimal model $X_0$, we have
\[
 L^2\leq
 \begin{cases}
 K_{X_0}^2,&p\geq3,\\
 \max\{K_{X_0}^2,K_{X_0}^2-3\chi(\OO_X)+2\},&p=2.
 \end{cases}
\]
Moreover $X$ does not admit a smooth lifting over $W_2(k)$.
\end{proposition}
\begin{proof}
If $X$ is neither of general type nor quasi-elliptic of Kodaira
dimension one, \cite[Theorem 1.6(1), pp.~191--192]{Terakawa1999}
gives the required vanishing for every nef and big $L$, hence for
our very ample $L$. In the remaining quasi-elliptic case let $F$
be a general fibre. This is an integral Gorenstein curve of arithmetic
genus one with $\omega_F\simeq\OO_F$. A positive-degree invertible
sheaf $M$ on $F$ has $h^0(F,M)=\deg M$ by duality and Riemann--Roch.
In degrees one and two its complete system cannot embed $F$:
it maps to a point or to $\PP^1$. Thus very ampleness of $L|_F$
forces $L\cdot F\geq3$. The corrected vanishing theorem
\cite[Theorem 3.1, pp.~4--5]{Zhang2022} applies: its bounds are
$L\cdot F>2$ for $p=2$ and $L\cdot F>1$ for $p=3$.
This excludes the quasi-elliptic case. The two numerical bounds
are contrapositives of \cite[Theorem 1.6(2)--(3)]{Terakawa1999}.
Finally, in the presence of a smooth $W_2(k)$ lifting,
\cite[section 11, Corollary 11.2]{Langer2011} with the numerically
flat bundle $E=\OO_X$, $i=0$, and $j=1$ gives vanishing because
$1<\min(p,2)$ for every prime $p$. This proves the last assertion.
\end{proof}
We use Zhang's corrected bounds rather than the stronger earlier
quasi-elliptic assertions which that paper disproves in characteristic two.
Combining Propositions~\ref{a:surface}, \ref{a:nilpotence}, and
\ref{a:generaltype}, any A3 example must be a surface
of general type with a very ample $L$, the displayed numerical bounds,
$q_{\mathrm{nil}}\geq h^0(L)\geq5$, and no smooth $W_2$ lifting.
These were necessary restrictions during the search; the explicit
construction below now proves existence.

\subsection{Tested approaches and their exact barriers}
\begin{proposition}[A3f: the indicated Raynaud polarizations are not very ample;
Proved]\label{a:fibre}
Let $p$ be the characteristic and let $\ell\geq2$ be prime to $p$.
Let $F$ be the projective curve obtained from
$\Spec k[u^p,u^\ell]$ by adding its smooth point $\infty$, with
normalization $\PP^1_u$.  For $0\leq a<\ell$, the complete system
$|\OO_F(a\infty)|$ does not give an immersion on a dense open subset
of $F$.  Consequently a line bundle restricting to this sheaf on such
a fibre cannot be very ample on the ambient variety.
\end{proposition}
\begin{proof}
A section is a polynomial in the monomials of the semigroup
$\langle p,\ell\rangle$ whose pole order at infinity is at most $a$.
Because $a<\ell$, every such exponent is a multiple of $p$.
Thus every section is a polynomial in $u^p$.  On the smooth locus
$u\ne0$, any map defined by these sections has zero differential, or
is constant.  It cannot be an immersion.  Restriction of a closed
immersion to a closed fibre would still be a closed immersion, proving
the last assertion.
\end{proof}
This applies to the polarizations $\OO(a\widetilde E)\otimes\phi^*N^b$
in the nonvanishing range $a\leq\ell-1$ of
\cite[Theorem 3.7]{Zheng2015}: on a geometric fibre the pullback factor
is trivial and $\widetilde E$ is its smooth point at infinity.  It also
applies to $\widetilde D=(\ell-1)F_\infty+g^*D'$ in
\cite[Definition 2.1 and Proposition 2.3, pp.~7--8]{Mukai2013revised},
where the fibre is explicitly $T^\ell=S^p$ and is normalized at infinity.
We have therefore closed the direct very-ampleness route for these
particular stated polarizations, not merely failed to find an embedding.
The finite semigroup checks in \texttt{code/check\_raynaud\_fibre.py}
cover 518 pairs with primes at most 19 and $2\leq\ell\leq80$;
the preceding proof covers all pairs.

The finite $p$-group assertion of R4 is proved above; its vanishing is
not.  We have checked the following specific limitations.
\begin{enumerate}[leftmargin=*]
\item The classical Picard theorem needs the vanishings in
Corollary~\ref{a:vanishing}; the editorial note in
\cite[Expos\'e XII, section 4, note 5, p.~128]{SGA2} explicitly separates
the characteristic-$p$ issue.
\item \cite[Theorem 0.1]{BhattDeJong2014} gives local injectivity up
to $p$-power torsion, not an isomorphism up to that torsion.  Thus the
recollection in SA14-H3 requires correction.  Its local examples do not
give the smooth projective pair required here.
\item The results of Cotner--Zavyalov
\cite[Theorem 3.4.4 and Corollary 3.4.5]{CotnerZavyalov2024}
require strong ampleness or additional cohomological hypotheses.
Examples 3.3.6 and 3.4.6 distinguish failure for the Picard scheme
from failure for its group of line bundles.  Tangent noninjectivity alone
does not settle A1 or A2.
\item The very ample family in \cite[section 2]{LauritzenRao1997}
has first nonzero negative cohomology in degree $3r-4$ on a
$(3r-3)$-fold, $r\geq3$; all lower degrees vanish by the same
calculation.  It does not supply the degree-one nonvanishing in A3.
\item Proposition~\ref{a:fibre} shows that the indicated Raynaud-type
nonvanishing polarizations are not very ample.  Passing to an unspecified
large power can remove the nonvanishing and is not a valid solution.
\end{enumerate}
For A1 the current gate is a smooth fourfold with a nonzero finite
obstruction in Statement~\ref{a:finite}, or a proof that every such
obstruction vanishes. For A2 it is the surviving $p$-torsion kernel and
the universal dimension-three case. These two gates remain unfinished.
The following construction settles the separate existence question A3.

\subsection{A3R: exclusion of every positive two-parameter Raynaud polarization}
\begin{lemma}[Connection bound; Proved]\label{a:connection-bound}
Let $C$ be a smooth projective connected curve over an algebraically closed
field of characteristic $p>0$. Let $L$ have positive degree, and let
$0\to\OO_C\to E\to L\to0$ be nonsplit with split Frobenius pullback.
For $0\leq n<p$ and any line bundle $M$ with
$\deg M<n\deg L$, we have $\Hom_C(\operatorname{Sym}^n E,M)=0$.
\end{lemma}
\begin{proof}
The splitting $F^*E=\OO_C\oplus L^p$ is unique since
$\Hom(L^p,\OO_C)=0$. Its canonical connection preserves $\OO_C$.
The second fundamental map $\beta:L^p\to\Omega_C^1$ is nonzero:
otherwise the splitting is horizontal and Cartier descent splits the
original extension. The precise descent source is Katz, \emph{Nilpotent
connections and the monodromy theorem}, Theorem~5.1, printed pp.~190--191 (opened in this run).
For a homomorphism $\varphi:\operatorname{Sym}^nE\to M$, its Frobenius
pullback is horizontal for the canonical connections. Its component on
$L^{pn}$ is zero by the degree inequality. On the symmetric power, the
connection carries the summand of exponent $j$ to the summand of exponent
$j-1$ with coefficient $j\beta$, in addition to terms of the same exponent.
If its components of exponents at least $j$ are zero, horizontality forces
$j\beta$ times its component of exponent $j-1$ to vanish. Since
$1\leq j\leq n<p$ and $\beta$ is nonzero, that component vanishes too.
Induction gives $F^*\varphi=0$, hence $\varphi=0$ by faithful flatness of
Frobenius on the smooth curve. For $n=0$ the assertion is simply that a
negative-degree line bundle has no sections.
\end{proof}

\begin{proposition}[A3R; Proved]\label{a:raynaud-all}
For the smooth Raynaud cyclic-cover surface of Zheng's construction,
with $\ell\mid p+1$, $L=N^\ell$ on its Tango base curve and
$\deg N>0$, let $\widetilde E$ be the positive ramification section.
For every $a,b>0$, the line bundle
$Z_{a,b}=\OO_X(a\widetilde E)\otimes\phi^*N^b$ cannot be both very
ample and have $H^1(X,Z_{a,b}^{-1})\ne0$.
\end{proposition}
\begin{proof}
Set $d=(p+1)/\ell$. For $a=\ell$, the finite-cover pushforward in
Zheng's Proposition~3.3 is
$\bigoplus_{i=0}^{\ell-1}\OO_{\PP(E)}(-id-1)\otimes\pi^*N^{ip-b}$.
Its direct image to $C$ is zero. Relative duality for the quotient
convention $\pi_*\OO_{\PP(E)}(1)=E$ therefore gives
\[
 H^1(X,Z_{\ell,b}^{-1})=
 \bigoplus_{i=1}^{\ell-1}
 H^0\!\left(C,(\operatorname{Sym}^{id-1}E)^\vee
                         \otimes N^{ip-b-\ell}\right).
\]
Here $n=id-1$ lies between zero and $p-d$, hence below $p$, and
\[
 \deg(N^{ip-b-\ell})-n\deg L=-(i+b)\deg N<0.
\]
The lemma makes every summand zero.
Next $\widetilde E^2=\deg N>0$, and $\phi|_{\widetilde E}$ is the
identity of $C$. Thus $Z_{a,b}^{-1}|_{\widetilde E}$ has negative
degree $-(a+b)\deg N$. The divisor exact sequence consequently
injects $H^1(X,Z_{a+1,b}^{-1})$ into $H^1(X,Z_{a,b}^{-1})$.
We obtain vanishing for every $a\geq\ell$ and $b>0$.
For $a<\ell$, restriction to the cuspidal fibre is
$\OO_F(a\infty)$ and Proposition~\ref{a:fibre} rules out very
ampleness. These two cases exhaust all positive $a,b$.
\end{proof}

This excludes the entire displayed two-parameter family. It is not a
nonexistence theorem for other line bundles or other surfaces.
The connection argument and all inequalities received an independent
targeted check in this run. The exact Cech computation covers1008 cases;
it serves as a check on this proof, not as its substitute.

\subsection{A3: an explicit very ample surface with negative first cohomology}

\begin{theorem}[A3; Proved]\label{a:very-ample-example}
There is a smooth geometrically integral projective surface $Y$ over an
algebraically closed field of characteristic two, and a very ample
invertible sheaf $L$ on $Y$, for which
\[
 h^0(Y,L)=562,\qquad h^1(Y,L^{-1})=1,\qquad h^1(Y,\OO_Y)=565.
\]
Thus the least dimension in the range requested in A3, namely
$\dim X\geq2$, is two. For the geometric generic hyperplane of this
embedding the kernel on $\mathbf{Pic}^{\tau}$ is nontrivial and has
one-dimensional tangent space. This conclusion concerns the group
scheme kernel; it does not give a counterexample on geometric points
to A1 or A2.
\end{theorem}

We give all the defining data, including a finite algorithm specifying
every coefficient of the sections used in the construction. The final
ground field is an explicitly specified algebraic closure of a purely
transcendental extension of $\mathbb F_2$. No numerical search for smooth
branch divisors is substituted for the proof of smoothness.

\paragraph{1. The curve and the extension.}
Let $C$ be the smooth projective model over $\mathbb F_2$ of
\[
 y^2+y=x^{15}.
\]
The Artin--Schreier polynomial is irreducible over
$\overline{\mathbb F}_2(x)$: a rational function whose Artin--Schreier
difference has a pole has an even pole order there, whereas $x^{15}$
has an odd pole order at infinity. There is one point $\infty$ over
infinity, with $v_\infty(x)=-2$ and $v_\infty(y)=-15$.
Write $P_0=(0,0)$, $P_1=(0,1)$, and put $t=y/x^8$. Then
\[
 \operatorname{div}(t)=7P_0+\infty-8P_1,\qquad
 x+t^{-2}=t^{-1}x^{-7},\qquad dt=x^6\,dx.
\]
In particular $t$ is a uniformizer at $\infty$, and
$v_\infty(x+t^{-2})=13$. The affine curve is smooth and $dx$ is
nowhere zero there, since the derivative of its equation with respect
to $y$ is one. The displayed formula for $dt$ therefore gives
\[
 \operatorname{div}(dx)=12\infty,\qquad g(C)=7.
\]
These formulas also verify the extension and smoothness calculations
below directly, without an appeal to an unspecified Tango curve.

Set $N=\OO_C(2\infty)$ and $D=6\infty$, so that
$\OO_C(D)=N^3$. Cover $C$ by
$U=C\setminus\{\infty\}$ and $V=C\setminus\{P_0,P_1\}$.
Define a rank-two bundle $E$ by frames $(e_0,e_1)$ on $U$ and
$(v_0,v_1)$ on $V$, with transition
\begin{equation}\label{a:explicit-transition}
 v_0=e_0,\qquad v_1=t^{-6}(e_1-t^{-1}e_0).
\end{equation}
It fits into $0\to\OO_C\to E\to N^3\to0$.
The extension is nonsplit. Indeed a splitting would express $t^{-1}$
as a function regular on $U$ plus a function divisible by $t^6$ and
regular at $\infty$. The function on $U$ would have its sole pole,
of order one, at $\infty$. This is impossible: the ring of $U$ is
$\mathbb F_2[x,y]/(y^2+y-x^{15})$, and its possible pole orders are
$2a$ and $15+2a$, $a\geq0$.

The Frobenius pullback splits. On $U$ its second summand is generated
by $e_1^2+xe_0^2$; on $V$ the corresponding expression is
\[
 v_1^2+t^{-12}(x+t^{-2})v_0^2.
\]
The coefficient in this expression is regular on $V$ and has order
one at $\infty$. This gives a splitting of the quotient $N^6$ and
also a map $N^6\to\operatorname{Sym}^2E$.

\paragraph{2. The smooth cyclic-cover surface.}
Use the quotient convention $P=\mathbf P_C(E)$, so that
$\pi_*\OO_P(j)=\operatorname{Sym}^jE$ for $j\geq0$.
Let $E_\infty$ be the section $e_0=0$; its divisor sheaf is
$\OO_P(1)$ and $E_\infty^2=6$. The map just constructed defines
a divisor
\[
 C''\in\big|\OO_P(2)\otimes\pi^*N^{-6}\big|.
\]
On $U$, in the affine chart $e_0\ne0$ with $w=e_1/e_0$, its
equation is $w^2+x=0$. It is disjoint from $E_\infty$.
It is smooth on this chart since $dx$ never vanishes on $U$.
Over $\infty$ its equation is
$w_V^2+t^{-12}(x+t^{-2})=0$; the second term has order one,
so the same smoothness assertion holds there. Thus $C''$ and
$E_\infty$ are disjoint smooth divisors.

Put $M=\OO_P(-1)\otimes\pi^*N^2$. Then
$M^{-3}=\OO_P(E_\infty+C'')$. Let $\sigma$ be the product of
the two defining sections, and form the cyclic cover
\[
 \psi:S=\operatorname{Spec}_P
       (\OO_P\oplus M\oplus M^2)\longrightarrow P,
 \qquad M^3\xrightarrow{\sigma}\OO_P.
\]
This is finite flat of degree three. Locally at the branch locus it
has equation $z^3=u$ with $u$ a regular parameter. Since three is
invertible in characteristic two and the branch components are smooth
and disjoint, $S$ is smooth. It is projective because it is finite over
$P$. Moreover $H^0(S,\OO_S)=\mathbb F_2$: the two nonconstant
summands have negative degree on the fibres of $\pi$. This remains
true after extension of the ground field, and smoothness implies
geometric integrality. Write $f=\pi\psi$ and let $T$ be the
ramification section over $E_\infty$. Then
\[
 \psi^*E_\infty=3T,\qquad T\simeq C,\qquad T^2=2.
\]
On the dense affine chart used above all equations are simply
\begin{equation}\label{a:surface-affine}
 y^2+y=x^{15},\qquad z^3=w^2+x.
\end{equation}

\paragraph{3. An ample generated bundle with the required cohomology.}
Define
\[
 Q=\OO_S(2T)\otimes f^*N.
\]
The bundle $E$ is nef: after pullback to any smooth projective curve,
a quotient line bundle receives either a nonzero map from $\OO$ or
factors through the positive-degree quotient $N^3$; in either case
its degree is nonnegative. Thus $\OO_P(1)$ is nef. The line bundle
$\OO_P(2)\otimes\pi^*N^3$ has positive intersection with every
curve and positive self-intersection: on a fibre its degree is two,
and on a horizontal curve its base factor has positive degree.
The Nakai criterion makes it ample. Since
\[
 Q^3=\psi^*(\OO_P(2)\otimes\pi^*N^3),
\]
$Q$ is ample as well.

For an integer $a$, the local equation $z^3=u$ at $T$ gives
\begin{equation}\label{a:cyclic-floor}
 \psi_*\OO_S(aT)=
 \bigoplus_{i=0}^2 M^i\bigl(\lfloor(a+i)/3\rfloor E_\infty\bigr).
\end{equation}
Indeed the $i$th eigensummand consists locally of
$z^i u^{-\lfloor(a+i)/3\rfloor}$ times regular functions.
For $a=2$, this formula and projection to $C$ give
\[
 f_*Q=N\oplus N^3.
\]
Both summands are generated by global sections: their bases are
$(1,x)$ and $(1,x,x^2,x^3)$, respectively. The first summand generates
$Q$ away from $T$. The second has valuation $1-3+2=0$ along $T$
and generates there. Hence $Q$ is globally generated and
\[
 h^0(S,Q)=6;\qquad
 1,x,z,zx,zx^2,zx^3
\]
are its sections in the affine trivialization of
\eqref{a:surface-affine}.

For $a=-2$, formula~\eqref{a:cyclic-floor} instead gives
\[
 \psi_*Q^{-1}=
 \bigl(\OO_P(-1)\otimes\pi^*N^{-1}\bigr)
 \oplus\bigl(\OO_P(-2)\otimes\pi^*N\bigr)
 \oplus\bigl(\OO_P(-2)\otimes\pi^*N^3\bigr).
\]
Relative duality, with the convention stated above, says
$R^1\pi_*\OO_P(-2)=(\det E)^{-1}=N^{-3}$.
All zeroth direct images in the preceding display vanish, so Leray
gives
\begin{equation}\label{a:Q-cohomology}
 H^1(S,Q^{-1})=H^0(C,N^{-2}\oplus\OO_C)=\mathbb F_2.
\end{equation}
Proposition~\ref{a:raynaud-all}, applied with $a=2n$, $b=n$,
$p=2$ and $\ell=3$, also gives
\[
 H^1(S,Q^{-n})=0\quad(n\geq2).
\]
For later dimension checks, $R^1f_*\OO_S=N$, whence
$h^1(S,\OO_S)=7+h^0(C,N)=9$. Intersection and duality give
$Q^2=16$, $K_S=f^*N^5$, $K_SQ=20$, and $\chi(\OO_S)=-2$.
The canonical formula follows as well from
$K_P=\OO_P(-2)\otimes\pi^*(K_C\otimes\det E)$ and
$K_S=\psi^*(K_P\otimes M^{-2})$.

\paragraph{4. A specified very ample power and its complete basis.}
Set $B=\OO_P(6)\otimes\pi^*N^9$. Then
\[
 Q^9=\psi^*B,\qquad
 \psi_*Q^9=B\oplus BM\oplus BM^2
   =\bigoplus_{i=0}^2\OO_P(6-i)\otimes\pi^*N^{9+2i}.
\]
The bundle $\operatorname{Sym}^a E\otimes N^b$ has a filtration
with quotients $N^{b+3j}$, $0\leq j\leq a$.
For $(a,b)=(6,9),(5,11),(4,13)$ all these quotients, even after
subtracting any effective divisor of degree two on $C$, have degree
at least sixteen, greater than $2g(C)-2=12$.
Thus their first cohomology vanishes, as does that of the filtered
bundle. Restriction of sections is consequently surjective over every
base subscheme of length at most two.

It follows that $B$ is very ample on $P$. To see every infinitesimal
case explicitly, a length-two subscheme of $P$ either lies over a
length-two subscheme of $C$, or lies in one fibre. In the first case
use the just-proved surjectivity and relative very ampleness of
$\OO_P(6)$; in the second, use $\OO_{\mathbf P^1}(6)$.
The same argument with length one proves that $BM$ is globally
generated. These facts show that $\psi^*B$ is very ample on $S$:
the $B$ sections recover the embedded base $P$, while, on an open
where a section of $BM$ is nonzero, its corresponding section of
$\psi^*B$ is a unit multiple of the cyclic generator $z$.
They generate the local algebra of the degree-three cover and hence
give a closed immersion. Equivalently they separate length-two
subschemes whose image is a point as well as those separated on $P$.

Riemann--Roch and the same vanishing give the exact dimensions
\[
 h^0(P,B)=210,\quad h^0(P,BM)=186,\quad
 h^0(P,BM^2)=160,\quad h^0(S,Q^9)=556.
\]
For complete explicit defining coefficients, the ordered basis
$v_1,\ldots,v_{556}$ is the basis over $\mathbb F_2$ in
\texttt{code/A3\_Q9\_basis.json}, generated by
\texttt{code/certify\_A3.py}. Here is its finite mathematical
specification. For $i=0,1,2$ put $a=6-i$, $b=9+2i$, and order the
monomials $w^jx^ry^sz^i$ lexicographically by $(j,s,r)$, where
\[
 0\leq j\leq a,\quad s\in\{0,1\},\quad r\geq0,
 \qquad 2r+15s\leq2b+7a-j.
\]
Use the expansion $x=t^{-2}u$, $y=t^{-15}u^8$, where
$u=1+t^{15}u^{-7}$. Impose that, for each $0\leq h\leq a$,
every negative power of $t$ in
\[
 \sum_{j\geq h,r,s}c_{jrs}{j\choose h}
       t^{2b+7h-j-2r-15s}u^{r+8s}
\]
has coefficient zero. This is precisely the regularity condition
after $w_U=t^6w_V+t^{-1}$ and $z_U=t^4z_V$.
Order the equations by their first occurrence when the ordered
monomials are processed, and take the reduced free-variable kernel
basis with free variables in increasing order, as specified by the
included elimination routine. Concatenate the three bases in
increasing $i$. All sections are regular on the affine curve, on
the chart at $\infty$, and along the section $T$ by their indicated
projective powers. Independence and the Riemann--Roch dimensions prove
that they form the complete basis.

The exact calculation has respectively $357,291,230$ variables,
ranks $147,105,70$, and nullities $210,186,160$. Only four, four,
and three terms beyond the constant in the $t^{15}$ expansion can
contribute to the negative powers, so the calculation is finite and
exact. Its output and time are recorded in
\texttt{code/certify\_A3\_output.json}; the basis SHA256 is recorded
there too. These checks verify coefficient and transition equations,
not a sampled collection of points on the surface.

\paragraph{5. The cover making the first power very ample.}
Let
\[
 k=\overline{\mathbb F_2(\lambda_{ij}\mid1\leq i,j\leq556)},
 \qquad s_i=\sum_{j=1}^{556}\lambda_{ij}v_j.
\]
All varieties and bundles above are now extended to $k$. The matrix
$\lambda$ is invertible, so the $s_i$ are a basis of $H^0(S,Q^9)$.
Their zero divisors are smooth, pairwise transverse, and no three
meet. Here is the exact genericity argument. The smooth-divisor
Bertini theorem applies because $Q^9$ is very ample (Stacks,
Lemma~33.47.3, opened)~\cite{StacksBertini}. After selecting any number of such divisors,
apply the same theorem to their smooth components to make the next
divisor transverse, and avoid their finitely many pairwise
intersection points. Also avoid the linear span of the already
chosen sections. At each stage these are finitely many nonempty
open conditions. Their intersection is nonempty over an infinite
field. They define a nonempty open set of ordered bases; algebraic
independence of all $\lambda_{ij}$ places our basis in this set.

For each $s_i$ take the ninth-root algebra
$\bigoplus_{r=0}^{8}Q^{-r}$ with multiplication determined by
$s_i:Q^{-9}\to\OO_S$, and set
\begin{equation}\label{a:root-algebra}
 \mathcal A=\bigotimes_{i=1}^{556}
       \left(\bigoplus_{r=0}^{8}Q^{-r}\right),\qquad
 \rho:Y=\operatorname{Spec}_S\mathcal A\longrightarrow S.
\end{equation}
This describes the whole finite cover, not a choice of a connected
component. It has degree $9^{556}$. There are tautological sections
$t_i$ of $L=\rho^*Q$ with $t_i^9=\rho^*s_i$.

The surface $Y$ is smooth. At any point of $S$ at most two branch
divisors pass, with independent regular parameters as equations.
After an étale extension extracting ninth roots of units, the cover
is obtained by replacing each such parameter by its ninth power.
The other equations define étale extensions, since nine is
invertible. The resulting local rings are regular of dimension two.
Projectivity follows from finiteness. Moreover
\begin{equation}\label{a:root-directsum}
 \rho_*\OO_Y=
 \bigoplus_{0\leq r_1,\ldots,r_{556}\leq8}
                   Q^{-(r_1+\cdots+r_{556})}.
\end{equation}
Ampleness of $Q$ makes all negative-power global sections zero, so
$H^0(Y,\OO_Y)=k$. Thus $Y$ is connected, and its smoothness over the
algebraically closed field makes it integral.

We verify very ampleness of $L$ without conflating this cover with
the inverse image under a coordinate power map. Use the six pulled
back sections of $Q$ and the $556$ sections $t_i$. They generate $L$
because the old six already do. Let $u$ be one of the old sections.
The open $S_u=S\setminus V(u)$ is affine: it is the intersection of
the $Q^9$ embedding with the affine hyperplane chart $u^9\ne0$.
Since the $s_i$ form a complete basis, its coordinate algebra is
generated over $k$ by $s_i/u^9$. Over this open put $z_i=t_i/u$.
Then
\[
 \OO(Y_u)=\OO(S_u)[z_1,\ldots,z_{556}]
                   /(z_i^9-s_i/u^9)_{i=1}^{556}.
\]
This algebra is generated over $k$ by the $z_i$ alone, because their
ninth powers include all the displayed generators of $\OO(S_u)$.
Consequently the projective morphism defined by the old sections and
the $t_i$ is a closed immersion on each target chart $u\ne0$.
Those charts cover $Y$. The morphism is therefore an immersion,
and properness makes it a closed immersion globally. Hence $L$ is
very ample. Retaining the old sections in this argument also
separates simultaneous multiplication of all $t_i$ by a ninth root
of unity; no torsion ambiguity in the identity $L=\rho^*Q$ remains.

Finally tensor~\eqref{a:root-directsum} by $Q^{-1}$ and use
\eqref{a:Q-cohomology} and the higher negative-power vanishing.
Only the multi-index $(0,\ldots,0)$ contributes to $H^1$, giving
$h^1(Y,L^{-1})=1$. The same decomposition gives
$h^0(Y,L)=6+556=562$ and $h^1(Y,\OO_Y)=9+556=565$.
This completes all geometric and cohomological assertions in
Theorem~\ref{a:very-ample-example}.

\paragraph{Verification and attribution.}
Three other AI agents independently checked the root-cover argument,
the smoothness and connectedness assertions, the cohomology summands,
and the very-ampleness chart proof; their audit notes are retained.
The finite certificate is independent of those prose checks. The
construction uses the classical Raynaud mechanism, but the equations,
specified polarization and simultaneous-root proof above are supplied
in full here. Searches for a published theorem with this exact
very-ample negative-$H^1$ conclusion did not locate one during this
run. No priority claim is made.

\subsection{The stronger generic-kernel assertion of SA14-H1}
\begin{proposition}[R3u/SA14-H1: generic kernel is infinitesimal unipotent; Proved]\label{a:generic-unipotent}
Let $X$ be a normal connected projective variety of dimension at least
two over an algebraically closed field $k$ of characteristic $p>0$.
For a specified very ample sheaf $L$, let $Y$ be the geometric generic
hyperplane section. Then the kernel
\[
 G=\ker\bigl(\mathbf{Pic}^{\tau}_{X/k}\otimes_k\overline{k(|L|)}
          \longrightarrow\mathbf{Pic}^{\tau}_{Y/\overline{k(|L|)}}\bigr)
\]
is finite, connected, and unipotent.
\end{proposition}
\begin{proof}
The finiteness proof in Proposition~\ref{a:generic} applies to normal
$X$: $\mathbf{Pic}^{\tau}$ is proper of finite type, and a
positive-dimensional reduced connected kernel would contain
prime-to-$p$ torsion, contrary to the fundamental-group surjection.

The argument excluding geometric points also works for normal $X$,
with one precision concerning differentials. Associate to a
$p$-torsion line bundle its rational logarithmic differential in
$\Omega^1_{k(X)/k}$, viewed as a regular form on $X_{\mathrm{sm}}$.
If it is zero, the local units describing a trivialization of the
$p$th tensor power have $p$th roots in $k(X)$. Normality implies that
each root and its inverse are regular on the relevant trivializing
open, so the roots trivialize the line bundle. Thus this rational
$d\log$ is injective on $\Pic(X)[p]$.

The incidence argument is carried out over $X_{\mathrm{sm}}$ and the
open locus of hyperplanes smooth at the specified point. Its relative
cotangent sheaf is locally free. A form defined over $k$ that vanishes
on the geometric generic hyperplane therefore vanishes throughout
this integral open incidence locus. At a smooth point where the form
is nonzero, very ampleness lets the tangent hyperplane vary over all
codimension-one subspaces. Some such subspace is not annihilated by
the form, a contradiction. Triviality of the restricted line bundle
implies vanishing of its relative logarithmic differential on the
hyperplane: its global units are constants. As in
Proposition~\ref{a:generic}, every finite-order geometric point of the
Picard scheme is defined over $k$. The argument excludes $p$-torsion,
and taking powers excludes all $p$-power torsion. Consequently $G$
has just its identity as a geometric point and is connected.

It remains to prove unipotence, which does not follow from
connectedness alone. Raynaud's Proposition~6.2.1, p.~50, and
Proposition~6.2.3(1), p.~51, give a natural isomorphism
\[
 \Hom(\mu_p,\mathbf{Pic}_{Z/\Omega})
       \simeq H^1_{\text{ét}}(Z,\ZZ/p)
\]
for a proper scheme $Z$ over an algebraically closed field $\Omega$
\cite{Raynaud1970PicardDuality}. Apply it to $X_\Omega$ and $Y$.
The surjection $\pi_1(Y)\to\pi_1(X_\Omega)$ makes restriction
injective on these cohomology groups. Hence $\Hom(\mu_p,G)=0$.

Here is an elementary finite-group argument completing the proof.
Let $E$ be the maximal étale quotient of the Cartier dual $G^D$.
Its dual $E^D$ is a closed subgroup of the finite connected scheme
$G$, and hence is connected. Over $\Omega$, $E$ is a constant finite
abelian group; connectedness of its diagonalizable dual implies
that it is a $p$-group. If $E\ne0$, a quotient $E\to\ZZ/p$ dualizes
to a subgroup $\mu_p\subset E^D\subset G$, a contradiction.
Thus $G^D$ is connected, so its finite coordinate algebra is local
with residue field $\Omega$.

Finite-dimensional representations of $G$ are modules over
$\mathcal O(G)^*=\mathcal O(G^D)$. That local algebra has only the
trivial simple module. Every such representation therefore has a
finite filtration with trivial one-dimensional factors. Applying
this to the faithful regular representation embeds $G$ in a group
of upper unitriangular matrices. This is unipotence.
\end{proof}

For normal $X$, the tangent-kernel computation also remains valid.
The geometric generic hyperplane is geometrically integral, so
$H^0(X_\Omega,\mathcal O)=H^0(Y,\mathcal O)=\Omega$; the sequence
$0\to L^{-1}\to\mathcal O_X\to\mathcal O_Y\to0$ identifies the tangent
space of $G$ with $H^1(X,L^{-1})\otimes_k\Omega$. A finite connected
group scheme is nontrivial exactly when its tangent space is nonzero,
by Nakayama applied to its local coordinate algebra.

\subsection{A conditional source example has incompatible hypotheses}
\begin{proposition}[A-CZ346ii: incompatible source hypotheses; Proved]\label{a:cz-vacuity}
Let $i:X\to Y$ be a morphism of smooth schemes over a perfect field
$k$ of characteristic $p>0$. If
\[
 H^0(Y,\Omega^1_{Y/k})\longrightarrow H^0(X,\Omega^1_{X/k})
\]
is injective, then
$H^1_{\mathrm{fppf}}(Y,\alpha_p)\to
 H^1_{\mathrm{fppf}}(X,\alpha_p)$ is injective.
\end{proposition}
\begin{proof}
There is a natural injective additive map
\[
 H^1_{\mathrm{fppf}}(Z,\alpha_p)
       \longrightarrow H^0(Z,\Omega^1_{Z/k})
\]
for every smooth $k$-scheme $Z$; its scalar convention is Frobenius
semilinear, which does not affect injectivity. On affine opens,
the sequence $0\to\alpha_p\to\mathbf G_a\xrightarrow{F}\mathbf G_a\to0$
writes an $\alpha_p$-torsor as $z^p=a_j$. On overlaps the functions
$a_j$ differ by $p$th powers, so $da_j$ glue. If the glued form is
zero, smoothness over the perfect field implies locally
$a_j=b_j^p$. For instance, this follows first in the function field
from $\ker(d)=k(Z)^p$, and then in local rings from normality.
The torsor is therefore Zariski locally trivial. Its local
trivializations glue uniquely, since $\alpha_p$ has no nonzero
sections on a reduced scheme. This proves injectivity.
The construction commutes with pullback. Applying it to $X$ and
$Y$ proves the assertion.
\end{proof}

Cotner--Zavyalov, Example~3.3.6, supplies a nonzero kernel on
$H^1(\alpha_p)$, while the additional hypothesis of their
Example~3.4.6(ii), printed pp.~1867--1868, asks that $X$ also be
smooth and that restriction on global one-forms be an isomorphism
\cite{CotnerZavyalov2024}. The proposition shows that these
hypotheses cannot hold simultaneously. Thus that conditional
example gives no counterexample, or realizable remaining gate,
for injectivity on $\Pic[p]$. This observation concerns the
auxiliary example only, not their main Lefschetz theorems.

\subsection{A finite determinantal test for special hyperplanes}

\begin{proposition}[A2D; Proved]\label{a:determinantal}
Let $X/k$ be smooth projective and integral of dimension at least two,
where $k$ is algebraically closed of characteristic $p>0$.
Let $L$ be very ample, let $M\ne\OO_X$ satisfy $M^p\simeq\OO_X$,
and put
\[
 a=h^1(X,M\otimes L^{-1}),\qquad b=h^1(X,M),\qquad
 r=h^0(X,L)-1.
\]
On the open subset of $|L|=\mathbb P^r$ parametrizing smooth divisors,
the locus where $M|_Y$ is trivial is finite. It is the restriction of
the maximal-minor locus of the linear multiplication matrix
\[
 H^1(X,M\otimes L^{-1})\otimes\OO_{\mathbb P^r}(-1)
 \longrightarrow H^1(X,M)\otimes\OO_{\mathbb P^r}.
\]
We have $a\leq b$. If this finite locus is nonempty, then $a>0$ and
\begin{equation}\label{a:determinantal-bound}
 b-a\ \geq\ h^0(X,L)-2.
\end{equation}
In dimension at least three, it follows that restriction on geometric
Picard groups is injective for every smooth member of $|L|$ if, for
each nontrivial $M\in\Pic(X)[p]$, either $a=0$ or the reverse strict
inequality to \eqref{a:determinantal-bound} holds.
\end{proposition}

\begin{proof}
The injective logarithmic differential construction of
Proposition~\ref{a:generic} gives a nonzero form
$\eta=d\log(M)\in H^0(X,\Omega_X^1)$.
If $M|_Y$ is trivial, its restriction to $\Omega_Y^1$ is zero.
The conormal sequence and
$H^0(Y,\OO_Y(-Y))=0$ therefore show that
$\eta|_Y=0$ already in $\Omega_X^1|_Y$.
Every such $Y$ is consequently a divisorial component of the zero
scheme of $\eta$. A nonzero section of a vector bundle on an integral
noetherian scheme has only finitely many divisorial components in its
zero scheme. A smooth ample divisor here is connected and hence
integral. There are therefore only finitely many possible $Y$.

A nontrivial numerically trivial line bundle on a projective integral
scheme has no nonzero section: any nonzero effective divisor has
positive intersection with a suitable power of an ample divisor.
The same intersection argument gives $H^0(X,M\otimes L^{-1})=0$.
For a section $h$ cutting out a smooth $Y$, the divisor sequence gives
\[
 H^0(Y,M|_Y)=\ker\bigl(H^1(X,M\otimes L^{-1})
                  \xrightarrow{\,h\,}H^1(X,M)\bigr).
\]
Since $M|_Y$ is numerically trivial, its space of sections has dimension
one exactly when it is trivial, and dimension zero otherwise.
This proves the asserted determinantal description. There are smooth
members outside the finite exceptional locus, so a multiplication map
is injective and $a\leq b$.

At a bad smooth member, the matrix has rank exactly $a-1$.
If $a=0$ there is no such member. Otherwise choose an invertible
$(a-1)$-minor locally, allowing the empty minor when $a=1$.
Elementary row and column operations reduce the maximal-minor ideal
locally to $b-a+1$ functions, the entries of its remaining column.
The bad locus has dimension zero in the smooth-member open subset
of the regular variety $\mathbb P^r$. The height theorem for an ideal
generated by $b-a+1$ elements therefore gives $r\leq b-a+1$,
which is \eqref{a:determinantal-bound}.

In dimension at least three the point kernel is a finite $p$-group
by the already proved A2 reduction. Any nonzero such group contains
an element of order $p$. Applying the preceding exclusion to every
nonzero $p$-torsion line bundle proves the last assertion.
\end{proof}

This supplies a finite linear-algebra test for a fixed pair $(X,L)$
once its $p$-torsion classes and multiplication maps are computed.
The entries depend linearly on the coefficients of the hyperplane.
The full special-hyperplane problem is not settled: the inequality
can hold, and we have not proved that the resulting determinantal
points must lie in the discriminant of singular divisors.

\subsection{The leading Cartier term of a possible point kernel}
\paragraph{A2V (Proved necessary condition; not a solution of A2).}
Let $X/k$ be smooth, projective and integral over an algebraically
closed field of characteristic $p>0$, with $\dim X\geq2$.
Let $Y\subset X$ be a smooth ample divisor and put
$L=\mathcal O_Y(Y)$. Suppose a nonzero closed global one-form
$\eta$ vanishes along $Y$. Let $m\geq1$ be its exact vanishing
order along the integral divisor $Y$. Then $p\mid m$.
If in addition $C\eta=\eta$, write $m=pa$. There is a nonzero
element
\[
 H^0(Y,B_Y^1\otimes L^{-a})
  \simeq
 \ker\bigl(F:H^1(Y,L^{-a})\longrightarrow
                   H^1(Y,L^{-pa})\bigr),                 \tag{A2V.1}
\]
where $B_Y^1=\operatorname{im}(d:F_*\mathcal O_Y\to
F_*\Omega_Y^1)$ has its Frobenius-module structure.
The cohomology group in the Frobenius target is understood
with the corresponding Frobenius scalar structure; equivalently,
the map is $p$-semilinear on the usual $k$-vector spaces.
In particular a nonzero $p$-torsion line bundle on $X$ that
restricts trivially to $Y$ forces this obstruction for some $a\geq1$.

Here is the proof. The leading coefficient of $\eta$
is a nonzero section
\[
 \sigma\in H^0\bigl(Y,\Omega_X^1|_Y\otimes L^{-m}\bigr).
\]
Locally write $\eta=t^m\alpha$, where $t$ cuts out $Y$.
The equation $d\eta=0$ implies
$m\,dt\wedge\alpha=0$ modulo $t$. If $p\nmid m$, the
leading coefficient is purely conormal. The conormal sequence
would put it in
\[
 H^0(Y,L^{-m-1})=0,
\]
contrary to its definition. Thus $m=pa$.

Since $d(t^{pa})=0$, the local form $\alpha$ is closed.
Cartier's rule $C(f^p\alpha)=fC(\alpha)$ and $C\eta=\eta$
give
\[
 C\alpha=t^{(p-1)a}\alpha.
\]
Restriction to the smooth divisor commutes with Cartier, so
the tangential leading coefficient $\beta$ has $C_Y\beta=0$.
It is nonzero: otherwise $\sigma$ would again be a section of
$L^{-pa-1}$, which has no nonzero global section. When local
equations change by $t_i=u_{ij}t_j$, the coefficients change
by $\alpha_i=u_{ij}^{-pa}\alpha_j$. Because the multiplier is
a $p$th power, closedness and exactness are preserved, and the
Frobenius-module interpretation makes these precisely the
transition functions of $B_Y^1\otimes L^{-a}$. The exact
Cartier sequence identifies $\ker C_Y$ on closed forms with
$B_Y^1$.

Finally, tensoring the Frobenius sequence with $L^{-a}$ gives
\[
 0\longrightarrow L^{-a}\longrightarrow F_*L^{-pa}
  \longrightarrow B_Y^1\otimes L^{-a}\longrightarrow0.
\]
Both negative line bundles have no global section, so its
cohomology sequence proves (A2V.1). The two exact sequences
and projection-formula conventions can also be checked in
\cite[Lemmas 9.2 and 9.5, pp.~2106--2107]{DGJ_Langer2011}.
For a point kernel take $\eta=d\log(M)$. It is nonzero and
Cartier-fixed; triviality of $M|_Y$ first makes its tangential
restriction zero, and the conormal sequence with
$H^0(Y,L^{-1})=0$ makes its full restriction zero.

There is also a numerical exclusion that can be checked on the
ambient variety. Since $m\geq p$, this same nonzero form defines
an injection of sheaves
\[
 \mathcal O_X(pY)\hookrightarrow\Omega_X^1.
\]
It is injective because $X$ is integral and the map is nonzero
at its generic point. Taking global sections gives
\[
 h^0(X,\Omega_X^1)\ \geq\ h^0(X,\mathcal O_X(pY)).
                                                        \tag{A2V.2}
\]
Consequently, if $\dim X\geq3$, the reverse strict inequality
excludes every nontrivial point in the restriction kernel by
the already established finite-$p$-group reduction. This bound
does not assert that such an inequality holds for all $X$.

\paragraph{Why Cartier iteration alone stops here.}
For every prime $p$, the affine local model
\[
 U=\operatorname{Spec}k[t,x,(1+t^px)^{-1}],\qquad
 \eta=d\log(1+t^px)=\frac{t^p\,dx}{1+t^px}
\]
is closed and Cartier-fixed and has exact order $p$ along
$t=0$. Its tangential leading term is the nonzero exact form
$dx$. Explicitly,
\[
 \eta=\sum_{n\geq0}(-1)^n t^{p(n+1)}x^n\,dx.
\]
Cartier retains $n=pr+p-1$ and transforms those terms into
$(-1)^r t^{p(r+1)}x^r\,dx$, giving the same series.
Thus one cannot infer divisibility of the vanishing order by
$p^2$, or infinite-order vanishing, from Cartier fixedness
alone. This is an affine local model, not a projective ample
counterexample. Its exact finite coefficient identities for
$p=2,3,5,7$ are recorded in
\texttt{check\_cartier\_leading.py} and the saved JSON output.

\paragraph{A necessary singularity of the associated cover.}
Any actual point-kernel counterexample has a nonnormal
$\mu_p$-torsor along $Y$. Indeed, the torsor has local equation
$z^p=u$, with $u$ a unit and $du/u=\eta$. It is integral:
if $u$ were a $p$th power in the function field, the generic
torsor section would extend over the normal scheme $X$,
trivializing $M$. Since $\eta$ vanishes along $Y$, the
Jacobian criterion makes this hypersurface singular at the
generic point over $Y$. That is a codimension-one point.
The integral hypersurface is Cohen--Macaulay, but fails
regularity in codimension one, so it is not normal. In
particular constructions whose $\mu_p$-cover stays normal
cannot yield the required point kernel. This observation
does not rule out the nonnormal-cover route.

\paragraph{Exact remaining gate.}
The necessary Frobenius kernel on the smooth ample divisor
need not vanish in positive characteristic. Its existence
does not construct a Cartier-fixed form on an ambient smooth
projective variety, and no converse is proved here. The
remaining task is either to exclude that global extension or
to realize it with a fully specified smooth very ample pair.

\subsection{An exclusion from closed global one-forms}
\paragraph{A2W (Proved under the stated differential hypothesis).}
Let $X/k$ be smooth, projective and integral of dimension $n\geq2$,
where $k$ is algebraically closed of characteristic $p>0$.
Suppose every global regular one-form on $X$ is closed:
\[
 d:H^0(X,\Omega_X^1)\longrightarrow H^0(X,\Omega_X^2)
 \quad\hbox{is zero}.
\]
If $L$ is very ample, then
\[
 H^0(X,\Omega_X^1\otimes L^{-1})=0.                 \tag{A2W.1}
\]
Consequently, if $n\geq3$, restriction on geometric Picard
groups is injective for every smooth divisor in $|L|$, using
the established reduction of its kernel to a finite $p$-group.
This is a conditional theorem: the closedness hypothesis is
not asserted for all smooth projective varieties in positive
characteristic.

To prove (A2W.1), suppose that
$0\ne\zeta\in H^0(X,\Omega_X^1\otimes L^{-1})$.
Choose a nonzero $s_0\in H^0(X,L)$ and put $\eta=s_0\zeta$.
For every $s\in H^0(X,L)$, both $s\zeta$ and $\eta$ are
global one-forms and hence closed. On the open set $s_0\ne0$
we therefore have
\[
 0=d\bigl((s/s_0)\eta\bigr)=d(s/s_0)\wedge\eta.
\]
Because $L$ gives a closed embedding, its affine coordinate
ratios $s/s_0$ generate the function field $K=k(X)$ over $k$.
Their differentials consequently span the $n$-dimensional
$K$-vector space $\Omega_{K/k}^1$. For a nonzero one-form
$\eta_K$ on a vector space of dimension at least two, the
linear map $v\mapsto v\wedge\eta_K$ is not zero. The
displayed equality for a spanning set is a contradiction.

For the Picard consequence, a nonzero finite $p$-group kernel
would contain a line bundle $M$ of order $p$. The nonzero
logarithmic form $d\log(M)$ restricts to zero tangentially on
the smooth divisor $Y$. The conormal exact sequence and
$H^0(Y,\mathcal O_Y(-Y))=0$ then make its full restriction
zero, putting it in the group (A2W.1). This proves injectivity.

There is a quantitative necessary condition even without the
closedness hypothesis. If a nonzero closed global form $\eta$
vanishes along a very ample divisor $Y$, write $\eta=s_0\zeta$
as above. The differentials of the global forms $s\zeta$ are
$d(s/s_0)\wedge\eta$. Their values over $K$ span the image
of the wedge map by $\eta$, which has rank $n-1$. It follows
that
\[
 \dim_k\operatorname{im}\bigl(
  d:H^0(X,\Omega_X^1)\to H^0(X,\Omega_X^2)\bigr)
  \geq n-1.                                         \tag{A2W.2}
\]
Thus any hypothetical point-kernel example in dimension
$n\geq3$ must have at least $n-1$ independent nonzero
differentials of global one-forms. In particular it must
have a nonzero Hodge-to-de Rham differential $d_1^{1,0}$.
The argument uses only closedness of the original logarithmic
form for this last necessary condition.

\subsection{Two target-side exclusions for a point kernel}
\paragraph{A2Y (Proved under either stated target hypothesis).}
Let $X/k$ be smooth, projective and integral of dimension
$n\geq3$ over an algebraically closed field of characteristic
$p>0$, and let $Y\subset X$ be a smooth very ample divisor.
Restriction
\[
 \operatorname{Pic}(X)\longrightarrow\operatorname{Pic}(Y)
\]
is injective on geometric line-bundle classes if either
\begin{enumerate}
\item Frobenius is injective on $H^1(Y,\mathcal O_Y)$; or
\item the smooth variety $Y$ has a smooth lifting to $W_2(k)$.
\end{enumerate}
No lifting of $X$ is assumed in the second assertion, and no
lifting of the very ample line bundle on $Y$ is required.
The first assertion includes $H^1(Y,\mathcal O_Y)=0$.
Frobenius injectivity here is the precise cohomological
hypothesis; it is not replaced by a claim about the reduced
Picard variety alone.

Put $L=\mathcal O_Y(Y)$, which is very ample on $Y$.
If the point kernel were nonzero, the finite-$p$-group
reduction gives a nontrivial class of order $p$. By A2V,
there would then be an integer $a\geq1$ and a nonzero class
\[
 \alpha\in H^1(Y,L^{-a}),\qquad F\alpha=0
       \quad\hbox{in }H^1(Y,L^{-pa}).                  \tag{A2Y.1}
\]
We exclude this class under each hypothesis.

First suppose Frobenius is injective on $H^1(Y,\mathcal O_Y)$.
Since $\dim Y\geq2$ and $L^a$ is very ample, Bertini gives
a smooth connected divisor $E\in|L^a|$, cut out by a section
$s$. Projectivity and geometric connectedness give
$H^0(Y,\mathcal O_Y)=H^0(E,\mathcal O_E)=k$. The exact
sequence
\[
 0\longrightarrow L^{-a}\xrightarrow{\,s\,}\mathcal O_Y
  \longrightarrow\mathcal O_E\longrightarrow0
\]
therefore makes multiplication by $s$ an injection
\[
 j_s:H^1(Y,L^{-a})\hookrightarrow H^1(Y,\mathcal O_Y).
\]
The Frobenius square commutes:
\[
\begin{array}{ccc}
 H^1(Y,L^{-a})&\xrightarrow{\ F\ }&H^1(Y,L^{-pa})\\
 {\scriptstyle j_s}\downarrow&&\downarrow{\scriptstyle j_{s^p}}\\
 H^1(Y,\mathcal O_Y)&\xrightarrow{\ F\ }&H^1(Y,\mathcal O_Y).
\end{array}
\]
Its horizontal arrows are $p$-semilinear, or linear when
their targets are given the Frobenius scalar structure.
Equation (A2Y.1) implies $F(j_s\alpha)=0$. The assumed
injectivity of the bottom arrow and then the injectivity of
$j_s$ give $\alpha=0$, a contradiction. No assertion that
$j_{s^p}$ is injective is needed.

Now suppose that $Y$ lifts smoothly to $W_2(k)$. Apply
\cite[Section 11 hypothesis and Corollary 11.2, p.~2116]{DGJ_Langer2011}
with variety $Y$, numerically flat bundle $\mathcal O_Y$,
ample divisor corresponding to $L^a$, and indices $i=0,j=1$.
Its required inequality is
\[
 1<\min(p,\dim Y),
\]
which holds for every prime $p$ and every $\dim Y\geq2$.
It yields $H^1(Y,L^{-a})=0$, again contradicting (A2Y.1).
The proof of that corollary uses the truncated
Deligne--Illusie inequality only in total degree one and
iterates Frobenius until Serre vanishing applies. It does
not require $\dim Y<p$.

\emph{Remaining gate.}
A counterexample to point injectivity must therefore have
noninjective Frobenius on $H^1(Y,\mathcal O_Y)$ and a target
$Y$ that admits no smooth lifting to $W_2(k)$, as well as
the independent ambient restrictions of A2V and A2W. These
are necessary conditions, not an existence construction.

\subsection{Targets covered by recent vanishing theorems}
\begin{proposition}[A1Q: quasi-F-split and Fano targets; Known deduction]
\label{a:known-targets}
Let $X$ be a smooth connected projective $n$-fold over an algebraically
closed field of characteristic $p>0$, and let $Y$ be a smooth effective
ample Cartier divisor. The following statements hold for this specified
divisor, and hence for a smooth hyperplane of a specified embedding.
\begin{enumerate}
\item If $Y$ is quasi-F-split and $n\geq3$, then
$\Pic(X)\to\Pic(Y)$ is injective. If also $n\geq4$, it is an isomorphism.
\item If $n=4$ and $Y$ is Fano, then
$\Pic(X)\to\Pic(Y)$ is an isomorphism, including in characteristic two.
\end{enumerate}
Here quasi-F-split means that, for some integer $r\geq1$, the natural
$\OO_{Y^{(1)}}$-linear map
$\OO_{Y^{(1)}}\to F_*W_r(\OO_Y)/p$ has an
$\OO_{Y^{(1)}}$-linear retraction, with the Frobenius module structure
as in \cite[\S1, p.~1]{RootPetrov2025}.
\end{proposition}
\begin{proof}
Put $d=\dim Y=n-1$ and $A=\OO_Y(Y)$, which is ample.
Petrov's Corollary~4.10(2), printed p.~16 of the opened version,
gives
\[
 H^{d-i}(Y,\omega_Y\otimes A^m)=0
       \qquad (m>0,\ 0<i<d)
\]
in the first case: use differential-form degree $d$, cohomological
degree $d-i$, and ample bundle $A^m$. The sum of the two degrees is
$2d-i>d$. Serre duality identifies the dual of the displayed group
with $H^i(Y,A^{-m})$. Thus all the negative first cohomology groups
vanish for $n\geq3$, and all the negative second groups vanish for
$n\geq4$. Apply Corollary~\ref{a:vanishing}.

For the second case, Kawakami--Tanaka prove Akizuki--Nakano vanishing
for every smooth Fano threefold in every positive characteristic:
\cite[Theorem~7.2, p.~56]{RootKawakamiTanakaII2025}, together with
\cite[Theorem~6.3, p.~50]{RootKawakamiTanakaI2026} for the case
$\Pic(Y)=\mathbf ZK_Y$ with very ample anticanonical bundle.
Take differential-form degree zero, cohomological degree $i=1,2$,
and ample divisor corresponding to $A^m$. Both sums of degrees are
strictly less than three, so $H^i(Y,A^{-m})=0$ for every $m>0$.
Corollary~\ref{a:vanishing} again applies. Neither proof requires
that the ambient variety lift, that the given divisor lift, or that
$p$ exceed its dimension.
\end{proof}

These are target hypotheses. They are not properties established for
an arbitrary smooth ample divisor. In particular this proposition
does not settle A1 or A2. The cited papers also distinguish Fano
threefolds from quasi-F-split varieties: some smooth Fano threefolds
are not quasi-F-split, so the second assertion is not redundant.

\subsection{An adjunction obstruction to a separable-cover construction}
\paragraph{A2Z (Proved exclusion of the stated construction).}
Let $k$ be algebraically closed of characteristic $p>0$, and let
$\pi:X\to B$ be a finite, flat, generically separable morphism of
degree $p$ between smooth projective integral varieties of dimension
$n\geq2$. Let $C\subset B$ be a smooth ample integral divisor, and
assume that its scheme-theoretic inverse image
\[
 Y=\pi^*C
\]
is smooth and integral. Then
\[
 \operatorname{Pic}(C)[p](k)\longrightarrow
 \operatorname{Pic}(Y)[p](k)
\]
is injective. In particular this construction cannot make a
nontrivial $M|_C$, for $M\in\operatorname{Pic}(B)[p](k)$, disappear
on $Y$. This is a restriction on the proposed construction; it does
not prove the general point-injectivity assertion A2.

Suppose that $N$ is a nontrivial order-$p$ line bundle on $C$ and
$\rho^*N$ is trivial, where $\rho=\pi|_Y$. Choose $N^p\simeq
\mathcal O_C$, and let $\tau:T\to C$ be the associated
$\mu_p$-torsor. It is integral. Indeed, if its generic equation
$z^p=u$ split in $k(C)$, the rational root would give a rational
trivialization of $N$ with zero divisor, hence a global
trivialization because $C$ is normal. Its degree is $p$.
A trivialization of $\rho^*N$ can be made compatible with the
chosen $p$th-power trivialization by rescaling by a constant:
$H^0(Y,\mathcal O_Y^*)=k^*$ and $k$ is algebraically closed.
It therefore gives a factorization
\[
 Y\xrightarrow{\nu}T\xrightarrow{\tau}C.
\]
Since both $\rho$ and $\tau$ have degree $p$, $\nu$ is finite
birational. Thus $Y$ is the normalization of $T$. In particular
$\rho$ is purely inseparable.

The torsor $T$ is Gorenstein and
\[
 \omega_T\simeq\tau^*\omega_C.                         \tag{A2Z.1}
\]
For an explicit check, realize it inside the total space of the
line bundle $N$ by its equation $z^p=u$. Adjunction gives
$\omega_T\simeq\tau^*(\omega_C\otimes N^{p-1})$; the
nowhere-zero tautological root trivializes $\tau^*N$, giving
(A2Z.1). This argument applies when $T$ is nonnormal.

Write $\mathfrak c$ for the conductor, viewed as an ideal of
$\mathcal O_Y$. Finite duality identifies
\[
 \omega_Y\simeq\mathfrak c\otimes\nu^*\omega_T.
\]
Here the conductor identification is literal: on an affine
birational extension $A\subset A'$ inside a common fraction field,
\[
 \operatorname{Hom}_A(A',A)
   =\{q:qA'\subset A\}
   =\operatorname{Ann}_A(A'/A).
\]
Both canonical sheaves are invertible, so $\mathfrak c$ is an
invertible nonzero ideal on $Y$. Consequently it is
$\mathcal O_Y(-D)$ for an effective Cartier divisor $D$, and
\[
 \omega_Y\simeq\rho^*\omega_C\otimes\mathcal O_Y(-D).
                                                               \tag{A2Z.2}
\]
The finite-duality input is the $R\operatorname{Hom}$ description
of the dualizing complex under a finite map
\cite[Lemma 47.15.8, with the finite-module description in Section 49.5]{DGJ_StacksFiniteDuality};
taking the common lowest cohomological degree gives the displayed
module identity. The divisor adjunction used here also follows
from \cite[Lemma 48.14.1, Tag 0B4B]{DGJ_StacksAdjunction}.

On the other hand, the determinant of the generically invertible
map $d\pi:\pi^*\Omega_B^1\to\Omega_X^1$ is a nonzero regular
section of $\omega_X\otimes(\pi^*\omega_B)^{-1}$. Let $R$ be
its effective divisor. Hence
\[
 \omega_X\simeq\pi^*\omega_B\otimes\mathcal O_X(R).
                                                               \tag{A2Z.3}
\]
The divisor $R$ contains $Y$. At its generic point the normal
map in the conormal sequences for $Y\subset X$ and $C\subset B$
is an isomorphism, because $\pi^*C=Y$ scheme-theoretically,
whereas the tangential determinant vanishes since $\rho$ is
purely inseparable of degree $p$. Thus the full determinant
vanishes at that generic point.
Applying smooth-divisor adjunction to (A2Z.3) gives
\[
 \omega_Y\simeq\rho^*\omega_C\otimes
                  \mathcal O_X(R)|_Y.                         \tag{A2Z.4}
\]
Comparison with (A2Z.2) says
$\mathcal O_X(R)|_Y\otimes\mathcal O_Y(D)\simeq\mathcal O_Y$.
Intersecting with $\mathcal O_X(Y)|_Y^{\,n-2}$ gives
\[
 0=R\cdot Y^{n-1}+D\cdot (Y|_Y)^{n-2}>0,
\]
a contradiction: $Y$ is ample because it is a finite pullback
of $C$, $R$ is a nonzero effective divisor (indeed $R\geq Y$),
and $D$ is effective. The first summand is strictly positive and
the second is nonnegative. This completes the proof.

\emph{Scope of the obstruction.}
A local separable equation such as
$z^p-s^{p-1}z=b$ can have a smooth total space and purely
inseparable residue extension at $s=0$. The proposition explains
why these local equations cannot be glued into the stated
projective construction with a smooth ample full inverse image
that kills a nontrivial $p$-torsion class. Allowing the
$\mu_p$-torsor on $C$ to be nonnormal does not evade the argument:
its conductor contributes the additional nonnegative term above.
The degree-$p$ hypothesis and the equality $\pi^*C=Y$ are retained;
no assertion about arbitrary higher-degree covers or a single
component of a reducible inverse image is made.

\subsection{An inseparable multisection construction is also excluded}
\begin{proposition}[A2P: the relative Frobenius multisection obstruction]
Let $S$ be a smooth projective integral variety of positive dimension
over an algebraically closed field of characteristic $p>0$.
Use the quotient convention for $P=\mathbf P_S(E)$, where $E$ has
rank two. Let $F_S^*E\twoheadrightarrow N$ be a line quotient.
The relative Frobenius
$P\to\mathbf P_S(F_S^*E)$ has degree $p$ on each projective-line
fibre. Let $Y\subset P$ be the inverse image of the section
defined by that quotient, and suppose $Y$ is smooth and integral.
If $Y$ is ample, then
\[
       \Pic(S)[p](k)\longrightarrow\Pic(Y)[p](k)
\]
is injective. Thus this construction cannot give an A2 example
by normalizing a nontrivial $\mu_p$-torsor on $S$.
\end{proposition}
\begin{proof}
Write $\rho:Y\to S$ and $K=\OO_P(1)|_Y$. The morphism $\rho$
is finite flat and purely inseparable of degree $p$. The section
of $\mathbf P_S(F_S^*E)$ is cut out by its kernel line bundle
$(\det E)^p\otimes N^{-1}$. Pulling back its divisor gives
\[
 \OO_P(Y)=\OO_P(p)\otimes
                 \pi^*((\det E)^{-p}\otimes N),
 \qquad K^p=\rho^*N.                                  \tag{A2P.1}
\]
Put $Q=K^2\otimes\rho^*(\det E)^{-1}$. The projective-bundle
canonical formula and smooth-divisor adjunction give
\[
 \OO_P(Y)|_Y=Q^p,\qquad
 \omega_Y\otimes(\rho^*\omega_S)^{-1}=Q^{p-1}.          \tag{A2P.2}
\]
Indeed $\omega_P=\OO_P(-2)\otimes\pi^*(\omega_S\otimes\det E)$;
substituting (A2P.1) gives exponents $2p-2$ on $K$ and $1-p$
on $\rho^*\det E$, exactly as in (A2P.2).

Suppose a nontrivial $M\in\Pic(S)[p](k)$ becomes trivial on $Y$.
As in A2Z, its integral $\mu_p$-torsor $T\to S$ has degree $p$.
A compatible trivialization of $\rho^*M$ gives a finite birational
map $Y\to T$, so $Y$ is its normalization. The Gorenstein torsor
has $\omega_T\simeq\omega_S|_T$, and finite duality gives
\[
        \omega_Y\otimes(\rho^*\omega_S)^{-1}
                    \simeq\OO_Y(-D)                         \tag{A2P.3}
\]
for an effective conductor divisor $D$. The precise conductor
argument, including nonnormal $T$, was proved in A2Z using
\cite[Lemma 47.15.8]{DGJ_StacksFiniteDuality}.
Equations (A2P.2)--(A2P.3) imply $Q^{p-1}\simeq\OO_Y(-D)$.
If $Y$ is ample in $P$, its normal line $Q^p$ is ample, hence
$Q$ is ample. An ample positive power cannot be the negative of
an effective divisor on the positive-dimensional projective
variety $Y$: intersection with any ample class of complementary
degree is positive on the former and nonpositive on the latter.
This contradiction proves injectivity.
\end{proof}

This is an exclusion for the specified relative Frobenius
multisection. It neither assumes a disjoint infinity section nor
classifies arbitrary smooth ample divisors in projective bundles.
The unrestricted A2 question remains undecided.


% Delegated research by GPT-6 Astra (OpenAI), 20 November 2026.
% This is an INCOMPLETE contribution. Statuses are explicit in results.csv.
\section{Seshadri, henselian comparison, and the structure of the Picard functor}
\label{sec:bce}

\subsection{The distinction between the manuscript and the supplied claims}

At 57:110 the question retains the hypothesis of essential properness of the
\emph{Picard functor}; replacing the representable Picard scheme by its functor
does not discard that properness hypothesis. The supplied Part B1 and Lead H2
discard it. The counterexample below disproves this broader formulation.
At 57:111 the text is unambiguous:
\begin{quote}
``une intersection compl\`ete relative de dim relative $\geq3$''.
\end{quote}
The complete-intersection hypothesis must likewise be retained in the positive
$R_2$ criterion. Finally, 57:314 says
\begin{quote}
``Prouver, si $X\to S$ est normal, que $\underline{\mathrm{Pic}}^0_{X/S}$ est
`ferm\'e' dans $\underline{\mathrm{Pic}}_{X/S}$ mieux, que
$\underline{\mathrm{Pic}}^0_{X/S}$ est propre sur $S$ ??''
\end{quote}
Here normality is a property of the \emph{morphism}: flatness with geometrically
normal fibres. The shorter wording on 57:315 must not silently be substituted
by mere normality of the total space.

\subsection{B1: a proper counterexample in every characteristic}

Let $k$ be algebraically closed, let $R=k[[t]]$, and put
\[
 X=\operatorname{Proj}R[x,y,z,w]/(xy-z(z-tw)),\qquad D=V(x,z).
\]
The equation is primitive over the DVR, and its quotient ring is torsion-free;
thus $X/R$ is flat. It is projective. For $t\ne0$, the change of homogeneous
coordinate $v=z-tw$ identifies its fibre with the smooth quadric $xy=zv$.
The special fibre is the integral quadric cone $xy=z^2$ in $\mathbf P^3$.
It is a hypersurface and hence Cohen--Macaulay, and its only singular point is
the vertex $[0:0:0:1]$, also in characteristic two. It is therefore normal and
$S_2$. These assertions persist after extending $k$.

The divisor $D$ is $\operatorname{Proj}R[y,w]\simeq\mathbf P^1_R$, so it has no
vertical component, and its generic fibre is Cartier on the smooth quadric.
On the chart $w=1$, at the origin, its prime ideal is $(x,z)$ in
\[
 B=\bigl(R[x,y,z]/(xy-z^2+tz)\bigr)_{(t,x,y,z)}.
\]
The ring $B$ is normal: it is a hypersurface with an isolated singularity.
All its defining relations have order at least two, whereas $x,z$ have
linearly independent images in $\mathfrak m/\mathfrak m^2$. Since
$\mathfrak m(x,z)\subset\mathfrak m^2$, the ideal $(x,z)$ needs two generators.
It is not invertible in this local ring. Consequently $D$ is not Cartier.
Taking the product with $\mathbf P^1_R$ also gives relative dimension three.

The sign in the supplied R1 example is correct when $\operatorname{char}k\ne2$:
if $i^2=-1$, then $(z-it)(z+it)=z^2+t^2$, so $(x,z-it)$ is indeed a prime
divisor on $xy-z^2-t^2=0$. The split projective version above removes the
restriction on the characteristic.

\subsection{B2: sufficiency, nonnecessity, and the exact obstruction}

Suppose $S$ is regular and integral, $f:X\to S$ is flat of finite type and a
relative local complete intersection, and its geometric fibres satisfy $R_2$.
The fibres are Cohen--Macaulay, hence normal, and $X$ is a normal absolute
local complete intersection. Let $D$ be a Weil divisor Cartier on $X_\eta$.
At any point of $X$ of codimension at most three lying over a nongeneric point
of $S$, the fibre-local dimension is at most two. The fibre is regular there,
so flatness and regularity of the base make the local ring of $X$ regular.
At points over $\eta$, $D$ is Cartier by hypothesis. Thus the non-Cartier locus
of $D$ has codimension at least four.

If it is nonempty, choose a generic point $x$ of that locus. On
$\operatorname{Spec}\mathcal O_{X,x}$ the divisor is Cartier away from the
closed point. This local ring is a complete intersection of dimension at least
four. Parafactoriality (SGA~2, XI, Theorem~3.13(ii)) makes the corresponding
line bundle on its punctured spectrum trivial. Normality then extends a
principal Weil divisor from the punctured spectrum, forcing $D$ to be Cartier
at $x$, a contradiction.

The $R_2$ condition is not necessary, even within the proper scope of B1.
Over $R=k[[t]]$ in characteristic different from two, compare the projective
families in $\mathbf P^3_R$
\[
 X_1: xy-z^2-tw^2=0,\qquad X_2: xy-z^2-t^2w^2=0.
\]
The only singular point of either special fibre is its cone vertex $w=1$.
At that point the derivative of the first equation with respect to $t$ is
$-1$, so the total space $X_1$ is regular everywhere and locally factorial.
Its special fibre is a normal surface node, not $R_2$, but every divisor on
$X_1$ is Cartier. The family $X_2$ has the non-Cartier horizontal divisor
$(x,z-itw)$, as in R1. Both special fibres are the same quadric cone, and
both geometric generic fibres are smooth projective quadrics. Taking the
product of both families with $\mathbf P^1_R$ gives relative dimension three
and preserves all these properties. Thus even proper relative complete
intersections with the same base and abstract geometric fibre types can
have different answers: those data alone do not determine the Cartier
extension property.

Without the relative complete-intersection assumption, even $R_2$ is not
sufficient. In $\mathbf P^6_R$ with coordinates $[a:b:c:d:e:f:w]$ impose
\[
 \operatorname{rank}\begin{pmatrix}a&b&c\\d&e&f\end{pmatrix}\leq1,
 \qquad a+e=tw.
\]
The generic fibre is the Segre variety $\mathbf P^1\times\mathbf P^2$.
The special fibre is the projective cone over the smooth hyperplane section
$u_0v_0+u_1v_1=0$ of that Segre variety. It is a normal Cohen--Macaulay
threefold with just its vertex singular, hence $R_2$. Flatness follows from
the constant Hilbert polynomial of the hyperplane sections of the projective
cone over the Segre variety (equivalently from the regular sequence defined
by $t$ and the hyperplane equation). At the vertex $w=1$, the completed local
ring of the total space, after eliminating $t=a+e$, is the rank-one
$2\times3$ determinantal cone. This ring is normal. The divisor
$D=V(a,b,c)$ is $\mathbf P^2_R$, with $e=tw$, and is generically Cartier.
Its ideal at the vertex requires three generators, since the determinantal
relations are quadratic. It is not Cartier. The same construction with a
$2\times(m+1)$ matrix, $m\geq2$, has fibres of dimension $m+1$ whose isolated
singularity satisfies $R_m$. No fixed finite Serre regularity condition alone
replaces the complete-intersection hypothesis.

For a fixed family there is an exact, deformation-sensitive criterion. Suppose
$S$ is regular integral and $X/S$ is flat with geometrically integral normal
fibres. Every vertical prime divisor on $X$ is the pullback, with multiplicity
one, of a prime divisor of $S$, and hence is Cartier. Therefore every horizontal
Weil divisor whose generic restriction is Cartier is Cartier if and only if
\[
 \ker\left(\operatorname{Cl}(X)/\operatorname{Pic}(X)
 \longrightarrow \operatorname{Cl}(X_\eta)/\operatorname{Pic}(X_\eta)\right)=0.
\tag{B.1}
\]
Indeed split any Weil divisor into its horizontal and vertical parts; the
vertical part does not change its Cartier obstruction. If $X_\eta$ is locally
factorial, (B.1) is simply local factoriality of $X$. This is an exact criterion
for the given family, not an assertion that the requested classification by
abstract fibres has been produced.


\section{Henselian comparison questions}
\subsection{C1 and C3: precise approximation deductions}

Let $A$ be the henselization at a prime of an algebra of finite type over a
field or an excellent DVR, and let $X$ be proper over $A$.
Artin's Theorem~3.5 \cite{Artin1969}, pp.~31--32, says that
\[
 \operatorname{Pic}(X)\longrightarrow\operatorname{Pic}(X_{\widehat A})
\]
is injective with dense image for the topology of restriction to the closed
fibre thickenings. If $X_{\widehat A}$ is projective, approximate an ample
line bundle on it modulo the maximal ideal. The resulting line bundle on
$X$ is ample on the closed fibre. Openness of relative ampleness makes it
relatively ample on a neighbourhood of the closed point of $\operatorname{Spec}A$;
that neighbourhood is the entire local spectrum. Thus $X/A$ is projective.
This proves C1 under the stated approximation hypotheses in all characteristics.

Elkik \cite{Elkik1973}, III~\S1, Theorem~3, p.~569 and its explicit Picard
corollary p.~573, proves the following stronger version of C3. For every
noetherian henselian pair $(A,J)$, with $\widehat A$ its $J$-adic completion,
\[
 \operatorname{Pic}(\operatorname{Spec}A-V(J))
 \xrightarrow{\sim}
 \operatorname{Pic}(\operatorname{Spec}\widehat A-V(J\widehat A)).
\tag{C.1}
\]
Taking $J$ maximal proves C3, without an isolated-singularity or characteristic
zero assumption.

For C2 in characteristic zero, write $\mathfrak m$ for the maximal ideal and
suppose $H^1(X,\mathcal O_X)$ has finite length. The proof of
\cite[Proposition~1.2]{PerlingSchroer2017} shows that
$H^1(X,\mathfrak m^n\mathcal O_X)=0$ for $n\gg0$: the graded direct sum is
finite over the Rees algebra, its sufficiently high graded terms are obtained
by multiplication by $\mathfrak m$, and each such term has finite length.
In characteristic zero formal logarithm identifies
$1+\mathfrak m^n\mathcal O_{\widehat X}$ with
$\mathfrak m^n\mathcal O_{\widehat X}$. Formal functions therefore gives
$H^1(\widehat X,1+\mathfrak m^n\mathcal O_{\widehat X})=0$ for large $n$.
Hence $\operatorname{Pic}(\widehat X)\to\operatorname{Pic}(X_n)$ is injective.
Approximate a given formal line bundle at this one finite level using Artin's
theorem; injectivity then identifies its approximation with the original
formal bundle. This proves C2 in characteristic zero.
Formal logarithm does not apply to an arbitrary maximal ideal in mixed
characteristic. The arguments below replace it first for a DVR and then
in arbitrary dimension. Perling--Schr\"oer's two-cohomology finiteness
theorem will be used only on a suitable auxiliary modification.

\subsection{C2 over a henselian DVR, in every characteristic}

\paragraph{Theorem C2.D.}
Let $R$ be an excellent henselian DVR for which Artin approximation applies,
and let $X$ be any proper $R$-scheme. If $H^1(X,\mathcal O_X)$ has finite
length, then $\operatorname{Pic}(X)\to\operatorname{Pic}(X_{\widehat R})$
is an isomorphism. No flatness of $X$ and no finite-length hypothesis on
$H^2(X,\mathcal O_X)$ are needed.

First suppose $X/R$ flat, and write $\pi$ for a uniformizer. Work after
completion to prove finite determination. Choose $c,d\ge0$ such that
$\pi^c H^1(X,\mathcal O_X)=0$ and $\pi^d$ kills the torsion submodule of
$H^2(X,\mathcal O_X)$; these exist by proper coherent cohomology finiteness.
Set $e=c+d$. The cohomology sequence for multiplication by $\pi^n$ gives
\[
 0\longrightarrow H^1(\mathcal O_X)/\pi^n
 \longrightarrow H^1(\mathcal O_{X_n})
 \longrightarrow H^2(\mathcal O_X)[\pi^n]\longrightarrow0,
 \qquad X_n=V(\pi^n).
\]
Consequently $\pi^e H^1(\mathcal O_{X_n})=0$ for every $n$.
If $L|_{X_n}\simeq\mathcal O_{X_n}$ with $n>e$, the analogous sequence
injects $H^1(L)/\pi^n H^1(L)$ into $H^1(\mathcal O_{X_n})$. Therefore
$\pi^e H^1(L)=\pi^nH^1(L)$, and Nakayama applied to the finite module
$\pi^eH^1(L)$ gives $\pi^eH^1(L)=0$.

Let $a_n=1$ be the trivializing section of $L|_{X_n}$. In the multiplication
exact sequences, the obstruction to lifting its restriction $a_{n-e}$ to
$H^0(X,L)$ is $\pi^e$ times the obstruction for $a_n$; it is zero. Thus
$a_{n-e}$ lifts. Its reduction on the closed fibre is a generator, so the
lift is a nowhere-vanishing section on all of $X$: its failure locus is
closed, and its proper image in the local base would meet the closed point.
Hence $L\simeq\mathcal O_X$. This proves injectivity of
$\operatorname{Pic}(X_{\widehat R})\to\operatorname{Pic}(X_n)$ for $n>e$.
Artin's density theorem now proves surjectivity of completion on Picard
groups, as in C1; its injectivity was already proved there.

For a general proper $X/R$, let $I\subset\mathcal O_X$ be its coherent
$\pi$-power torsion ideal and put $X'=V(I)$. Then $X'$ is flat over $R$.
For $N$ sufficiently large, Artin--Rees and bounded torsion give
$I\cap\pi^N\mathcal O_X=0$. Thus
\[
 \mathcal O_X\simeq
 \mathcal O_{X'}\times_{\mathcal O_{X'_N}}\mathcal O_{X_N}.
\]
The generic fibres of $X$ and $X'$ agree, so $H^1(X',\mathcal O_{X'})$
is a finite torsion $R$-module and the flat case applies to $X'$.
An invertible sheaf on the completed $X$ restricts to one on completed
$X'$, which algebraizes, and one on $X_N$, which is already defined over
$R$. Their gluing isomorphism on $X'_N$ is also defined over $R$, since
$R/\pi^N=\widehat R/\pi^N$. Invertible modules patch across this fibre
product (one of the two maps is surjective), giving a line bundle on $X$.
This completes the proof.

The one-parameter proof uses multiplication by a uniformizer. The following
argument turns an arbitrary maximal ideal into an invertible ideal by a
blowup and supplies the additional descent step.

\subsection{C2 in arbitrary dimension and characteristic}

\paragraph{Theorem C2.G.}
Let $A$ be an excellent henselian local ring, $f:X\to\operatorname{Spec}A$
a proper morphism, and suppose $H^1(X,\mathcal O_X)$ has finite length.
Then
\[
 \operatorname{Pic}(X)\xrightarrow{\sim}
 \operatorname{Pic}(X_{\widehat A}).                 \tag{C2.G.1}
\]
In particular this includes every henselization in Part C, in positive,
zero, and mixed characteristic, without assuming $X/A$ flat.

\paragraph{Proof.}
Artin approximation gives injectivity and density as above; for the stated
henselizations this is Artin's Theorem~3.5, and for general excellent
henselian rings the same argument uses Popescu approximation, as explained
in \cite[Remark~1.3]{PerlingSchroer2017}. We prove the required finite
determination after completion. We use formal functions in the precise
form of Stacks tag~02OB(3): for a proper $A$-scheme $W$ and every fixed
$k$, there is $N_0$ such that for $N\geq N_0$,
\[
 \operatorname{Im}\bigl(H^0(W_N,\mathcal O)\to H^0(W_k,\mathcal O)\bigr)
 =\operatorname{Im}\bigl(H^0(W,\mathcal O)\to H^0(W_k,\mathcal O)\bigr),
 \quad W_n=V(\mathfrak m^n\mathcal O_W).             \tag{C2.G.2}
\]
An element on the right lifting a unit on the closed fibre is a global unit:
its nonunit locus is closed with proper image in the local base, hence
would meet the closed fibre if it were nonempty.

First remove the coherent ideal $I\subset\mathcal O_X$ of sections killed
by a power of $\mathfrak m$. A uniform power kills $I$, and Artin--Rees
gives $I\cap\mathfrak m^N\mathcal O_X=0$ for $N$ sufficiently large.
With $X^\circ=V(I)$ we have the fibre-product algebra
\[
 \mathcal O_X=
 \mathcal O_{X^\circ}\times_{\mathcal O_{X^\circ_N}}\mathcal O_{X_N}.
\tag{C2.G.3}
\]
The cohomology sequence and the fact that all cohomology of $I$ is killed
by a power of $\mathfrak m$ show that $H^1(X^\circ,\mathcal O)$ still has
finite length. If (C2.G.1) holds for $X^\circ$, an arbitrary completed line
bundle patches from its algebraized restriction to $X^\circ$ and its
restriction to $X_N$. Their gluing isomorphism is defined on
$X^\circ_N$, which is unchanged by completion. Invertible modules patch
over a fibre-product ring with a surjective side map (Stacks
Lemma~15.6.9, also used in tag~0DVP). It suffices therefore to handle $X$
with no $\mathfrak m$-power torsion in its structure sheaf.

From now on $A$ is complete and $X$ has this property. Put
\[
 Z=\operatorname{Bl}_{\mathfrak m}(\operatorname{Spec}A),\qquad
 X'=\operatorname{Bl}_{\mathfrak m\mathcal O_X}(X),\qquad
 \pi:X'\to X.
\]
The blowup makes $\mathfrak m\mathcal O_{X'}$ an invertible ideal defining
an effective Cartier divisor (Stacks tag~02OS), and its universal property
(tag~0806) gives a proper map $X'\to Z$. Take its Stein factorization
\[
 X'\xrightarrow{h}T\longrightarrow Z.
\]
Here $T\to Z$ is finite, $T/A$ is proper, and
$h_*\mathcal O_{X'}=\mathcal O_T$ (tag~03H0). The ideal
$J=\mathfrak m\mathcal O_T$ is invertible and locally generated by a
nonzerodivisor: a generator acts injectively on $\mathcal O_{X'}$, and
therefore on its pushforward.

Away from $V(J)$, the map $X'\to Z$ is the original $X$ over the punctured
base. Write $g:T\to Z$ for the finite morphism. Since $g_*$ is exact and
faithful on quasi-coherent modules, the identity
$g_*R^1h_*\mathcal O_{X'}=R^1(g\circ h)_*\mathcal O_{X'}$ and the
finite-length hypothesis imply that $R^1h_*\mathcal O_{X'}$ is supported
on $V(J)$, and hence is killed
by $J^c$ for some $c$. The $J$-power torsion in
$R^2h_*\mathcal O_{X'}$ is killed by $J^d$ for some $d$.
Set $e=c+d$.

The argument of Theorem C2.D applies on every affine chart of $T$ where
$J=(a)$, after $a$-adic completion. It only requires that $a$ be a
nonzerodivisor on $X'$ and that $a^cH^1(\mathcal O)=0$; the base need not
be a DVR. The use of Nakayama is valid because $a$ belongs to the Jacobson
radical of its adic completion. Thus, if a line bundle $L$ on $X$ has a
chosen trivialization on $X_N$ with $N>e$, its pullback has, locally on
the formal completion of $T$, generators agreeing with that trivialization
modulo $J^{N-e}$.

The ratios of these generators descend to units on formal $T$: formal
functions and $h_*\mathcal O_{X'}=\mathcal O_T$ identify the completed
pushforward of the structure sheaf with $\mathcal O_{\widehat T}$.
They define a formal line bundle $M$ on $T$ with
$h^*M=(\pi^*L)^{\wedge}$. Moreover $M$ is trivial modulo $J^{N-e}$.
Indeed the local multiplication sequence gives an injection
\[
 \mathcal O_T/J^r\longrightarrow h_*\mathcal O_{X'_r};
\]
so a descended ratio equal to one on $X'_r$ is already one on $T_r$.

Off the closed point of $\operatorname{Spec}A$, the scheme $T$ is the
finite Stein space $\operatorname{Spec}(f_*\mathcal O_X)$. Hence
$H^1(T,\mathcal O_T)$ and $H^2(T,\mathcal O_T)$ both have finite length.
Formal GAGA first algebraizes $M$ on the proper $A$-scheme $T$.
Apply \cite[Proposition~1.2 and the proof of Theorem~1.1]{PerlingSchroer2017}
with the structure sheaf: for some $q$, a formal line bundle on $T$ trivial
on $T_q$ is trivial. Choose $N\geq q+e$. Then $M$ is trivial, and formal
GAGA over complete $A$ implies $\pi^*L\simeq\mathcal O_{X'}$.
Notice that finite length of $H^2$ has been used only on this auxiliary
$T$, where it follows from the construction, not on the original $X$.

It remains to descend triviality at finite order. Let
$Y=\operatorname{Spec}_X(\pi_*\mathcal O_{X'})$ be the finite Stein space
of $\pi$. The inclusion $\mathcal O_X\hookrightarrow\mathcal O_Y$ is
injective: its kernel is zero off the closed fibre, and $X$ has no
$\mathfrak m$-power torsion. Also $Y\to X$ is an isomorphism off that
fibre. The map $\operatorname{Pic}(Y)\to\operatorname{Pic}(X')$ is
injective by the projection formula and the defining Stein equality.
Therefore $L|_Y$ is trivial.

Let $C\subset X$ and $C'\subset Y$ be the conductor subschemes. Their
ideals contain $\mathfrak m^k$ for some $k$, and
$\mathcal O_X=\mathcal O_Y\times_{\mathcal O_{C'}}\mathcal O_C$.
A trivialization of $L|_Y$ and the chosen trivialization of $L|_{X_N}$
differ on $Y_N$ by a unit $u_N$. Increase $N$ to satisfy (C2.G.2) for
$W=Y$ and this $k$. Then $u_N|_{Y_k}$ lifts to a global function on $Y$,
which is a unit by the observation following (C2.G.2). Adjusting the
trivialization on $Y$ by that unit makes it agree with the one on $C$
after pullback to $C'$. Fibre-product patching gives $L\simeq\mathcal O_X$.

All bounds depend only on $X/A$, not on $L$. This proves finite
determination in the case with no closed-fibre torsion. Artin's density
then gives algebraization in that case, and (C2.G.3) gives it in general.
Together with injectivity, this proves (C2.G.1).

\subsection{The characteristic-$p$ local NS comparison: a discrete obstruction}

The original setting of 57:184--194 is the local ring at the vertex of an
affine cone over a smooth projective variety $V$, with its given very ample
bundle $\mathcal O_V(1)$. In that setting the characteristic-$p$ assertion
has a direct proof in every dimension. The blowup is the total space of
$\mathcal O_V(-1)$, and its completion along the zero section has a projection
back to $V$. Thus restriction to $V$ on Picard groups is split surjective.
The kernel has the filtration arising from $1+I^n$, with successive quotients
subquotients of $H^1(V,\mathcal O_V(n))$. Both $H^1$ and $H^2$ of these sheaves
vanish for $n\gg0$, so the formal Picard group is determined at a finite level.
In characteristic $p$ these finitely many successive additive groups are killed
by $p$. The comparison cokernel is consequently killed by some $p^N$.
Removing the exceptional divisor quotients both groups by the same class
$\mathcal O_V(1)$, and the conclusion persists. Elkik then identifies the
henselian and completed punctured Picard groups. This proves the cone-context
claim even at the level of Picard groups modulo bounded $p$-power torsion;
it is not a proof of the extrapolation to all local schemes.

Under the literal comparison of the Picard groups of the three punctured
spectra, H12 is false, already in characteristic three and dimension three.
Let $k$ be algebraically closed of characteristic three and put
\[
 A=\left(k[x,y,z,t]/(xy-z^2-t^2(1+t))\right)_{(x,y,z,t)},
 \qquad U=\operatorname{Spec}A-\{\mathfrak m\}.
\tag{C.2}
\]
This is a normal isolated hypersurface singularity of dimension three.
The polynomial $z^2+t^2(1+t)$ is irreducible in $k[z,t]$: a factorization
would make $1+t$ a square in $k(t)$, impossible because it has an odd valuation
at $t=-1$. Thus $(x)$ is a prime divisor in $A$. The ring $A[1/x]$ is a
localization of $k[x,x^{-1},z,t]$, hence factorial. Nagata's class-group
localization sequence then says that $\operatorname{Cl}(A)$ is generated by
the principal prime $(x)$, so $\operatorname{Cl}(A)=0$. Since $U$ is regular,
$\operatorname{Pic}(U)=0$.

Both the henselization and the completion contain the unique square root
$s=\sqrt{1+t}$ congruent to one. For $i^2=-1$, set $u=z-its$ and $v=z+its$.
Their linear terms are independent because $2i\ne0$. The formal inverse
function theorem gives
\[
 \widehat A\simeq k[[x,y,u,v]]/(xy-uv).
\]
The divisor class group of this completed threefold node is $\mathbf Z$,
generated by $(x,u)$. One direct computation is to blow up its closed point:
the exceptional divisor is $Q=\mathbf P^1\times\mathbf P^1$, its normal bundle
is $\mathcal O_Q(-1,-1)$, and the successive conormal powers are
$\mathcal O_Q(n,n)$. Their $H^1$ and $H^2$ vanish for $n\geq1$, so formal
functions and existence identify the Picard group of the blowup with
$\operatorname{Pic}(Q)=\mathbf Z^2$. Removing the exceptional divisor quotients
by its class $(-1,-1)$, leaving $\mathbf Z$. Formula (C.1) gives the same
punctured Picard group for $A^h$.

There is no connected algebraic part accounting for this extra integer:
the node is Cohen--Macaulay of dimension three, so the tangent space of its
local Picard functor is $H^1(U,\mathcal O_U)=H^2_{\mathfrak m}(A)=0$.
Equivalently the above exceptional-quadric computation has zero connected
Picard part. Thus the comparison is $0\to\mathbf Z\to\mathbf Z$ on the
resulting discrete local divisor groups, and its cokernel is not a $3$-group.
Some modern definitions define the local Picard/NS group using a strict
henselization from the outset; with that convention the notation
$\mathrm{NS}_{\mathrm{loc}}(S)$ already incorporates the second ring, and
one must not interpret it as $\operatorname{Pic}(U)/\operatorname{Pic}^0(U)$.
The ambiguity must be resolved before identifying this with a differently
defined local N\'eron--Severi invariant.


\section{Quotient pathologies for Picard schemes}
\subsection{Quotients: the scope of the product formula}
\paragraph{D0 (Refutation of R6 as stated; correction of scope).}
At 57:286 the transcription explicitly assumes
``supposons que $A_{\Gamma}=0$ \dots et $\Gamma$ opère trivial[ement]
sur $\operatorname{Pic}^{\tau}_{B/S}$'' before writing the full Picard product.
Those hypotheses were omitted from R6 in the research prompt.  Without them
the full product formula is false: take $\Gamma=1$ and $A=B=E$, an elliptic
curve.  The Poincaré line bundle on $E\times E$ is not an exterior product:
its restriction to $E\times\{b\}$ varies with $b$, whereas that of an
exterior product has constant class.  Thus the missing correspondence
factor is essential.  This is a correction of the prompt, not of the
transcribed manuscript.

We use the following form, sufficient for all examples below.  Let $A,B$
be abelian schemes over a locally noetherian base $S$, and let a constant
finite group $\Gamma$ act on $A$ by group automorphisms and on $B$ by
translations through a subgroup $\Gamma_S\hookrightarrow B$.  Set
$Y=B/\Gamma$ and $X=(A\times_S B)/\Gamma$, with $f:X\to Y$.  Then, as fppf
group sheaves, and hence as their representing group algebraic spaces,
\begin{equation}\label{dgj:pic-tau}
 \operatorname{Pic}^{\tau}_{X/S}
 \simeq Y^\vee\times_S(A^\vee)^\Gamma.
\end{equation}
Here $\operatorname{Pic}^{\tau}$ means the union of geometric components
whose classes in the Néron--Severi group are torsion.  The action on the
dual is pullback; replacing all actions by inverse pullback gives the
same fixed subgroup.

For clarity we give the descent argument.  The zero of $A$ induces a
section $j:Y\to X$, so $j^*f^*=1$ splits off $\operatorname{Pic}_{Y/S}$.
An algebraically trivial line bundle on an abelian scheme, rigidified at
zero, has no nontrivial rigidified automorphism.  Consequently a fixed
point of $A^\vee$ has a unique rigidified $\Gamma$-linearization and its
pullback to $A\times B$ descends to $X$, trivially on $j(Y)$.  This
constructs a map to the relative Picard sheaf.  It lands in
$\operatorname{Pic}^{\tau}$: on each geometric fibre, its pullback
under the finite étale quotient $q:A\times B\to X$ is algebraically
trivial, while the norm identity
$\operatorname{Nm}_q(q^*N)=N^{\otimes|\Gamma|}$ shows that its
Néron--Severi class is torsion.  Norm preserves algebraic triviality,
since it carries connected families of line bundles to connected
families and sends the trivial bundle to the trivial bundle.
Conversely, a $\tau$-class pulls back to
an algebraically trivial class on $A\times B$: the Néron--Severi group of
an abelian variety is torsion-free.  The product duality
$(A\times B)^\vee=A^\vee\times B^\vee$ then shows that, after subtracting
the pullback of its restriction to $j(Y)$, it is a constant class from
$A^\vee$.  Descent makes this class fixed.  A discrepancy in descent is
a character of $\Gamma$.  Its restriction along $j$ is the line
bundle descended along $B\to Y$ by that character.  The resulting
map $\Gamma^D\hookrightarrow Y^\vee$ is the injective kernel map in
the dual isogeny sequence for $B\to Y$; thus restriction along $j$
removes the discrepancy even on nonreduced test schemes.  These
constructions commute with base change and
are inverse.  The same argument applies fppf locally to classes rather
than line bundles, so no unmentioned universal line bundle is required.

The quotient map $A\times B\to X$ is finite étale and surjective.  Thus
$X/S$ is smooth and proper.  It is projective: the tensor product of all
translates of an ample bundle, with its permutation linearization,
descends to an ample bundle.  In particular these examples are schemes,
not merely algebraic spaces.

\subsection{Both signs of the Hodge asymmetry}
\paragraph{D1 (Proved: both signs, and variation in a smooth family).}
At 57:287 Grothendieck asks for examples where the difference is
``$>0$, ou $<0$'' and where its terms vary with the parameter.
Let $k$ be algebraically closed of characteristic $2$, let $C,D$ be
elliptic curves, and let $u:D\to C$ be an isogeny.  On $A=C\times D$ put
\[
 g(x,y)=(-x+u(y),y),\qquad \Gamma=\langle g\rangle\simeq\mathbb Z/2.
\]
This is an involution because $g^2=1$.  Its tangent matrix and that of
its pullback action on $A^\vee=C^\vee\times D^\vee$ give
\begin{align*}
 \dim (t_A)_\Gamma&=2-\operatorname{rank}(du),\\
 \dim(t_{A^\vee})^\Gamma&=2-\operatorname{rank}(du^\vee),\\
 \dim(t_{A^\vee})^\Gamma-\dim(t_A)_\Gamma
   &=\operatorname{rank}(du)-\operatorname{rank}(du^\vee).
\end{align*}
Indeed the matrices of $g-1$ and $g^\vee-1$ in characteristic $2$ are
$\left(\begin{smallmatrix}0&du\\0&0\end{smallmatrix}\right)$ and
$\left(\begin{smallmatrix}0&0\\du^\vee&0\end{smallmatrix}\right)$.

Choose an ordinary elliptic curve $E/k$.  For
$u=F_{E/k}:E\to E^{(2)}$ its dual is Verschiebung $V$, with $dF=0$ and
$dV\ne0$.  The displayed difference is $-1$.  For $u=V$ the difference
is $+1$.  An explicit auxiliary elliptic curve is
\[
 Q:\quad y^2+xy=x^3+1,
 \qquad P=(0,1)\in Q[2](k)\setminus\{0\}.
\]
Take $B=Q$ and let the generator act by translation by $P$.  This is a
free action, so the construction above gives a smooth projective
threefold $X$.  Since $Y=Q/\langle P\rangle$ is elliptic,
\[
 h^1(X,\mathcal O_X)=1+\dim(t_{A^\vee})^\Gamma,
 \qquad
 h^0(X,\Omega_X^1)=1+\dim(t_A)_\Gamma.
\]
The first equality is the tangent space of \eqref{dgj:pic-tau}; the
second is étale descent of differential forms and the product formula
for $\Omega^1_{A\times Q}$.  It does not use the false assertion that
wild descent of $H^1(\mathcal O)$ is simply taking invariants.
The two examples have respectively $(h^{01},h^{10})=(2,3)$ and $(3,2)$.

For variation take $R=k[[t]]$ and the elliptic scheme
\begin{equation}\label{dgj:elliptic-family}
 E:\quad y^2+txy+y=x^3.
\end{equation}
Its discriminant is $1+t^3$, a unit.  The generic fibre is ordinary and
the special fibre is supersingular.  Repeating the construction with
$u=F$ or $V$ gives respectively $-1\to0$ and $+1\to0$ under
specialization.  The assertions about ordinarity may also be read
directly from the formal parameter $z=x/y$: multiplication by $2$ has
initial term $tz^2$, while at $t=0$ its first nonzero term has degree
$4$.  The symbolic certificate supplied with this paper checks the
discriminant and $[2](z)=tz^2\pmod {z^4}$.

\subsection{A nonflat fixed subgroup and nonflat $\operatorname{Pic}^{\tau}$}
\paragraph{D2 (Proved: the pathologies requested at 57:288).}
Keep \eqref{dgj:elliptic-family}, put $C=E^{(4)}$, $u=F^2:E\to C$, and
$A_0=C\times E$.  Let $g(x,y)=(-x+F^2y,y)$ and $h=-1$.  They commute,
are distinct involutions, and generate a faithful constant group
$\Gamma\simeq(\mathbb Z/2)^2$.  Scheme-theoretically
\[
 A_0^\Gamma=C[2]\times H,
 \qquad H=\ker(F^2:E\to C)\cap E[2].
\]
There is also a short independent check of the formal calculation.
For $z=x/y$ and $w=1/y$, the Weierstrass equation gives
$w+tzw+w^2=z^3$, hence $w\in z^3R[[z]]$.
Negation sends $(x,y)$ to $(x,y+tx+1)$, so
\[
 [-1](z)=\frac{z}{1+tz+w}=z+tz^2+O(z^3).
\]
Write the commutative formal group law as
$F(U,V)=U+V+cUV+d(U^2V+UV^2)+O((U,V)^4)$.
The identity $F(z,[-1](z))=0$ forces $c=t$.
Its diagonal has no cubic term in characteristic two, so
$[2](z)=tz^2\pmod {z^4}$.  This reasoning is independent of
rational-function duplication and the symbolic certificate.

The formal coordinate computation gives the exact local algebra
\[
 \mathcal O_H=R[z]/(z^4,tz^2).
\]
Here $\ker F^2$ is supported on zero with algebra $R[z]/(z^4)$, and
$[2](z)\equiv tz^2\pmod {z^4}$ supplies the additional equation.
Thus $H$ has generic rank $2$ and special rank $4$; equivalently the
nonzero classes $z^2,z^3$ are killed by $t$.  Since $C[2]$ has constant
rank $4$, $A_0^\Gamma$ has fibre lengths $8$ and $16$ and is not flat.

Take $A=A_0^\vee$ with the dual action, and $B=Q^2$, where the two
generators translate by $(P,0)$ and $(0,P)$.  The resulting $X/S$ is
smooth projective of relative dimension $4$, and
\[
 \operatorname{Pic}^{\tau}_{X/S}=Y^\vee\times A_0^\Gamma
\]
is not flat.  Nonflatness follows from the displayed local algebra:
if the product were flat over $S$, flatness would descend along its
smooth surjective projection to $A_0^\Gamma$.  This answers the existence question within the specified
quotient family itself.

\subsection{No abelian subscheme with the fibrewise Picard varieties}
\paragraph{D3 (Proved: failure of the stated specialization compatibility).}
At 57:288 the precise geometric request is the absence of a
``bon sous-schéma abélien'', followed by
``la formation de Pic variété n'est pas compatible avec la spécialisation''.
Use the same $A_0=C\times E$ and $g$, but only the cyclic group
$\langle g\rangle$, and take $A=A_0^\vee$, $B=Q$.  The fixed group
$G=A_0^g$ is the kernel of the $R$-homomorphism
\[
 C\times E\longrightarrow C,\qquad (x,y)\longmapsto[2]x-F^2y.
\]
On every geometric fibre its reduced neutral component is the image
of
\[
 n:E\longrightarrow C\times E,\qquad y\longmapsto(F^2y,[2]y).
\]
Indeed this is an elliptic image contained in $G$, and $G$ has
dimension one because its projection to $C$ is finite, a base change
of $F^2$.  But $\ker n=H$ has lengths $2$ and $4$ on the generic and
special fibres.

There cannot exist a closed abelian $R$-subscheme $P\subset G$ with
$P_s=(G_s^0)_{\mathrm{red}}$ for every geometric $s$.  If it existed,
$n$ would factor through $P$: this holds generically and extends
because $E$ is $R$-flat and $P$ is closed.  The resulting morphism
$E\to P$ would be a fibrewise isogeny of elliptic schemes and hence a
finite locally free morphism of constant degree.  Its kernel $H$
contradicts that constancy.

The same obstruction applies to
$\operatorname{Pic}^{\tau}_{X/S}=Y^\vee\times G$.  If a closed abelian
subscheme with the fibrewise reduced neutral components existed,
the morphism $Y^\vee\times E\to Y^\vee\times G$ induced by $(1,n)$
would factor through it and be a fibrewise isogeny of abelian schemes.
Its kernel again has nonconstant length.  This is an explicit, precise
failure of formation of the Picard variety to commute with this
specialization; it does not assert that its dimension changes.

\subsection{The product Igusa family and its limitations}
\paragraph{D4 (Product family and modular deformation analysis).}
For arbitrary elliptic schemes $E/S$ in characteristic $2$ and ordinary
$Q/S$ with a chosen nonzero étale $2$-torsion section, the usual Igusa
quotient $X=(E\times Q)/\langle(-1,+P)\rangle$ satisfies
\[
 \operatorname{Pic}^{\tau}_{X/S}=Y^\vee\times E^\vee[2].
\]
It is flat, and the fibrewise Picard variety is $Y^\vee$, compatibly
with base change.  Both first Hodge numbers are $2$.
The torsion Néron--Severi group is $\mathbb Z/2$ on an ordinary fibre
of $E$ and zero on a supersingular fibre, since it is the component
group of $E[2]$.  In the family \eqref{dgj:elliptic-family}, torsion
therefore disappears under specialization.  This product family is
not the full deformation family of the surface.  A first-order
nonflatness assertion is stated in accessible published literature,
but the proof audit below identifies a gap, and D4o subsequently refutes the assertion by computing the full ordinary deformation.  For the ordinary
Igusa surface $X_0$ put $I=H^1(X_0,T_{X_0})^\vee$ and
$R_1=k\oplus I$ with $I^2=0$, and take the first-order deformation
whose Kodaira--Spencer class is the identity of $H^1(X_0,T_{X_0})$.
Here $T_{X_0}=\mathcal O_{X_0}^{\oplus2}$ and
$h^1(\mathcal O_{X_0})=2$, so $\dim_k I=4$.
Chai--Oort, \S3.9(iv), printed p.~35
\cite{DGJ_ChaiOort2017}, states that
$\operatorname{Pic}^{\tau}_{X_1/R_1}$ is not flat.
The source identifies this assertion with Mumford's example.
The surface deformation
can be taken projective: the obstruction to lifting the square of
an ample line bundle is twice its square-zero obstruction, hence
zero in characteristic two, and ampleness over an Artinian base
is detected on the special fibre.

The annotated March 31, 1962 letter, p.~27, footnote~1
\cite{DGJ_GrothMumfordLetters}, gives the more detailed description
that the other connected component factors through a proper closed
subscheme of $\operatorname{Spec}R_1$.  Its auxiliary assertion
about \emph{every} nontrivial line bundle killed by two having
nonzero Hodge Chern class is too broad as printed: pull back a
nontrivial degree-zero two-torsion line bundle from the elliptic
Albanese $Y$.  The pullback remains nontrivial because the zero
section splits pullback, whereas its Hodge Chern class vanishes,
as degree zero line bundles on an elliptic curve have zero Chern
class in $H^1(\Omega^1)$.  We do not copy that overbroad auxiliary
claim into a new proof.  Direct Chai PDF downloads still return 403,
but the complete seven-page v2 has now been opened in web retrieval.
Theorem~2.4, pp.~6--7, and its Remark~(1), p.~7, have been inspected
in full.  The following caution applies to their final inference;
it does not by itself refute the actual Igusa statement; the independent proof D4o does so.

\paragraph{D4-gap (Proved caution on the inspected tangent argument).}
The proof argues that the off-neutral component cannot be a
flat torsor under the neutral component because their absolute
tangent dimensions are five and six, respectively, and their
tangent maps to the base have different ranks.  That inference
is false for torsors under a nonsmooth group.  Here is a complete
example with exactly those dimensions.  Over an algebraically
closed field $k$ of characteristic two put
\[
 R=k[e_1,e_2,e_3,e_4]/(e_i e_j:1\leq i,j\leq4),
 \qquad P=\operatorname{Spec}R[z]/(z^2-e_1).
\]
It is finite free of rank two over $R$, with basis $1,z$.
The unit $1+z$ satisfies $(1+z)^2=1+e_1$; multiplication on this
unit makes $P$ a $\mu_2$-torsor.  Explicitly, with $y^2=0$ the
coordinate of $\mu_2$, the torsor map is
$z_1\mapsto z$, $z_2\mapsto z+y+zy$.
On the two free rank-four coordinate modules over $R$ its
determinant is $1+e_1$, a unit.  This proves the torsor condition,
including faithful flatness.

The maximal ideal of $P$ has cotangent basis $z,e_2,e_3,e_4$,
because $e_1=z^2$.  Thus its tangent dimension is four, and its
tangent map to $\operatorname{Spec}R$ has image the hyperplane
$de_1=0$.  In contrast $\mu_2\times\operatorname{Spec}R$ has
tangent dimension five and a surjective tangent map to the base.
After multiplying both spaces by any smooth elliptic curve,
these dimensions become five and six, exactly as in the inspected
argument.  The first space remains a flat torsor under the
second group scheme.  Moreover $P(R)$ is empty, since every
principal unit of $R$ has square one, while $1+e_1\ne1$.
Consequently absence of an $R$-valued lift also does not imply
that $P$ factors through a proper closed subscheme of the base.
The finite algebra calculation is checked in
\texttt{check\_flat\_torsor.py} and its saved output.

The tangent inference is invalid. The complete ordinary calculation D4o below proves flatness for every Artinian deformation and therefore refutes the literal published first-order claim. This supersedes both the earlier Known classification and the subsequent unresolved verification status. The independent nonflat family D2 is unaffected.

\paragraph{A further source discrepancy and its independent replacement.}
Partsch's published classification \cite{DGJ_Partsch2013} gives, for the ordinary Igusa
case $p=d=2$, $h^1(T_X)=4$ in Proposition~6.1, p.~95,
whereas the Jacobian-deformation hull in Proposition~6.2, p.~99,
is $W[[t_E,q-1]][\sqrt q]$, whose relative tangent dimension is
two.  Theorem~4.4, p.~89, adds $h^1(C,\operatorname{Lie})=1$.
The claimed equality of the resulting tangent dimension with
four on p.~100 therefore needs correction or an additional
qualification.  We have not imported the classification to
close the calculation. The independent calculations D4o and D4s below
supply both miniversal hulls, including the quotient by changes of origin,
and prove flatness of the torsion Picard scheme for every Artinian
deformation. They were separately audited by root and fhi.
The explicit first-order deformation D4a also gives the direct
counterexample D4b to the universal lifting assertions on
pp.~102--106 of that article; the issue is therefore not only a
tangent-count discrepancy.

% BEGIN D4m independent modular degeneration
\paragraph{D4m (Proved: a modular Igusa degeneration with flat
$\operatorname{Pic}^{\tau}$ but incompatible Picard varieties).}
This example is independent of the full versal-ring calculation.
Let $k$ be algebraically closed of characteristic two, let $R=W(k)$,
and take both $E$ and $F$ to be the elliptic curve
\[
 Y^2Z+XYZ=X^3+XZ^2
\]
over $R$, with origin $O=[0:1:0]$ and $P=[0:0:1]$.
The discriminant is $-63$, a unit of $R$, so these are smooth elliptic
schemes. The point $P$ has order two and specializes to a nonzero
point; hence the special fibres are ordinary. Use subscripts to
distinguish $P_E$ and $P_F$.

Let $w\in\mathbb Z_2\subset R$ be the unique root congruent to zero
modulo two of
\[
 2w^2-w+8=0.
\]
It exists uniquely by Hensel's lemma, because the derivative is odd;
also $w\in8R$. The point
\[
 Q_F=[-2:1:w]\in F(R)
\]
lies on $F$, is fixed by elliptic negation
$[X:Y:Z]\mapsto[X:-Y-X:Z]$, and is nonzero on the generic fibre.
Thus it has order two generically and as an $R$-section, whereas
$(Q_F)_k=O$. It is the nontrivial generic section of the multiplicative
two-torsion subgroup. No choice of a root of a nonunit in $R$ is
being assumed here; the displayed Hensel equation defines $Q_F$.

Define two commuting involutions on $E\times_R F$ by
\[
 h(x,y)=(x+P_E,y+Q_F),\qquad
 \tau(x,y)=(-x,y+P_F).
\]
They generate the constant group $\Gamma=(\mathbb Z/2)^2$.
The action is free. For $h$ this follows from its nonzero $E$-translation;
for $\tau$ and $h\tau$ it follows from their $F$-translations
$P_F$ and $P_F+Q_F$, both of which specialize to a nonzero point.
The zero locus of either section is closed in $\operatorname{Spec}R$
and misses its closed point, so it is empty. Consequently
\[
 \mathcal X=(E\times_R F)/\Gamma
\]
is smooth projective over $R$, and its quotient map is finite etale.
For quotient existence one may use Stacks, Lemma~66.14.2, tag07S7;
an invariant ample tensor product descends to show projectivity.
Its special fibre is the ordinary Igusa surface
\[
 \mathcal X_k\simeq
 \bigl((E_k/\langle P_E\rangle)\times F_k\bigr)
 /\langle(-1,+P_F)\rangle.                              \tag{M1}
\]

\emph{The torsion Picard scheme is flat.}
First divide by $h$. Since $h$ is translation by the etale subgroup
$\langle(P_E,Q_F)\rangle$, the quotient
$A=(E\times F)/\langle h\rangle$ is an abelian scheme.
The copy of $F$ embeds in $A$, because that subgroup meets
$0\times F$ trivially. Its quotient is $\overline E=E/\langle P_E\rangle$.
Duality gives
\[
 0\longrightarrow\overline E^{\vee}\longrightarrow A^{\vee}
 \longrightarrow F^{\vee}\longrightarrow0.              \tag{M2}
\]
The linear part $\sigma$ of the remaining deck involution acts by
$-1$ on the first term and by $+1$ on the last term. If
$G=(A^{\vee})^{\sigma^{\vee}}$, the same elementary lifting argument
works over every base change: lift a point of $F^{\vee}$ to $A^{\vee}$,
then subtract half of its failure of invariance in
$\overline E^{\vee}$, taking the half fppf locally. Therefore
\[
 0\longrightarrow\overline E^{\vee}[2]\longrightarrow G
 \longrightarrow F^{\vee}\longrightarrow0               \tag{M3}
\]
is fppf exact. The map $G\to F^{\vee}$ is a finite locally free
torsor of degree four.

For the remaining double cover $A\to\mathcal X$, cyclic descent
gives
\[
 0\longrightarrow\mu_2\longrightarrow
 \operatorname{Pic}^{\tau}_{\mathcal X/R}
 \longrightarrow G\longrightarrow0.                    \tag{M4}
\]
Indeed an invariant algebraically trivial line bundle can be
linearized fppf locally: its chosen isomorphism with the deck
pullback has square a base unit, and adjoining a square root of
that unit corrects the cocycle. The ambiguity is exactly a
character of the cyclic group, namely $\mu_2$. The descended
class lies in $\operatorname{Pic}^{\tau}$ because its square is
the norm of an algebraically trivial bundle, hence is in the
neutral Picard component on each fibre. Conversely every torsion
Neron--Severi class pulls back into $A^{\vee}$, since the
Neron--Severi group of an abelian variety is torsion free.
These are fppf-sheaf statements on arbitrary $R$-schemes.
They also prove representability of the torsion Picard functor
here by finite-flat torsor descent from $G$.
Combining (M3) and (M4),
$\operatorname{Pic}^{\tau}_{\mathcal X/R}\to F^{\vee}$
is finite locally free of degree eight. In particular the full
torsion Picard scheme is flat over $R$.

\emph{The Picard varieties cannot specialize as an abelian subscheme.}
Let $f:F\to\mathcal X$ send $y$ to the class of $(0,y)$ and use
$f(0)$ as base point. For either geometric fibre $s$, every pointed
map $\mathcal X_s\to Z$ to an abelian variety pulls back to a
homomorphism $u_E+u_F:E_s\times F_s\to Z$. Invariance under $\tau$
forces $2u_E=0$ and $u_F(P_F)=0$. Since homomorphisms of abelian
varieties have no additive torsion, $u_E=0$. Invariance under $h$
then forces $u_F(Q_F)=0$. Conversely these conditions suffice for
descent. The Albanese variety and the map induced by $f_s$ are thus
\[
 \operatorname{Alb}(\mathcal X_s)=
 F_s/\langle P_F,Q_F\rangle.
\]
On the generic fibre the two distinct nonzero rational two-torsion
points $P_F,Q_F$ are independent, so the subgroup is all of $F[2]$
and the map has degree four. On the special fibre $Q_F=0$, so its
degree is two. Dualizing, restriction along $f$ on the fibrewise
Picard varieties has degrees
\[
 \deg f_\eta^*=4,\qquad \deg f_k^*=2.                    \tag{M5}
\]
If a closed abelian subscheme of
$\operatorname{Pic}^{\tau}_{\mathcal X/R}$ had these Picard varieties
as all its fibres, restriction along $f$ would be an isogeny of
abelian schemes with the two degrees in (M5). An isogeny of abelian
schemes is finite locally free, so its degree is locally constant.
This contradiction proves the asserted specialization failure
inside an actual modular Igusa degeneration, even though
$\operatorname{Pic}^{\tau}$ itself is flat.

The finite calculation records the discriminant, the projective
point and two-torsion identities, and Hensel approximations. See
\texttt{check\_igusa\_trait.py} and its saved output. The calculation
does not infer flatness from tangent dimensions and does not use
the full-hull assertion.

% END D4m independent modular degeneration

% BEGIN D4o completed ordinary calculation
% Completed ordinary calculation; independently audited by root and fhi.
\subsection*{D4o: the complete ordinary Igusa deformation calculation}
\textbf{Status: Proved.} The calculation was independently audited by root and fhi, including the opened Katz and Mumford--Fogarty sources. It covers ordinary first factors; the supersingular case is completed separately in D4s.

Let $k$ be algebraically closed of characteristic two, let $E_0,F_0$ be ordinary
elliptic curves, and let $P\in F_0[2](k)$ be nonzero. Put
\[
 A_0=E_0\times F_0,\quad \sigma_0=(-1,1),\quad
 \tau_0=t_{(0,P)}\sigma_0,\quad X_0=A_0/\langle\tau_0\rangle.
\]
The base category consists of local Artinian $W(k)$-algebras with residue
field $k$, and all deformations are marked on the closed fibre.

\paragraph{Step 1: lifting the cover and choosing its origin.}
The finite etale cover $A_0\to X_0$, together with its deck involution,
lifts uniquely to every deformation $X/R$, by invariance of the etale site
under a nilpotent immersion. This equivalence also applies to marked
isomorphisms: specify the identity lift on $A_0$ to remove the alternative
deck transformation. The lifted cover $A/R$ is smooth proper. A section
lifting $0\in A_0(k)$ exists by smoothness. It gives a unique abelian-scheme
structure, by Mumford--Fogarty, Proposition~6.15, pp.~124--126
\cite{DGJ_MumfordFogarty1982}, applied successively to small extensions.
This proposition assumes properness, not projectivity.

Write the deck map uniquely as $\tau=t_c\sigma$, where
$c=\tau(0)$ and $\sigma$ fixes the origin. Rigidity makes $\sigma$ a group
automorphism lifting $\sigma_0$, and $\sigma^2=1$. Thus the only additional
equation for the affine map is $(1+\sigma)c=0$.

Conversely every such triple $(A,\sigma,c)$ gives a free involution: its
fixed locus is closed in the proper scheme $A$ and has empty special
fibre, hence is empty. Its quotient is a smooth proper deformation of
$X_0$. At each Artinian stage $A$ is projective: a sufficiently high
power of an ample bundle on $A_0$ lifts through successive small
extensions, because their obstruction groups are killed by two.
An invariant ample bundle is obtained by tensoring a bundle with its
translate, so the quotient is projective too.

\paragraph{Step 2: Serre--Tate parameters and exact point coordinates.}
Choose bases $e_1,e_2$ of $T_2A_0(k)$ and $e_1^\vee,e_2^\vee$ of
$T_2A_0^\vee(k)$, with $e_2/2=P$. Use Katz's convention
\[
 q_{ij}=q(A/R;e_i,e_j^\vee)\in1+\mathfrak m_R.
\]
Katz, Theorem~2.1, pp.~148--149, identifies deformations with this matrix,
identifies duality with transposition, and gives the homomorphism-lifting
criterion. Therefore $\sigma_0$ lifts exactly when
\[
 q_{12}^2=q_{21}^2=1.                                      \tag{I1}
\]

Here is an $R$-point interpretation that avoids any unstated surjectivity
of a torsion subgroup. Choose a compatible system
$P_n\in F_0[2^{n+1}](k)$ with $2^nP_n=P$. For $n$ sufficiently large,
$[2^n]$ kills the formal group $\widehat A(\mathfrak m_R)$.
Choose an arbitrary lift $\widetilde P_n\in A(R)$ and put
\[
 c_0=[2^n]\widetilde P_n.
\]
It is independent of the lift and of sufficiently large $n$.
Indeed changes of lift lie in the killed formal group; replacing
$n$ by $n+1$ gives the same result because
$2\widetilde P_{n+1}-\widetilde P_n$ lies in that formal group.
It is functorial in $R$, reduces to $(0,P)$, and is fixed by $\sigma$,
because $\sigma_0P_n=P_n$. By Katz's explicit recovery of the
pushout homomorphism, pp.~150--152,
\[
 [2]c_0=(q_{21},q_{22})\in
 \widehat A(\mathfrak m_R)\simeq(1+\mathfrak m_R)^2.        \tag{I2}
\]
The isomorphism on the right is the one furnished by the chosen dual
Tate basis. It also identifies $\sigma$ on the formal group with
$(z_1,z_2)\mapsto(z_1^{-1},z_2)$.

Every lift $c$ of $(0,P)$ is uniquely $c_0+z$, with
$z=(z_1,z_2)\in(1+\mathfrak m_R)^2$. Consequently
\[
 (1+\sigma)c=(q_{21},q_{22}z_2^2).
\]
The affine involution condition is therefore exactly
\[
 q_{21}=1,\qquad q_{22}z_2^2=1.                           \tag{I3}
\]
Together with (I1), the functor of triples is represented by
\[
 R_{\rm tr}=W(k)[[t,s,u,v]]/(s^2+2s),                     \tag{I4}
\]
where
\[
 q_{11}=1+t,\quad q_{12}=1+s,\quad q_{21}=1,\quad
 q_{22}=(1+v)^{-2},\qquad z_1=1+u,\ z_2=1+v.
\]
This inverse in $q_{22}$ fixes Katz's convention; replacing $q$ by
$q^{-1}$ yields the earlier equivalent coordinate draft.

\paragraph{Step 3: the ring is a miniversal hull for the unpointed surface.}
Forgetting the chosen origin gives a formally smooth transformation
from the triple functor to $\operatorname{Def}_{X_0}$: over a small
extension, lift the cover uniquely and lift its prescribed origin by
smoothness. At first order this transformation is injective as well
as surjective. In fact any marked isomorphism of the unpointed covers
is a translation followed by a marked abelian isomorphism; rigidity
forces the latter to preserve the same $q$ parameters. Translation
by a formal point $d$ changes $c$ by $(1-\sigma)d$. In formal torus
coordinates this changes $z_1$ by a square and leaves $z_2$ fixed.
On $k[\epsilon]/(\epsilon^2)$ every principal unit has square one.
Thus all four tangent parameters in (I4) remain distinct. This proves
formal smoothness and tangent bijectivity, the defining hull property.
It also explains why the ring represents triples but need not represent
the surface deformation functor: higher-order changes of origin can
identify different triples.

The compatible formal family can be algebraized over
$\operatorname{Spec}R_{\rm tr}$. Let $\lambda_0$ be the product principal
polarization on $A_0$. The homomorphism $2\lambda_0$ lifts through all
stages, since its Serre--Tate condition is $q_{ij}^2=q_{ji}^2$, which
holds by (I1) and (I3). The canonical diagonal pullback of the
Poincare bundle attached to this symmetric homomorphism gives a
compatible relatively ample line bundle whose special-fibre class is
$4\lambda_0$. Formal existence algebraizes the abelian scheme and
the involution, its translation point, and the quotient. The hull is
unpolarized; the auxiliary polarization is used only for algebraization.

\paragraph{Step 4: the fixed dual group is flat.}
Let $\mathcal E^\vee,\mathcal F^\vee$ denote the elliptic deformations
with Serre--Tate coordinates $q_{11},q_{22}$. The condition $q_{21}=1$
and the transposition rule lift both the original inclusion and projection
to an exact sequence of abelian schemes
\[
 0\longrightarrow\mathcal E^\vee\longrightarrow A^\vee
 \stackrel{\rho}{\longrightarrow}\mathcal F^\vee
 \longrightarrow0.                                      \tag{I5}
\]
For precision, the inclusion condition is the $(1,2)$ entry of the
dual matrix, namely $q_{21}=1$; the projection has the same condition.
The compositions are zero by rigidity. The lifted inclusion is a closed
immersion and the projection is smooth with connected elliptic kernel,
as can be checked on the unique geometric fibre and by the smoothness
criterion; their kernel is therefore the stated elliptic scheme.
Rigidity shows that $\sigma^\vee$ acts by $-1$ on the kernel and by
$+1$ on the quotient.

Put $G=\ker(1-\sigma^\vee:A^\vee\to A^\vee)$.
The restricted map $G\to\mathcal F^\vee$ is an epimorphism of fppf
sheaves. To lift $b\in\mathcal F^\vee(T)$, first lift it fppf locally
to $a\in A^\vee$. The element $(1-\sigma^\vee)a$ lies in
$\mathcal E^\vee$. Divide it by two fppf locally in that elliptic
scheme and subtract the chosen half from $a$. The result is fixed.
Its kernel is exactly $\mathcal E^\vee[2]$. Hence
\[
 0\longrightarrow\mathcal E^\vee[2]\longrightarrow G
 \longrightarrow\mathcal F^\vee\longrightarrow0          \tag{I6}
\]
is fppf exact, and $G\to\mathcal F^\vee$ is a finite locally free
torsor of degree four. In particular $G/R$ is flat.

\paragraph{Step 5: the relative Picard scheme is flat.}
For the finite etale double cover $\pi:A\to X$, translation acts
trivially on $\operatorname{Pic}^0_{A/R}=A^\vee$. Pullback gives an
exact sequence of fppf sheaves
\[
 0\longrightarrow\mu_2\longrightarrow
 \operatorname{Pic}^{\tau}_{X/R}\stackrel{\pi^*}{\longrightarrow}
 G\longrightarrow0.                                    \tag{I7}
\]
Here are the descent details. An invariant relative class is represented
fppf locally by a line bundle $L$. Invariance of its relative class gives an isomorphism up to a base line bundle; trivialize that line bundle fppf locally to obtain
$\alpha:\tau^*L\simeq L$. The square of this isomorphism is a base
unit, since $\Gamma(A_T,\mathcal O)=\Gamma(T,\mathcal O)$ universally.
Adjoin a square root of that unit fppf locally and rescale $\alpha$;
it then obeys the cocycle condition, so $L$ descends to $X$.
The descended class $D$ satisfies
$D^{\otimes2}=\operatorname{Nm}_{A/X}(L)$ in the relative Picard
functor. Fibrewise the norm of a class in $A^\vee$ belongs to the
neutral Picard component, because the norm maps the connected
abelian variety $A^\vee$ to that component. Thus $D$ belongs to
$\operatorname{Pic}^{\tau}$.
Conversely a numerically trivial class pulls back into $A^\vee$, as
the Neron--Severi group of an abelian variety is torsion free.
The choices of linearization on the trivial bundle form precisely
$\operatorname{Hom}(\mathbb Z/2,\mathbb G_m)=\mu_2$, proving the
kernel assertion. These arguments work over arbitrary base changes
$T\to\operatorname{Spec}R$, not only at geometric points.

Thus $\operatorname{Pic}^{\tau}_{X/R}\to G$ is a finite locally free
$\mu_2$-torsor of degree two. Composed with (I6), it is finite locally
free of degree eight over $\mathcal F^\vee$, and is flat over $R$.
In particular the universal first-order ordinary Igusa deformation has
flat $\operatorname{Pic}^{\tau}$.
This refutes the literal conclusion of Chai's Theorem~2.4(iii),
while leaving its line-bundle obstruction and differing tangent-space
computations entirely compatible with flatness.

\paragraph{Opened source locators and independent review.}
Katz \cite{DGJ_Katz1981}, pp.~148--152, contains the coordinate,
duality, lift criterion, and pushout recovery used above. The finite
etale lifting statement is Stacks, Theorem~59.45.2, tag04DY,
\url{https://stacks.math.columbia.edu/tag/04DY}; quotient existence is
Lemma~66.14.2, tag07S7,
\url{https://stacks.math.columbia.edu/tag/07S7}. Both were opened.
Mumford--Fogarty, Proposition~6.15, pp.~124--126, was opened through
the primary-book OCR mirror listed in the bibliography.
The full seven-page Chai v2 was opened in web retrieval, and its final
flatness inference has no additional qualification beyond the tangent
comparison already tested by the explicit torsor example.

No classification beyond ordinary $E_0$ is claimed in this proof. Root and fhi independently checked (I2)--(I7), the origin quotient, and algebraization; fhi's audit is recorded separately. The flatness assertion for every Artinian ordinary deformation uses Steps1,2,4,5 alone and does not depend on miniversality.

\paragraph{Two concrete first-order checks.}
Over $k[\epsilon]/(\epsilon^2)$ the parameters $q_{12}=1+\epsilon$, all other $q_{ij}=1$, and $z_1=z_2=1$ give an actual nonproduct direction of the surface deformation, with flat torsion Picard scheme by (I6)--(I7). There is also a completely explicit product direction: take $E_\epsilon:y^2+xy=x^3+1+\epsilon$ and keep the ordinary translation factor fixed. Its nonidentity two-torsion component is $\operatorname{Spec}R[a]/(a^2-\epsilon)$, obtained by putting $x=0,y=1+a$. It is finite flat despite lacking an $R$-point. Its absolute tangent dimension is one smaller than that of the neutral $\mu_2$ component. Thus the nonsmooth-torsor behavior used to diagnose the published inference occurs in an Igusa family itself.

% END D4o completed ordinary calculation

% BEGIN D4s completed supersingular calculation
% Completed supersingular calculation; independently audited by root and fhi.
\subsection*{D4s: the supersingular Igusa hull and the canonical elliptic subgroup}
\textbf{Status: Proved.} Root and fhi independently audited the complete argument.
This section treats the remaining case in D4o. All bases in the
deformation calculation are local Artinian $W(k)$-algebras with residue
field an algebraically closed field $k$ of characteristic two; all
deformations and isomorphisms are marked on the closed fibre. Let
$E_0$ be supersingular, let $F_0$ be ordinary, and choose
$0\ne P\in F_0[2](k)$. Put
\[
 A_0=E_0\times F_0,\qquad \sigma_0=(-1,1),\qquad
 \tau_0=t_{(0,P)}\sigma_0,\qquad X_0=A_0/\langle\tau_0\rangle.
\]
The conclusions are
\[
 R_{X_0}^{\mathrm{hull}}
   =W(k)[[t,s,u,v]]/(s^4-ts^2-2s),                       \tag{S1}
\]
and flatness of $\operatorname{Pic}^{\tau}$ for every Artinian
deformation of $X_0$. The ring in (S1) is finite free of rank four over
$W(k)[[t,u,v]]$. Together with D4o, this gives the two cases of the full
Igusa calculation; the ordinary ring is
$W(k)[[t,s,u,v]]/(s^2+2s)$.

\paragraph{Step 1: the cover and the multiplicative subgroup.}
As in D4o Step1, every deformation $X/R$ uniquely lifts the marked
finite etale cover $A_0\to X_0$ and its deck involution. A chosen lift
of its origin makes $A/R$ an abelian scheme; write its deck involution
as $\tau=t_c\sigma$, where $\sigma$ is a group automorphism lifting
$\sigma_0$. Thus $\sigma^2=1$ and $(1+\sigma)c=0$.
The sources are Stacks tag04DY and Mumford--Fogarty,
Proposition6.15, pp.124--126, already opened in D4o.

Write $G=A[2^\infty]$ and $D=G^0$. The connected--etale sequence,
and Cartier duality applied to that sequence for $G^\vee$, give
\[
 G/D\simeq (\mathbb Q_2/\mathbb Z_2)_R,
 \qquad M:=((G^\vee)^{\mathrm{et}})^\vee
       \simeq\mu_{2^\infty}\ \subset D.                 \tag{S2}
\]
The marked special fibre identifies $M$ with the multiplicative
part of $F_0[2^\infty]$. These are canonical subgroups: the only
etale part of $G_0$ and of $G_0^\vee$ comes from $F_0$.
The quotient $G/M$ is a Barsotti--Tate group, being Cartier dual to
$(G^\vee)^0$. Its connected part is $D/M$, so $D/M$ is a connected
Barsotti--Tate group deforming $E_0[2^\infty]$.
For these facts one can use Tate, \S1.3--1.4, pp.160--161,
\S2.2, Proposition1 and equation(4), pp.162--164, and
\S2.3, p.166 \cite{DGJ_Tate1967}; the exact sequence and duality
are also obtained level by level from their finite flat versions.

The action of $\sigma$ on $M$ is $+1$, and its action on $D/M$ is
$-1$. Here no eigenspace splitting in characteristic two is being
assumed. A homomorphism between Barsotti--Tate groups over $R$ that
reduces to zero is zero: its image is killed by a fixed power of two
in the infinitesimal formal kernel, whereas its source is two-divisible.
This is the rigidity argument in Katz, Lemmas1.1.1--1.1.3,
pp.141--143, using the equivalence between connected Barsotti--Tate
groups and divisible formal groups. Applying it to $\sigma-1$ on
$M$ and $\sigma+1$ on $D/M$ proves the assertions.

\paragraph{Step 2: the canonical divisible system of actual points.}
Choose a compatible generator of $T_2F_0(k)$, or equivalently points
$e_r\in F_0[2^r](k)$ for $r\geq1$, with $e_1=P$ and
$2e_{r+1}=e_r$; put $e_0=0$. The formal group
$\widehat A(\mathfrak m_R)$ is killed by $2^N$ for sufficiently large
$N$ (Katz, Lemma1.1.1). For $r\geq0$ define
\[
 j(2^{-r})=[2^N]\widetilde e_{N+r}\in A(R),             \tag{S3}
\]
where $\widetilde e_{N+r}$ is any lift to $A(R)$, and take $N$ large.
Such lifts exist by smoothness. Changing the lift or increasing $N$
does not change (S3): the difference of two choices, or of
$2\widetilde e_{N+r+1}$ and $\widetilde e_{N+r}$, belongs to the
killed formal kernel. These identities also give
$2j(2^{-r-1})=j(2^{-r})$. Thus (S3) defines a functorial additive map
\[
 j:\mathbb Z[1/2]\longrightarrow G(R).
\]
Its values really are two-power torsion: their reductions are torsion
and a further uniform power of two kills the formal kernel. Every
value of $j$ is fixed by $\sigma$, since the chosen special-fibre
points are fixed and the same independence argument applies to their
lifts. Set $c_0=j(1/2)$ and $\kappa=j(1)=2c_0\in D(R)$.

Write $c=c_0+z$, where $z\in\widehat A(\mathfrak m_R)=D(R)$.
The last equality follows either from the formal-group description
of $D$ or directly by killing the formal kernel by a power of two.
The affine involution condition is
\[
 0=(1+\sigma)c=\kappa+(1+\sigma)z.
\]
Since $\sigma=-1$ on $D/M$, this forces
\[
                         \kappa\in M(R).                \tag{S4}
\]
This is the step that uses the existence of the affine deck map, not
merely a lift of its linear part.

\paragraph{Step 3: construction and uniqueness of the elliptic subgroup.}
As fppf sheaves over $R$ there is a presentation
\[
 G\simeq
 (D\oplus\underline{\mathbb Z[1/2]})/
       \{(-n\kappa,n):n\in\underline{\mathbb Z}\},       \tag{S5}
\]
whose map to $G$ is $(d,r)\mapsto d+j(r)$. This is a statement on
arbitrary test schemes $T/R$, not just on $R$-points. The image of
an element of $G(T)$ in the constant etale quotient is locally a
constant rational class in $\mathbb Z[1/2]/\mathbb Z$; choose a
representative $r$ locally and subtract the base section $j(r)_T$.
The difference lies in $D(T)$. The kernel consists exactly of the
pairs $(-n\kappa,n)$, which proves (S5) as a sheaf quotient.

By (S4) the analogous subquotient
\[
 H=(M\oplus\underline{\mathbb Z[1/2]})/
       \{(-n\kappa,n):n\in\underline{\mathbb Z}\}
       \ \subset G                                      \tag{S6}
\]
is defined. It fits in an exact sequence
$0\to M\to H\to\mathbb Q_2/\mathbb Z_2\to0$.
It is a Barsotti--Tate group, as can be checked without invoking an
extension-closure assertion: multiplication by two is surjective on
both end terms and therefore on $H$; for every $n$ the snake lemma
and $M/2^nM=0$ give
\[
 0\longrightarrow M[2^n]\longrightarrow H[2^n]
 \longrightarrow (\mathbb Z/2^n)_R\longrightarrow0.     \tag{S7}
\]
Thus $H[2^n]$ is a torsor under the finite flat group $M[2^n]$ over
the constant finite etale scheme $(\mathbb Z/2^n)_R$; it is finite
locally free of rank $2^{2n}$. These finite levels give a height-two,
dimension-one ordinary Barsotti--Tate group. Their monomorphisms into
$G[2^n]$ are closed immersions. Its special fibre is exactly the
given subgroup $F_0[2^\infty]\subset A_0[2^\infty]$.

Serre--Tate, in the general form of Katz Theorem1.2.1,
pp.143--144, now produces an elliptic scheme $F/R$ and a
homomorphism $F\to A$ lifting $F_0\hookrightarrow A_0$ and
realizing $H\hookrightarrow G$. This map is a closed immersion:
it is proper and quasi-finite, since its only geometric fibre is a
closed immersion, hence finite; the map on structure sheaves is
surjective by reduction and Nakayama. We identify $F$ with its image.
Both $M$ and $j$ are fixed by $\sigma$, so $\sigma|_F=1$ by
Serre--Tate full faithfulness. The quotient $E=A/F$ is an elliptic
scheme, and $\sigma$ induces $-1$ on it by rigidity.

This subgroup is unique among lifts of $F_0$ fixed pointwise by
$\sigma$. Indeed such a lift $F'$ has multiplicative subgroup $M$
and the same etale quotient. In (S3) one may choose all the lifts
inside the smooth scheme $F'$, so every $j(r)$ belongs to $F'$.
Consequently its Barsotti--Tate subgroup contains $H$; (S7) and
the ranks at every finite level force equality. Full faithfulness
then gives $F'=F$. In particular $F$, $E$, and the subsequent
classification are functorial in the deformation.

\paragraph{Step 4: an elementary graph description of $(A,\sigma)$.}
Let $\rho:A\to E=A/F$. The homomorphism $1-\sigma$ kills $F$, so
it factors uniquely as $a\rho$, where $a:E\to A$. Since
$\sigma=-1$ on $E$, one has $\rho a=[2]_E$.
The homomorphism
\[
 \Phi:E\times F\longrightarrow A,\qquad(e,f)\longmapsto a(e)+f
\]
is an fppf epimorphism with kernel
\[
 K=\{(e,h(e)):e\in E[2]\},\qquad
 h=-a|_{E[2]}:E[2]\longrightarrow F[2].                 \tag{S8}
\]
Indeed one can lift $\rho(x)$ by dividing it by two fppf locally in
$E$, and then the residual difference lies in $F$; the kernel formula
is immediate. Thus $\Phi$ is an isogeny of degree four.
On the special fibre $a(e)=(2e,0)$, so $h$ reduces to zero.
Its composite with $F[2]^{\mathrm{et}}\simeq\mathbb Z/2$ is
zero, by invariance of maps into an etale scheme under a nilpotent
immersion. Hence $h$ factors through $F[2]^0\simeq\mu_2$.
Cartier duality gives
\[
 \operatorname{Hom}(E[2],\mu_2)
  =\operatorname{Hom}(\mathbb Z/2,E^\vee[2])
  =E^\vee[2](R).                                       \tag{S9}
\]
Let $Q$ denote the corresponding point; it reduces to zero.

Conversely, elliptic deformations $E,F$ and such a point $Q$ define
$h_Q$ by (S9) and the abelian quotient
\[
 A_Q=(E\times F)/\operatorname{graph}(h_Q).               \tag{S10}
\]
The involution $(-1,+1)$ preserves the graph, because $-1=+1$ on
$E[2]$, and descends to $\sigma$. The factor $F$ embeds in $A_Q$,
and the quotient by it is $E$ via $[2]$ on the first factor.
The marked special-fibre identification is induced by
$(e,f)\mapsto(2e,f)$. The homomorphism $a$ is $e\mapsto[(e,0)]$,
so this construction recovers exactly (S8) and $Q$. By Step3 the
elliptic subgroup used in the reconstruction is unique. We have
therefore classified all the required $(A,\sigma)$, not just those
which already come with an elliptic fibration.

\paragraph{Step 5: the affine norm equation removes the $F$ parameter.}
The marked deformation space of $E_0$ has a parameter $t$.
Choose a coordinate $s$ at zero on the dual elliptic deformation,
and let $g_E(t,s)=0$ be its two-torsion equation. The ordinary
elliptic deformation $F$ has Serre--Tate parameter $q\in1+\mathfrak m$.
Thus the pairs $(A,\sigma)$ of Step4 have complete local base
\[
 B=W(k)[[t,q-1,s]]/(g_E(t,s)).                           \tag{S11}
\]
The marking identifies the formal group of $F$ with
$\widehat{\mathbb G}_m$. The canonical point $c_0$ in Step2 lies
in $F$, by its construction there, and Katz's pushout recovery,
pp.150--152, gives $[2]c_0=q$ in this formal-group coordinate.

Choose a formal-variety section $b(u)$ of
$\widehat A\to\widehat E$ over $B$, through the identity; this
exists by formal smoothness and a choice of relative formal
coordinates. It is not required to be a homomorphism. Every formal
point of $A$ can then be written uniquely as $b(u)+(1+v)$, with
$1+v$ in the formal group of $F$. Since $\sigma=-1$ on $E$,
write $(1+\sigma)b(u)=n(t,q,s,u)$ in the formal group of $F$.
It is a principal unit and satisfies $n(t,q,s,0)=1$ identically.
The affine condition becomes exactly
\[
                 q\,n(t,q,s,u)(1+v)^2=1.               \tag{S12}
\]
The derivative of its left side minus one with respect to $q$ is
one at the closed origin. The formal implicit-function theorem
therefore eliminates $q$ uniquely. This elimination is valid over
the possibly nonreduced base (S11): lift $n-1$, which is divisible
by $u$, to a power series before imposing $g_E$, eliminate $q$
there, and then impose $g_E(t,s)$, which is independent of $q$.
The functor of triples $(A,\sigma,c)$ is consequently represented
by
\[
                    W(k)[[t,s,u,v]]/(g_E(t,s)).          \tag{S13}
\]

Forgetting the chosen origin is formally smooth onto
$\operatorname{Def}_{X_0}$, by the cover-lifting and smooth-origin
argument in D4o. It is tangent bijective: a marked isomorphism of
the covers is a translation by a formal point $d$, followed by a
marked group isomorphism. The group isomorphism preserves the
classified data; the translation changes $c$ by $(1-\sigma)d$.
On $k[\epsilon]/(\epsilon^2)$ this last expression is zero,
because the derivative of $\sigma_0=(-1,1)$ is the identity in
characteristic two. Hence no tangent parameters in (S13) are
identified. This proves the hull property. As before, the assertion
is miniversality for unpointed surfaces, not representability of that
functor; changes of origin can act at higher orders.

\paragraph{Step 6: the exact supersingular equation.}
Choose $E_0:y^2+y=x^3$, which represents the unique supersingular
elliptic isomorphism class over $k$. A universal marked deformation
is
\[
                  E_t:\ y^2+txy+y=x^3,
                  \qquad\Delta=t^3-27.                 \tag{S14}
\]
For completeness this normal form follows directly from general
Weierstrass equations. At the displayed special equation, changes
$x=\lambda^2x'+r$, $y=\lambda^3y'+\lambda^2 b x'+d$ have an
invertible Jacobian on the four coefficients
$(a_2,a_3-1,a_4,a_6)$: modulo two its diagonal entries, in the
order $(r,\lambda-1,b,d)$, are $(3,-3,-1,-1)$.
Thus these four coefficients can be normalized formally to
$(0,0,0,0)$, leaving $a_1=t$. There are no infinitesimal pointed
automorphisms of an elliptic curve. This proves the asserted
one-parameter marked deformation chart. The canonical principal
polarization identifies $E_t^\vee$ with $E_t$.

Near $O=[0:1:0]$, use $s=X/Y=x/y$ and $w=Z/Y=1/y$.
The curve equation and inversion are
\[
 w+t s w+w^2=s^3,\qquad
 (s,w)\longmapsto
 \left(-\frac{s}{1+ts+w},-\frac{w}{1+ts+w}\right).       \tag{S15}
\]
There is a unique formal solution $w=w_t(s)\in(s^3)$.
Because the denominator in (S15) is a unit and $w\in(s)$, the
fixed-point ideal of inversion is generated by
$s(2+ts+w)$. The following exact identity in the curve ring removes
the auxiliary coordinate:
\[
 (w-1)(w+2+ts)=s^3-ts-2.
\]
Since $w-1$ is a unit at the origin, the two-torsion ideal is
\[
 (s(2+ts+w))=(s^4-ts^2-2s).                             \tag{S16}
\]
Equations (S13) and (S16) prove (S1). All of the supersingular
two-torsion lies in this formal neighborhood. The polynomial is
monic in $s$ and has its lower coefficients in the maximal ideal of
$W(k)[[t,u,v]]$; Weierstrass division makes the complete quotient
free with basis $1,s,s^2,s^3$. At $2=t=0$ its two-torsion factor is
$k[s]/(s^4)$, displaying the fourth-order thickening explicitly.

\paragraph{Step 7: algebraization and flatness of the torsion Picard scheme.}
The hull has an algebraic realization over its complete local base.
The elliptic deformation $E_t$ is given by (S14); the ordinary
elliptic family $F_q$ is effective using its Weierstrass deformation
chart and the Serre--Tate coordinate isomorphism. The point $Q$ is
an actual section of its finite two-torsion scheme. Hence (S10) is
an abelian scheme over the complete ring after substituting the
solution of (S12). The compatible point $c$ is effective as well:
choose an affine neighborhood of its closed-fibre point in $A$;
the compatible Artinian sections give its coordinate values in the
complete base ring. The resulting affine involution is free, because
its proper fixed locus has empty closed fibre.

There is a useful explicit ample bundle, so no unpolarized formal
existence assumption is needed. Since the kernel of
$\Phi:E\times F\to A$ lies in $(E\times F)[2]$, multiplication
by two factors as $g\Phi$ for an isogeny $g:A\to E\times F$.
Pull back the product ample bundle by $g$, then tensor it with its
$\tau$-pullback. Swapping the factors supplies a $\tau$-linearization.
It descends to an ample bundle on the finite etale quotient $X$.
This gives the claimed algebraized projective hull family.

Finally the canonical subgroup constructed in Step3 gives
\[
 0\longrightarrow E^\vee\longrightarrow A^\vee
 \longrightarrow F^\vee\longrightarrow0,
\]
with $\sigma^\vee=-1$ on the kernel and $+1$ on the quotient.
Put $G_\sigma=\ker(1-\sigma^\vee)$. Exactly as in D4o Steps4--5,
fppf division by two in $E^\vee$ and cyclic line-bundle descent give
\[
 0\longrightarrow E^\vee[2]\longrightarrow G_\sigma
 \longrightarrow F^\vee\longrightarrow0,
 \qquad
 0\longrightarrow\mu_2\longrightarrow
 \operatorname{Pic}^{\tau}_{X/R}\longrightarrow G_\sigma
 \longrightarrow0.                                    \tag{S17}
\]
Both surjections are torsors in the fppf topology. Their degrees
are four and two. In the second sequence an invariant relative
class is represented fppf locally by a line bundle; trivialize the
base factor in its invariance isomorphism, then adjoin a square root
of the cocycle unit to linearize it. The descended class is torsion
modulo the neutral component by the norm identity, just as in (I7).
These steps do not require any section of a nonsmooth torsor over
the original base. Thus $\operatorname{Pic}^{\tau}_{X/R}$ is finite
locally free of degree eight over $F^\vee$ and is flat over $R$.
This conclusion holds for every Artinian deformation already by
Steps1--3 and (S17), independently of the explicit hull calculation.

\paragraph{Finite certificate and scope.}
The code \texttt{check\_igusa\_supersingular.py} records the exact
Weierstrass discriminant, the inversion identities, the polynomial
identity in (S16), the normal-form Jacobian, and the special-fibre
length four. Its output does not purport to mechanize the
Barsotti--Tate construction, fppf descent, or the hull argument.
Those are supplied in full above. Root and fhi independently checked the subgroup construction, the
classification, the norm elimination, the hull and effectivity arguments,
and the two torsor sequences. Their audits are recorded in the
check-event ledger.

% END D4s completed supersingular calculation

\subsection{No compatible Picard variety over dual numbers}
\paragraph{D4a (Explicit counterexample to Artinian compatibility).}
Let $k=\overline{\mathbb F}_2$, $R=k[\epsilon]/(\epsilon^2)$,
and take two labelled copies $E_0,F_0$ of the ordinary elliptic
curve
\[
 y^2+xy=x^3+1,
 \qquad P=(0,1)\in F_0[2](k).
\]
Its discriminant is one and the coefficient of $xy$ is one, so it
is smooth ordinary; inversion sends $(x,y)$ to $(x,y+x)$,
showing that the displayed nonzero point has order two.
Set $A_0=E_0\times F_0$ and
$X_0=A_0/\langle(-1,t_P)\rangle$, as in D4o.
Choose the marked abelian deformation $A/R$ with Katz coordinates
\[
 (q_{ij})=\begin{pmatrix}1&1+\epsilon\\1&1\end{pmatrix}.
                                                        \tag{D4a.1}
\]
The involution and its translation lift: in D4o, equations
(I1)--(I3), choose $z_1=z_2=1$. Indeed
$(1+\epsilon)^2=1$ and the second row is $(1,1)$.
This gives an actual smooth projective quotient deformation
$X/R$ and a finite étale double cover $\pi:A\to X$.
Its torsion Picard scheme is flat, by D4o.

Nevertheless there is \emph{no} abelian $R$-scheme $B$ with a
homomorphism
\[
 \iota:B\longrightarrow\operatorname{Pic}^{\tau}_{X/R}
                                                        \tag{D4a.2}
\]
whose special fibre identifies $B_k$ isomorphically with
$\operatorname{Pic}^{\rm var}_{X_0/k}$ and whose special-fibre
map is its natural inclusion. In particular no abelian closed
subscheme of the relative Picard scheme can have the specified
reduced identity component on the closed fibre. The assertion
rules out existence, not merely a preferred construction.

Here is the proof. Put $Y_0=F_0/\langle P\rangle$ and write
$\phi:F_0\to Y_0$ for the étale degree-two isogeny. The product
Igusa calculation gives
\[
 \operatorname{Pic}^{\tau}_{X_0/k}
     \simeq Y_0^\vee\times E_0^\vee[2],
 \qquad \operatorname{Pic}^{\rm var}_{X_0/k}=Y_0^\vee.
\]
Under pullback to $A_0$, the latter inclusion becomes
\[
 f_0=(0,\phi^\vee):Y_0^\vee
                         \longrightarrow E_0^\vee\times F_0^\vee.
                                                        \tag{D4a.3}
\]
This also follows directly by pulling back a degree-zero line
bundle from the quotient map $X_0\to Y_0$.
If (D4a.2) existed, compose it with $\pi^*$ to the relative
Picard scheme of $A$. This composite lands in its open neutral
component $A^\vee$: it does so on the closed fibre, and an
Artinian nilpotent thickening has the same underlying topological
space, so the inverse image of the complementary open is empty.
We therefore obtain a homomorphism of ordinary abelian $R$-schemes
\[
 f:B\longrightarrow A^\vee
\]
lifting $f_0$. Both source and target are abelian schemes, so its
dual $f^\vee:A\to B^\vee$ is a homomorphism as well; no image,
kernel-flatness, or splitting assertion is needed.

Choose the ordinary étale Tate bases for $A_0$ and $A_0^\vee$
as in D4o, with indices $E,F$. Since $\ker\phi^\vee=\mu_2$,
its map on étale $2$-adic Tate modules is an isomorphism. Choose
the generator $x\in T_2Y_0^\vee(k)$ so that
\[
 T_2(f_0)(x)=e_F^\vee.
\]
The dual of (D4a.3) is $(0,\phi):E_0\times F_0\to Y_0$,
so
\[
 T_2(f_0^\vee)(e_E)=0.
\]
Katz's necessary and sufficient morphism-lifting criterion
\cite[Theorem 2.1(4), p.~149]{DGJ_Katz1981}, applied to
$f:B\to A^\vee$, would imply
\[
 q_B\bigl(x,T_2(f_0^\vee)e_E\bigr)
    =q_{A^\vee}\bigl(T_2(f_0)x,e_E\bigr).
\]
The left side equals one. By the transposition formula for dual
abelian schemes, Katz's Theorem~2.1(3), pp.~148--149, the right
side equals
\[
 q_{A^\vee}(e_F^\vee,e_E)=q_A(e_E,e_F^\vee)
      =q_{12}=1+\epsilon\ne1.
\]
This contradiction proves the assertion. It uses the exact
Serre--Tate lifting criterion on the special-fibre Tate modules,
not a claim that an arbitrary point lifts through a torsion
isogeny. It remains valid for every possible elliptic deformation
$B$ of $Y_0^\vee$: its first pairing entry is evaluated at zero
in the second slot, independently of its own parameter.

The leaf 57:311 asks ``Existence de'' $\operatorname{Pic}^{00}$
(the superscript is transcribed uncertainly as $00$ or $\infty$),
``dans le cas Artinien'' with ``local'' uncertain. Under the
compatible-abelian-subscheme interpretation specified above,
the answer is negative already on an Artinian base of length two. The full torsion Picard scheme
of this same example is flat. The two properties must not be
identified. The calculation of the matrix and the nonzero
obstruction in $R$ is certified in
\texttt{check\_artin\_picvar.py}; existence and the contradiction
are the geometric proof above.

The targeted independent audits by root and bce checked the étale Tate
isomorphism for $\phi^\vee$, the zero dual entry, the exact
transposition convention, and the existence of the quotient from
D4o. No gap was found. This explicit Serre--Tate proof is consistent
with Grothendieck's published first-order Igusa remark
\cite[Remark 3.4, p.~238]{DGJ_FGAVI1962}, not a claim of priority.

% BEGIN D4b
\paragraph{D4b (Refuted universal lifting of the Albanese fibration).}
\emph{Existing-source check.}
Grothendieck's earlier first-order Igusa obstruction to a compatible
abelian Picard subgroup is already recorded in
\cite[Remark 3.4, p.~238]{DGJ_FGAVI1962}.
The exact lifting assertions addressed here were read in the published
2013 text of \cite{DGJ_Partsch2013}, pp.~102--106.
No correction to these assertions was found in the retrieved
publisher and author-source search; this limited search is not a
claim that no correction exists. No priority claim is made for the
obstruction. Its fully specified deformation is D4a.

The explicit deformation $X/R$ of D4a, with
$R=\overline{\mathbb F}_2[\epsilon]/(\epsilon^2)$, does not
admit a lifting of the Albanese elliptic fibration
\[
 f_0:X_0\longrightarrow Y_0=F_0/\langle P\rangle.
\]
Here a lifting means a morphism $f:X\to C$ over $R$, where
$C/R$ is a smooth proper deformation of $Y_0$, and with
special fibre identified with $f_0$. The special fibration is
Jacobian: its section is induced by $\{0\}\times F_0$ in the
cover $E_0\times F_0$.

Suppose such $f$ exists. The relative Picard scheme
$B=\operatorname{Pic}^0_{C/R}$ is an elliptic abelian scheme,
and pullback is a homomorphism of $R$-group schemes
\[
 f^*:B\longrightarrow\operatorname{Pic}_{X/R}.
\]
Its image lies in $\operatorname{Pic}^{\tau}_{X/R}$, since
pullback preserves algebraic triviality on geometric fibres.
Its special fibre is precisely the canonical inclusion
$Y_0^\vee=\operatorname{Pic}^{\rm var}(X_0)
\hookrightarrow\operatorname{Pic}^{\tau}(X_0)$, because $f_0$
is the Albanese map. This is the homomorphism ruled out by
D4a, for every possible deformation $B$. Thus $f$ cannot
exist. No assertion about the scheme-theoretic image of
$f^*$, or its flatness, enters this contradiction.

This refutes the literal universal lifting assertions in
\cite[Proposition 6.6, pp.~102--103; Theorem 6.7, p.~104]{DGJ_Partsch2013}:
the opened published text states that every deformation of
a Jacobian bielliptic surface extends its elliptic fibration,
and then claims the same for the Albanese fibration of every
bielliptic surface. The special fibre above is exactly a
Jacobian ordinary Igusa surface, and D4a constructs an actual
smooth projective deformation. The stronger assertion that
both fibrations lift, in Proposition~6.9, p.~106, is refuted
by the same example. This does not make a claim about other
cases of that article's classification.

The last trace step of the proof on p.~103 illustrates a
specific danger. Killing a connected subgroup and inducing
zero on an etale quotient does not force the original map to
vanish: its remaining map may go from that etale quotient to
a connected infinitesimal target. Over this very $R$ there
is a nontrivial homomorphism
\[
 h:(\mathbb Z/2)_R\longrightarrow\mu_{2,R},\qquad
 h(0)=1,\quad h(1)=1+\epsilon,
\]
whose special fibre is zero. The relation
$(1+\epsilon)^2=1$ verifies the group law. Thus rigidity of
maps between etale group schemes cannot be applied when
only the source is etale. The complete geometric refutation
above uses D4a; this finite-group example identifies an
invalid inference without being substituted for that proof.
The Hopf-map and group-law identities are recorded by
\texttt{check\_partsch\_etale\_step.py} and its saved output.

% END D4b

\section{The relative Picard functor}
\subsection{E1: literal counterexample and the relative-normal result}

Over $R=\mathbf C[[t]]$ take the plane cubic
\[
 X:\ Y^2Z=X^3+X^2Z+t^2Z^3,
 \quad O=[0:1:0],\quad P=[0:t:1].
\]
Its generic fibre is smooth since the cubic discriminant is
$-t^2(4+27t^2)$. Its special fibre is an irreducible nodal cubic. Its total
space is normal, because it is a hypersurface with an isolated singularity
at $(x,y,t)=(0,0,0)$. It is proper and flat, has a section, and has
$f_*\mathcal O_X=\mathcal O_R$ universally by the plane-cubic exact sequence.
The generic line bundle $\mathcal L_\eta=\mathcal O(P_\eta-O_\eta)$ is
algebraically equivalent to zero.

The closure $P$ has local ideal $(x,y-t)$ at the node, with two independent
linear generators, so it is not Cartier. If $\mathcal L_\eta$ extended to
an invertible sheaf on $X$, choose a rational trivialization agreeing with
the generic divisor $P_\eta-O_\eta$. Its divisor differs from $P-O$ by a
vertical divisor. The only vertical prime is the reduced irreducible special
fibre, whose class is the principal divisor of $t$. As $O$ is Cartier, this
would make $P$ Cartier, a contradiction. Thus mere total-space normality does
not suffice, even with projectivity.

For the original relative-normal reading, assume instead that $f$ is proper
flat with geometrically normal, geometrically connected fibres over a trait,
and has a section. Universal $H^0=\mathcal O$ gives the algebraic Picard space
by Artin, and geometrically integral fibres give separatedness. Raynaud
\cite[Theorem~3.3.1, p.~40]{Raynaud1970} makes this separated group algebraic
space a scheme. Grothendieck \cite[Theorem~1.1(i), Corollary~2.3]{Grothendieck1962VI}
then says that the union of its fibrewise neutral components is a proper
\emph{subset}, in his precise topological/valuative sense. Thus a generic
algebraically trivial Picard point extends uniquely to a Picard point with
special fibre algebraically trivial; the section supplies an actual line
bundle representing it. No projectivity is required.

It is essential not to recast this as representability and properness of
the fibrewise $\operatorname{Pic}^0$ subfunctor. That subfunctor need not be an
algebraic space, even for smooth projective families, as
\cite[Proposition~2.2]{LaurentSchroer2024} demonstrates for Enriques surfaces.
The topological subset statement and the representability statement differ.

\subsection{E2/H9: separatedness of correspondences}

Let $S$ be locally noetherian and $X,Y$ proper flat and finitely presented
over $S$, both with universal $H^0=\mathcal O_S$. Define $\operatorname{Corr}_S(X,Y)$
as the fppf quotient of $\operatorname{Pic}_{X\times_SY/S}$ by the two factors'
Picard functors. It is a separated unramified $S$-scheme locally of finite type.

Here are the valuative argument and the exact representability deduction.
Given a trait $T=\operatorname{Spec}R$ over $S$ and a correspondence class
vanishing over its generic point, make a faithfully flat extension of traits
so that $X_T,Y_T$ have sections $x,y$. Such sections are obtained by taking
closed points of the generic fibres over a finite field extension and using
properness over a dominating valuation ring. A correspondence is now uniquely
represented by a line bundle bi-rigidified along $X\times y$ and $x\times Y$:
from a representative $L$ one replaces it by
\[
 L\otimes p_X^*(L|_{X\times y})^{-1}
  \otimes p_Y^*(L|_{x\times Y})^{-1}
  \otimes f^*(L|_{(x,y)}).
\]
This bundle is trivial on the generic fibre. It defines a map
$a:X_T\to P=\operatorname{Pic}_{Y_T/T}$ with $a(x)=0$ and $a_\eta=0$.

Let $E\subset P$ be the schematic closure of the generic identity.
Raynaud \cite[\S3.2(c), p.~40]{Raynaud1970} defines this subsheaf by precisely
the maps from flat $T$-schemes whose generic maps land in the identity.
Since $X_T$ is flat, $a$ factors through $E$. Universal $H^0=\mathcal O_T$
for $Y_T$ and Raynaud's Proposition~5.2, p.~46, imply that $E/T$ is an
\emph{\'etale algebraic space}. Each geometric fibre of $X_T$ is connected,
and a map from it to the discrete fibre of $E$ is constant. The section fixes
that value as zero. The equality locus of $a$ and zero is open, since $E/T$
is unramified, and contains every geometric point; it is therefore all of
$X_T$. The bi-rigidified line bundle is trivial. Faithfully flat descent
returns the original correspondence class to zero.

This proves the separatedness condition (F6) in Murre
\cite[p.~245]{Murre1966}. Murre's Theorem~4, p.~256, proves that, under exactly
these hypotheses, \emph{if} the correspondence functor is separated then it
is represented by an unramified scheme locally of finite type. The conditional
theorem becomes applicable after the valuative proof just supplied; it must
not be cited as if it already included unconditional separatedness.

\subsection{E3 and H6: flattening and the numerical quotient}

For $f:X\to S$ proper of finite presentation and $F$ finitely presented and
quasi-coherent on $X$, the functor taking $T/S$ to a singleton exactly when
$F_T$ is flat over $T$ is represented by a monomorphism $S_F\to S$ of finite
presentation. This is precisely the requested Problem A: the scheme version
is Stacks, Lemma~38.27.6(3), tag~05PS; the algebraic-space version is
Lemma~77.11.8(3), tag~0CWQ. No projectivity hypothesis is used.

Suppose $P=\operatorname{Pic}_{X/S}$ is an algebraic space locally of finite
presentation and $H=P^\tau$ is flat, of finite presentation, open and closed.
The free translation action gives an algebraic-space quotient $Q=P/H$ with
$P\to Q$ an fppf torsor (Stacks, Lemma~80.11.7, tag~06PH). Its relation
$P\times H\to P\times P$ is open and closed. Faithfully flat descent of the
diagonal makes $Q/S$ unramified and separated. A separated locally
quasi-finite algebraic space is a scheme (Stacks, Lemma~67.51.1, tag~0418).
This proves both readings of H6. The prompt summary asks for an algebraic-space
quotient under flatness: that already follows from tag~06PH, without closedness.
The fuller wording in \texttt{HUNCHES.md} assumes closedness as well, and obtains
the stronger separated scheme conclusion proved above. Compare
\cite[Theorem~3.1]{LaurentSchroer2024}.

\subsection{Some decisive leaf-by-leaf answers on 57:308--315}

The accompanying CSV retains every substantive leaf, including duplicates and
ambiguous readings; 57:308 is a heading and 57:312 is an unrelated tensor-product
insertion. The remaining broad entries are not silently marked settled.

For $X:\ xy=tz^2\subset\mathbf P^2_{\mathbf C[[t]]}$, its relative Picard
scheme is the disjoint union over total degrees $n$ of schemes $S_n$ obtained
by gluing copies of $\operatorname{Spec}\mathbf C[[t]]$, indexed by pairs
$(a,b)$ with $a+b=n$, along the common generic point. This is the explicit
split version of \cite[Example~4.14]{Kleiman2005}. Thus a connected component
need not be of finite type; the closure of a single generic point need not
be quasi-compact; and $P^\tau$ (the identity copy inside $S_0$) need not be
closed. These answer 57:309(a),(a'),(a'') in general and 57:310(8) negatively.

For the cuspidal cubic $C:y^2z=x^3$ in characteristic $p>0$, normalization
and the conductor exact sequence give
$\operatorname{Pic}^0_C=\operatorname{Pic}^\tau_C=\mathbf G_a$:
the relevant unit quotient is
$(k[\epsilon]/\epsilon^2)^\times/k^\times\simeq\mathbf G_a$.
Multiplication by $p$ is the constant zero map and is not open. Every
$k$-point is killed by $p$, so the set of points killed by a fixed integer
need not be finite. This answers 57:310(4),(6) in the generality actually
written there. For an integer invertible on the base, multiplication on a
commutative group scheme locally of finite type is \`etale
\cite[Theorem~2.5]{Grothendieck1962VI}; in particular it is universally open.

The tau-part is always an open algebraic subspace of finite presentation
under the proper-flat-finitely-presented, cohomologically-flat-in-degree-zero
hypotheses: \cite[Theorem~2.1]{LaurentSchroer2024}. Its finite-type and bounded
NS-rank assertions have the classical sources SGA~6 XIII, Theorems~4.7 and
5.1; see the exact discussion in \cite[Remark~6.19]{Kleiman2005}. Multiplication
by every nonzero integer on a represented relative Picard functor is of finite
type by SGA~6 XIII, Theorem~3.6; \cite[Remark~6.28]{Kleiman2005} records the
full proper version. These finiteness theorems do not imply flatness or
universal openness of the tau-part.

For geometrically integral locally factorial fibres the tau-part is proper.
This follows from \cite[Proposition~2.3(ii)]{LaurentSchroer2024}. Under mere geometrically
normal fibres, this particular theorem has an extra locally factorial
hypothesis and cannot be cited after deleting it. The properness of the
\emph{neutral subset} in E1 similarly does not itself prove properness of
the entire tau-part.

\subsection{E4: separation under reduced fibres, and a represented counterexample}

The two questions on 57:313 are:
\begin{quote}
``S\'epar\'e sur $S$ si $X_s$ sont g\'eom$^t$ loc. int\`egres ?''
and ``$\underline{\mathrm{Pic}}^\tau_{X/S}$ s\'epar\'e sur $S$ si les $X_s$
s\'eparables ?''
\end{quote}
Both have affirmative answers in the proper-flat setting. For the second,
test separatedness after base change to a strictly henselian trait. Geometric
reducedness of the special fibre implies Raynaud's property (N)
\cite[Example~6.1.6, p.~49]{Raynaud1970}. After the Stein-factor reduction
of \S6.1.9 one has property (N)$^*$, and every special-fibre multiplicity is
one. Thus $d=1$ and $d'=1$ in Proposition~6.1.11. Theorem~6.4.1(3), p.~56,
or equivalently Corollary~6.4.5, p.~57, makes $P^\tau$ separated. This holds
after every trait base change, proving the valuative criterion over $S$.
Every $n$-torsion Picard point belongs to $P^\tau$, so the $n$-torsion
subfunctor is separated as well. For the first question, a proper geometrically
locally integral fibre is the disjoint union of integral connected components.
The finite \'etale Stein-factor reduction therefore reduces to the usual
separatedness theorem for geometrically integral fibres. Separatedness
descends along that finite \'etale cover.

Representability alone does not make $P^\tau$ separated, answering the more
general question on 57:309(e) negatively. Here is a complete example over
$\mathbf C$. Put $\lambda_i=i$ for $0\leq i\leq5$ and define
\[
 Y=\left\{\sum_{i=0}^5\lambda_i^jX_i^2=0\quad(j=0,1,2)\right\}
 \subset\mathbf P^5_{\mathbf C}.
\]
This is a smooth connected complete intersection of three quadrics. To check
smoothness, a nonzero dependence among the three Jacobian rows gives a nonzero
polynomial of degree at most two vanishing at every $\lambda_i$ for which
$X_i\ne0$. There would be at most two such coordinates; the corresponding
Vandermonde columns are independent, and the three defining equations force
all their squares to vanish, a contradiction. The Koszul complex computes
$H^0(Y,\mathcal O_Y)=\mathbf C$ and $H^1(Y,\mathcal O_Y)=0$; adjunction gives
$\omega_Y\simeq\mathcal O_Y$. The involution
\[
 g[X_0:X_1:X_2:X_3:X_4:X_5]
 =[X_0:X_1:X_2:-X_3:-X_4:-X_5]
\]
is free: its fixed locus consists of the two coordinate planes, and on each
the relevant three Vandermonde columns are independent, so it misses $Y$.
All twenty three-column determinants, including the two fixed-plane ones,
are computed exactly in the accompanying code.

Set $R=\mathbf C[[t]]$, $R'=\mathbf C[[u]]$ with $t=u^2$, and form the
projective quotient
\[
 \mathcal X=(Y\times\operatorname{Spec}R')/
 \langle(g,u\mapsto-u)\rangle\longrightarrow\operatorname{Spec}R.
\]
The action is free, so the quotient map is finite \'etale. Consequently
$\mathcal X$ is regular and flat over $R$. Write $Z=Y/\langle g\rangle$.
The special fibre is $2Z$. If $T$ is the nontrivial order-two line bundle
given by the cover, then
$\nu_*\mathcal O_Y=\mathcal O_Z\oplus T$ and the square-zero ideal in
$\mathcal X_s$ is $T$. The cover is connected, so $T$ is nontrivial and
$H^0(Z,T)=0$. The vanishing of $H^1(Y,\mathcal O_Y)$ gives
$H^1(Z,\mathcal O_Z)=H^1(Z,T)=0$. Hence
$H^0(\mathcal X_s,\mathcal O)=\mathbf C$ and $H^1(\mathcal X_s,\mathcal O)=0$.
The same statements hold geometrically on the generic fibre, a form of $Y$.
Thus $f_*\mathcal O=\mathcal O_R$ universally and the Picard algebraic space
is unramified. Its identity is an open subgroup represented by $\operatorname{Spec}R$;
Raynaud's Proposition~5.3, since the residue field is separably closed, makes
the entire Picard space a scheme.

Nevertheless $\mathcal L=\mathcal O_{\mathcal X}(Z)$ satisfies
$\mathcal L^{\otimes2}\simeq\mathcal O_{\mathcal X}$ and
$\mathcal L_\eta\simeq\mathcal O$, while $\mathcal L|_Z\simeq T\ne\mathcal O_Z$.
Here $Z$ is Cartier because $\mathcal X$ is regular, and $2Z=\operatorname{div}(t)$.
So $\mathcal L$ and the identity are two sections of the represented tau-part
that agree generically and disagree on the special fibre. The tau-part is not
separated. The special fibre in this example is nonreduced; it does not
contradict the reduced-fibre theorem.

\subsection{E4: universal openness and constant dimension fail in general}

The repeated question ``$\underline{\mathrm{Pic}}^\tau_{X/S}$ est-il
universellement ouvert sur $S$ ?'' occurs on 57:309--311. The following
example disproves it in the unrestricted proper-flat setting, already in
characteristic zero. It also disproves unrestricted constancy of
$\dim\operatorname{Pic}^0_{X_s}$ in 57:313(i). It has nonreduced fibres, so
the separately written normal-morphism version remains a distinct question.

Let $E/\mathbf C$ be the elliptic curve $y^2z=x^3-xz^2$, with origin $o$, let
$S=E$, and let
$\mathcal P=\mathcal O_{E\times S}(\Delta-E\times o-o\times S)$ be the
normalized Poincar\'e bundle. On $V=E\times\mathbf P^1\times S$ put
\[
 M=\operatorname{pr}_{E,S}^*\mathcal P\otimes
       \operatorname{pr}_{\mathbf P^1}^*\mathcal O_{\mathbf P^1}(-2),
 \qquad \mathcal X=\operatorname{Spec}_V(\mathcal O_V\oplus M),\quad M^2=0.
\]
The split square-zero algebra is finite locally free of rank two over $V$.
Thus $\mathcal X/S$ is projective and flat. The vanishing
$H^0(\mathbf P^1,\mathcal O(-2))=0$ gives
$f_*\mathcal O_{\mathcal X}=\mathcal O_S$ after every base change.
Over a geometric point $s$, the splitting of units
$\mathcal O_{\mathcal X_s}^{\times}=\mathcal O_{V_s}^{\times}\times(1+M_s)$
gives
\[
 \operatorname{Pic}_{\mathcal X_s}
 \simeq\operatorname{Pic}_{E\times\mathbf P^1}\times H^1(V_s,M_s)
\]
as group schemes, with the last factor additive. A degree-zero line bundle
on $E$ has both $H^0$ and $H^1$ zero unless trivial, and both one-dimensional
when trivial. K\"unneth therefore gives
\[
 (h^0(M_s),h^1(M_s),h^2(M_s))=
 \begin{cases}(0,0,0),&s\ne o,\\(0,1,1),&s=o.\end{cases}
\]
Consequently $\dim\operatorname{Pic}^0$ is one away from $o$ and two at $o$.
SGA~3 VI$_B$, Proposition~4.2, says that a group scheme locally of finite
presentation is universally open along the identity precisely when its fibre
dimension is locally constant there. Thus the tau-part is not universally
open. The same statement can be seen in the extra Picard factor over the
local DVR at $o$: its equation is $tz=0$ in $\mathbf G_{a,R}$, and the open
$z\ne0$ maps to the closed point. The dimension argument alone suffices and
does not require this local equation.

There is a positive statement where the manuscript specifies normal fibres:
when the fibrewise Picard identity components have no additive subgroup,
FGA VI Corollary~2.7 makes their dimension locally constant, and SGA~3
VI$_B$~4.2 gives universal openness along the neutral part. Multiplication by
an integer invertible on the base is \'etale, so the same holds along components
whose order in the component group is prime to the residue characteristic.
In characteristic zero this covers the whole tau-part. For normal fibres
in characteristic $p$, the question remains open here on the components
whose classes have $p$-power order.

\subsection{E4: the product formula, with the correct hypotheses}

For $X,Y$ proper flat and finitely presented over $S$, each with universal
$H^0=\mathcal O_S$, the canonical map
\[
 \operatorname{Pic}^\tau_{X/S}\times_S\operatorname{Pic}^\tau_{Y/S}
 \xrightarrow{\sim}\operatorname{Pic}^\tau_{X\times_SY/S}
\]
is an isomorphism of fppf sheaves. Here is the argument completing the
question on 57:309(g). E2 represents the correspondence quotient by a
separated unramified group scheme. Its geometric fibres are torsion-free.
Indeed over an algebraically closed field, choose base points and normalize
a correspondence as in E2. An $n$-torsion class gives a based map
$a:X\to\operatorname{Pic}^0_Y[n]$. This target is affine. To see this, apply
Chevalley's theorem to the reduced connected Picard group: its affine part
maps to a finite $n$-torsion subscheme of its abelian quotient, and a nilpotent
thickening of an affine scheme is affine. Since $X$ is proper with
$H^0(X,\mathcal O_X)=k$, every map from $X$ to an affine scheme is constant.
The based map is therefore zero.

An algebraically trivial Picard class maps to zero in the discrete
correspondence scheme, and a tau-class maps to torsion, hence to zero by
the preceding paragraph. Over any base, its correspondence section is zero
on all geometric fibres. The equalizer with the identity is open because
the correspondence scheme is unramified, and contains every point, so the
section is zero. Fppf locally one can choose base points, and the normalized
bundle is then an external product. Restriction to the two axes says that
both factors are tau-trivial; uniqueness and fppf descent finish the proof.

Geometric connectedness alone, one of the uncertain marginal readings, is
insufficient. Take $X=\operatorname{Spec}k[\epsilon]/(\epsilon^2)$ and an
elliptic curve $Y/k$. Then $\operatorname{Pic}^\tau_X=0$, while the Picard
functor of $X\times Y$ has an extra additive factor $H^1(Y,\epsilon\mathcal O_Y)$
of dimension one. Its neutral dimension is two, whereas that of
$\operatorname{Pic}^\tau_Y$ is one. The universal $H^0=\mathcal O$ condition,
or geometric reducedness together with connectedness, removes this defect.

\subsection{E4: a proper normal family with nonproper tau-part}

\paragraph{Theorem E4.N (explicit counterexample).}
There is a projective flat morphism $X\to\operatorname{Spec}\mathbf C[[t]]$
with geometrically integral normal fibres and a section, for which the
represented separated finite-type group scheme $\operatorname{Pic}^\tau_{X/R}$
is not proper. In particular the normal-fibre branch of 57:310(5) is false,
and normal projectivity of $X/R$ does not imply projectivity of its tau-part.
Since the base has characteristic zero, its tau-part is universally open;
thus the implication proposed in 57:310(5') is false as well.

\paragraph{Construction.}
Put $R=\mathbf C[[t]]$, $K=\mathbf C((t))$. In $\mathbf P^5_R$ take the
three quadrics $\sum_i a_{ji}x_i^2=0$ with coefficient matrix
\[
 A(t)=\begin{pmatrix}
 1&1&1&1&1&1\\
 0&1&2&3&5&7\\
 0&1&2+t&9&25&49
 \end{pmatrix}.
\]
Call their intersection $Y$, and let $g$ fix the first three homogeneous
coordinates and negate the last three. Every three-column determinant of
$A(t)$ is nonzero in $K$; at $t=0$ exactly the determinant of columns
$0,1,2$ vanishes, and it equals $t$. These twenty exact determinants are
listed in the accompanying computation output. Any two columns are
independent. The generic surface is therefore a smooth complete intersection
of three quadrics: a deficient Jacobian would be supported on at most two
coordinates, and then the equations force all coordinates to vanish.
The same argument on the special fibre shows that its only singular points
are the four points supported on $0,1,2$ with squared coordinates
\[
 (x_0^2,x_1^2,x_2^2)=(1,-2,1).
\]
Indeed the kernel of the first three columns at zero is spanned by
$(1,-2,1)$, whose three entries are nonzero. After projectivization there
are precisely four choices of signs.

On the affine chart $x_2=1$ at each of these points, the first two equations
solve for $x_0,x_1$ in the completed local ring: their $2\times2$ Jacobian
has nonzero determinant since $x_0x_1\ne0$. Subtracting the second quadric
from the third leaves exactly
\[
 t+6x_3^2+20x_4^2+42x_5^2=0.                         \tag{E4.N.1}
\]
Thus $Y$ is regular at each of these points and the special fibre has an
ordinary double point there. Away from them, fibre smoothness makes the
total space regular. The squaring map $[x_i]\mapsto[x_i^2]$ on $\mathbf P^5_R$ is finite
locally free. The three linear equations with matrix $A(t)$ cut out a
$\mathbf P^2_R$, because the last three columns have unit determinant.
Thus $Y$, their inverse image under squaring, is finite flat over that
$\mathbf P^2_R$ and is flat over $R$; it is also a relative complete
intersection.
The fibres are connected by their Koszul resolution, are Cohen--Macaulay,
and are regular in codimension one; consequently they are normal and
geometrically integral. The same resolution gives $H^0=\mathbf C$ and
$H^1(\mathcal O)=0$ on both geometric fibres. Adjunction gives trivial
canonical bundle. In particular the generic fibre is a K3 surface.

The fixed locus of $g$ in projective space consists of the two coordinate
planes. The corresponding three-column matrices show that $g$ has no
fixed point on $Y_K$ and exactly the four displayed fixed points on $Y_0$.
The finite quotient $q:Y\to X=Y/\langle g\rangle$ exists as a projective
scheme: the linear action supplies an invariant Veronese homogeneous
coordinate ring. Since $2$ is invertible, invariants commute with every
base change. The quotient is flat over $R$ (its local invariant rings are
$R$-torsion-free) and its geometric fibres are normal integral schemes,
because invariant subrings of normal domains are normal. Reynolds
averaging and the preceding Koszul calculation give universal
$f_*\mathcal O_X=\mathcal O_R$ and $H^1(X_s,\mathcal O)=0$.
Choose a smooth closed point of $X_0$ away from the four quotient singularities.
Henselian lifting through the smooth locus gives an $R$-section of $X$.
This ensures that every $R$-point of the relative Picard sheaf is represented
by an actual line bundle.

\paragraph{Failure of extension.}
The connected finite \'etale double cover $q_K:Y_K\to X_K$ has a nontrivial
anti-character line bundle $T_K$, with $T_K^{\otimes2}\simeq\mathcal O$.
It is nontrivial even geometrically: otherwise the geometrically connected
K3 cover would split after extending $K$. Suppose it extended to a line
bundle $L$ on $X$. Then $q^*L$ is trivial on $Y_K$. Since $Y$ is regular,
a rational trivialization gives a Cartier divisor supported on $Y_0$.
There is only one reduced irreducible vertical divisor, $Y_0$, and it is
$\operatorname{div}(t)$. Thus $q^*L\simeq\mathcal O_Y$.

The natural $g$-linearization on a pullback line bundle acts trivially on
its fibre at any fixed point of $g$. On the other hand, any linearization
on $\mathcal O_Y$ is given by a unit $c\in H^0(Y,\mathcal O_Y)^*=R^*$ with
$c^2=1$; hence $c=+1$ or $-1$. Its generic restriction must be the sign
linearization defining $T_K$, so $c=-1$. This contradicts triviality on
fibres at the four fixed points. Therefore $T_K$ does not extend to $X$.

The projective geometrically integral family has a represented separated
Picard scheme, and its tau-part is open of finite presentation by the
standard finiteness theorem. Since $T_K$ is a tau-point and cannot extend,
the tau-part fails the valuative criterion for properness. There is no
conflict with E1: $T_K$ is a nonzero N\'eron--Severi torsion class; it is not
in the neutral component. Indeed $H^1(X_K,\mathcal O)=0$ and characteristic
zero make the neutral component trivial. The universally open assertion
follows from FGA VI Theorem~2.5 and its corollary: normal geometric fibres
have proper neutral Picard varieties, and every component order is
invertible over this base.

\subsection{E4: dimension jump even with geometrically reduced fibres}

\paragraph{Theorem E4.R.}
The constancy question on 57:313(i) is false even with geometrically reduced
fibres, the marginal hypothesis ``du moins pour fibres s\'eparables''.
Consequently the universal-openness questions 57:309(c), 310(3), and 311
also have negative answers under this weaker hypothesis than normality.

Keep the elliptic curve $E$, base $S=E$, origin $o$, and normalized
Poincar\'e bundle $\mathcal P$ on $C=E\times S$ used above. Let $V\to C$ be
the projective-line bundle compactifying the total space of
$\mathcal P^\vee$, and write $C\subset V$ for its zero section. Its conormal
bundle is $\mathcal P$. The doubled zero section $D=2C$ is the split ribbon
\[
 \mathcal O_D=\mathcal O_C\oplus\mathcal P,\qquad\mathcal P^2=0,
\]
with the splitting induced by $V\to C$. Define a finite $V$-scheme
\[
 X=\operatorname{Spec}_V\!\left(
 \mathcal O_V\times_{\mathcal O_D}\mathcal O_V\right).       \tag{E4.R.1}
\]
Writing $I=\mathcal O_V(-D)$, the fibre-product algebra is
$\mathcal O_V\oplus I$ as an $\mathcal O_V$-module. Locally, if $d$ is an
equation for $D$, its equation is $z(z-d)=0$. In particular $X\to V$ is
finite locally free of rank two, so $X/S$ is projective and flat.

Formation of this algebra commutes with arbitrary base change on $S$:
in the exact sequence
\[
 0\longrightarrow I\longrightarrow\mathcal O_V
     \longrightarrow\mathcal O_D\longrightarrow0
\]
all three terms are flat over $S$; $D$ is a square-zero extension of the
smooth $C/S$ by an invertible module. Thus $I_T$ remains the kernel after
any base change $T\to S$, and the fibre-product algebra remains the same
fibre product. Each geometric fibre of $V$ is integral and $d$ is a
nonzerodivisor. Hence the algebra in (E4.R.1), on that fibre, embeds in
$\mathcal O_{V_s}\times\mathcal O_{V_s}$ and is reduced. The two components
meet along the nonempty ribbon $D_s$, so the fibre is connected. Moreover
\[
 0\longrightarrow\mathcal O_X\longrightarrow
 \mathcal O_V\oplus\mathcal O_V\longrightarrow\mathcal O_D\longrightarrow0
\]
shows $H^0(X_T,\mathcal O)=H^0(T,\mathcal O_T)$ for every $T$: the map on
the two constant summands is their difference in the copy of
$\mathcal O_T\subset H^0(D_T,\mathcal O)$, so its kernel is the diagonal.
This proves universal cohomological flatness in degree zero directly.

For a geometric point $s$, the units Mayer--Vietoris sequence gives
\[
 H^0(V_s,\mathcal O)^*\times H^0(V_s,\mathcal O)^*
 \longrightarrow H^0(D_s,\mathcal O)^*
 \longrightarrow\operatorname{Pic}(X_s)
 \longrightarrow\operatorname{Pic}(V_s)^2
 \longrightarrow\operatorname{Pic}(D_s).
\]
Here $H^0(D_s,\mathcal O)^*=k^*\times H^0(E,\mathcal P_s)$ with the
second factor additive; the first map has image exactly the scalar $k^*$.
Also $\operatorname{Pic}^0(V_s)=\operatorname{Pic}^0(E)$, and restriction to
$D_s$ is the split pullback from $E$. The kernel on neutral components of
the last map is therefore the diagonal elliptic curve. Pullback along the
folding map $X_s\to V_s\to E$ splits this diagonal, giving an isomorphism
of group schemes
\[
 \operatorname{Pic}^0_{X_s/k}
  \simeq E\times\mathbf G_a^{\,h^0(E,\mathcal P_s)}.
\]
The units sequence and the splittings used here are functorial after a
$k$-algebra extension, so this is a group-scheme computation, not only a
calculation of rational points. Since $\mathcal P_s$ is a nontrivial
degree-zero bundle for $s\ne o$, the exponent is zero there and is one
at $o$. The neutral dimension jumps from one to two. SGA~3 VI$_B$~4.2
then rules out universal openness at the identity. Although the fibres
are reduced, they are not normal: the two smooth components have been
glued along a doubled divisor. This is the precise scope of the example.

\subsection{E4: the abelian and additive rank inequalities}

FGA VI Remark~2.8, pp.234--235, fixes the convention in 57:313(ii):
unprimed invariants belong to the special point $s$, primed invariants to
its generization $\eta$. Thus the manuscript asks for
$\alpha_s\le\alpha_\eta$ and $\lambda_s\ge\lambda_\eta$.
Here $\alpha,\mu,\lambda$ are the dimensions of the abelian quotient,
torus, and unipotent part of the reduced connected geometric Picard group.
Both inequalities hold in the represented proper-flat Picard setting;
in particular they hold under the geometrically reduced hypothesis in
the margin. They do not imply constancy of their sum, as E4.R shows.

We use Theorems~1.1 and~4.1 of Gille--Guralnick
(arXiv:2211.09876, pp.~1, 7--8), \emph{Unipotent Subgroups of
Stabilizers}. For separated group schemes of finite presentation,
they prove upper semicontinuity of two dimensions: the maximal
unipotent-subgroup dimension $d_u$, and the maximal dimension $d_c$
of an affine commutative subgroup. No flatness hypothesis is required. For a commutative algebraic group,
\[
 d_u=\lambda,\qquad d_c=\mu+\lambda,\qquad
 \alpha=\dim G-d_c.                                  \tag{E4.rank.1}
\]
The elementary Chevalley decomposition proves these identifications.
The cited proof takes the closure of a generic subgroup over a trait,
uses Raynaud's affineness theorem (proved in Prasad--Yu, Proposition~3.1,
p.5), and the finite smooth normalization after a finite extension
(Conrad's Appendix, Corollary~A.4, p.29). Hence no unproved extension
of SGA~3's smooth-group announcement is being used.

First work over a strictly henselian trait and suppose $P^\tau$ is
separated; it is of finite presentation. The $d_u$ theorem immediately
gives $\lambda_s\ge\lambda_\eta$. Let $G$ be the schematic closure of
$P^\tau_\eta$ in $P^\tau$. This is a closed finite-type flat subgroup:
flatness is torsion-freeness over the DVR, and multiplication and inversion
preserve the closure because products of flat schemes have schematically
dense generic fibre. Put $B=(P_s)^0_{\mathrm{red}}$ and
$H=(G_s)^0_{\mathrm{red}}$, over an algebraic closure of the residue field;
these are smooth connected groups.

The quotient $B/H$ is unipotent.
Here is a direct verification of the relevant assertion, also recorded
in Raynaud 1970, Proposition~6.3.1, p.52. A line bundle of order prime to
the residue characteristic on $X_s$ is represented by a finite \'etale
Kummer torsor. The equivalence between finite \'etale covers of a proper
scheme over a henselian local base and covers of its closed fibre
(Artin 1969, Theorem~3.1, p.30) lifts that torsor, with its group action,
to $X$. It therefore lifts the torsion Picard point. Its generic value is
in $P^\tau_\eta$, so its section lies in $G$. For every sufficiently large
prime $\ell$ invertible in the residue field, the finite group of
components of $G_s$ has no $\ell$-torsion, and hence
$(P_s)^0[\ell]\subset H$. Multiplication by $\ell$ is surjective on the
connected algebraic group $H$, so the quotient has no $\ell$-torsion.
A positive-dimensional torus or abelian quotient would have nonzero
$\ell$-torsion for every such $\ell$. The quotient is therefore unipotent.
It follows that $\alpha(P_s)=\alpha(G_s)$.

The flat group $G$ has constant fibre dimension. Apply upper
semicontinuity of $d_c$ and subtract in (E4.rank.1):
\[
 \alpha_s=\alpha(G_s)
 =\dim G_\eta-d_c(G_s)
 \le\dim G_\eta-d_c(G_\eta)=\alpha_\eta.
\]
This proves the two claimed inequalities over traits; constructibility
and the usual trait reduction give the specialization statement over a
locally noetherian base.

For completeness, separation of the whole Picard scheme is unnecessary.
If the proper-flat Picard functor $P$ is represented over the trait,
Raynaud's Proposition~5.2 makes the closure $E$ of the generic identity
\'etale, and his \S4.1 constructs a separated group scheme $Q=P/E$.
Here the passage from a separated group algebraic space to a scheme uses
precisely Raynaud Theorem~3.3.1 over the trait.
The image $Q'$ of the finite-presentation open subgroup $P^\tau$ is a
quasi-compact open subgroup of $Q$, hence a separated finite-presentation
group scheme. On every geometric fibre the kernel is zero-dimensional;
thus $P^\tau\to Q'$ preserves $\alpha,\mu,\lambda$ on the reduced connected
groups. Apply the preceding argument to $Q'$; prime-to-characteristic
torsion lifts after choosing $\ell$ also prime to the finite special
kernel. This proves the same conclusion when the represented tau-part
is nonseparated. The statement is not extended here to an arbitrary
nonrepresentable Picard functor without checking the same quotient data.

\paragraph{The prime-to-characteristic rank inequality.}
The original printed formulas in FGA VI Remark~2.8, pp.234--235,
are unambiguous: with $s$ special and $\eta$ generic,
\[
 2\alpha_s+\mu_s\le 2\alpha_\eta+\mu_\eta.       \tag{E4.rank.2}
\]
In the geometrically reduced setting of 57:313, this follows directly
from separation of $P^\tau$ proved above. Over a strictly henselian
trait, choose a prime $\ell$ invertible on the trait and not dividing
the orders of the torsion in either geometric N\'eron--Severi group.
The kernel $P[\ell]$ is \'etale by FGA VI Theorem~2.5 and is separated
because it lies in $P^\tau$. Each special point therefore extends to
a section, and distinct sections have distinct generic values.
The geometric cardinalities are respectively
$\ell^{2\alpha_s+\mu_s}$ and
$\ell^{2\alpha_\eta+\mu_\eta}$: prime-to-characteristic
multiplication on the unipotent part is an isomorphism, and on the
torus and abelian variety it has these familiar orders.
This proves (E4.rank.2), including its direction.
More generally the same proof applies whenever the relevant
prime-to-characteristic torsion subspace is separated.
This assertion is distinct from constancy of Picard dimension,
which E4.R refutes even with geometrically reduced fibres.

\subsection{Flatness and openness: a Heisenberg degeneration}

\paragraph{Theorem E4.H.}
There exists a smooth projective family of geometrically connected surfaces
$X/R$, with $R=W(k)$ and $k$ an algebraically closed field of characteristic
$2$, such that
\[
 \operatorname{Pic}^{\tau}_{X/R}
 \simeq\mu_{2,R}^{\,2}\times_R K,
 \qquad K=\operatorname{Spec}R\amalg\operatorname{Spec}k.       \tag{E4.H.1}
\]
The second summand of $K$ has the group law inherited from the nonzero
component of the constant group $(\mathbf Z/2)_R$; its product with itself
lands in the identity summand. Consequently $\operatorname{Pic}^{\tau}_{X/R}$
is not open over $R$, even though $X/R$ is smooth. Over
$A=R/4R$, its fibrewise neutral subgroup is represented by
$\mu_{2,A}^{\,2}$ and is finite flat, whereas its tau-part is not flat.
This refutes both the normal-family universal-openness question on 57:310(3)
and the implication written on 57:311, even with represented neutral part.

\paragraph{The finite group and its characters.}
Define a finite flat group $G/R$ by
\[
 G=\mu_{2,R}\times (\mathbf Z/2)_R^2,
\quad
 (z,a,b)(z',a',b')=
 \bigl(zz'(-1)^{ba'},a+a',b+b'\bigr).                        \tag{E4.H.2}
\]
Here the displayed product decomposition is a decomposition of schemes;
the group is not commutative. The exponent is taken modulo two.
The elementary bilinear cocycle identity proves associativity. Its central
$\mu_2$ has quotient $(\mathbf Z/2)^2$. The generic geometric group is the
dihedral group of order eight, with commutator subgroup its central order-two
subgroup. The special group is the commutative group
$\mu_{2,k}\times(\mathbf Z/2)_k^2$, since $-1=1$ in $k$.
This is the case $p=2$, $e=p-1$ of Raynaud's construction
\cite[Lemma~4.2.1, pp.~135--136]{BCERaynaud1979}.

A character of $G_T$ restricts on the central $\mu_{2,T}$ to the locally
constant exponent $n\in\mathbf Z/2$. Its values on $(1,1,0)$ and
$(1,0,1)$ are arbitrary elements $u,v\in\mu_2(T)$. The commutator relation
imposes exactly $(-1)^n=1$: the component $n=0$ is defined over all of $R$,
and $n=1$ exists precisely over $R/2R=k$. There are no further relations.
Thus the character functor, including nonflat tests, is
\[
 \underline{\operatorname{Hom}}(G,\mathbf G_m)
 =\mu_{2,R}^{\,2}\times_R K.                               \tag{E4.H.3}
\]
In particular its coordinate module is $R^4\oplus k^4$.
The code \texttt{check\_raynaud\_group.py} checks all 64 cocycle identities,
all 512 associativity instances of the generic eight-element group, its
nontrivial commutator, and its four characters. The Hopf-scheme calculation
(E4.H.3), rather than point counting, supplies the special length eight.

\paragraph{A completely specified complete-intersection realization.}
The following finite formula fixes the degree, every section, and every
coefficient. Use the 24 coordinates $x_{r,n,a,b}$, where $1\leq r\leq3$
and $n,a,b\in\{0,1\}$. The central $\mu_2$ has weight $n$, and
\[
 A(x_{r,n,a,b})=x_{r,n,a+1,b},\qquad
 B(x_{r,n,a,b})=(-1)^{na}x_{r,n,a,b+1}.                    \tag{E4.H.4}
\]
The indices $a,b$ are modulo two. These are three copies of the regular
coordinate representation, so $P=\mathbf P^{23}$ is the projectivization
of the threefold direct sum of the contragredient regular representation
in Raynaud's Lemma~4.2.2. Its free-action open $U$ has complement of
fibrewise codimension at least three.

Order the coordinates lexicographically. Let $\mathcal A$ be the ordered
finite set of degree-192 monomials with even total $\mu_2$ weight. Form
\[
 F=\sum_{\alpha\in\mathcal A}u_\alpha x^\alpha,
 \qquad F\,A(F)\,B(F)\,AB(F)=\sum_{|\beta|=4}u^\beta p_\beta(x).
                                                               \tag{E4.H.5}
\]
The actions fix the $u$ variables. On even-weight forms $A,B$ commute,
so every integral polynomial $p_\beta$ is $G$-invariant of degree 768.
Retain zero polynomials and repetitions; no basis choice is involved.
Let $\mathcal B$ be the ordered set of the multiindices $\beta$ above.
Define
\[
 k=\overline{\mathbf F_2(c_{i\beta}:1\leq i\leq21,
                                  \ \beta\in\mathcal B)},
 \quad R=W(k),\qquad
 s_i=\sum_{\beta\in\mathcal B}[c_{i\beta}]p_\beta.          \tag{E4.H.6}
\]
Here all variables are algebraically independent and brackets are
Teichm\"uller representatives. These are finite definitions, with
\[
 h=|\mathcal A|=\tfrac12\{\binom{215}{23}+\binom{107}{11}\},
 \qquad |\mathcal B|=\binom{h+3}{4}.
\]

To type the descent precisely, put $M=P/\mu_2$ and $Q=M/H=P/G$, where
$H=(\mathbf Z/2)^2$. The central norm of $\mathcal O_P(1)$ is a bundle
$\mathcal N$ on $M$ with pullback $\mathcal O_P(2)$, naturally
$H$-linearized. Put $\mathcal L=\operatorname{Norm}_H(\mathcal N)$ on
$Q$. It is ample with pullback $\mathcal O_P(8)$. These are Raynaud's
norm bundles with their factorization made explicit
\cite[proof of Proposition~4.2.3, p.~137]{BCERaynaud1979}.
Even-weight degree-192 forms are sections of $\mathcal N^{96}$.
Their $H$-norms, hence all $p_\beta$, are sections of $\mathcal L^{96}$.
These integral formulas avoid any assumption about reduction of a chosen
basis of sections. Define
\[
 X=V(s_1,\ldots,s_{21})\subset Q,
 \qquad Y=q^{-1}(X)\subset P.                              \tag{E4.H.7}
\]
Thus $Y$ is specified by 21 equations of the fixed degree 768.

\paragraph{Why this fixed degree suffices.}
Write $V=U/G$. The maps $U\to U/\mu_2\to V$ are respectively a degree-two
$\mu_2$-torsor and a degree-four finite \emph{\'etale} $H$-torsor.
The scheme $V$ is smooth of relative dimension 23: flatness descends
from $U$, and geometric regularity descends through the faithfully flat
finite maps on fibres. The complement $Q_k\setminus V_k$ has dimension
at most 20.

The sections $p_\beta$ generate $\mathcal L^{96}$ on $Q_k$. Indeed even
coordinate powers generate $\mathcal N^{96}$ on $M_k$. Since $k$ is
infinite, a linear combination can be chosen nonzero at all points of
any finite $H$-orbit. Its norm is nonzero at the corresponding point of
$Q_k$, and every norm is a linear combination of the coefficient
sections $p_\beta$.

They also span every first jet at $v\in V(k)$. The first infinitesimal
neighborhood of $v$ has length 24; its inverse image $Z$ in $P_k$ has
length $8\cdot24=192$. A length-$\ell$ subscheme of projective space
over an infinite field is interpolated by $\mathcal O(d)$ for
$d\geq\ell-1$: choose a hyperplane avoiding its support and let $W_d$
be the span of degree-at-most-$d$ monomials in its length-$\ell$
coordinate algebra. If $W_{d+1}=W_d$, this space is stable under every
coordinate and is the whole algebra. Otherwise its dimension strictly
increases. Since $\dim W_0=1$, it reaches the algebra by $d=\ell-1$.
Thus degree-192 forms interpolate all sections on $Z$.

Taking $\mu_2$ invariants preserves that surjection: projection onto
weight zero is exact even in characteristic two. Sections of
$\mathcal N^{96}$ therefore interpolate arbitrary first jets at the
four points above $v$ in $U/\mu_2$. Prescribe the desired jet at one
point and unit jets at the other three, in $H$-compatible
trivializations. The norm has the desired jet at $v$. Consequently
the span of the $p_\beta$ surjects onto every first jet on $V_k$.

The incidence proof of Bertini now makes 21 generic sections smooth on
$V_k$, while global generation and the dimension bound 20 make them
avoid its complement. The independent coefficients in (E4.H.6) specify
such a generic tuple. Amplitude makes the intersection a nonempty
surface. Thus $X_k$ is smooth and contained in $V_k$. Since $Q$ is smooth
near $X_k$ and the 21 equations cut its smooth special fibre in codimension
21, they form a fibrewise regular sequence there. The local flatness
criterion makes $X/R$ flat near $X_k$, and the smooth special fibre then
makes $X/R$ smooth on that neighborhood. These properties hold throughout
$X$: their bad loci are closed in the projective
$R$-scheme $X$, and a nonempty closed image in the local base must meet
the closed point. Therefore $X/R$ is smooth projective, $Y\to X$ is a
$G$-torsor, and the 21 pullback equations are a relative complete
intersection in $P^{23}_R$. This verifies all hypotheses of Raynaud's
Proposition~4.2.3 with an explicit power $m=96$.

A complete intersection of dimension at least two has
$H^0(\mathcal O)=k$, $H^1(\mathcal O)=0$, and trivial Picard tau-group,
including when singular or nonreduced; the last statement is
\cite[Theorem~1.1.8(iii), proved in \S3.4]{BCECotnerZavyalov2024}.
Hence $Y/R$ has universal $H^0=\mathcal O_R$ and
$\operatorname{Pic}^\tau_{Y/R}=0$, by the fibrewise assertion and the
unramified identity-section criterion. Geometric connectedness of $Y$
and its surjection onto $X$ make the smooth fibres of $X$ geometrically
integral and normal. The coordinate-check script verifies the signed
representation and the finite index cardinalities. It does not expand
the enormous norm polynomial or claim a computational smoothness test;
the interpolation and incidence arguments supply that proof.

\paragraph{Picard descent and the two failures.}
The kernel of
$\operatorname{Pic}_{X/R}\to\operatorname{Pic}_{Y/R}$ is the character
functor of $G$. Indeed a line bundle trivialized after pullback to $Y$ is a
$G$-linearization of $\mathcal O_Y$; universal $H^0(Y,\mathcal O_Y)=R$
identifies the linearization with a unit on $G$ satisfying exactly the
character identity. This argument is fppf-local on every test scheme and
therefore identifies the full sheaves, including their nilpotents.
Since $\operatorname{Pic}^\tau_{Y/R}=0$, every tau-class lies in this
kernel. Conversely the character group in (E4.H.3) is killed by two and
hence lies in tau. This proves (E4.H.1) and its compatibility with arbitrary
base change.

The open-and-closed component
$\mu_{2,k}^{\,2}\subset\operatorname{Pic}^\tau_{X/R}$ has image the closed
point of the trait. The map is therefore not open, and a fortiori not
universally open. This closes the positive-residue-characteristic gate;
no implication from nonflatness alone is being used.
After base change to $A=R/4R$, the neutral component of the only geometric
fibre lies in the first summand, so the fibrewise neutral subfunctor is
exactly $\mu_{2,A}^{\,2}$, finite locally free of rank four. The other
summand has coordinate module $k^4$ and is not $A$-flat: a nonzero finite flat module over the
local ring $A$ is free, whereas this module is annihilated by $2\ne0$. Thus flatness of the represented
neutral part does not imply flatness of tau. The construction depends on
Raynaud's noncommutative finite-flat group, not on the disputed tangent-space
argument for ordinary Igusa deformations.

\subsection{The quotient question in its written generality}

\paragraph{Theorem E4.Q (refutation of 57:309(b)).}
The manuscript asks for ``existence du pr\'esch\'ema quotient''
$\underline{\mathrm{Pic}}_{X/S}/\underline{\mathrm{Pic}}^\tau_{X/S}$.
In this generality the answer is negative, even when the tau-part is flat.
Let
\[
 S=\operatorname{Spec}\mathbf R[[t]],\qquad
 X=V(x^2+y^2-tz^2)\subset\mathbf P^2_S .                  \tag{E4.Q.1}
\]
This is Mumford's example as proved in
\cite[Example~4.14, pp.~31--33]{Kleiman2005}. Its Picard fppf sheaf
$P$ is an algebraic space, not a scheme. After the finite \'etale base change
$S'=\operatorname{Spec}\mathbf C[[t]]$, its component of total degree $n$
is obtained by gluing copies $S'_{a,b}$, $a+b=n$, along their generic point.
Complex conjugation interchanges the origins $(a,b)$ and $(b,a)$.
For $a\ne b$ these origins have no common affine neighborhood: their two
trait sections agree generically, contrary to separatedness of such an
affine neighborhood. If $P$ were a scheme, an affine neighborhood of their
common descended point would pull back to just such an affine neighborhood.

In this example $P^\tau$ is exactly the identity subgroup $S$. Indeed the
geometric generic conic has Picard group $\mathbf Z$, and the geometric
special fibre, two projective lines meeting in one point, has Picard group
$\mathbf Z^2$ by multidegree. Both have trivial neutral and torsion parts.
In the displayed scheme $P_{S'}$, the tau-subfunctor is consequently the
identity copy $S'_{0,0}$; the explicit representation also rules out any
additional infinitesimal structure. This identification descends to $S$.
Therefore, as fppf sheaves of groups,
\[
 P/P^\tau=P/S\simeq P,
\]
which is not a scheme. Here $S$ denotes the trivial group object, not a
nontrivial quotient of the underlying base. Its identity immersion into
$P$ is open and not closed. Thus this example is consistent with the flat-only algebraic-space assertion
in the prompt summary of H6. The stronger scheme conclusion in the fuller
\texttt{HUNCHES.md} formulation additionally assumes closedness, which fails here.

The stronger variant that presupposes $P$ itself schematic, or imposes
closedness of $P^\tau$ while dropping flatness, is a distinct question.
No answer to those extra variants is inferred from (E4.Q.1). In particular,
this example has a geometrically reducible special fibre and gives no
counterexample to the geometrically integral or normal question on 57:315.

\paragraph{Two further consequences of E4.H.}
The character group in (E4.H.1) is killed by two. Its identity section is
not open: in the component $\mu_{2,R}^{\,2}$, the four geometric generic
points specialize to the same special identity. Thus the image of the
whole group under $[2]$ is not open. Question 57:310(4) is therefore false
even for a smooth projective family, extending the simpler cuspidal
field example.

For the margin question about ``Multiplicit\'e de
$\underline{\mathrm{Pic}}^\tau$'' and specialization on 57:310, the same
example makes several precise natural interpretations explicit. The
generic fibre has four reduced points, each of multiplicity one. The
special fibre consists of two copies of $\mu_{2,k}^{\,2}$, each of length
four. Total length changes from four to eight, the number of geometric
components from four to two, and the length of the neutral geometric
component from one to four. Thus none of these numerical multiplicities
is specialization invariant without further hypotheses. The manuscript's
additional words about existence and descent do not specify an object;
we do not replace them by a different theorem.

\subsection{E4: closedness of the numerical locus without projectivity}

\paragraph{Theorem E4.T.}
Let $f:X\to S$ be a proper flat morphism of finite presentation with
geometrically integral fibres. The base scheme $S$ is arbitrary. For every
invertible sheaf $L$ on $X$, the set
\[
 S^\tau(L)=\{s\in S:L_{\bar s}\text{ is numerically trivial}\}
                                                        \tag{E4.T.1}
\]
is open and closed. Equivalently, the open subgroup algebraic space
$\operatorname{Pic}^{\tau}_{X/S}$ is closed in
$\operatorname{Pic}_{X/S}$. Thus projectivity is unnecessary in the
question ``projectif est-il n\'ecessaire ?'' on 57:314. Here numerical
triviality and tau-equivalence agree for proper schemes over a field;
see \cite[Remark~6.14, pp.~57--58]{Kleiman2005}.

\paragraph{1. The trait assertion.}
Let $R$ be a discrete valuation ring, with fraction field $K$, and let
$X/R$ satisfy the hypotheses. Suppose $L_K$ is numerically trivial.
We prove that $L_k$ is numerically trivial. Replacing $R$ by its completion
preserves the special fibre and the hypotheses; numerical triviality of
the generic bundle is preserved under field extension. We may therefore
assume $R$ complete, hence excellent. Flatness and integrality of the generic
fibre make $X$ integral. If $\xi$ is the generic point of its reduced
integral special fibre, $\mathcal O_{X,\xi}$ is a one-dimensional local
domain with maximal ideal generated by a uniformizer of $R$. It is a DVR.

Zheng's Galois alteration lemma
\cite[Lemma~3.8, pp.~313--314]{BCEZheng2009} gives a finite extension of
complete traits $R\subset R'$, a regular integral scheme $Y$ projective
over $R'$, and a proper generically finite map
\[
 \pi:Y\longrightarrow X
\]
equivariant for a finite group $G$ acting trivially on $X$. Moreover
$k(Y)^G/k(X)$ is purely inseparable. This is precisely the definition of
Galois alteration used on p.~313 of that source. To apply the lemma take
the input group trivial and the dense open $X_K$. Its compactification
$Y$ is the projective end of a plurinodal fibration over $R'$. The source
$X'$ of the open immersion in the general statement is proper over $X$,
hence over $R$. Since $R'/R$ is separated, $X'$ is proper over $R'$ as
well. Its map into the separated projective end is therefore proper.
The dense open immersion is consequently an isomorphism. The lemma also makes the
boundary a strict normal-crossings divisor. Only regularity of $Y$ and
the intersection theory below are needed: regularity of a geometric
special fibre is not assumed.

Put $M=\pi^*L$ and write
\[
 Y_{k'}=\sum_{i=1}^r m_iY_i,\qquad m_i>0.                \tag{E4.T.2}
\]
The special fibre is connected. Indeed $H^0(Y,\mathcal O_Y)$ is finite
over the complete local ring $R'$ and is a domain, hence a local ring;
Stein factorization then gives connectedness. Each $Y_i$ is a Cartier
divisor on the regular $Y$. Let $d=\dim Y_{K'}$. The case $d=0$ is
immediate, since every line bundle on a proper zero-dimensional scheme
is numerically trivial. Assume $d\geq1$, and choose a relatively ample
line bundle $H$ on $Y$.

\paragraph{2. The vertical correction.}
Define the rational symmetric matrix $A$ and vector $b$ by intersection
numbers on the special fibre:
\[
 A_{ij}=(Y_i\cdot Y_j\cdot H^{d-1}),\qquad
 b_i=(M\cdot Y_i\cdot H^{d-1}).                          \tag{E4.T.3}
\]
For $i\ne j$ one has $A_{ij}\geq0$, with strict inequality exactly when
$Y_i$ and $Y_j$ meet. This follows because their proper intersection is
an effective cycle and $H$ is ample on the projective special fibre.
The relation $\sum_jm_jY_j=\operatorname{div}(\varpi')$ gives $Am=0$.
The intersection graph is connected. Consequently
\[
 v^tAv=-\sum_{i<j}A_{ij}m_im_j
       \left(\frac{v_i}{m_i}-\frac{v_j}{m_j}\right)^2,
 \qquad \ker A=\mathbf Qm.                              \tag{E4.T.4}
\]
Constancy of intersection numbers in a flat projective family gives
\[
 \sum_i m_i b_i=(M_{K'}\cdot H_{K'}^{d-1})=0.
\]
Thus there is $a\in\mathbf Q^r$ with $Aa=-b$. Set
\[
 D=\sum_i a_iY_i,\qquad N=M+D\quad\text{in }\operatorname{Pic}(Y)_\mathbf Q.
                                                               \tag{E4.T.5}
\]
By construction $(N|_{Y_i}\cdot H|_{Y_i}^{d-1})=0$ for every $i$.
For $d=1$ this already says that $N|_{Y_i}$ is numerically trivial.
For $d\geq2$, the Hodge index theorem gives
\[
 (N|_{Y_i}^{\,2}\cdot H|_{Y_i}^{d-2})\leq0.             \tag{E4.T.6}
\]
Equality implies numerical triviality. The higher-dimensional version
follows by cutting with sufficiently ample integral hyperplanes and the
surface Hodge theorem; the equality implication is
\cite[Theorem~6.3(h), proof pp.~53--54, and Theorem~B.27, p.~79]{Kleiman2005}.
These statements apply to integral projective schemes, with no smoothness
hypothesis. Over a nonclosed field, pass to an algebraic closure: the
geometric components of an integral scheme are conjugate, and the
intersection numbers of descended bundles are equal on these components
up to their common positive multiplicity. The degree-zero condition,
the inequality, and its equality conclusion therefore give the same
statement over the original field. Denominators in $N$ are cleared before
applying the line-bundle theorem.

The same flat-family intersection calculation gives
\[
 \sum_i m_i(N|_{Y_i}^{\,2}\cdot H|_{Y_i}^{d-2})
       =(M_{K'}^{\,2}\cdot H_{K'}^{d-2})=0.              \tag{E4.T.7}
\]
Every summand is nonpositive and every $m_i$ is positive. Each is therefore
zero, and $N$ is numerically trivial on every special component, hence
on the whole special fibre. It is also numerically trivial on the generic
fibre. This construction gives a rational vertical correction; it does
not assert that an arbitrary generic line bundle extends integrally.

\paragraph{3. Galois symmetry removes the correction.}
Replace $D$ by $|G|^{-1}\sum_{g\in G}g^*D$. This is a rational average,
so no invertibility assumption on $|G|$ in $R$ is used. Since $M$ is pulled
back from $X$, it is $G$-invariant, and numerical triviality is preserved
by each (possibly semilinear over $R'$) automorphism. Thus $M+D$ remains
numerically trivial on both fibres.

Let $Z$ be the normalization of $X$ in $k(Y)$. It is finite over $X$ by
excellence. Normality of $Y$ factors $\pi$ as
\[
 Y\xrightarrow{h}Z\longrightarrow X,
\]
where $h$ is proper birational between normal schemes. The finite group
$G$ acts transitively on the prime divisors of $Z$ above $\xi$.
Here is the precise valuation argument. The DVR $\mathcal O_{X,\xi}$ has
a unique extension to the purely inseparable field $k(Y)^G$; the field
extension $k(Y)/k(Y)^G$ is Galois with group the faithful quotient of $G$,
and its Galois group permutes transitively the extensions of that valuation.
The same transitivity holds for their strict transforms on $Y$, since
a proper birational map to a normal scheme is an isomorphism at the
generic point of every target prime divisor.

These are exactly the components $Y_i$ that dominate $X_k$. Indeed a
vertical prime of the finite scheme $Z\to X$ has dimension $d$, hence
maps onto the integral $d$-dimensional special fibre. Their multiplicities
$m_i$ in (E4.T.2) are constant on the orbit: $G$ fixes the uniformizer of
$R$, whose divisor is the ramification index of $R'/R$ times (E4.T.2).
The averaged coefficients $a_i$ are constant on this orbit as well.
Consequently $a_i/m_i$ has a single value $c$ on every nonexceptional
vertical component. Subtract $c\operatorname{div}(\varpi')$ from $D$.
This is a principal rational divisor, so it preserves all numerical
conclusions and makes $D$ supported on $h$-exceptional prime divisors.

For any curve contracted by $h$, the degrees of $M$ and of $N=M+D$ are
zero. This includes curves over generic-fibre points, since there $D=0$
and $N=M$ is numerically trivial, as well as the special-fibre curves
treated above. Hence both $D$ and $-D$ are $h$-nef. The negativity lemma for proper
birational morphisms of normal schemes
\cite[Lemma~11.60, p.~436]{KollarFamilies} applied to both signs makes
$D\geq0$ and $D\leq0$, because $h_*D=0$. Thus $D=0$.
We conclude that $M|_{Y_{k'}}=\pi^*L|_{Y_{k'}}$ is numerically trivial.
Numerical triviality is detected by proper surjective pullback: for an
integral curve $C\subset X_k$, take a closed point of the generic fibre
of $Y_{k'}\times_{X_k}C\to C$ and close it to an integral curve $C'$.
The map $C'\to C$ is generically finite of positive degree, and the
projection formula gives
$0=\deg(M|_{C'})=[k(C'):k(C)]\deg(L|_C)$.
This proves the trait assertion.

\paragraph{4. From traits to an arbitrary base.}
First suppose $S$ is locally noetherian. The set (E4.T.1) is open by
the proper-flat Picard tau theorem
\cite[Theorem~2.1]{LaurentSchroer2024}; proper geometrically integral
fibres give universal $f_*\mathcal O_X=\mathcal O_S$.
Every specialization in a locally noetherian scheme is realized by a
map from the spectrum of a DVR (Stacks, Lemma~28.5.10, tag~054F).
The trait assertion shows that (E4.T.1) is stable under specialization.
Locally on the noetherian base it is constructible, being open; a
constructible specialization-stable subset is closed. Thus it is open
and closed.

For arbitrary $S$ the assertion is local, so take $S=\operatorname{Spec}A$.
Noetherian approximation descends $X/S$ and $L$ to some scheme
$S_i=\operatorname{Spec}A_i$ of finite type over $\mathbf Z$, with $X_i/S_i$
proper, flat and finitely presented. The needed exact limit results are
Stacks, Lemmas~32.10.1, 32.13.1, and 32.8.7 (tag~04AI), together with
finite-presentation descent of an invertible module. To ensure geometric
integrality of every fibre, let $E_i\subset S_i$ be the locus where the
fibre is geometrically reduced and geometrically irreducible. It is
constructible by Lemmas~37.26.5 and 37.27.7 (the latter tag~055B), and
$S\to S_i$ has image in $E_i$. Lemma~32.4.10, tag~05F4, shows that for
a sufficiently large index the entire $S_j\to S_i$ has image in $E_i$.
Replacing the model by $S_j$ therefore gives proper flat geometrically
integral fibres, while retaining $L_j$.
The cited limit and constructibility sources are collected in
\cite{BCEStacksTauDescent}.

The already proved noetherian assertion says $S_j^\tau(L_j)$ is open and
closed. Numerical triviality is invariant under field extension, so
$S^\tau(L)$ is exactly its inverse image. This proves (E4.T.1) for every
base. Finally, an fppf-local representative of a Picard section is an
actual invertible sheaf; applying (E4.T.1) on such a cover and descending
the open-and-closed locus proves the asserted closed immersion of the
tau subfunctor into the full Picard algebraic space. No schematicity
of the full Picard functor over $S$ is assumed or deduced.

\subsection{Fibrewise smoothness can improve at specialization}

\paragraph{Theorem E4.S.}
There is a projective flat family of geometrically integral reduced surfaces
over $R=\overline{\mathbf F}_2[[t]]$ whose geometric generic Picard
tau-part is nonsmooth, whereas its special Picard tau-part is smooth.
Thus the locus of fibrewise smoothness is not open. This is a precise
negative answer to the smoothness interpretation of 57:309(b$'$);
it does not identify fibrewise smoothness with relative smoothness.

Let $C\subset\mathbb P^2_R$ be the cubic
\[
 y^2z+txyz+tyz^2=x^3,\qquad o=[0:1:0].                 \tag{E4.S.1}
\]
Its discriminant and $c_4$ are
$\Delta=t^4(1+t^2)$ and $c_4=t^4$. Hence its generic fibre is a
smooth ordinary elliptic curve, with $j=t^8/(1+t^2)\ne0$.
The special fibre is the integral cuspidal cubic $y^2z=x^3$.
The coefficient of $y^2z$ makes $o$ smooth over $R$.
The involution
\[
 \sigma[x:y:z]=[x:y+tx+tz:z]                           \tag{E4.S.2}
\]
is inversion generically and fixes $o$.
Let $B/R$ be the constant ordinary elliptic curve
$y^2z+xyz=x^3+z^3$, let $P=(0,1)$ on its affine chart,
and let $\Gamma=(\mathbb Z/2\mathbb Z)_R$ act on $C\times_R B$
through $(\sigma,t_P)$. The nonzero point $P$ has order two, so this
action is free. Set
\[
 X=(C\times_R B)/\Gamma,\qquad Y=B/\langle P\rangle.
\]
The finite-group quotient is projective: an orbit tensor product of
an ample line bundle supplies an ample linearized line bundle.
Its quotient map is finite \'etale, so flatness descends, and all
geometric fibres of $X/R$ are integral and reduced. The invariant
section $o\times B$ gives a section $i:Y\to X$.

Write $G=\operatorname{Pic}^0_{C/R}$. It is a smooth separated group
scheme of finite presentation: the projective integral curve family
has its Picard scheme, and $H^2(C_s,\mathcal O_{C_s})=0$ removes
all infinitesimal obstructions. The degree-zero locus is its neutral
part and also its tau-part. The equality $\sigma^*=[-1]$ on $G$
holds generically; it holds everywhere because $G$ is flat over
the trait and the target is separated.
There is a base-change compatible isomorphism
\[
 \operatorname{Pic}^\tau_{X/R}
       \simeq Y^\vee\times_R G[2].                  \tag{E4.S.3}
\]
Here is the descent argument, including the singular special fibre.
Rigidification along $o$ makes a representing line bundle on $C$
unique up to a unique rigidified isomorphism. For a point of $G[2]$
the unique isomorphism with its $\sigma$-pullback respects this
rigidification, and its square is the identity. It therefore supplies
a canonical $\Gamma$-linearization. Pulling to $C\times B$ and
descending gives a line bundle on $X$ trivial on $i(Y)$.
Together with pullback from $Y$, this defines (E4.S.3).
Conversely, subtract the restriction to $i(Y)$ from a tau-trivial
line bundle on $X$. Its pullback to $C\times B$ is a bundle from
$C$ alone, by the product formula for tau and its trivialization on
$o\times B$. Descent makes its class $\sigma$-invariant. Its
linearization agrees with the canonical one: their ratio is a unit
on the proper connected product and is one along $o\times B$.
These constructions are inverse and commute with arbitrary base
change. Numerical triviality is detected by the finite surjective
cover, by the degree formula on curves. This proves (E4.S.3) as a
group-scheme assertion.

Over a geometric generic point,
$G[2]\simeq\mu_2\times(\mathbb Z/2\mathbb Z)$, so (E4.S.3)
is nonsmooth and has dimension one. On the special fibre,
$G_s\simeq\mathbb G_a$. Indeed, the cusp has normalization
$x=u^2,y=u^3$, conductor $(u^2)$, and the units quotient
\[
 (k[u]/(u^2))^\times/k^\times\simeq\mathbb G_a.
\]
The normalization has no degree-zero Picard group, giving the
assertion. Since $[2]$ is zero on $\mathbb G_a$,
\[
 (\operatorname{Pic}^\tau_{X/R})_s
       \simeq Y_s^\vee\times\mathbb G_a
\]
is smooth of dimension two. Its tangent dimension is two, as is
the generic tangent dimension, whereas the defect
$\dim\operatorname{Lie}-\dim$ drops from one to zero.
Consequently that defect is not upper semicontinuous either.
Relative smoothness remains an open condition; in this example it
fails at every special point because the relative dimension jumps.
The separate dimension function is upper semicontinuous, in agreement
with FGA VI Corollary~2.7, p.234.

The finite polynomial identities, Jacobian checks, discriminant,
and cusp parameterization are recorded in
\texttt{code/check\_cuspidal\_igusa.output.json}.
The singular-cubic extension of the rigidified descent proof received
a separate targeted AI audit from /root/dgj, which found no gap.

\subsection*{E4.Curve: schematicity for every smooth proper curve family}

\paragraph{Theorem E4.Curve.}
Let $C\to S$ be a proper smooth finitely presented morphism of schemes
with geometrically connected one-dimensional fibres, where $S$ is an
arbitrary scheme. Then the full fppf Picard functor $\operatorname{Pic}_{C/S}$
is represented by a smooth separated $S$-scheme. This settles the
geometrically normal curve case of the last question at 57:315. It does
not settle higher-dimensional families, or geometrically integral
singular curves without further hypotheses.

\paragraph{Proof.}
The arithmetic genus is locally constant, so we work on a locus where
it is a fixed integer $g$. If $g=0$, the dual of the relative canonical
bundle is relatively ample. If $g\geq2$, the relative canonical bundle
is relatively ample. In either case the family is locally projective
on $S$, and the projective Picard theorem gives a smooth separated
Picard scheme. These are exactly the arguments of
\cite[\S9.3, Theorem~1 and Remark~2, p.~252]{BCEBLR1990}.

Suppose $g=1$. Put $E=\operatorname{Pic}^0_{C/S}$; this is an abelian
scheme by \cite[\S9.4, Proposition~4, p.~260]{BCEBLR1990}.
The degree map decomposes the Picard algebraic space as
\[
 \operatorname{Pic}_{C/S}=\coprod_{n\in\mathbb Z}\operatorname{Pic}^n_{C/S}.
\]
The diagonal divisor defines $C\to\operatorname{Pic}^1_{C/S}$, and
this map is an isomorphism: after an etale cover giving a section of
$C/S$, it is the usual genus-one Abel map. Thus $C$ is an $E$-torsor.
For $n\neq0$, the map
\[
 C\longrightarrow\operatorname{Pic}^n_{C/S},\qquad
        x\longmapsto[\mathcal O_C(nx)]                  \tag{E4.Curve.1}
\]
is equivariant for multiplication by $n$ on $E$. After an etale
cover trivializing $C$, it is a translate of $[n]:E\to E$.
Consequently it is a finite locally free $E[n]$-torsor, and
\[
        \operatorname{Pic}^n_{C/S}=C/E[n]               \tag{E4.Curve.2}
\]
as fppf sheaves. This conclusion also holds when $n$ is divisible by
a residue characteristic; etaleness of $E[n]$ is unnecessary.

We show that (E4.Curve.2) is a scheme. First suppose $S$ is locally
noetherian. The assertion is Zariski local on $S$. Fix $s\in S$ and
choose an affine open $V\subset C$ meeting $C_s$. The image of $V$
in $S$ is open because $C/S$ is smooth. Replace $S$ by an affine open
neighborhood of $s$ in that image, and replace $V$ by its inverse
image. The latter is a quasi-affine scheme. Since $C$ is an $E$-torsor
and $V$ meets every fibre, the image of $E\times_S V\to C$ is all of
$C$. Apply \cite[\S6.6, Theorem~2(c), pp.~164--166]{BCEBLR1990}:
every finite set of points of $C$ lies in one affine open. In
particular, each finite $E[n]$-orbit does. The finite locally free
quotient theorem, \cite[Proposition~66.14.1, tag~07S6]{BCEStacksFiniteQuotient},
therefore represents $C/E[n]$ by a scheme. The degree-zero piece is
$E$ itself; hence every piece and their disjoint union are schemes.
Smoothness and separatedness follow etale locally from the projective
curve theorem, or from the torsor description.

Finally, the arbitrary-base assertion follows by approximation,
without assuming that schematicity descends along arbitrary fpqc
covers. Locally take $S$ affine. The proper smooth finitely presented
scheme $C/S$ descends to a proper smooth finitely presented scheme
$C_0/S_0$ with $S_0$ noetherian. Geometric connectedness of proper
smooth fibres is an open condition and the arithmetic genus is
locally constant. Enlarging the model and restricting to the
corresponding open locus therefore arranges geometrically connected
genus-one fibres. The noetherian argument makes
$\operatorname{Pic}_{C_0/S_0}$ a scheme; its base change to $S$ represents
$\operatorname{Pic}_{C/S}$ by the definition of the relative fppf
Picard functor. This completes the proof.

\paragraph{Provenance.}
The smooth-curve argument above was already given by K\k{e}stutis
\v{C}esnavi\v{c}ius in his answer of 28 April 2015 to Poonen's question
\cite{BCECesnavicius2015}. We independently checked the cited BLR and
finite-quotient statements before locating that answer. No novelty is
claimed for Theorem E4.Curve.

\paragraph{Why the genus-one torsor counterexample route fails.}
An arbitrary torsor under an abelian scheme can be a nonschematic
algebraic space. Such a torsor alone does not give the requested
counterexample, whose original family must be a scheme. If a genus-one
scheme family $C$ is used, its multiples in torsor cohomology are
exactly (E4.Curve.2), and the affine-neighborhood theorem above makes
all of them schemes. Thus taking a scheme genus-one torsor and hoping
that $\operatorname{Pic}^n$ becomes a nonschematic multiple cannot work.

\subsection{The two marginal notes on 57:310: mathematics and reading}

The multiplication margin is diplomatically transcribed as
``La multiplication par un multiple de $p$ est [surjective?]
[two illegible words] [ouverte?] (?) [comp. irr\'eductibles?]''.
The original image confirms ``Non, exemple'' and the branch diagram,
but does not justify removing the marked uncertainties. The mathematical
universal-openness question is nevertheless answered negatively by
Theorem E4.H, including for a smooth projective family. Its status is
therefore separate from the unresolved transcription of the qualifier.

The bracketed finite-type criterion has the alternatives ``au dessus''
and ``en dehors'' of an open set, with an uncertain dense/constructible
qualification. In the proper-flat-finitely-presented,
cohomologically-flat-in-degree-zero Picard setting over a noetherian
base, finite presentation of $\operatorname{Pic}^{\tau}$ already proves
that every such base change, open restriction, and closed complement has
finitely many irreducible components. Thus either indicated
finite-component condition holds in this Picard setting. This does not
prove a proposed general criterion for arbitrary group schemes.

\paragraph{An explicit warning about the abstract generalization.}
Over $S=\mathbf A^1_{\mathbf Q}$, enumerate the distinct points
$s_i=(t-i)$, $i\geq1$. Glue a copy $S_0$ of $S$ and copies $S_i$,
$i\geq1$, by identifying $S_i\setminus\{s_i\}$ with the corresponding
open of $S_0$. Call the resulting scheme $G$. Its fibre over $s_i$ has
two points, and every other fibre has one. It is an \'etale commutative
group scheme with the two-point fibre given the group law
$\mathbf Z/2$: on two different charts $S_i,S_j$, their sum is the
section that uses the extra origin at $s_i$ and $s_j$, obtained by gluing
on $S\setminus\{s_j\}$ and $S\setminus\{s_i\}$. On the same chart
the sum is the zero section. These formulas agree on chart intersections;
the group laws hold fibrewise and hence hold as morphisms into the
\'etale space.

The scheme $G$ is irreducible, since all its charts have the same dense
generic point, but is not quasi-compact: any finite collection of these
charts misses infinitely many extra origins. For every closed irreducible
$S'\subset S$, its restriction over a nonempty open of $S'$ has finitely
many irreducible components. Consequently this generic-open
component-count condition does not characterize finite type for arbitrary
locally finite-type group schemes. This is deliberately not presented as
a Picard realization. It only rules out silently replacing the uncertain
margin by that unrestricted abstract assertion.


\subsection*{E4.Pinching: an excluded route to nonschematic Picard spaces}

\paragraph{Proposition.}
Suppose $X'\to X$ is a finite schematically dominant surjection
of proper $S$-schemes, with its conductor square, and both conductor schemes are finite over
$S$. If $X'\to S$ is locally projective, so is $X\to S$.
Consequently, when $X/S$ is also flat, finitely presented and has
geometrically integral fibres, its full Picard functor is a scheme.

\paragraph{Proof.}
Work near any point $s\in S$ and choose a relatively ample invertible
sheaf $L'$ on $X'$. Its restriction to the inverse-image conductor is
trivial after shrinking $S$ around $s$: over the local scheme
$\operatorname{Spec}\mathcal O_{S,s}$, that conductor is finite and
hence semilocal, so every invertible module is free. The trivialization
spreads to a neighborhood by finite presentation. The conductor
fibre-product square therefore glues $L'$ to the trivial bundle on
the conductor of $X$, producing an invertible sheaf $L$ on $X$.
Since its pullback under the finite surjection is ample, $L$ is
relatively ample. This is the ampleness argument in
\cite[Proposition~2.4(i), p.~6]{BCEBrion2015}; for this argument only
finiteness of the conductors is needed. The final assertion is the
projective Picard theorem.

Thus finite-conductor pinching of a locally projective simultaneous
normalization, in particular of $\mathbf P^1_S$, cannot realize a
nonschematic Picard space. This does not assert that an arbitrary
pinching is flat or that normalization commutes with base change.
The general integral-curve schematicity question remains when no
such normalization over the original base is available. For normal
families of higher dimension, the arbitrary-base question likewise
remains unresolved here.



\subsection{H8: the exceptional divisor lattice, with its dimension-one exception}

On 57:7 Grothendieck asks that $D\to\mathrm{NS}(X_0)$ be injective; on 57:8
he writes ``I expect that this is always the case'' and proposes intersection
with curves on the exceptional components. These leaves are in English, not
French, so a French original must not be invented.

Let $A$ be a complete noetherian local ring, let
$f:X\to\operatorname{Spec}A$ be proper, regular, and the closure of the
regular punctured spectrum, and let $D$ denote the free group of prime
divisors supported on the closed fibre. If the irreducible components of
$X$ have dimension at least two, then
\[
 D\longrightarrow\mathrm{NS}(X_0)
\]
is injective, even after replacing algebraic equivalence by numerical
equivalence. Indeed first suppose $A$ normal and integral. A divisor $B$
whose restriction to $X_0$ is algebraically trivial has degree zero on every
proper curve contracted by $f$: all those curves lie in $X_0$. Both $B$ and
$-B$ are therefore $f$-nef, and $f_*B=0$ because the closed point has
codimension at least two. The negativity lemma, applied to $B$ and to $-B$,
gives $B\geq0$ and $B\leq0$, whence $B=0$.
The exact source opened is Koll\'ar, \emph{Families of varieties of general
type}, Lemma~11.60, p.~436, citing Koll\'ar--Mori~3.39. Section~11.4 explicitly
treats arbitrary schemes; the characteristic-zero restriction stated at the
start of Chapter~11 is for Sections~11.1--11.3.

For general $A$, replace it by the normalization of $A_{\mathrm{red}}$.
It is finite since $A$ is excellent. The normal scheme $X$ factors through
it, and the punctured spectrum is unchanged since it was regular. The finite
normalization is a disjoint union of complete normal local schemes. Apply the
preceding argument on each component. All contracted curves remain in the
original closed fibre, so numerical triviality is preserved.

There is an indispensable qualification in dimension one. For
$A=k[[t]]$ and $X=\operatorname{Spec}A$, the divisor group is $\mathbf Z$,
whereas $\mathrm{NS}(X_0)=\mathrm{NS}(\operatorname{Spec}k)=0$. Thus the
unqualified assertion for every dimension is false. More generally, in the
setting just stated, the whole kernel is freely generated by the closed points
of the one-dimensional components of $X$: after normalization those components
are spectra of DVRs and their proper birational maps are isomorphisms; the
dimension-at-least-two components have zero kernel by negativity. Components
of dimension zero contribute no divisors. This gives the complete correction.

\subsection{H10: smoothness in characteristic zero over nonreduced bases}

The supplementary question on 57:64 reads:
\begin{quote}
``Si $G$ est un sch\'em en groupes sur $S$, $S$ \`a corps r\'esiduels de car.
$0$, et si $\underline\omega_{G/S}$ [faisceau conormal \`a la section unit\'e]
est loc. libre, alors $G$ est simple sur $S$ le long de la section unit\'e''.
\end{quote}
With the locally-finite-presentation hypothesis used throughout the discussion,
this is exactly SGA~3, VI$_B$, Proposition~1.6 (original p.~327, p.~6 of the
opened retyped expos\'e). No flatness hypothesis is needed. The proof uses
$\Omega^1_{G/S}=\pi^*\omega_{G/S}$ and EGA IV$_4$, 16.12.2 and 17.12.5.

For $X/S$ smooth projective in characteristic zero, Deligne's
\emph{Th\'eor\`eme de Lefschetz et crit\`eres de d\'eg\'en\'erescence de suites
spectrales}, Theorem~5.5(i), p.~123, gives local freeness and arbitrary
base-change compatibility of every $R^qf_*\Omega^p_{X/S}$, with no reducedness
condition on $S$. The infinitesimal Picard tangent functor is
$M\mapsto R^1f_*(\mathcal O_X\otimes f^*M)$, which, by this base-change theorem,
equals $R^1f_*\mathcal O_X\otimes M$. Consequently the conormal module of Pic
at the identity is $(R^1f_*\mathcal O_X)^\vee$ and is locally free.
The preceding SGA~3 criterion makes Pic smooth along the identity. Its neutral
part is therefore an open smooth group scheme. Its fibres are abelian varieties
(smooth proper varieties have proper reduced neutral Picard schemes), and its
relative properness follows from the normal-fibre proper-neutral-subset theorem
once that part is open and represented. Thus $\operatorname{Pic}^0_{X/S}$ is an
abelian scheme over arbitrary, possibly nonreduced, $S$.
If the geometric fibres of $X/S$ are disconnected, make the finite \'etale
Stein-factor reduction and take the corresponding product/Weil restriction.
Kleiman's Remark~5.21, pp.~48--49, provides an independent presentation of the
same deduction using exponential lifting; this is a source comparison, not a
claim of an independent research verification.

\subsection{H11: affineness after shrinking for normal families over a curve}

The question at 57:74, repeated on 57:78, is in English in the manuscript:
``one could replace `quasi-affine' by `affine'.'' It specifically proposes
finite maps between projective normal families, and birational maps between
smooth projective families. The finite-normal possibility has an affirmative
answer after the permitted shrinking; the argument below also allows proper
surjections between projective normal families. It does not assert the same
result for every nonnormal proper family.

\paragraph{Theorem H11.N.}
Let $S$ be integral noetherian of dimension at most one, and let
$f:X'\to X$ be a surjective morphism of projective $S$-schemes whose geometric
generic fibres are normal. After replacing $S$ by a nonempty open, the map
$f^*:\operatorname{Pic}_{X/S}\to\operatorname{Pic}_{X'/S}$ is affine.
In particular the finite normal-family counterexample suggested on 57:78
cannot exist after further shrinking of the base.

\paragraph{Proof.}
Shrink first so both families are flat with geometrically normal fibres,
and so both Picard functors are separated group schemes locally of finite
presentation and $u=f^*$ is quasi-affine of finite presentation, as in
SGA~6 XII Theorem~1.1 and Corollary~1.2 (also Kleiman, Theorem~4.18.5).
Write $P=\operatorname{Pic}_{X/S}$, $P'=\operatorname{Pic}_{X'/S}$,
$H=P^\tau$, $H'=P'^\tau$. These are open-and-closed group subschemes of
finite presentation. Their generic fibres are proper, since the generic
geometric fibres of $X,X'$ are normal. Generic flatness and spreading out
properness for a finitely presented morphism allow one more shrinking so
that $H,H'$ are proper and flat over $S$.

By H6 the quotients $Q=P/H$, $Q'=P'/H'$ are separated unramified schemes
locally of finite presentation over $S$. The projections $P\to Q$ and
$P'\to Q'$ are fppf torsors under $H,H'$, respectively, and hence are proper,
flat, surjective, and of finite presentation. Pullback induces $v:Q\to Q'$.

First $v$ is a monomorphism. On a geometric fibre, a line bundle whose
pullback is numerically trivial is itself numerically trivial. In the finite
case, for every integral curve $C\subset X_s$ one takes a component
$C'\subset X'_s$ finite and surjective onto $C$ and uses
\[
 \deg_{C'}f^*L=[k(C'):k(C)]\deg_C L.
\]
For a proper surjection one can use Chow's lemma and successive hyperplane
sections over $C$ to obtain such a dominating proper curve, and the same
degree formula applies. Numerical triviality is equivalent to membership
in $\operatorname{Pic}^\tau$ for a proper scheme over an algebraically closed
field. Thus $v$ is injective on geometric points. If two $T$-points of $Q$
have the same image, their values agree on every geometric fibre of $T$.
Their equalizer is open, since $Q/S$ is unramified, and contains all points
of $T$; therefore the two $T$-points are equal.

Next $v$ is quasi-compact. For any quasi-compact open $V\subset Q'$, the
scheme $P'\times_{Q'}V$ is quasi-compact because $P'\to Q'$ is proper.
Its inverse image under the quasi-compact map $u$ is quasi-compact, and maps
surjectively to $v^{-1}(V)$. Thus $v^{-1}(V)$ is quasi-compact. The morphism
$v$ is consequently separated and quasi-finite. Every local ring of $Q'$
has dimension at most one, since $Q'/S$ is unramified and $\dim S\le1$.
Zariski's Main Theorem embeds $v^{-1}(V)$, for affine $V$, as a quasi-compact
open in a finite $V$-scheme. Such a finite scheme is affine and noetherian
of dimension at most one, and every quasi-compact open in it is affine
(Stacks, Lemma~33.38.2, tag~09N9). Thus $v$ is affine.

Finally consider
\[
 P\xrightarrow{\psi} P'\times_{Q'}Q\longrightarrow P'.
\]
Both $P\to Q$ and $P'\times_{Q'}Q\to Q$ are proper, and the latter is
separated, so $\psi$ is proper. It is quasi-affine: this follows by taking
the graph over $P'$ and using that $u$ is quasi-affine and the second map
is separated. A proper quasi-affine morphism is finite. The second map
is affine as a base change of $v$, so their composite $u$ is affine.

The proof keeps the finite-type hypothesis on $H$ visible: normality alone
does not make $H$ proper over the original unshrunk base, as Theorem E4.N
shows. The smooth birational case on 57:78 is covered by the same proof.
One must not claim that the full Picard functor is unchanged by a blowup:
exceptional divisor classes enlarge the N\'eron--Severi group. The familiar
birational invariance concerns the neutral Picard variety, and pullback on
the full Picard functor is injective when $f_*\mathcal O_{X'}=\mathcal O_X$.
The argument above does not address nonnormal fibres. The following theorem
settles that branch by a different quotient, using the kernel itself.

\subsection{H11: the full one-dimensional affineness statement}

\paragraph{Theorem H11.G.}
Let $S$ be an integral noetherian scheme of dimension at most one, and let
$f:X'\to X$ be a surjective morphism of proper $S$-schemes. After replacing
$S$ by a nonempty open subscheme, the pullback
\[
 f^*:\operatorname{Pic}_{X/S}\longrightarrow\operatorname{Pic}_{X'/S}
\]
is affine. No normality or projectivity assumption is necessary.
Thus the general one-dimensional question on 57:74 and 57:78 has an
 affirmative answer, extending Theorem H11.N above.

\paragraph{Proof.}
SGA~6 XII Theorem~1.1 and Corollary~1.2,
\cite[original p.~596]{BCESGA6XII}, allow a first shrinking so that
$P=\operatorname{Pic}_{X/S}$ and $P'=\operatorname{Pic}_{X'/S}$ are separated
group schemes locally of finite presentation, and $u=f^*:P\to P'$ is
quasi-affine of finite presentation. Its kernel $K$ is a quasi-affine group
scheme of finite presentation over $S$. The generic fibre $K_\eta$ is an
affine algebraic group, by SGA~3 VI$_B$~11.11. Affineness spreads out for a
finitely presented scheme over the generic point: take an affine model of
$K_\eta$ over a nonempty open and spread out the generic isomorphism and
its inverse. After further shrinking, $K$ is therefore affine. Generic
flatness permits a final shrinking making $K/S$ flat.

The free translation action of the flat finitely presented group $K$ on
$P$ has a quotient algebraic space $Q=P/K$ (Stacks tag~06PH). The quotient
map $p:P\to Q$ is an fppf $K$-torsor, hence affine and of finite presentation
by descent. The map $u$ factors through a monomorphism $v:Q\to P'$.
This map is locally of finite presentation, by fppf descent from $u$ along
$p$, and is quasi-compact: over a quasi-compact open $V\subset P'$, the
quasi-compact scheme $u^{-1}(V)$ maps surjectively to $v^{-1}(V)$.
Thus $v$ is a quasi-finite separated morphism of finite presentation.
Stacks tag~0418 shows that it is representable by schemes, so $Q$ is a
scheme as well.

For every point $s\in S$, the morphism $v_s$ is a quasi-compact
monomorphism of group schemes locally of finite type over the field
$\kappa(s)$. SGA~3 VI$_B$, Corollary~1.4.2, original pp.325--326, says
precisely that such a monomorphism is a closed immersion. In particular
$v_s$ is finite. Now apply SGA~3 XVI, Lemma~4.2, original pp.521--522
(p.~19 of the opened retyped expos\'e): over a locally noetherian base
of dimension at most one, a quasi-finite separated morphism between
schemes locally of finite type, whose morphism on every fibre is finite,
is affine. Therefore $v$ is affine and so is $u=v\circ p$.

For clarity, the quoted lemma imposes no dimension bound on $P'$ itself.
Its proof passes to complete local bases and then a DVR. Zariski's Main
Theorem splits the source over a complete local target into a finite part
and a part missing its closed point. Fibrewise finiteness forces the latter
part to lie in the generic fibre; that generic fibre is a principal affine
open over the DVR. Both parts are affine over the target. This is the
one-dimensional input that removes the earlier normality restriction.
The statement does not assert the corresponding affineness over an
arbitrary higher-dimensional base.


\section{The local fundamental-group conjecture}\label{sec:F}
\subsection{The actual published boundary}
\begin{proposition}[F1: checked known cases; Known]\label{f:known}
Under the assumptions of Part F, the maximal prime-to-residue-characteristic
quotient is topologically finitely presented in the prime-to-characteristic
category. In residue characteristic zero this is the whole group.
In all residue characteristics the maximal pro-nilpotent quotient is
topologically finitely generated. In equal characteristic $p>0$, strong
$F$-regularity implies that the whole group is finite of order prime to $p$.
\end{proposition}
\begin{proof}[Verification of hypotheses and locators]
Grothendieck--Murre, \S9.1 and Theorem 9.2, p.~122, assume a complete
noetherian local ring of dimension two with algebraically closed residue
field, a closed subset $E$, and connected $S=\Spec(A)\setminus E$.
Their conclusion concerns the maximal prime-to-$p$ quotient; the
presentation convention on pp.~122--123 uses a free pro-prime-to-$p$
group. They allow $p=0$ \cite{GrothendieckMurre1971}.
The original two-page preview containing the theorem and definition was
opened directly. No analytic algebraization assumption is needed here.

Suzuki's Theorem 1.2, p.~2, proves the pro-nilpotent statement in equal
and mixed characteristic \cite{Suzuki2023}. Its hypotheses, printed as
Conjecture 1.1 on p.~1, require every irreducible component to have
dimension at least two and the punctured spectrum to be connected.
A normal two-dimensional local domain satisfies both conditions.
This theorem does not assert finite generation of the whole group.

Theorem A=Theorem 5.1 of \cite[p.~994]{CRST2018} assumes normality,
$F$-finiteness, strong $F$-regularity, strict henselianity, and dimension
at least two. It bounds the group order by $1/s(A)$ and makes it prime
to $p$. In the complete equicharacteristic situation here, strict
henselianity follows from the algebraically closed residue field.
A Cohen presentation makes $A$ a quotient of $k[[x_1,\ldots,x_e]]$;
Frobenius is finite on this power-series ring because $k$ is perfect,
so $A$ is $F$-finite. Strong $F$-regularity remains an additional
hypothesis and has not been deduced from normality.
\end{proof}

\begin{proposition}[F1b: all log terminal surfaces, including small and mixed characteristic; Known]
If the ring in Part F is log terminal, then its entire punctured
\'etale fundamental group is finite, in every equal or mixed
characteristic.
\end{proposition}
\begin{proof}[Verification of the exact published theorem]
This is Carvajal-Rojas--Yasuda, Theorem~B, p.~4 of the opened
27 March 2025 author version, equivalently Corollary~6.4,
pp.~29--30 \cite{CarvajalYasuda2025}. Its statement assumes exactly
a two-dimensional complete log terminal local ring with
algebraically closed residue field. The mixed-characteristic
setup on p.~28 specifies normality, a torsion canonical class
($\QQ$-Gorenstein), and positive log discrepancies on a minimal
resolution. The theorem includes characteristics $2$ and $3$ and
does not impose strong $F$-regularity. It asserts finiteness, not
that the order is prime to the residue characteristic.
\end{proof}

\begin{proposition}[F1c: the mixed-characteristic perfectoid-signature theorem; Known]
Let $A$ satisfy Part F in mixed characteristic $(0,p)$. If there is an
effective divisor $\Delta$ such that $K_A+\Delta$ is $\QQ$-Cartier
and $(A,\Delta)$ is BCM-regular, then the full punctured fundamental
group is finite, of order prime to $p$, and its order is at most the
inverse of the perfectoid signature of $A$.
\end{proposition}
\begin{proof}[Verification of the current primary source]
Cai, Lee, Ma, Schwede, and Tucker, Theorem~D on pp.~5--6 and
Theorem~7.0.2 on pp.~73--74 of the opened 3 July 2025 version,
give this assertion for any closed subset of codimension at least
two in a complete normal local domain of mixed characteristic with
algebraically closed residue field \cite{CaiEtAl2026}.
The journal article appeared online on 16 December 2025 and in
\emph{Inventiones mathematicae} 244 (2026), pp.~357--454.
Normality is not the BCM-regular-pair hypothesis. The latter is
used to make the perfectoid signature positive; the displayed
inverse-signature bound alone is vacuous when it is zero.
\end{proof}

The primary-source search through 20 November 2026 found no later
resolution for arbitrary normal complete surfaces. A negative search is not a
nonexistence proof. The full positive-residue-characteristic case
outside these additional regularity hypotheses remains
undecided in this work.

\begin{proposition}[F2: missing dimension qualification; Refuted wording]
\label{f:dimension}
The prompt's unqualified consequence for \emph{every} complete local
ring with algebraically closed residue field is false in dimension one.
\end{proposition}
\begin{proof}
Take $A=k[[t]]$ in characteristic $p>0$, with $k$ algebraically closed.
Artin--Schreier theory gives
\[
 H^1_{\mathrm{et}}(\Spec k((t)),\ZZ/p)
       =k((t))/(F-1)k((t)).
\]
The classes $t^{-n}$, $n>0$ and $p\nmid n$, are linearly independent
over $\mathbf F_p$. A nonzero finite combination has a largest pole
order prime to $p$, whereas a Laurent series $g$ with pole order $m$
has $g^p-g$ of pole order $pm$; an integral $g$ has no negative pole.
Thus this cohomology has infinite dimension. A group with $d$
topological generators has at most $p^d$ homomorphisms to $\ZZ/p$,
so its fundamental group is not finitely generated.
This does not refute the normal surface case. It explains why the
dimension-at-least-two qualification in SGA~2 XIII, Conjecture 3.1,
must remain in the consequence quoted from 58:19.
\end{proof}

\subsection{Why even prosolvable d\'evissage is insufficient}
\begin{proposition}[F3: an explicit metabelian obstruction; Proved]
\label{f:metabelian}
There is a metabelian profinite group $G$ whose maximal pro-nilpotent
quotient is $\ZZ_3$, such that every open subgroup has finitely generated
abelianization, but $G$ is not topologically finitely generated.
\end{proposition}
\begin{proof}
Let $f_n$ be the multiplicative order of $2$ modulo $3^n$ and
$L_n=\mathbf F_{2^{f_n}}$. Choose $\zeta_n\in L_n^\times$ of order
$3^n$. Put
\[
 V_n=L_n^n,\qquad V=\prod_{n\geq1}V_n,\qquad
 \Gamma=\ZZ_3,
\]
and let $\gamma\in\Gamma$ act on $V_n$ by multiplication by
$\zeta_n^{\gamma\bmod 3^n}$. This continuous action defines the
metabelian profinite group $G=V\rtimes\Gamma$.

The commutator operator with $1\in\Gamma$ is multiplication by
$\zeta_n-1$ on each coordinate and hence maps $V$ onto itself.
In any finite nilpotent quotient, the image of $V$ is a $2$-group
and that of $\Gamma$ a $3$-group. They commute, so the quotient kills
$[\Gamma,V]=V$. Conversely finite quotients of $\Gamma$ are
nilpotent. This proves the first assertion.

If $H$ is open, it contains $\prod_{n>M}V_n$ and $3^N\Gamma$ for
some $M,N$. Put $r=\max(M,N)$ and $T=\prod_{n>r}V_n$.
The commutator with $3^N\in\Gamma$ is invertible on every factor
of $T$, with continuous coordinatewise inverse. Thus
$T\subset[H,H]$. The quotient $H/T$ is a closed subgroup of
$(\prod_{n\leq r}V_n)\rtimes\Gamma$, with finite kernel and
procyclic image under projection to $\Gamma$. Finitely many
generators of the kernel and one lift of a generator of the image
generate it topologically. Its quotient $H^{\mathrm{ab}}$ is
therefore finitely generated. In particular
$H^1_{\mathrm{cont}}(H,\mathbf F_\ell)$ is finite for every prime
$\ell$.

Suppose $d$ elements generated $G$ densely. Their vector coordinates
in $V_n$ span an $L_n$-subspace $W_n$ of dimension at most $d$.
All words in their images in $V_n\rtimes(\Gamma/3^n\Gamma)$ lie
in $W_n\rtimes(\Gamma/3^n\Gamma)$, since the action is scalar.
For $n>d$ this is proper, contradicting generation of that finite
quotient. Thus $G$ is not finitely generated.
\end{proof}

This is a counterexample to an abstract inference, with no claimed
realization as a local fundamental group. It strengthens the
$\prod A_5$ warning in the supplied reading: the gap already occurs
for metabelian groups. The finite checks verify the relevant field
orders through $n=8$ and the perfection of $A_5$, while the preceding
proof handles the infinite product.

\paragraph{Approaches and exact gates.}
Rank-one Artin--Schreier/Frobenius estimates and their application to
each finite cover establish qualitative constant-coefficient
finiteness. Statement~\ref{f:metabelian} proves that these estimates
alone cannot give finite generation, even after restricting to
prosolvable quotients. Unbounded multiplicities of nontrivial finite
coefficient modules are an explicit remaining obstruction.
The resolution/intersection argument in Suzuki would need a uniform
extension to nonconstant local systems with controlled ramification;
no such estimate is proved here. Conversely wild covers of an affine
curve or cone do not automatically survive completion: regular
Artin--Schreier functions are killed by henselianity and the
algebraically closed residue field. No normal complete surface
counterexample has been constructed. The smallest undecided scope
recorded is a normal complete surface in positive residue
characteristic outside the log terminal and strongly $F$-regular
cases, with its full wild fundamental group.

\paragraph{The newer birational approach and its limitation.}
Carvajal-Rojas--Yasuda prove strict decrease of stringy motives
under ramified quasi-\'etale Galois covers and a descending chain
condition for the relevant log terminal surface invariants
\cite[Theorem 6.3 and Corollary 6.4, pp.~29--30]{CarvajalYasuda2025}.
This controls nonabelian wild covers in their stated class.
Positive discrepancies and the resulting convergence are
part of the construction. For an arbitrary normal non-log-terminal
surface, reusing that invariant and its descending chain condition
would be an additional unproved step; the theorem cannot be
applied after dropping log terminality.

\begin{corollary}[F3b: an explicit failure of uniform twisted-cohomology bounds; Proved]
For the group $G$ of Proposition~\ref{f:metabelian}, let $M_n=L_n$
with $V$ acting trivially and $\Gamma$ acting by $\zeta_n$.
Then
\[
 H^1_{\mathrm{cont}}(G,M_n)\simeq L_n^n,
 \qquad
 \frac{\dim_{\mathbf F_2}H^1(G,M_n)}{\dim_{\mathbf F_2}M_n}=n.
\]
By contrast, any profinite group generated by $d$ elements satisfies
$\dim_{\mathbf F_\ell}H^1(G,M)\leq d\dim_{\mathbf F_\ell}M$
for every finite continuous $\mathbf F_\ell[G]$-module $M$.
\end{corollary}
\begin{proof}
A continuous crossed homomorphism is determined by its values on a
topological generating set: the cocycle identity determines it on
words and continuity determines it on their closure. This proves
the stated necessary bound.

For our group, the restriction of a cocycle to $\Gamma$ is a
coboundary, since $\zeta_n-1$ is invertible on $L_n$: subtract
the coboundary whose value at $1\in\ZZ_3$ is the prescribed
value of the cocycle there, then use the density of the integers
in $\ZZ_3$. A normalized cocycle vanishing on $\Gamma$ is exactly
a continuous $\Gamma$-equivariant additive homomorphism $V\to L_n$.
Such a homomorphism factors through finitely many factors of $V$.
For $m<n$, the element $\gamma^{3^m}-1$ is zero on $L_m$
and invertible on $L_n$, forcing an equivariant homomorphism to
vanish. For $m>n$, use $\gamma^{3^n}-1$ in the reverse way. On $V_n=L_n^n$,
equivariance means $L_n$-linearity, because
$\mathbf F_2[\zeta_n]=L_n$. Thus the normalized cocycles form
$\Hom_{L_n}(L_n^n,L_n)=L_n^n$. No nonzero normalized cocycle
is a coboundary, since a coboundary vanishing on $\Gamma$ would
come from an element fixed by $\zeta_n$, necessarily zero.
This proves the formula and the unbounded ratio.
\end{proof}

\begin{proposition}[F4: a quantitative sufficient criterion for prosolvable quotients; Proved]
Let $G$ be a prosolvable profinite group. Suppose a real constant $C\geq0$
satisfies
\[
 \dim_{\mathbf F_\ell}H^1_{\mathrm{cont}}(G,M)
 \leq C\dim_{\mathbf F_\ell}M
\]
for every prime $\ell$ and every finite simple continuous
$\mathbf F_\ell[G]$-module $M$. Then $G$ has at most
$\lceil C\rceil+1$ topological generators. Conversely, $d$ generators
imply this bound with $C=d$.
\end{proposition}
\begin{proof}
Only sufficiency is new here. Put $D=\lceil C\rceil+1$.
For every finite quotient $F$ of $G$, inflation injects
$H^1(F,M)$ into $H^1(G,M)$ for an $F$-module $M$.
We show by induction on $|F|$ that every finite solvable quotient
$F$ is generated by $D$ elements, allowing identities in the tuple.
Choose a minimal nontrivial normal subgroup $M$ of $F$.
Solvability and minimality imply that $M$ is an elementary abelian
$\ell$-group and a simple module for $H=F/M$ under conjugation:
its derived subgroup, primary parts, and subgroup $\ell M$ are
characteristic, so minimality successively forces these assertions.
By induction choose a generating $D$-tuple of $H$.

Any lift of this tuple generates a subgroup $J$ with $JM=F$.
The intersection $J\cap M$ is an $H$-submodule, hence is either
$0$ or $M$. In the latter case $J=F$; in the former case $J$ is
a complement to $M$. Thus a nonsplit extension has a generating
lift immediately. In the split case identify $F=M\rtimes H$.
The complements are precisely the graphs of the crossed
homomorphisms $H\to M$. For the chosen tuple each complement
contains exactly one lifting tuple, so the number of nongenerating
lifts is $|Z^1(H,M)|$.

Restriction to $M$ and the splitting give, directly on cocycles,
\[
 H^1(F,M)\simeq H^1(H,M)\oplus\operatorname{End}_H(M).
\]
Indeed a cocycle vanishing on $H$ is exactly an $H$-equivariant
additive map $M\to M$, and a coboundary restricts to zero on $M$.
The endomorphism group contains its nonzero identity, so has
cardinality at least $\ell$. Consequently
\[
 |Z^1(H,M)|
 =|H^1(H,M)|\frac{|M|}{|M^H|}
 \leq \frac{|M|^{C+1}}{\ell}
 <|M|^D.
\]
There are $|M|^D$ lifting tuples, so a generating one exists.
This proves the finite induction.

For each open normal subgroup $N$ of $G$, the tuples in $G^D$
generating $G/N$ form a nonempty closed subset. Any finite
intersection contains the subset corresponding to the
intersection of the relevant normal subgroups, which is nonempty
by the finite result. Compactness therefore supplies one tuple
generating every finite quotient; its closed generated subgroup
is $G$. Necessity was proved in F3b by evaluating cocycles on a
generating tuple.
\end{proof}

\paragraph{The resulting geometric target.}
For the maximal prosolvable quotient of the local fundamental group,
it would suffice to prove a constant $C_A$, independent of the
finite simple local system $\mathcal M$ and its rank, with
\[
 \dim_{\mathbf F_\ell}H^1_{\mathrm{et}}
       (\Spec(A)\setminus\{\mathfrak m\},\mathcal M)
 \leq C_A\operatorname{rank}(\mathcal M).
\]
One only needs local systems whose monodromy is solvable; inflation
then bounds the cohomology of the prosolvable quotient. The argument
does not provide this uniform estimate. In particular, separate
finiteness on each finite cover has no uniform constant, as F3
demonstrates. Nonabelian simple chief factors of the full fundamental
group require an additional argument even after this target is met.

\subsection{A uniform bound for abelian prime-to-$p$ monodromy}

\begin{proposition}[F5: uniform twisted cohomology in a geometric subclass; Proved]
Let $A$ be a complete normal two-dimensional local domain of
characteristic $p>0$ with algebraically closed residue field $k$,
and put $U=\Spec A\setminus\{\mathfrak m\}$. There is a constant
$C_A$ such that, for every finite locally constant
$\mathbf F_p$-sheaf $M$ on $U$ whose monodromy group is abelian
of order prime to $p$,
\[
 \dim_{\mathbf F_p}H^1_{\mathrm{et}}(U,M)
 \leq C_A\operatorname{rank}_{\mathbf F_p}M.
\]
Consequently the group in F3 cannot be a quotient of such a local
fundamental group when $p=2$.
\end{proposition}

\begin{proof}
We first produce one finite cover, independent of $M$, on whose
fixed resolution all these coefficient systems extend without
ramification. We then make Suzuki's equal-characteristic
nearby-cycle argument uniform for the resulting torsion line bundles.
The argument uses the explicit filtration and Frobenius-cokernel
calculation of \cite[Propositions 2.1 and 3.1--3.4, pp.~4--7]{Suzuki2023};
the uniformity and the preliminary fixed-cover reduction are proved below.

\paragraph{A finite bound for abelian tame inertia.}
Choose a coefficient field $k\hookrightarrow A$ by the equal-characteristic
Cohen structure theorem. Since $k$ is perfect, a power-series presentation
makes $A$, and every scheme of finite type over it used below, $F$-finite.
Choose a resolution $f:X\to\Spec A$ with reduced exceptional divisor
$D=\bigcup_{i=1}^s D_i$ a simple normal crossing divisor, each $D_i$
smooth and $D$ nonempty; if $A$ is regular, first blow up its closed point.
Its intersection matrix $B=(D_i\cdot D_j)$ is negative definite,
so $b=|\det B|$ is a positive integer. Write
$Q=\pi_1(U)^{\mathrm{ab},(p')}$ and let $I_i\subset Q$ be the image
of tame inertia at $D_i$. These are procyclic groups.

For every character $\chi:Q\to\mu_n$ with $(n,p)=1$, Kummer theory
represents the torsor by a line bundle $L$ on $U$ with a trivialization
of $L^{\otimes n}$. Extend $L$ to a line bundle $\widetilde L$ on the
regular scheme $X$: extend a representing divisor and use regularity.
The trivialization identifies
\[
 \widetilde L^{\otimes n}\simeq
 \mathcal O_X\Bigl(\sum_i a_iD_i\Bigr)
\]
for integers $a_i$. Its inertia character at $D_i$ is $a_i\bmod n$,
with the usual common choice of roots of unity. Restriction to $D_j$
gives
\[
 \sum_i a_i(D_i\cdot D_j)
       =n\deg(\widetilde L|_{D_j}).
\]
The adjugate matrix therefore gives $ba_i=0\bmod n$ for every $i$.
Characters separate points of the abelian pro-prime-to-$p$ group $Q$,
so $bI_i=0$. The closed subgroup $I=\sum_i I_i$ is finite.

The same subgroup contains the image of tame inertia for every
exceptional divisorial valuation over $\Spec A$. Indeed, at the center
of that valuation on $X$, trivialize $\widetilde L$. The Kummer equation
has the form
\[
 z^n=u\prod_{i\in S}t_i^{a_i},
\]
where $u$ is a unit and the $t_i$ are the local equations of the
components through the center. Its valuation is
$\sum_{i\in S}v(t_i)a_i$. The unit contributes only unramified
monodromy. In particular, if $\chi$ is trivial on $I$, then all $a_i$ are zero
modulo $n$, so this valuation inertia also has trivial character.
Characters of $Q/I$ separate points; this places its image in $I$
without making compatible choices of inertia generators.
There is no wild inertia in an abelian prime-to-$p$ quotient.

Choose an open subgroup $Q'\subset Q$ with $Q'\cap I=0$; a finite
quotient separating the finitely many nonzero elements of $I$ does
this. Let $U'\to U$ be the connected finite abelian prime-to-$p$
cover defined by $Q'$. Normalize $\Spec A$ in its function field to
obtain $\Spec A'$; excellence makes this normalization finite, and
completeness and the algebraically closed residue field make $A'$
a complete normal local domain with residue field $k$.
Choose an SNC resolution $X'\to\Spec A'$ with reduced exceptional
curve $D'$. For every $M$ in the proposition, its pullback to $U'$
has trivial inertia at every exceptional divisor of $X'$: its image
in $Q$ lies both in $I$ and in $Q'$. Here a divisorial valuation of
the finite field extension restricts to an exceptional divisorial
valuation of the base field, because residue transcendence degree one
is preserved under finite extension. Normalize $X'$ in its finite
monodromy cover. It is unramified in codimension one, hence is
\'etale everywhere by purity of the branch locus on the regular
scheme $X'$ \cite[Lemma 58.21.4, tag 0BMB]{StacksBranchPurity}.
Thus every such $M|_{U'}$ extends to a locally constant
sheaf on the same fixed $X'$.

Restriction
$H^1(U,M)\to H^1(U',M|_{U'})$ is injective, since restriction followed
by corestriction is multiplication by the degree, which is prime to
$p$. It suffices to obtain a uniform bound on $U'$. In the rest of
the proof rename $A',X',D',U'$ as $A,X,D,U$; these are now fixed.

\paragraph{The extended coefficient bundle.}
Let $M$ extend to $X$, set $r=\operatorname{rank}_{\mathbf F_p}M$, and
let $V=M\otimes_{\mathbf F_p}\mathcal O_X$. This is a vector bundle
with its natural Frobenius-semilinear map $F_V$, whose linearization
$F_X^*V\to V$ is an isomorphism. Because the monodromy is finite,
abelian, and of order prime to $p$, the representation over $k$
splits into characters. Hence
\[
 V\simeq\bigoplus_{a=1}^r L_a,
 \qquad L_a\text{ a finite-order line bundle on }X.
\]
In particular $\deg(L_a|_{D_i})=0$ for every $a,i$.
The summands need not be individually Frobenius-stable.

Choose a positive exceptional cycle $Z=\sum_i c_iD_i$ such that
$J=\mathcal O_X(-Z)$ has positive degree on every $D_i$, as in
\cite[p.~6]{Suzuki2023}. All choices below depend on this fixed
resolution and $J$, not on $M$.
With $j:U\hookrightarrow X$, $i:D\hookrightarrow X$, and
$\Psi=i^*j_*$ on \'etale sheaves, the Artin--Schreier sequence for
$V$ gives
\[
 0\longrightarrow V(*D)/V
 \xrightarrow{F_V-1} V(*D)/V
 \longrightarrow R^1\Psi M\longrightarrow0.
\]
Here we identify coherent quotients supported on $D$ with their
pullbacks by $i$. Exactness follows locally from the ordinary
Artin--Schreier sequence after an \'etale trivialization of $M$;
$j$ is affine and $F_V-1$ is surjective on the extended bundle in
the \'etale topology.

Let $U_mR^1\Psi M$ be the image of $J^{-pm}V/V$.
Valuations along the regular components of $D$, and the fact that
$F_X^*V\to V$ is an isomorphism, give the exact sequence
\[
 0\longrightarrow J^{-m}V/V
 \xrightarrow{F_V-1}J^{-pm}V/V
 \longrightarrow U_mR^1\Psi M\longrightarrow0.
 \tag{F5.1}
\]
In detail, a section with a pole of order $e>0$ has Frobenius pole
order $pe$, so its leading pole cannot cancel with its own order-$e$
term; regularity tests membership in these divisor lattices at
codimension-one points. This also proves the claimed kernel and
preimage assertions in (F5.1).

\paragraph{A pole cutoff independent of the coefficients.}
For $m\geq2$, the same quotient calculation as in Suzuki's
Proposition~3.1, using the projection formula for Frobenius and
$F_X^*V\simeq V$, gives
\[
 \operatorname{gr}_m R^1\Psi M
 \simeq B_Z\otimes_{\mathcal O_Z}V|_Z
                  \otimes_{\mathcal O_Z}J^{-m}|_Z,
 \qquad
 B_Z=(F_*\mathcal O_{pZ})/\mathcal O_Z.
 \tag{F5.2}
\]
The sheaf $B_Z$ is locally free on $Z$: the finite flat Frobenius
on the regular scheme $X$ makes $(F_*\mathcal O_X)/\mathcal O_X$
locally free, and the displayed sheaf is its base change to $Z$.

Fix a filtration of $\mathcal O_Z$ with successive terms line bundles
$P_\nu$ on components $D_{i(\nu)}$, obtained by removing one
component at a time from the effective cycle $Z$. Tensoring by the
locally free sheaves in (F5.2) gives successive terms
\[
 (B_Z|_{D_{i(\nu)}}\otimes P_\nu)
   \otimes L_a|_{D_{i(\nu)}}
   \otimes J^{-m}|_{D_{i(\nu)}}.
\]
For each of the finitely many fixed vector bundles
$B_Z|_{D_i}\otimes P_\nu$, choose an injection into a finite sum
of copies of a line bundle $H_\nu$ of positive degree. Such an
injection is obtained by globally generating a sufficiently ample
twist of the dual bundle and then dualizing. Since
$\deg L_a|_{D_i}=0$ and $\deg J|_{D_i}>0$, all the resulting line
bundles have negative degree for $m\geq m_0$, with one $m_0$
independent of $a,M,r$. Therefore
\[
 H^0(D,\operatorname{gr}_mR^1\Psi M)=0
 \quad(m\geq m_0).
\]
Filtered colimits and quasi-compactness now give
\[
 H^0(D,R^1\Psi M)=H^0(D,U_{m_0}R^1\Psi M),
 \tag{F5.3}
\]
increasing $m_0$ by one if necessary.

\paragraph{The remaining pole range has dimension at most a constant times rank.}
Suzuki's negative-degree filtration, Proposition~2.1(1), applies to
$J^{-pm}/\mathcal O_X$ and $J^{-pm}/J^{-m}$. Tensoring every layer
by a finite-order $L_a$ keeps its degree negative. Thus both
quotients tensored by $V$ have no global sections. Taking cohomology
of (F5.1) injects $H^0(D,U_mR^1\Psi M)$ into the kernel of
\[
 F_V-1:H^1(D,J^{-m}V/V)\longrightarrow
                    H^1(D,J^{-pm}V/V),
\]
and the map denoted $1$ here is injective. If a Frobenius-linear
map and an injective linear map have common domain of $k$-dimension
$N$, their difference has an $\mathbf F_p$-kernel of dimension at
most $N$: project to a linear splitting of the latter map and use
the fixed-space bound for a semilinear endomorphism. For completeness,
$\mathbf F_p$-linearly independent fixed vectors are $k$-linearly
independent: normalize a shortest $k$-linear relation to have first
coefficient $1$, apply Frobenius, and subtract; minimality forces
all coefficients into $\mathbf F_p$, a contradiction. This is the
quantitative form of Suzuki's Lemma~3.3.

A fixed filtration of $J^{-m_0}/\mathcal O_X$ has line-bundle layers
$T_\nu$ on smooth components $D_{i(\nu)}$. Their degrees are
negative. On a smooth projective curve $C$ of genus $g$, Riemann--Roch
therefore gives, for every degree-zero $L$,
\[
 h^1(C,T_\nu\otimes L)=g-1-\deg T_\nu.
\]
The finite filtration and the splitting of $V$ imply
\[
 \dim_k H^1(D,J^{-m_0}V/V)\leq C_1r
\]
for a constant $C_1$ fixed by $X,J,m_0$. Combining this with
(F5.3) gives
\[
 \dim_{\mathbf F_p}H^0(D,R^1\Psi M)\leq C_1r.
\]

Finally, the five-term nearby-cycle sequence and proper base change
bound $H^1(U,M)$ by $H^1(D,M)$ and $H^0(D,R^1\Psi M)$.
The Artin--Schreier sequence on the proper curve $D$ identifies
$H^1(D,M)$ with the Frobenius-fixed subgroup of $H^1(D,V|_D)$: on the finite-dimensional space of
global sections, $F_V-1$ is surjective over the algebraically closed
field. The normalization sequence of the nodal curve $D$ bounds
\[
 h^1(D,L_a|_D)\leq\sum_i g(D_i)+\#\{\text{nodes of }D\},
\]
since each restriction to a component has degree zero. Consequently
$\dim_{\mathbf F_p}H^1(D,M)\leq C_2r$. This yields the desired
bound with $C_A=C_1+C_2$ for the fixed cover.

For the last assertion, suppose $\pi_1(U)$ surjected onto the group
of F3 in characteristic $2$. Its modules $M_n$ have cyclic odd-order
monodromy. Inflation and F3b would give
\[
 n=\frac{\dim H^1(G,M_n)}{\dim M_n}
 \leq\frac{\dim H^1(U,M_n)}{\dim M_n}\leq C_A
\]
for every $n$, a contradiction.
\end{proof}

\paragraph{Exact scope of this advance.}
F5 bounds coefficients with abelian prime-to-$p$ monodromy and
excludes the specific metabelian obstruction F3 as a geometric
quotient. It does not bound all simple modules of a prosolvable
fundamental group: their finite monodromy can be nonabelian and
can have order divisible by $p$. It therefore does not yet prove
finite generation of the maximal prosolvable quotient, much less
the full local fundamental group. The fixed-cover argument depends
on the finite abelian inertia subgroup and does not extend formally
to those other monodromy groups.

\begin{corollary}[F6: nonabelian pro-$p$ kernels over abelian tame quotients; Proved]
Under the equal-characteristic hypotheses of F5, every quotient $G$
of $\pi_1(U)$ which admits an exact sequence
\[
 1\longrightarrow P\longrightarrow G\longrightarrow\Gamma
 \longrightarrow1,
 \qquad P\text{ pro-}p,\quad
 \Gamma\text{ abelian pro-prime-to-}p,
\]
is topologically finitely generated. In fact one bound on the number
of generators works for all such quotients of the fixed $\pi_1(U)$.
The kernel $P$ is allowed to be nonabelian. The same finite-generation
conclusion holds when $\Gamma$ is a virtually abelian pro-prime-to-$p$
group, without asserting a uniform bound independent of its abelian
subgroup's index.
\end{corollary}
\begin{proof}
Let $d_0$ generate $\pi_1(U)^{\mathrm{ab},(p')}$, which is finitely
generated by F1. Then every $\Gamma$ in the first assertion has at
most $d_0$ generators. We bound $H^1(G,M)$ for every finite simple
$\mathbf F_\ell[G]$-module $M$.

For $\ell=p$, the normal pro-$p$ subgroup $P$ acts trivially on $M$.
Indeed its action factors through a finite $p$-group, whose every
nonzero characteristic-$p$ representation has a nonzero invariant
vector (induct on the group order using a central element of order
$p$ and the nonzero kernel of its nilpotent operator $z-1$).
The invariant subspace $M^P$ is $G$-stable, hence is all of the
simple module $M$. Thus $M$ has abelian prime-to-$p$ monodromy.
Inflation from $G$ to $\pi_1(U)$ is injective in degree one, and F5
therefore gives
\[
 \dim_{\mathbf F_p}H^1(G,M)\leq C_A\dim_{\mathbf F_p}M.
\]

For $\ell\ne p$, one has $H^1(P,M)=0$. A continuous cocycle factors
through a finite $p$-group together with its finite coefficient
action, and averaging by that group's invertible order makes it a
coboundary. Inflation--restriction now gives
\[
 H^1(G,M)\simeq H^1(\Gamma,M^P).
\]
Again $M^P$ is either $0$ or $M$ by simplicity and normality. In the
latter case evaluation on $d_0$ generators bounds the dimension of
this cohomology by $d_0\dim M$; in the former it is zero.

All finite quotients of $G$ are solvable, since they are extensions
of finite $p$-groups by finite abelian groups. Apply F4 with
$C=\max(C_A,d_0)$ to obtain at most $\lceil C\rceil+1$ generators,
uniformly in $G$.

If $\Gamma$ has an open abelian subgroup $\Gamma'$, its inverse
image $G'$ has finite index in $G$. The corresponding finite \'etale
cover $U'\to U$ has degree prime to $p$ and is the punctured spectrum
of a complete normal local surface $A'$ with the same algebraically
closed residue field. The first assertion applied to $A'$ shows
that $G'$ is finitely generated. Adding finitely many coset
representatives generates $G$.
\end{proof}

The class in F6 does not contain arbitrary solvable quotients:
alternating wild and tame normal factors can leave no normal
pro-$p$ subgroup with abelian tame quotient. Thus the full
prosolvable and full profinite problems remain separate gates.

\begin{proposition}[F7: arbitrary coefficients extending over a fixed resolution; Proved]
Let $A,U$ satisfy F5, and fix an SNC resolution $X\to\Spec A$,
with nonempty reduced exceptional curve $D$. There is a constant
$C_X$ such that every finite locally constant $\mathbf F_p$-sheaf
$M$ on $X$ satisfies
\[
 \dim_{\mathbf F_p}H^1(U,M|_U)
       \leq C_X\operatorname{rank}_{\mathbf F_p}M.
\]
The monodromy of $M$ may be nonabelian and may have order divisible
by $p$.
\end{proposition}
\begin{proof}
We give all changes needed to the uniform part of F5, retaining its
notation $V=M\otimes_{\mathbf F_p}\mathcal O_X$, $J$, $Z$, $B_Z$,
and the nearby-cycle filtration. The Artin--Schreier and Frobenius
identities (F5.1)--(F5.2) remain valid: they only use that $M$ is
locally constant on $X$, equivalently that $V$ has its unit
Frobenius structure. We must replace the decomposition into torsion
line bundles by rank-uniform estimates for $V$ itself.

For each smooth exceptional component $C=D_i$, the bundle $V|_C$
is trivialized by a finite \'etale cover of $C$. It is semistable of
degree zero. Indeed its determinant becomes trivial on that cover,
so has degree zero. The pullback of any subbundle becomes a
subbundle of a trivial bundle. Taking its determinant gives a
nonzero map to a finite sum of trivial line bundles, so its degree
is nonpositive. Degree multiplies by the positive covering degree,
so the original subbundle also has nonpositive degree. This proof
works equally for covers of degree divisible by $p$.

It follows that
\[
 H^0(C,V|_C\otimes T)=0\quad\text{if }\deg T<0,
 \qquad h^0(C,V|_C)\leq r,
 \quad r=\operatorname{rank}V.
\]
For the second inequality, evaluation of sections at any closed
point of $C$ is injective: a section vanishing there would give a
line subbundle of positive degree after saturation. Riemann--Roch
therefore gives
\[
 h^1(C,V|_C)\leq r g(C),\qquad
 h^1(C,V|_C\otimes T)=r(g(C)-1-\deg T)\quad(\deg T<0).
 \tag{F7.1}
\]

For the high-pole calculation in F5, replace a layer indexed by
$L_a$ by
\[
 V|_C\otimes(B_Z|_C\otimes P_\nu)\otimes J^{-m}|_C.
\]
The fixed vector bundle in parentheses embeds in a finite sum of
copies of the fixed line bundle $H_\nu$. Semistability and the first
vanishing above show that the displayed layer has no sections as
soon as $\deg H_\nu-m\deg J|_C<0$. Thus the same fixed $m_0$
works, independently of $V$ and $r$.

For the low-pole calculation, the negative-degree filtration of
$J^{-m}/\mathcal O_X$ and $J^{-pm}/J^{-m}$, tensored with $V$,
has no sections. Formula (F7.1) bounds its first cohomology by a
constant times $r$. The semilinear fixed-space lemma in F5 bounds
$H^0(D,R^1\Psi M)$ by the same type of constant. Finally the
normalization sequence for $D$ and (F7.1) give
\[
 h^1(D,V|_D)\leq r\left(\sum_i g(D_i)+\#\{\text{nodes}\}\right).
\]
As in F5, the Artin--Schreier sequence identifies $H^1(D,M)$ with
the Frobenius-fixed subgroup of $H^1(D,V|_D)$, so its
$\mathbf F_p$-dimension is at most the latter $k$-dimension.
The nearby-cycle five-term sequence finishes the proof.
\end{proof}

\begin{lemma}[F7a: a relative pro-$p$ generation criterion; Proved]
Let $1\to P\to G\to\Gamma\to1$ be an exact sequence of profinite
groups, with $P$ pro-$p$ and $\Gamma$ generated by $d$ elements.
Suppose that for every finite simple $\mathbf F_p[G]$-module $M$
there is a uniform bound
$\dim H^1(G,M)\leq C\dim M$. Then $G$ is generated by at most
$\max(d,\lceil C\rceil+1)$ elements. No solvability hypothesis on
$\Gamma$ is needed.
\end{lemma}
\begin{proof}
Put $D_0=\max(d,\lceil C\rceil+1)$. In any finite quotient $F$
of $G$, let $N$ be the image of $P$. Induct on $|N|$ to generate
$F$ by $D_0$ elements. If $N=1$, $F$ is a quotient of $\Gamma$
and the assertion holds. Otherwise choose a minimal nontrivial
$F$-normal subgroup $M\subseteq N$. It is an elementary abelian
$p$-group and a simple $\mathbf F_p[F/M]$-module, by the same
characteristic-subgroup argument as in F4. Induction generates
$F/M$ by $D_0$ elements.

The nonsplit case and the count of nongenerating lifts as
$|Z^1(F/M,M)|$ are exactly those in F4. In the split case the
nonzero $\operatorname{End}_{F/M}(M)$ summand in $H^1(F,M)$
gives the strict bound
\[
 |Z^1(F/M,M)|\leq |M|^{C+1}/p<|M|^{D_0}.
\]
Inflation to $G$ supplies the assumed bound on $H^1(F,M)$.
Thus a generating lift exists. Compactness of $G^{D_0}$, applied
to all finite quotients, completes the argument as in F4.
\end{proof}

\begin{corollary}[F7b: pro-$p$ extensions of quotients from the resolution; Proved]
Let $G$ be a quotient of $\pi_1(U)$ with normal pro-$p$ subgroup
$P$. Suppose the composite
\[
 \pi_1(U)\longrightarrow G\longrightarrow\Gamma=G/P
\]
factors through $\pi_1(X)$. Then $G$ is topologically finitely
generated, with a bound depending only on the fixed resolution $X$.
The group $\Gamma$ need not be abelian, solvable, or prime to $p$.
\end{corollary}
\begin{proof}
Every finite simple $\mathbf F_p[G]$-module is trivial on the normal
pro-$p$ subgroup, by the invariant-vector argument in F6. Its
pullback to $\pi_1(U)$ therefore extends to $X$ under the stated
factorization. Inflation and F7 give the uniform bound required
by F7a.

It remains to bound the generators of $\Gamma$. Proper
specialization over the complete henselian ring $A$ gives
$\pi_1(X)\simeq\pi_1(D)$; passing to the reduced special fibre does
not change finite \'etale covers
\cite[Lemma 58.9.1, tag 0A48]{StacksProperEtaleSpecialization}.
The connected curve $D$ is proper over the algebraically closed
field, so its fundamental group is finitely generated. For a
current stronger primary statement, see
\cite[Theorem 1.2, p.~1]{LaraSrinivasStix2024}, which proves finite
presentation for every connected proper scheme over an
algebraically closed field (and recalls SGA~1 X, Theorem~2.9 for
finite generation). Hence one fixed generator bound works for
all quotients $\Gamma$ of $\pi_1(X)$, and F7a applies.
\end{proof}

This does not assert that the kernel of $\pi_1(U)\to\pi_1(X)$ is
pro-$p$. Tame inertia and more complicated wild quotients may
remain in that kernel. The explicit normal-pro-$p$ hypothesis in
F7b is essential to the argument.

\subsection{Uniform tame coefficients: a root-stack argument}

\begin{proposition}[F8: uniform tame coefficient bound; Proved]
Let $A$ be a complete normal two-dimensional local domain of
characteristic $p>0$ with algebraically closed residue field $k$.
Put $U=\Spec A\setminus\{\mathfrak m\}$, and fix an SNC resolution
$X\to\Spec A$ with reduced exceptional divisor
$D=\bigcup_{i=1}^s C_i$. There is a constant $C_X$ such that
\[
 \dim_{\mathbf F_p}H^1(U,M)\leq C_X\operatorname{rank}M
 \tag{F8.0}
\]
for every finite locally constant $\mathbf F_p$-sheaf $M$ on $U$
whose monodromy is tamely ramified along $D$.
The finite monodromy group need not be abelian or of order prime to $p$.
\end{proposition}

\paragraph{Status.}
The complete proof passed an independent root audit, including the
primary sources and nearby-cycle formulas; BCE separately checked
the root-gerbe degree, coarse-descent, and character-count steps.
This is a proved result of the present work, without a historical
priority claim. The uniformity in root orders is part of the proof.

\begin{proof}
Choose a coefficient field $k\hookrightarrow A$. Cohen structure
expresses $A$ as a quotient of a power-series ring over perfect $k$,
which is finite over its Frobenius image; hence $A$ and the
essentially finite-type chart rings below are $F$-finite. If $A$ is regular, first blow up its closed point,
so $D$ is nonempty. Write
\[
 B_{ij}=C_i\cdot C_j,\qquad g_i=g(C_i),\qquad
 \nu_i=\#\{\text{crossings on }C_i\}.
\]
Choose positive integers $c_i$ so that
$J=\mathcal O_X(-Z)$, $Z=\sum_i c_iC_i$, has positive degree
\[
 d_i=-\sum_jc_j B_{ij}>0
\]
on every component. Such a cycle exists by the negative definiteness
used in Suzuki's proof. Set
\[
 L_i=\sum_jc_j|B_{ij}|,\qquad
 K_i=\max\{0,-B_{ii},2g_i-2+\nu_i\}.
 \tag{F8.1}
\]
These constants depend only on the chosen resolution.

\paragraph{1. Root stacks, extension, and degree conventions.}
Choose an integer $N$ prime to $p$ divisible by the tame inertia
orders of the finite monodromy of $M$ along all $C_i$. Let
\[
 \pi:\mathfrak X_N=\sqrt[N]{(X,C_1)}\mathbin{\times_X}\cdots
                 \mathbin{\times_X}\sqrt[N]{(X,C_s)}\longrightarrow X.
\]
Locally at a crossing this has the quotient chart
\[
 \left[\Spec R[u,v]/(u^N-x,v^N-y)\,/\,\mu_N^2\right],
\]
with the analogous one-variable chart away from crossings.
These charts are regular, and the stack is a regular tame
Deligne--Mumford stack, proper over $A$, with coarse space $X$.
Let $\mathfrak C_i$ be its reduced divisor above $C_i$. Then
$N\mathfrak C_i=\pi^*C_i$.

The sheaf $M$ extends to a finite locally constant sheaf on
$\mathfrak X_N$. Indeed the displayed root charts kill its
codimension-one tame inertia by the DVR form of Abhyankar's lemma.
Normalization on the regular charts is then \'etale by purity;
the canonical normalizations and their descent data agree on
chart overlaps. This proves the asserted extension, using only
Abhyankar's lemma and regular-scheme purity
\cite[tags 0BRM and 0BMB]{StacksAbhyankar,StacksBranchPurity}.
The same argument applies if $M$ has wild monodromy away from
these inertia groups.

The component $\mathfrak C_i$ is a $\mu_N$-gerbe over the root
curve obtained from $C_i$ by inserting $N$th roots at its $\nu_i$
crossing points. Degrees on $\mathfrak C_i$ will be normalized
so that the pullback of a line bundle on $C_i$ has its ordinary
degree. Concretely every line bundle has its $N$th power descending
to $C_i$, and we divide that descended degree by $N$. This convention
is additive and gives
\[
 \deg\bigl(\mathcal O(\mathfrak C_j)|_{\mathfrak C_i}\bigr)
        =B_{ij}/N,
 \qquad \deg(J|_{\mathfrak C_i})=d_i.
 \tag{F8.2}
\]
For vector bundles use the degree of the determinant. Nonzero
maps of line bundles do not increase degree, since their zeros
are effective divisors of nonnegative rational degree.
The degree multiplies by $p$ under absolute Frobenius, since
Frobenius pulls a line bundle back to its $p$th power.
The local charts and degree convention agree with the usual
root-stack/parabolic description
\cite[\S\S2.3--2.4]{Cadman2007}
\cite[Theorems 4--5, pp.~12,17]{Borne2007}.

Set $V=M\otimes_{\mathbf F_p}\mathcal O_{\mathfrak X_N}$, of rank $r$.
It has a unit Frobenius structure $F^*V\simeq V$ and is trivialized
by a finite \'etale cover. Consequently $V|_{\mathfrak C_i}$ is
semistable of degree zero: pull a subbundle back to that cover,
take its determinant inside an exterior power of a trivial bundle,
and use the preceding degree inequality. Under a finite representable
cover the normalized degrees scale by a positive rational factor:
ordinary stack degree scales by the degree of the cover, while
normalizing multiplies by the generic stabilizer order. Thus degree
signs are preserved. Its determinant has degree zero because it
becomes trivial on the cover. No restriction
on the degree of the trivializing cover is used here. For root curves with trivial generic stabilizer this is also
the finite-bundle semistability assertion of
\cite[Proposition 6, p.~22]{Borne2007}.

\paragraph{2. Coarse cohomology and a useful curve estimate.}
For a vector bundle $W$ on $\mathfrak C_i$, decompose it by the
characters of the generic $\mu_N$. All nontrivial generic-character
summands have zero coherent cohomology: taking invariants along
that gerbe is exact and kills them. The character-zero summand
$W_0$ descends to the root curve of $C_i$ at the crossings.
Its pushforward $E_0$ to $C_i$ is a vector bundle of the same rank.
On a local root chart a graded basis gives
\[
 0\leq \deg W_0-\deg E_0\leq\nu_i\operatorname{rank}W_0.
 \tag{F8.3}
\]
For clarity, adjunction from the pullback of $E_0$ is a generically
invertible map to $W_0$. In a character basis at a stacky point,
its matrix has diagonal entries $z^{a_j}$ with $0\leq a_j<N$.
The determinant degree lost in pushforward is $\sum_j a_j/N$.
Summing over the crossings proves (F8.3). Graded bases are obtained
by lifting a character basis modulo the maximal ideal and using
Nakayama's lemma; averaging makes the lifts equivariant.

Thus if $W_0$ is semistable of negative slope $t$, then it has no
sections and ordinary Riemann--Roch for $E_0$ gives
\[
 h^1(\mathfrak C_i,W_0)
 \leq \operatorname{rank}(W_0)(g_i-1-t+\nu_i).
 \tag{F8.4}
\]
If its slope is zero, evaluation at a nonstacky point is injective
on sections, since a section vanishing there gives a positive-degree
line subbundle after saturation. Hence
\[
 h^1(\mathfrak C_i,W_0)
 \leq \operatorname{rank}(W_0)(g_i+\nu_i).
 \tag{F8.5}
\]
These statements also apply to an isotypic summand tensored by a
line bundle that makes its generic character zero. Each isotypic
summand of a degree-zero semistable bundle is again semistable of
degree zero: it is a direct summand, every summand has degree at
most zero, and their degrees sum to zero.

Coarse pushforward on these tame stacks is exact on coherent
sheaves and on sheaves of $\mathbf F_p$-vector spaces. Locally it
is invariants for a finite group of order prime to $p$, followed
by finite affine pushforward; the invariants functor is exact by
averaging. For the \'etale assertion, take a strict henselian
neighborhood on the coarse space: its finite root cover is a finite
union of strictly henselian local schemes, whose \'etale global
section functors are exact. The quotient atlas then computes the
pushforward by finite-group cohomology, which vanishes in positive
degrees. The same description commutes with restriction to a closed
fibre. More explicitly, write $j:U\hookrightarrow\mathfrak X_N$ and
$j_X:U\hookrightarrow X$. Exact coarse pushforward gives
$R\pi_*Rj_*M=Rj_{X*}M$, and its closed-fibre compatibility gives
$i_X^*R\pi_*Rj_*M=R\pi_{D*}i^*Rj_*M$. These identities apply
to the bounded-below torsion complex $Rj_*M$ degree by degree.
Proper base change for the proper scheme $X$ over the strictly
henselian base $\Spec A$ therefore identifies
$H^1(U,M)$ with $\mathbb H^1(\mathfrak D_N,i^*Rj_*M)$.
Its spectral sequence yields the required nearby-cycle five-term
sequence. Coherent cohomology is computed by the same exact
coarse pushforward and the preceding coarse-space estimates.

\paragraph{3. A Frobenius-cokernel slope bound independent of $N$.}
Let
\[
 B_N=(F_*\mathcal O_{\mathfrak X_N})/\mathcal O_{\mathfrak X_N}.
\]
It is locally free: absolute Frobenius induces the automorphism
$\zeta\mapsto\zeta^p$ on each $\mu_N$ stabilizer, so it is
representable and is finite flat on the regular root charts. Its
unit section extends locally to a basis.
The pullback $F^*B_N$ is canonically isomorphic, as a module, to
the augmentation ideal in $F^*F_*\mathcal O$: the unit provides a
section of the augmentation. The powers of that ideal give a
filtration with locally free graded pieces
\[
 T^a(\Omega^1_{\mathfrak X_N/k}),\qquad 1\leq a\leq2(p-1),
 \tag{F8.6}
\]
where $T^a$ is the degree-$a$ truncated symmetric power.
This is Sun's canonical filtration
\cite[Theorem 3.7, pp.~12--13; local proof pp.~6--8]{Sun2008}.
In the present complete-local setting it can also be checked
directly after faithfully flat completion of a regular root chart:
with parameters $u,v$ and coefficient field $k$, the algebra is
\[
 k[[u,v]]\otimes_{k[[u^p,v^p]]}k[[u,v]]
 \simeq k[[u,v]][\alpha,\beta]/(\alpha^p,\beta^p),
\]
with augmentation ideal $(\alpha,\beta)$. Its graded generators
are $du,dv$. This calculation is invariant under changes of
parameters and therefore descends. All modules and maps involved are finite. Every closed point on
the proper exceptional divisor has residue field $k$, so its
completed regular chart is $k[[u,v]]$. Faithfully flat completion
checks these identities at every such point; any nonzero coherent
error supported on that divisor would have a closed point in its
support. Thus the calculation suffices on a neighborhood of the
exceptional divisor, which is all that the proof uses.

The conormal sequence on $\mathfrak C_i$ is an extension of two
line bundles, of degrees
\[
 -B_{ii}/N,\qquad
 2g_i-2+\nu_i(1-1/N).
 \tag{F8.7}
\]
The second formula follows from $dx=Nz^{N-1}dz$ at each root point;
the generic gerbe is \'etale and contributes no cotangent direction.
Both numbers are at most $K_i$. Each $T^a(\Omega^1)|_{\mathfrak C_i}$
has a filtration by the characters of this two-line flag: its
weights are $u,v$ with $0\leq u,v\leq p-1$ and $u+v=a$.
Thus its line-bundle layers have degrees at most $2(p-1)K_i$.
Tensoring with the semistable degree-zero bundle $V|_{\mathfrak C_i}$
shows that every subbundle of
$F^*(V\otimes B_N)|_{\mathfrak C_i}$ has slope at most
$2(p-1)K_i$. Here we used $F^*V\simeq V$ and the filtration (F8.6).
The assertion about subbundles follows by intersecting with the
filtration and adding degrees and ranks; each graded subbundle has
the indicated slope bound.

Pulling any subbundle of $(V\otimes B_N)|_{\mathfrak C_i}$ back by
Frobenius and dividing degrees by $p$ therefore proves the uniform
upper slope bound
\[
 \mu(\text{any subbundle of }(V\otimes B_N)|_{\mathfrak C_i})
       \leq b_i:=2(p-1)K_i/p.
 \tag{F8.8}
\]
This is independent both of $N$ and of $r$.

\paragraph{4. Nearby-cycle filtration and a uniform cutoff.}
Write $\mathfrak Z_N=\pi^*Z=\sum_iNc_i\mathfrak C_i$ and continue
to denote its ideal by $J$. The sheaf $M$ is locally constant on
$\mathfrak X_N$, so the Artin--Schreier calculation of F5/F7 applies
on its \'etale charts and descends. In particular, with
$\Psi=i^*j_*$ for the reduced exceptional stack,
\[
 0\to J^{-m}V/V\xrightarrow{F_V-1}J^{-pm}V/V
       \to U_mR^1\Psi M\to0,
 \tag{F8.9}
\]
and, for $m\geq2$,
\[
 \operatorname{gr}_mR^1\Psi M
  \simeq B_N|_{\mathfrak Z_N}\otimes V|_{\mathfrak Z_N}
                          \otimes J^{-m}|_{\mathfrak Z_N}.
 \tag{F8.10}
\]
The projection formula proves the latter identity, as in Suzuki's
Proposition 3.1; $F^{-1}\mathfrak Z_N=p\mathfrak Z_N$.

Filter $\mathcal O_{\mathfrak Z_N}$ by removing one irreducible
component at a time. A line-bundle layer $P$ on $\mathfrak C_i$
has degree bounded above by $L_i$, by (F8.2): all cycle coefficients
are between $0$ and $Nc_j$. The corresponding layer of (F8.10)
therefore has every subbundle of slope at most
\[
 b_i+L_i-md_i.
\]
Choose an integer $m_0\geq2$ strictly greater than
$\max_i(b_i+L_i)/d_i$. The last expression is negative for every
$m\geq m_0$, so there are no sections on any layer and hence
\[
 H^0(\operatorname{gr}_mR^1\Psi M)=0\quad(m\geq m_0).
 \tag{F8.11}
\]
The number of layers grows with $N$, but their common vanishing
bound does not. Quasi-compactness and the increasing union of the
pole lattices give
$H^0(R^1\Psi M)=H^0(U_{m_0}R^1\Psi M)$, increasing $m_0$ by one
if desired.

\paragraph{5. Low poles: cancellation of the growing number of layers.}
The negative-degree filtration of Suzuki's Proposition 2.1 works
on the root surface as well. Its proof only uses the negative
definiteness of the intersection matrix and the nonnegativity of
off-diagonal intersections; here the normalized matrix is $B/N$.
Thus $J^{-pm}/\mathcal O$ and $J^{-pm}/J^{-m}$ admit filtrations
by line bundles of strictly negative degree on the $\mathfrak C_i$.
Tensoring with $V$ gives no sections, because $V|_{\mathfrak C_i}$
is semistable of degree zero. Equation (F8.9) consequently injects
$H^0(U_mR^1\Psi M)$ into the kernel of a Frobenius-linear map minus
an injective linear map on
$H^1(J^{-m}V/V)$. The fixed-space argument from F5 bounds its
$\mathbf F_p$-dimension by $h^1(J^{-m}V/V)$.

We bound the latter dimension uniformly for the fixed $m=m_0$.
Choose the negative filtration of $J^{-m}/\mathcal O$ by adding
components to the cycle $m\mathfrak Z_N$. For a layer $P$ on
$\mathfrak C_i$ we have
\[
 -mL_i\leq\deg P<0.
 \tag{F8.12}
\]
Decompose $V|_{\mathfrak C_i}=\bigoplus_{\chi}V_{i,\chi}$ by its
generic $\mu_N$-characters, and let $r_{i,\chi}$ be their ranks.
Only the summand whose character cancels that of $P$ contributes
to cohomology. As the coefficient of $\mathfrak C_i$ increases
from $0$ to $mNc_i$, the generic character of the successive
layers runs through every residue modulo $N$ exactly $mc_i$
times. Contributions from other components have trivial generic
character at $\mathfrak C_i$. Therefore the sum of the ranks of
all contributing summands for this component is exactly
\[
 mc_i\sum_\chi r_{i,\chi}=mc_i r.
 \tag{F8.13}
\]
This is the needed cancellation of the extra $N$ layers.

Each contributing bundle $V_{i,\chi}\otimes P$ is semistable of
negative slope $\deg P$ and descends after killing the generic
gerbe. By (F8.4) and (F8.12), its first cohomology has dimension
at most its rank times $g_i+\nu_i+mL_i+1$. Summing the exact
filtration and using (F8.13) gives
\[
 h^1(J^{-m}V/V)
 \leq r\,m\sum_i c_i(g_i+\nu_i+mL_i+1).
 \tag{F8.14}
\]
The right side is independent of $N$. This bounds
$\dim_{\mathbf F_p}H^0(R^1\Psi M)$ by a constant times $r$.

\paragraph{6. The proper exceptional term.}
Let $\mathfrak D_N=\bigcup_i\mathfrak C_i$ with its reduced
structure. Its normalization sequence has one term at each crossing,
with stabilizer $\mu_N^2$. Invariants there have dimension at most
$r$ and no higher coherent cohomology. On a component, only the
generic-character-zero summand of $V$ contributes; it is semistable
of degree zero. Estimate (F8.5) thus yields
\[
 h^1(\mathfrak D_N,V)
 \leq r\left(\sum_i(g_i+\nu_i)+\#\{\text{crossings}\}\right).
 \tag{F8.15}
\]
The Artin--Schreier sequence identifies
$H^1(\mathfrak D_N,M)$ with the Frobenius-fixed subgroup of this
finite-dimensional coherent cohomology. Surjectivity of $F_V-1$
on global sections follows from the same finite-dimensional
semilinear argument as in F5. Its $\mathbf F_p$-dimension is
therefore at most (F8.15).

The nearby-cycle five-term sequence, justified by coarse descent
in step 2, combines (F8.14) and (F8.15) to prove (F8.0). One possible
constant is
\[
 C_X=m_0\sum_i c_i(g_i+\nu_i+m_0L_i+1)
            +\sum_i(g_i+\nu_i)+\#\{\text{crossings}\}.
\]
It depends on $X,J,p$ and not on the coefficient rank, the monodromy
group, or any tame inertia order.
\end{proof}

\begin{corollary}[F8b: all pro-$p$-by-pro-prime-to-$p$ quotients; Proved]
Every quotient $G$ of $\pi_1(U)$ admitting an
exact sequence $1\to P\to G\to\Gamma\to1$ with $P$ pro-$p$ and
$\Gamma$ pro-prime-to-$p$ is uniformly topologically finitely
generated. No abelian hypothesis on $\Gamma$ is necessary.
\end{corollary}
\begin{proof}
Every finite simple $\mathbf F_p[G]$-module kills $P$ and hence has
prime-to-$p$ monodromy. Such a system is tamely ramified on the
resolution. F8 and inflation give the uniform degree-one bound.
The quotient $\Gamma$ has at most a fixed number of generators,
because it is a quotient of $\pi_1(U)^{(p')}$, which is finitely
generated by Grothendieck--Murre. The relative criterion F7a gives
the conclusion. Alternating wild and tame normal factors remain
outside this assertion, as do general nonsolvable wild kernels.
\end{proof}

\subsection{Pro-$p$ extensions of the full tame quotient}

The tame finite-generation statement below follows by the standard
curve and gluing methods; the uniform extension consequence F9b
uses the coefficient estimate proved in F8. Both arguments and the
SGA~1 source passed an independent root audit. No historical
priority assertion is made.

\begin{lemma}[F9: a uniform generator bound for the tame quotient; Known method, proved here]
Keep the notation of F8, and let $\pi_1^t(U;X,D)$ denote the quotient
classifying finite covers tamely ramified along the fixed SNC divisor
$D$ on $X$. Let $\Delta$ be the connected dual graph of $D$. Then
\[
 d\bigl(\pi_1^t(U;X,D)\bigr)
 \leq d_X:=\sum_i(2g_i+\nu_i+1)+b_1(\Delta).
 \tag{F9.1}
\]
In particular this is a bound for the full tame group, not only its
maximal prime-to-$p$ quotient.
\end{lemma}
\begin{proof}
Let $N$ be prime to $p$, and use the root stack $\mathfrak X_N$
and reduced exceptional stack $\mathfrak D_N$ of F8. We prove
that every finite quotient of $\pi_1(\mathfrak X_N)$ satisfies
(F9.1), independently of $N$.

\emph{The closed fibre surjects on fundamental groups.}
Let $\mathfrak Y\to\mathfrak X_N$ be a connected finite \'etale
cover. Its coarse space $Y$ is a finite scheme over $X$.
Indeed, locally $\mathfrak X_N=[V/H]$ with $V$ finite over an
affine open $\Spec R\subset X$ and $H=\mu_N$ or $\mu_N^2$.
The pullback of the cover is $W\to V$ finite \'etale, with
$H$-equivariant descent, and its coarse space is
$\Spec\Gamma(W,\mathcal O_W)^H$. This invariant ring is finite
over $R$: it is an $R$-submodule of a finite $R$-module and $R$
is noetherian. These finite coarse schemes glue. Thus $Y$ is
proper over $A$, and it is connected because a coarse-space map
identifies the underlying topological spaces.

For a proper scheme over a henselian local ring, connectedness
is equivalent to connectedness of the closed fibre. For example,
apply the full faithfulness in
\cite[Lemma 58.9.1, tag 0A48]{StacksProperEtaleSpecialization}
to the scheme and its trivial two-sheeted cover: sections of that
cover correspond exactly to open-and-closed decompositions.
Consequently $Y_k$ is connected. Locally in the above quotient
charts, taking invariants commutes with the residue-field base
change by averaging, since $|H|$ is invertible. Therefore $Y_k$
is the coarse space of $\mathfrak Y_k$, which is connected.
Replacing the closed fibre by its reduced support does not change
its underlying topological space. It follows that the restriction
of every connected finite \'etale cover to $\mathfrak D_N$ stays
connected. This is precisely the surjectivity of
\[
 \pi_1(\mathfrak D_N)\longrightarrow\pi_1(\mathfrak X_N).
 \tag{F9.2}
\]
This argument requires only this surjectivity; it does not assume
a full stack version of the proper specialization equivalence.

\emph{A bound for each normalization component.}
The component $\mathfrak C_i$ is the banded $\mu_N$-gerbe over
the root curve $\mathcal C_i$ obtained from $C_i$ by taking roots
at its $\nu_i$ crossing points. In every finite quotient $G$ of
$\pi_1(\mathfrak C_i)$, the image $I$ of the generic relative
$\mu_N$ is cyclic and central. To see centrality, the band supplies
natural automorphisms of every fibre functor; naturality makes
them commute with its automorphisms. Quotienting the corresponding
torsor by $I$ kills the relative gerbe inertia, so the torsor
descends to a finite \'etale $G/I$-torsor on $\mathcal C_i$.
This descent can be checked on a neutralizing \'etale cover of the
gerbe, where it is exactly descent of an equivariant finite \'etale
algebra with trivial $\mu_N$-action.

Every finite \'etale cover of $\mathcal C_i$ restricts to a cover
of $C_i$ minus its crossing points, tamely ramified at those
points. The local chart is $[\Spec k[[z]]/\mu_N]$ with coarse
parameter $z^N$, so the inertia order divides $N$. A connected
cover remains connected on this dense open subset because the
root curve is normal. Hence
$\pi_1^t(C_i\setminus\{\text{crossings}\})$ surjects onto
$\pi_1(\mathcal C_i)$.
SGA~1, Expos\'e~XIII, Corollaire~2.12 gives $2g_i+\nu_i$
generators for this full tame curve group
\cite[original p.~392; retyped pp.~290--292]{SGA1TameCurves}.
It follows that every such $G$ needs at most
\[
                              2g_i+\nu_i+1               \tag{F9.3}
\]
generators: lift generators of $G/I$ and adjoin one generator of
$I$. The compactness criterion for generators gives the same bound
for $\pi_1(\mathfrak C_i)$ itself.

\emph{Gluing the components.}
Finite \'etale covers of $\mathfrak D_N$ are obtained by gluing
finite \'etale covers of its normalization components, with
identifications over the two preimages of each node gerbe.
The local assertion follows on the node chart
\[
 [\Spec k[[u,v]]/(uv)\,/\,\mu_N^2].
\]
The strictly henselian nodal local ring and its two strictly
henselian branches have only trivial finite \'etale covers. Thus
on the quotient chart all three cover categories are finite
$\mu_N^2$-sets, and the gluing statement is the identification
of these sets at the two branches. The assertion descends from
these charts and the smooth locus. Equivalently, it is the usual
normalization gluing of finite \'etale algebras, with equivariant
descent along a finite tame atlas.

For completeness, the resulting generator bound does not require
injectivity of any edge-group map. Choose a spanning tree of
$\Delta$, base points on the components, and paths to the incident
node points. The gluing identifications along the tree allow one
to use a single fibre. The monodromy is then generated by the
images of the component fundamental groups and one gluing
element for every edge outside the tree. If the subgroup they
generated in a finite connected torsor were proper, its orbits
would give a nontrivial decomposition of all component covers
compatible with every gluing identification, hence a decomposition
of that torsor. This contradicts connectedness. Thus the number
of generators is bounded by the sum of (F9.3) and $b_1(\Delta)$.
Together with (F9.2), this proves the asserted uniform bound for
every $\pi_1(\mathfrak X_N)$.

Every finite tame quotient of $\pi_1(U)$ extends to one of these
root stacks by the Abhyankar--purity argument of F8. Consequently
every finite quotient of $\pi_1^t(U;X,D)$ obeys the same bound.
The compactness argument for generating tuples now proves (F9.1).
\end{proof}

\begin{corollary}[F9b: a normal pro-$p$ kernel over a tame quotient; Proved]
Let $G$ be a quotient of $\pi_1(U)$, let $P\triangleleft G$ be
pro-$p$, and assume that the composite
\[
 \pi_1(U)\longrightarrow\Gamma:=G/P
\]
factors through $\pi_1^t(U;X,D)$. Then
\[
 d(G)\leq\max\{d_X,\lceil C_X\rceil+1\},
 \tag{F9.4}
\]
where $C_X$ is the uniform constant in F8. No solvability
hypothesis is imposed on $\Gamma$ or on its finite quotients.
\end{corollary}
\begin{proof}
Every finite simple $\mathbf F_p[G]$-module is killed by $P$,
because a normal finite $p$-group has nonzero invariants on every
nonzero representation in characteristic $p$, and the invariant
space is $G$-stable. The associated sheaf on $U$ therefore has
tame boundary monodromy. Inflation and F8 give the uniform
rank-linear bound on $H^1(G,M)$ for all these simple modules.
F9 bounds the number of generators of $\Gamma$. The relative
criterion F7a proves (F9.4).
\end{proof}

\paragraph{Remaining wild obstruction.}
The kernel of $\pi_1(U)\to\pi_1^t(U;X,D)$ is normally generated
by wild inertia groups. This does not make it a pro-$p$ group.
For example, a finite group generated by $p$-subgroups can have
minimal normal subgroups of another prime order: $S_3$ is generated
by its order-two subgroups and has normal subgroup $C_3$.
Thus F9b does not close the alternating tame/wild case or the full
local fundamental-group conjecture. No assertion of historical
priority is attached to F9 or its proof.

\subsection{Abelian wild inertia and metabelian quotients}

\begin{proposition}[F10: the abelian inertia kernel is finite; Proved]
Let $A,U,X,D=\bigcup_iC_i$ be as in F8. The kernel of
\[
 \pi_1(U)^{\mathrm{ab}}\longrightarrow\pi_1(X)^{\mathrm{ab}}
 \tag{F10.1}
\]
is finite. More explicitly, write $B=(C_i\cdot C_j)$ and set
\[
 k_i=\max\{0,2g_i-2+\nu_i\},\qquad
 t=(-B)^{-1}k.
 \tag{F10.2}
\]
Every $p$-primary rank-one character $\chi$ of $\pi_1(U)$ has
Swan conductor $s_i(\chi)\leq t_i$ at $C_i$.
Choose $e\geq0$ with $p^e>\max_i t_i$. Then the kernel in
(F10.1) is annihilated by $p^e|\det B|$.
\end{proposition}

\paragraph{Status and source.}
The complete argument passed an independent root proof and primary-source
audit; BCE separately passed F10 and F10b. The input used here is the
global refined Swan conductor on regular $F$-finite
schemes in Krishna--Majumder, arXiv:2601.12177v2, 8 November 2026,
Corollary 9.7 p.~56 (Theorem 9.4 pp.~55--56), with the hypotheses
on p.~2. The source is a preprint. Its formulation includes
complete-local resolutions and does not require algebraization
or finite type over a field. No priority assertion is made.

\begin{proof}
A coefficient field identifies $A$ as a quotient of a power-series
ring over perfect $k$, so $A$ and the resolution charts are
$F$-finite. The resolution is an equidimensional noetherian regular
$\mathbf F_p$-scheme with SNC boundary. These are precisely the
hypotheses of the cited global refined-conductor theorem.
All differential forms over $\mathbf F_p$ below may also be viewed
over $k$, since $k$ is perfect and $\Omega^1_{k/\mathbf F_p}=0$.

Fix a character $\chi\in H^1(U,\mathbf Q_p/\mathbf Z_p)$.
It has finite order, say dividing $p^m$. Let
\[
                         S=\sum_i s_i C_i
\]
be its logarithmic Swan-conductor divisor. Each $s_i$ is a
nonnegative integer, defined at the henselian discrete valuation
field at the generic point of $C_i$. If $s_i>0$, the character is
in $\operatorname{fil}^{\log}_S H^1(U,\mathbf Z/p^m)$ and is not
in $\operatorname{fil}^{\log}_{S-C_i}H^1(U,\mathbf Z/p^m)$.
Since $S-C_i\geq\lfloor S/p\rfloor$, Corollary 9.7 of
\cite{KrishnaMajumder2026} gives a nonzero section
\[
 0\ne\operatorname{Rsw}_i(\chi)\in
 H^0\!\left(C_i,
   \Omega_X^1(\log D)|_{C_i}\otimes\mathcal O_X(S)|_{C_i}\right).
 \tag{F10.3}
\]
The map is a levelwise monomorphism, not only an injection after
passing to an unspecified colimit; see Theorem 9.4(1).
Its generic stalk is the local refined Swan conductor by
Corollary 9.6. Thus the nonvanishing in (F10.3) has exactly the
required meaning.

The residue sequence on $C_i$ is
\[
 0\longrightarrow\Omega^1_{C_i/k}(\log\{\text{crossings}\})
 \longrightarrow\Omega_X^1(\log D)|_{C_i}
 \longrightarrow\mathcal O_{C_i}\longrightarrow0.
 \tag{F10.4}
\]
It follows either from the local coordinates $d\log x,dy$ and
$d\log x,d\log y$, or from Lemma 9.8 of the same source.
The two line bundles have degrees $2g_i-2+\nu_i$ and $0$.
Consequently every line subbundle of the middle term has degree
at most $k_i$. A nonzero section in (F10.3), after saturation,
therefore implies
\[
                              S\cdot C_i\geq-k_i.
 \tag{F10.5}
\]
When $s_i=0$, this inequality holds as well, because
$S\cdot C_i=\sum_{j\ne i}s_j(C_j\cdot C_i)\geq0$.
Thus, as a vector inequality, $Bs\geq-k$.

We include the elementary matrix step. Put $C=-B$. This matrix
is symmetric positive definite and has nonpositive off-diagonal
entries. Its inverse is entrywise nonnegative. Indeed, if
$Cx=v\geq0$, write $x=x_+-x_-$ with $x_+,x_-$ nonnegative and
of disjoint support. If $x_-\ne0$, then
\[
 0\leq x_-^Tv=x_-^TCx_+-x_-^TCx_-<0,
\]
because the mixed term is nonpositive and the quadratic term is
positive. This contradiction proves $x\geq0$.
Applying $C^{-1}$ to $Cs\leq k$ now gives $s\leq C^{-1}k=t$.
This proves the uniform conductor bound.

\emph{A fixed $p$-power kills all abelian wild inertia.}
We recall explicitly why multiplication by $p$ reduces the bound
by a factor $p$. For a henselian characteristic-$p$ DVR $R$ with
fraction field $L$, the Brylinski filtration on Witt vectors
$a=(a_{m-1},\ldots,a_0)\in W_m(L)$ is
\[
 \operatorname{fil}_nW_m(L)
   =\{a:p^jv(a_j)\geq-n\text{ for all }j\}.
\]
Its image under Artin--Schreier--Witt is the logarithmic conductor
filtration. Witt restriction drops the last displayed coordinate
and gives
\[
 R:\operatorname{fil}_nW_m(L)
       \longrightarrow\operatorname{fil}_{\lfloor n/p\rfloor}
                         W_{m-1}(L).
 \tag{F10.6}
\]
This is Lemma 3.2, p.~12, of \cite{KrishnaMajumder2026}; the
classical characteristic-$p$ description is also given in
\cite[\S3.2.2--3.2.4, pp.~14--15]{KatoLealSaito2019}.
Under the standard embeddings $\mathbf Z/p^m\hookrightarrow
\mathbf Q_p/\mathbf Z_p$, restriction represents multiplication
by $p$. Hence
\[
 s_i(p^e\chi)\leq\lfloor s_i(\chi)/p^e\rfloor=0.
\]
Conductor zero here means unramified: its Witt representative is
integral, and the Artin--Schreier--Witt equation then defines an
\'etale cover of the DVR. Thus $p^e\chi$ is unramified along every
$C_i$. Purity on the regular $X$ extends it to $X$.

\emph{The prime-to-$p$ part.}
The Kummer/intersection calculation in F5 shows that multiplying
any prime-to-$p$ character by $b=|\det B|$ kills all its inertia
along the $C_i$. Purity again extends the resulting character
to $X$. One can see the calculation directly: extend its torsion
line bundle to $X$, write the divisor of the trivialization of
its $n$th power as $\sum a_i C_i$, and intersect with $C_j$.
Then $Ba\equiv0\pmod n$, so $ba_i\equiv0\pmod n$ by the
adjugate matrix. These residues describe the tame inertia.

Set $Q=\pi_1(U)^{\mathrm{ab}}$ and let $I$ be the kernel in
(F10.1). The map is surjective, since a connected finite \'etale
cover of normal $X$ stays connected on its dense open subset $U$.
The previous two paragraphs show that all characters of $I$
obtained by restriction from $Q$ vanish on $bp^eI$. Finite
characters of the profinite abelian group $Q$ separate its
points, so $bp^eI=0$. By Grothendieck--Murre and Suzuki, $Q$ is
topologically finitely generated. Its subgroup of elements
annihilated by the fixed integer $bp^e$ is finite. Therefore $I$
is finite, as asserted.
\end{proof}

\begin{corollary}[F10b: one fixed cover removes every abelian inertia action; Proved]
There is a finite connected abelian cover $U'\to U$ such that,
for every finite locally constant sheaf on $U$ with abelian
monodromy, its pullback to $U'$ extends to a finite locally constant
sheaf on a fixed SNC resolution $X'$ of the normalization of $A$
in $U'$. The same cover and resolution work for all coefficient
primes and all ranks.
\end{corollary}
\begin{proof}
Choose an open subgroup $Q'\subset Q$ with $Q'\cap I=0$; this
is possible because $I$ is finite and $Q$ is profinite. Let $U'$
be the cover corresponding to the inverse image of $Q'$.
Its normalization $A'$ is finite over $A$, excellent, normal,
complete and local with residue field $k$. Fix an SNC resolution
$X'$ of $A'$. The pulled-back abelian monodromy factors through
$Q'$.

At a height-one exceptional point of $X'$, let $R$ be its DVR.
The morphism $\Spec\operatorname{Frac}(R)\to X$ extends to
$\Spec R\to X$ by properness of $X$ over $A$. Every finite
character of $Q/I$, being a cover of $X$, is unramified after this
pullback. Hence the image in $Q$ of the inertia at $R$ lies in
$I$. It also lies in $Q'$, so it is trivial. Thus the given pulled-back
monodromy is unramified at every such DVR. It is already \'etale
elsewhere. Normalization and regular-scheme purity now extend the
corresponding finite cover to $X'$. This proof uses one fixed
$U'$ and $X'$, independently of the sheaf.
\end{proof}

\begin{proposition}[F10c: all-prime bounds for abelian coefficient monodromy; Proved]
There is a constant $C_A$ such that for every prime $\ell$ and
every finite locally constant $\mathbf F_\ell$-sheaf $M$ on $U$
with abelian monodromy,
\[
                       \dim_{\mathbf F_\ell}H^1(U,M)
                       \leq C_A\operatorname{rank}M.
 \tag{F10.7}
\]
\end{proposition}
\begin{proof}
Use the fixed cover and resolution of F10b. If $\ell=p$, F7
bounds the cohomology after pullback by $C_{X'}\operatorname{rank}M$.
If $\ell\ne p$, every degree-one cocycle after pullback already
factors through the tame quotient. Indeed, it gives an affine
monodromy representation into $M\rtimes\operatorname{im}\rho$.
The original representation $\rho$ is unramified on $X'$, so its
restriction to boundary inertia is trivial. On inertia the cocycle
is a homomorphism to the additive $\ell$-group $M$. It kills wild
pro-$p$ inertia, and the affine cover is therefore tame.
Consequently
\[
 H^1(U',M)=H^1(\pi_1^t(U';X',D'),M),
\]
whose dimension is at most $d_{X'}\operatorname{rank}M$ by F9
and the values of a cocycle on generators.

Let $F=\operatorname{Gal}(U'/U)$. Inflation--restriction gives
\[
 0\to H^1(F,M^{\pi_1(U')})\to H^1(U,M)
          \to H^1(U',M)^F.
\]
The first term has dimension at most $d(F)\operatorname{rank}M$,
with no assumption that $|F|$ is prime to $\ell$.
Thus one can take
$C_A=d(F)+\max\{C_{X'},d_{X'}\}$.
\end{proof}

\begin{theorem}[F11: all metabelian local quotients are finitely generated; Proved]
Every metabelian quotient $G$ of $\pi_1(U)$ is topologically
finitely generated, with a bound depending only on $A$.
Both tame and wild normal factors are allowed.
\end{theorem}
\begin{proof}
Let $d_0=d(\pi_1(U)^{\mathrm{ab}})$ and let $C_A$ be as in
F10c. Put $D=\max\{d_0,\lceil C_A\rceil+1\}$.
We show that every finite metabelian quotient $F$ is generated
by $D$ elements. If $F'=1$, the bound $d_0$ applies. Otherwise
choose a minimal nonzero $F$-normal subgroup $M\subset F'$.
It is an elementary abelian $\ell$-group for some prime $\ell$,
and is simple as an $\mathbf F_\ell[F]$-module. Since $F'$ is
abelian, conjugation on $M$ factors through the abelian group
$F/F'$. Inflation from $F$ to $\pi_1(U)$ and F10c give
\[
                  \dim H^1(F,M)\leq C_A\dim M.
 \tag{F11.1}
\]
Inductively choose $D$ generators of $F/M$ and consider their
$|M|^D$ possible lifts. A nongenerating lift spans a proper
subgroup $J$ with $JM=F$. The subgroup $J\cap M$ is $F$-stable:
$J$ preserves it, and $M$ centralizes it. Minimality therefore
forces $J\cap M=0$, so $J$ is a complement. If there is no
complement, all lifts generate.

In the split case write $H=F/M$. Complements are indexed by
$Z^1(H,M)$. Since $M$ acts trivially on itself, splitting gives
\[
 H^1(F,M)\simeq H^1(H,M)\oplus\End_H(M).
\]
The endomorphism ring contains the scalar field $\mathbf F_\ell$.
Moreover $|B^1(H,M)|\leq|M|$. Hence (F11.1) implies
\[
 |Z^1(H,M)|\leq |M|^{C_A+1}/\ell<|M|^D.
\]
Each complement contains at most one lift of the chosen generating
tuple of $H$, so some lift generates $F$. This proves the bound
by induction. Finally the compactness of generating-tuple conditions
on $G^D$ gives $D$ topological generators of $G$.
\end{proof}

\paragraph{Scope remaining after these results.}
F11 closes derived length two, including the metabelian
abstract obstructions of F3 in every residue characteristic.
It does not treat arbitrary solvable quotients of larger derived
length: their minimal normal modules can have nonabelian
coefficient monodromy. Nor does it settle arbitrary nonsolvable
quotients. These are genuine remaining gates.

\subsection{The coefficient-$p$ half at derived length three}

Throughout this subsection, $A$ is the complete normal two-dimensional
local domain of equal characteristic $p>0$ with algebraically closed
residue field used in F8--F11, $U$ is its punctured spectrum, and
all resolutions are fixed SNC resolutions. No mixed-characteristic
assertion is made here. The proof passed an independent root audit.

\begin{lemma}[F12a: normal $p$-subgroups on simple characteristic-$p$ modules]
Let $\Gamma$ be a finite group and let $V$ be a nonzero simple
$\mathbf F_p[\Gamma]$-module. Every normal $p$-subgroup of
$\Gamma$ acts trivially on $V$. Consequently, if $\Gamma$ is
metabelian and its action on $V$ is faithful, then $\Gamma'$ has
order prime to $p$.
\end{lemma}
\begin{proof}
Every nonzero finite-dimensional representation of a finite
$p$-group $P$ over $\mathbf F_p$ has nonzero invariants. Here is an
inductive proof. A central subgroup $C_p\subset P$ exists; its
generator $u$ satisfies $(u-1)^p=0$, so the kernel of $u-1$ is
nonzero and $P$-stable. The quotient group $P/C_p$ acts on it,
and induction gives a nonzero $P$-fixed vector. If $P$ is normal
in $\Gamma$, the invariant space $V^P$ is $\Gamma$-stable, so
simplicity gives $V^P=V$.
For metabelian $\Gamma$, the $p$-primary part of its abelian
commutator group is characteristic in $\Gamma'$ and therefore
normal in $\Gamma$. It acts trivially; faithfulness makes it
trivial.
\end{proof}

\begin{proposition}[F12: characteristic-$p$ coefficients with metabelian monodromy; Proved]
There is a constant $C_{A,2,p}$ such that every simple finite
$\mathbf F_p$-local system $M$ on $U$ whose monodromy group is
metabelian satisfies
\[
                    \dim_{\mathbf F_p}H^1(U,M)
                    \leq C_{A,2,p}\operatorname{rank}M.
 \tag{F12.1}
\]
The system itself may be wildly ramified on the original resolution.
\end{proposition}
\begin{proof}
Use once and for all the finite connected abelian cover
$U_0\to U$ and the fixed resolution $X_0$ supplied by F10b.
Write $F=\operatorname{Gal}(U_0/U)$ and let
$\rho:\pi_1(U)\twoheadrightarrow\Gamma\subset\operatorname{GL}(M)$ be the
faithful monodromy image. It is metabelian by hypothesis.
Lemma F12a makes $\Gamma'$ a prime-to-$p$ group.

The finite quotient $\pi_1(U)\to\Gamma^{\mathrm{ab}}$ is an
abelian local system. Its pullback to $U_0$ extends to $X_0$,
by the uniform assertion of F10b. Thus at every exceptional
DVR of $X_0$, the image under $\rho$ of boundary inertia lies
in $\Gamma'$. Wild inertia is pro-$p$, so its image in the
prime-to-$p$ group $\Gamma'$ is trivial. Therefore the entire
pulled-back system $M|_{U_0}$ is tame along the fixed exceptional
divisor of $X_0$.

No simplicity assumption on this restricted system is needed:
F8 applies to arbitrary finite $\mathbf F_p$-local systems with
tame boundary monodromy. Hence
\[
 \dim H^1(U_0,M)\leq C_{X_0}\operatorname{rank}M.
\]
Inflation--restriction bounds the kernel of restriction by
$\dim H^1(F,M^{\pi_1(U_0)})\leq d(F)\operatorname{rank}M$.
Thus $C_{A,2,p}=C_{X_0}+d(F)$ works. In particular this argument
does not assume that $[U_0:U]$ is prime to $p$ or that restriction
preserves simplicity.
\end{proof}

\begin{corollary}[F12b: a pro-$p$ second derived subgroup; Proved]
Every quotient $G$ of $\pi_1(U)$ whose closed second derived
subgroup $G''$ is pro-$p$ is topologically finitely generated,
with a bound depending only on $A$.
\end{corollary}
\begin{proof}
The quotient $\Gamma=G/G''$ is metabelian, so F11 bounds its
number of generators uniformly. Every finite simple
$\mathbf F_p[G]$-module kills the normal pro-$p$ group $G''$ by
F12a applied to its finite image. Its faithful monodromy group
is therefore metabelian, so inflation and F12 give the uniform
$H^1$ bound. The relative criterion F7a proves the assertion.
\end{proof}

\paragraph{Precise remaining derived-length-three gate.}
For a finite group of derived length at most three, a minimal
normal subgroup inside its second derived group is an elementary
abelian $\ell$-group. Its conjugation representation is simple
and factors through a metabelian group, because the second derived
group is abelian. F12 supplies the needed uniform estimate when
$\ell=p$. The cases $\ell\ne p$, with possibly wild metabelian
coefficient monodromy, remain open here. F12b is a consequence
for a specified class, not a claim that every second derived group
is pro-$p$ or that bounded-exponent inertia inside a finitely
generated metabelian group is finite.

\subsection{A fixed logarithmic conductor bound for metabelian monodromy}

Throughout, $A$ is a complete normal two-dimensional local domain
of equal characteristic $p>0$ with algebraically closed residue field,
and $U$ is its punctured spectrum. Fix the cover $U_0\to U$ and
SNC resolution $X_0$ of its normalization supplied by F10b. Write
$D_0=\bigcup_i C_i$, $B=(C_i\cdot C_j)$,
\[
 k_i=\max\{0,2g(C_i)-2+\nu_i\},\qquad
 t=(-B)^{-1}k,\qquad T_0=\max_i t_i,
 \tag{F13.1}
\]
where $\nu_i$ is the number of nodes on $C_i$. As in F10, choose
a further blowup if necessary so that $D_0$ is nonempty.
The logarithmic breaks below are computed over the completed DVRs
at the generic points of $D_0$. The complete proof passed an
independent root audit; F13b also passed a separate BCE geometry audit.

\begin{lemma}[F13a: modular monomial representations with a specified normal subgroup]
Let $\Gamma$ be a finite group, let $N\triangleleft\Gamma$ be
abelian with $\Gamma/N$ abelian, and let $\kappa$ be an algebraically
closed field of characteristic $\ell$. Every nonzero simple
$\kappa[\Gamma]$-module is induced from a one-dimensional character
of a subgroup $K$ containing $N$. In particular $K$ is normal in
$\Gamma$ and $\Gamma/K$ is abelian. No restriction on the
$\ell$-part of $|\Gamma|$ is imposed.
\end{lemma}
\begin{proof}
The $\ell$-primary subgroup $N_\ell$ is normal in $\Gamma$ and
acts trivially by the normal-primary invariants argument of F12a,
with $p$ replaced by $\ell$. We may quotient it out. Restriction
to the resulting prime-to-$\ell$ abelian group $N$ is a direct
sum of character spaces. Simplicity makes the characters occurring
in $V$ a single $\Gamma$-orbit. Choose one character $\chi$, let
$I$ be its stabilizer, and put $W=V_\chi$. The character-space
sum gives $V=\operatorname{Ind}_I^\Gamma W$. Moreover $W$ is
simple as an $I$-module: a proper nonzero submodule would induce
a proper nonzero submodule of $V$.

Put $A_1=I/N$. The pairing
\[
 b:A_1\times A_1\longrightarrow\kappa^*,\qquad
 b(\bar x,\bar y)=\chi([x,y])
 \tag{F13.2}
\]
is well-defined, alternating and bimultiplicative. Indeed,
$[I,I]\subset N$, and the $I$-invariance of $\chi$ kills
$\chi([x,n])$ for $n\in N$; the commutator identities then
prove independence of lifts and bimultiplicativity. Its radical
contains the $\ell$-primary part of $A_1$, since $\kappa^*$ has
no nontrivial $\ell$-power torsion. Choose a maximal isotropic
subgroup $L\subset A_1$. It contains the radical and satisfies
$L=L^\perp$: an element of $L^\perp\setminus L$ could otherwise
be adjoined to produce a larger isotropic subgroup.

Let $K$ be the inverse image of $L$ in $I$. Commutators of
$K$ act on $W$ through the scalar $\chi([K,K])=1$, so the
finite family of matrices from $K$ commutes. Over $\kappa$ it
has a common eigenvector, defining a character
$\theta:K\to\kappa^*$ extending $\chi$. Frobenius reciprocity
gives a nonzero map
$\operatorname{Ind}_K^I\theta\to W$.
The stabilizer of $\theta$ in $I$ is $K$: for $x\in I$, equality
of $\theta$ and its conjugate means precisely
$b(\bar x,L)=1$, hence $\bar x\in L^\perp=L$.
Thus the restriction of $\operatorname{Ind}_K^I\theta$ to
$K$ is a direct sum of pairwise distinct character lines.
Every submodule is a sum of those lines, and $I$ acts transitively
on them. The induced module is therefore simple and the nonzero
map is an isomorphism. Induction in stages gives
$V=\operatorname{Ind}_K^\Gamma\theta$.
The preimage of $K$ before quotienting $N_\ell$ contains the
original $N$, as required. Finally every subgroup containing $N$
is normal in $\Gamma$, because $\Gamma/N$ is abelian.
\end{proof}

\begin{lemma}[F13b: the Swan bound is unchanged by a finite \'etale resolution cover]
Let $f:X_1\to X_0$ be any connected finite \'etale cover, and put
$U_1=f^{-1}(U_0)$ and $D_1=f^{-1}(D_0)$. If $C'_a$ lies over
$C_i$ with degree $d_a$, every finite $p$-primary character on
$U_1$ has Kato Swan conductor at $C'_a$ at most $t_i$.
In particular the bound is independent of $\deg f$.
\end{lemma}
\begin{proof}
The surface $X_1$ is regular and proper over $A_0$, the
normalization underlying $U_0$. Its Stein factor is the spectrum
of the normalization of $A_0$ in the connected cover $U_1$.
This normalization is finite, complete and local with residue
field $k$; $X_1$ is a resolution of it, and $D_1$ is its reduced
exceptional SNC divisor. These assertions can also be checked
from the finite \'etale morphism over $X_0$: away from $D_0$
it is exactly the specified finite \'etale cover, regularity
and the SNC condition are preserved, and the connected proper
closed fiber follows from proper henselian specialization.
Thus the proof of F10 applies to $X_1$.

Let $B'$ be the intersection matrix of the components $C'_a$.
The finite \'etale curve map $C'_a\to C_i$ gives
\[
 2g(C'_a)-2=d_a(2g(C_i)-2),\qquad \nu'_a=d_a\nu_i.
 \tag{F13.3}
\]
Consequently $k'_a=d_a k_i$. The projection formula gives, for
any rational vector $v=(v_i)$,
\[
 \left(B'f^*v\right)_a
 =\left(f^*\sum_i v_iC_i\right)\cdot C'_a
 =d_a(Bv)_i,
 \qquad (f^*v)_a=v_i.
 \tag{F13.4}
\]
Taking $v=t$ shows $(-B')f^*t=k'$, so
$(-B')^{-1}k'=f^*t$. Apply F10's refined-Swan inequality
$B's'\geq-k'$ and its nonnegative inverse-matrix argument:
$s'_a\leq(f^*t)_a=t_i$. This uses the same
Krishna--Majumder levelwise refined-Swan input as F10; it does
not assert a bound on the total order of the character.
\end{proof}

\begin{proposition}[F13: a uniform logarithmic conductor after one fixed cover; Proved]
For every prime $\ell\ne p$ and every finite locally constant
$\mathbf F_\ell$-sheaf $M$ on $U$ with metabelian monodromy,
the highest logarithmic break of $M|_{U_0}$ at $C_i$ is at most
$t_i$. The fixed cover and all the bounds are independent of
$\ell$, the rank and the finite metabelian monodromy group.
\end{proposition}
\begin{proof}
Let $\Gamma$ be the monodromy group of $M$, and let
$\Gamma_0$ be the image of $\pi_1(U_0)$. Put
$N=\Gamma_0\cap\Gamma'$. It is abelian and normal, and
$\Gamma_0/N$ is abelian. By F10b, the pullback of the original
abelian quotient $\Gamma^{\mathrm{ab}}$ extends over $X_0$.
Therefore the image of every boundary inertia group in
$\Gamma_0$ lies in this specified subgroup $N$.

Extend scalars to $\overline{\mathbf F}_\ell$ and take a
composition series of the restricted representation. For a
simple constituent $V$, F13a gives
$V=\operatorname{Ind}_K^{\Gamma_0}\theta$ with $K\supset N$.
The associated cover $U_1\to U_0$ is unramified at every
exceptional DVR, because all boundary inertia lies in $N$.
Purity on the regular surface $X_0$ extends its normalization
to a finite \'etale cover $f:X_1\to X_0$.
The character $\theta$ is a rank-one local system on $U_1$.
All its values lie in a finite extension of $\mathbf F_\ell$.
Split $\theta$ into its $p$-primary part and its part of order
prime to $p$. The latter is tame and has logarithmic break zero.
The former has Kato Swan conductor at $C'_a$ at most $t_i$
by F13b.

Kato--Saito, Theorem~1.3, identifies the rank-one Kato Swan
conductor with the Abbes--Saito logarithmic conductor over the
completed DVR, including imperfect residue fields
\cite[Theorem~1.3, p.~2]{KatoSaito2019}.
Identifying the finite cyclic character values with complex
roots of unity does not change its kernel or either conductor.
Thus the logarithmic break of $\theta$ at $C'_a$ is at most
$t_i$.

For precision, induction through this finite \'etale cover does
not multiply the break by its degree. At the completed DVR of
$C_i$, each local field extension belonging to a component of
$f$ is unramified. For any positive rational $r$, Abbes--Saito
base-change compatibility identifies its logarithmic group
$G^r_{\log}$ with the same subgroup of the original absolute
Galois group \cite[Proposition~3.15(3)]{AbbesSaito2002}.
For $r>t_i$, this subgroup acts trivially on each conjugate
character in the restricted induced representation, and trivially
on the cosets of every unramified local extension. It therefore
acts trivially on the induced representation. This proves the
claim for each simple constituent.

Finally the image of any positive logarithmic ramification group
is a finite $p$-group. Its representations in characteristic
$\ell\ne p$ are semisimple. Hence triviality on every constituent
implies triviality on the whole representation. Scalar extension
is faithful, so the same assertion holds for $M$. This last step
also shows that no semisimplicity of $M$ was assumed.
\end{proof}

\paragraph{The remaining complete-local Betti gate.}
F13 is a conductor estimate, not yet an $H^1$ estimate. The precise
additional assertion that would close the coefficient-$\ell\ne p$
case at derived length three is: for this fixed complete-local
pair $(X_0,D_0)$ and the fixed bounds $t_i$, there is a constant
$C$ such that
\[
 \dim_\Lambda H^1(U_0,L)\leq C\operatorname{rank}_\Lambda L
 \tag{F13.5}
\]
for all finite fields $\Lambda$ of characteristic different from
$p$ and all finite local systems $L$ whose logarithmic breaks
at $C_i$ are bounded by $t_i$. It would suffice to establish
(F13.5) just for the systems occurring in F13.
Inflation--restriction for the fixed $U_0/U$ and the complement
count of F11 would then give a uniform generator bound for all
quotients of derived length at most three. No such conclusion is
being made without (F13.5).

Hu--Teyssier prove the corresponding rank-linear estimate for a
proper smooth \emph{finite-type} $k$-scheme $\overline V$ and
$V=\overline V\setminus E$, uniformly in finite coefficient fields
\cite[Theorem~7.34, p.~41]{HuTeyssier2025Betti}. Their relative
Theorem~7.15 likewise requires finite-type schemes
(Definition~7.14 and Theorem~7.15, p.~33). The nearby-cycle
Theorem~4.7 requires a geometric henselian trait
(Definition~4.6 and Theorem~4.7, pp.~15--16), and its proof uses
that condition in the spreading-out step. Neither statement,
as printed, applies to the fixed complete-local pair above.
A comparison, approximation or local Betti theorem with constants
uniform for this pair remains necessary. This explicit scope
check prevents treating finite-type spreading out for one sheaf
as a uniform result for all ranks and all sheaves.


\subsection{Removing unramified order from the remaining rank-one problem}

The following reduction does not assert the complete-local Betti bound
(F13.5). It identifies a bounded-degree geometric problem sufficient for
that bound on the metabelian systems actually used in F13. The
complete reduction passed an independent root proof audit.

\begin{lemma}[F13c: bounded ramified order after an unramified twist]
Retain $X_0,D_0,U_0$ and $t_i$ from F13, and choose $e\geq0$
such that $p^e>\max_i t_i$. Let $f:X_1\to X_0$ be any connected
finite \'etale cover, with $U_1=f^{-1}(U_0)$. Every character
\[
 \chi\in H^1(U_1,\mathbf Q_p/\mathbf Z_p)
\]
can be written $\chi=\alpha+\psi|_{U_1}$, where
$p^e\alpha=0$ and $\psi\in H^1(X_1,\mathbf Q_p/\mathbf Z_p)$.
The same $e$ works for every such cover. The Swan conductor of
$\alpha$ at every exceptional component equals that of $\chi$.
\end{lemma}
\begin{proof}
First, $H^1(X_1,\mathbf Q_p/\mathbf Z_p)$ is divisible. Indeed,
proper henselian specialization identifies the finite \'etale cover
categories of $X_1$ and its reduced proper nodal closed fibre $D_1$
\cite[Lemma 58.9.1, tag 0A48]{StacksProperEtaleSpecialization}.
Thus their groups of finite continuous characters agree. To verify
divisibility without any assertion about the nonabelian fundamental
group, use the Artin--Schreier sequence on $D_1$. The group
$H^1(D_1,\mathcal O_{D_1})$ is a finite-dimensional $k$-vector space,
and $H^2(D_1,\mathcal O_{D_1})=0$. On this vector space, Frobenius
is additive and $p$-semilinear. The map $F-1$ is therefore a
homomorphism of vector algebraic groups with differential $-1$.
It is \'etale and its image is an open subgroup of a connected
vector group, hence is the whole group. Since $k$ is algebraically
closed, it is surjective on $k$-points. Artin--Schreier now gives
$H^2(D_1,\mathbf F_p)=0$.

The exact sequence
\[
 0\longrightarrow\mathbf F_p\longrightarrow
 \mathbf Q_p/\mathbf Z_p\xrightarrow{p}
 \mathbf Q_p/\mathbf Z_p\longrightarrow0
\]
shows that multiplication by $p$ is surjective on
$H^1(D_1,\mathbf Q_p/\mathbf Z_p)$. All of its elements are finite
characters, and the same is true on $X_1$. This proves the required
divisibility.

The F13b conductor bound applies to every rank-one character on
$U_1$, not only to the systems in the statement of F13.
By that bound and the Witt-restriction calculation (F10.6), the character
$p^e\chi$ has conductor zero along every exceptional component.
It extends to $X_1$ by purity. Choose an unramified character
$\psi$ whose $p^e$ multiple is this extension, using divisibility.
Then $\alpha=\chi-\psi|_{U_1}$ is killed by $p^e$. An unramified
twist does not change a positive logarithmic break and preserves
conductor zero; hence the Swan conductors are unchanged.
\end{proof}

\begin{proposition}[F13d: a bounded-degree sufficient criterion]
Keep the notation and $e$ of F13c. Suppose that there is a constant
$C$ with the following property. For every connected finite \'etale
cover $X_1\to X_0$ of degree $d$ and every connected cyclic cover
$V\to U_1$ of degree dividing $p^e$ corresponding to a character
whose Swan conductors satisfy the F13b bounds, the normalization
of the complete local surface underlying $X_1$ in $V$ admits an
SNC resolution $Y$ dominating $X_1$ such that
\[
 d\bigl(\pi_1^t(V;Y,D_Y)\bigr)\leq Cd.
 \tag{F13.6}
\]
Then there is a uniform rank-linear $H^1(U_0,-)$ bound for every
coefficient-$\ell\ne p$ system occurring in F13. Consequently
(F13.6), together with the already proved F12, would settle the
uniform finite-generation assertion for quotients of derived
length at most three.
\end{proposition}
\begin{proof}
First consider a rank-one character $\theta$ on $U_1$ with finite
values in a field of characteristic $\ell\ne p$. Split its finite
image into its $p$-primary and prime-to-$p$ parts. By F13c, after
a finite extension of the coefficient field if necessary, write
\[
                         \theta=\alpha\otimes\beta,
\]
where $\alpha$ has order dividing $p^e$ and $\beta$ is tame
along $D_1$. Here the $p$-primary factor removed from $\alpha$
extends over $X_1$, while every prime-to-$p$ rank-one character
is tame. Let $V\to U_1$ be the connected cyclic cover associated
to $\alpha$. Its degree is a power of $p$, hence invertible in
the coefficient field. Inflation--restriction gives an injection
\[
                H^1(U_1,\theta)\longrightarrow H^1(V,\beta|_V).
 \tag{F13.7}
\]

The pullback $\beta|_V$ is tame along $D_Y$. For its unramified
$p$-primary factor this follows by pulling back the cover of
$X_1$ along the proper morphism $Y\to X_1$; for its prime-to-$p$
factor it follows from its prime-to-$p$ image at each DVR.
Every degree-one cocycle for this rank-one system also factors
through the tame quotient. Indeed, the corresponding affine
representation has image in the semidirect product of the additive
coefficient group with the image of $\beta$. On wild inertia
$\beta$ is trivial, so the cocycle restricts to a continuous
homomorphism from a pro-$p$ group to an additive $\ell$-group,
which is zero. The affine cover is therefore tame. Thus (F13.6)
and values of a cocycle on generators give
\[
          \dim H^1(V,\beta|_V)\leq Cd.
 \tag{F13.8}
\]

Now let $M$ be a system occurring in F13. Extend scalars to a
finite field which splits its finite monodromy group. Each simple
constituent $W$ is, by F13a, induced from a rank-one character
$\theta$ on a finite \'etale resolution cover $X_1\to X_0$,
of degree $d=\operatorname{rank}W$. Shapiro's lemma and
(F13.7)--(F13.8) give
\[
 \dim H^1(U_0,W)=\dim H^1(U_1,\theta)
                   \leq C\operatorname{rank}W.
\]
The long exact sequence for a composition series adds these
bounds, proving $\dim H^1(U_0,M)\leq C\operatorname{rank}M$.
Finite extension of the coefficient field preserves these
dimensions; for a fixed finite local system this follows by
tensoring the continuous cochain complexes, or by their filtered
finite-quotient description. Finally apply the fixed-cover
inflation--restriction and complement count already explained
after (F13.5).
\end{proof}

\paragraph{Scope of the reduction.}
The exponent $e$ in F13c is proved uniform. The generator bound
(F13.6) is a separate, unresolved assertion. Bounded degree of a
ramified cover alone has not been shown here to bound the genera
and incidence graph of a resolution linearly in $d$. In particular,
this proposition neither substitutes a bounded-degree statement
for that missing geometric argument nor promotes F13.5 or the
derived-length-three case to Proved.


\subsection{A degree-linear bound for refined-Swan zeros}

The following estimate passed an independent root proof audit.

\begin{proposition}[F13e: zeros on the original exceptional curves]
Retain the fixed pair $X_0,D_0$, matrix $B$, and cutoff vector $t$
of F13. Let $f:X_1\to X_0$ be a connected finite \'etale cover
of degree $d$, and let $\chi$ be a finite $p$-primary character
on $U_1$. Write
\[
 S_\chi=\sum_a s_aC'_a
\]
for its Swan divisor. If $C'_a$ lies above $C_i$ with degree $d_a$
and $s_a>0$, let $Z_a$ be the zero divisor of its nonzero refined
Swan section
\[
 \omega_a\in H^0\!\left(C'_a,
 \Omega^1_{X_1}(\log D_1)|_{C'_a}
                  \otimes\mathcal O_{X_1}(S_\chi)|_{C'_a}\right).
\]
More precisely, $Z_a$ is defined by saturating the image of
$\mathcal O_{C'_a}\xrightarrow{\omega_a}$ in this rank-two
vector bundle. Then
\[
 0\leq\deg Z_a\leq d_a(-B_{ii})(t_i-s_a),\qquad
 \sum_{a:s_a>0}\deg Z_a
       \leq d\sum_i(-B_{ii})t_i.
 \tag{F13.9}
\]
In particular the total number of zeros, counted with their
vanishing multiplicities on the component normalizations, is
bounded linearly in the cover degree. If $s_a=t_i$, that section
has no zero on $C'_a$.
\end{proposition}
\begin{proof}
The section exists and is nonzero by the exact refined-Swan
input used in F10: Krishna--Majumder, Corollary~9.7 and
Theorem~9.4, with the regular $F$-finite hypotheses explicitly
checked there \cite{KrishnaMajumder2026}. Its saturated image
is a line subbundle of the displayed rank-two vector bundle.
The map from $\mathcal O_{C'_a}$ to that line bundle has divisor
$Z_a$, so the degree of the line bundle is $\deg Z_a$.

The residue sequence (F10.4) on $C'_a$ has line subbundle and
line quotient of degrees $2g(C'_a)-2+\nu'_a$ and $0$.
Every line subbundle of its middle term has degree at most
$k'_a=d_a k_i$. Indeed, its map to the line quotient is either
nonzero, in which case its degree is at most zero, or vanishes,
in which case it maps into the line subbundle and has degree at
most $2g(C'_a)-2+\nu'_a$. Tensoring by
$\mathcal O_{X_1}(S_\chi)|_{C'_a}$ therefore gives
\[
                 \deg Z_a\leq d_a k_i+S_\chi\cdot C'_a.
 \tag{F13.10}
\]

The components over a fixed $C_i$ are disjoint: their union is
the finite \'etale inverse image of the smooth curve $C_i$.
The projection formula consequently gives
\[
 (C'_a)^2=d_aB_{ii},\qquad
 \sum_{b:C'_b\mapsto C_j}(C'_b\cdot C'_a)=d_aB_{ij}
                     \quad(j\ne i).
\]
By F13b, $0\leq s_b\leq t_j$ for every component over $C_j$.
All intersection numbers between distinct components are
nonnegative. It follows that
\[
 S_\chi\cdot C'_a
 \leq d_a\left(s_aB_{ii}+\sum_{j\ne i}t_jB_{ij}\right).
\]
The identity $(-B)t=k$ transforms the right side of (F13.10)
into $d_a(-B_{ii})(t_i-s_a)$. This proves the first inequality.
For each $i$, the degrees of all components above $C_i$ sum to
$d$. Dropping the nonnegative $s_a$ terms and summing proves
the second inequality. The final assertion follows because an
effective divisor of degree zero on a proper smooth curve over
an algebraically closed field is zero.
\end{proof}

\paragraph{Exact limitation.}
This estimate concerns the refined-Swan sections on the fixed
resolution and its finite \'etale covers. It does not bound the
number of blowups in a clean-ramification resolution, nor the
cohomology created at those zeros. Those implications require
an additional argument, so (F13.5) and (F13.6) remain open.


\begin{corollary}[F13f: the cutoff below the characteristic]
Suppose the fixed cover and SNC resolution supplied by F10b can
be chosen so that the cutoff vector (F13.1) satisfies $t_i<p$
for every component. Then there is a uniform all-prime
rank-linear $H^1(U,M)$ bound for finite coefficient systems
$M$ with metabelian monodromy. In this subclass, all local
quotients of derived length at most three are uniformly finitely
generated.
\end{corollary}
\begin{proof}
Let $\ell\ne p$. After the fixed cover, each simple constituent
is induced as in F13a from a rank-one character on a finite
\'etale resolution cover. Its $p$-primary part has
Swan conductors bounded by $t_i<p$ under F13b; F13g then forces
all those conductors to be zero. It is unramified on that resolution by
purity. Its prime-to-$p$ part is tame. Induction through the
finite \'etale resolution cover preserves tameness. Thus all
simple constituents are tame. The image of wild inertia is a
finite $p$-group, whose representations in characteristic
$\ell\ne p$ are semisimple, so $M$ itself is tame after the
fixed cover.

Every degree-one cocycle for that pulled-back system factors
through the full tame quotient: its affine monodromy is trivial
on wild inertia, since the original representation is trivial
there and a pro-$p$ group has no nonzero homomorphism to an
additive $\ell$-group. F9 therefore bounds its $H^1$ by
$d_{X_0}\operatorname{rank}M$. Inflation--restriction for the
single fixed cover adds at most
$d(\operatorname{Gal}(U_0/U))\operatorname{rank}M$.

For coefficient $p$, F12 gives the required bound for each
simple constituent; the long exact cohomology sequences for a
composition series give the same rank-linear bound for $M$.
These uniform metabelian coefficient bounds supply precisely
the remaining hypothesis in the complement-count argument
following (F13.5), proving the final assertion.
\end{proof}

This numerical corollary does not assert that a cover satisfying
$t_i<p$ always exists. Nor does it extend the monomial argument
to representations of arbitrary solvable derived length.


\subsection{A residue obstruction and a stronger conductor gap}

The complete argument and exact primary-source formulas passed an
independent root audit.

\begin{lemma}[F13g: a negative Swan intersection forces a $p$-divisible coefficient]
Let $X\to\operatorname{Spec}A$ be an SNC resolution in the
equal-characteristic complete-local surface setting of F10, and
let $S=\sum_i s_iC_i$ be the Swan divisor of a finite
$p$-primary character on the punctured spectrum. If $s_i>0$
and $p\nmid s_i$, then $S\cdot C_i\geq0$. Consequently, if
$S\ne0$, there is a component $C_i$ such that
\[
                 S\cdot C_i<0\quad\hbox{and}\quad p\mid s_i.
 \tag{F13.11}
\]
In particular $S\ne0$ implies $\max_i s_i\geq p$.
\end{lemma}
\begin{proof}
We first check the local residue explicitly. At the completed
DVR of $C_i$, put $n=s_i$ and choose a uniformizer $z$.
The character has a Witt representative
$a=(a_{m-1},\ldots,a_0)$ in the filtration
$p^jv(a_j)\geq-n$. For $j\geq1$, the assumption $p\nmid n$
implies
\[
 p^jv(a_j)\geq-p^j\lfloor n/p^j\rfloor\geq-(n-1).
\]
The refined Swan conductor is represented by
\[
                  -\sum_{j=0}^{m-1}a_j^{p^j-1}\,da_j
 \tag{F13.12}
\]
modulo logarithmic pole order at most $n-1$; these coordinate
conventions and the formula are displayed in
\cite[Definition~2.1 and (2.5), pp.~8--9]{Yatagawa2017Cycle}.
All terms with $j\geq1$ have pole order at most $n-1$.
The remaining coordinate must have pole exactly $n$, since
otherwise the character would have Swan conductor at most
$n-1$. Write $a_0=uz^{-n}$ with $u$ a unit. The logarithmic
residue of $-da_0$ in pole degree $n$ is
$n\bar u$, which is nonzero.

The global refined-Swan section of F10 maps, under the residue
sequence (F10.4) tensored by $\mathcal O_X(S)|_{C_i}$, to a
section of $\mathcal O_X(S)|_{C_i}$. Its generic value is the
nonzero residue just calculated; this generic compatibility is
Krishna--Majumder, Corollary~9.6, and the global map is
Corollary~9.7. Their Lemma~9.11 identifies this local refined
conductor with Kato's refined Swan conductor used in (F13.12)
\cite[Corollaries~9.6--9.7, pp.~55--56, and Lemma~9.11,
p.~57]{KrishnaMajumder2026}.
Thus this line bundle on the proper smooth curve $C_i$ has a
nonzero section. Its degree $S\cdot C_i$ is nonnegative.

Suppose $S\ne0$. Negative definiteness gives
\[
                  0>S^2=\sum_i s_i(S\cdot C_i).
\]
There is therefore an $i$ with $s_i>0$ and $S\cdot C_i<0$.
The first assertion forces $p\mid s_i$, proving (F13.11) and
the final conclusion.
\end{proof}

\begin{corollary}[F13h: improved unramified exponent]
In F13b and F13c, one may choose any integer $e\geq0$ with
\[
                           p^{e+1}>T_0=\max_i t_i.
 \tag{F13.13}
\]
For every rank-one $p$-primary character on every finite \'etale
resolution cover, its $p^e$ multiple then extends unramified
over that cover's resolution. Hence F13c's bounded-order twist
decomposition holds with this improved $e$.
Likewise the $p$-primary inertia kernel in F10 is killed by
$p^e$ whenever $p^{e+1}>\max_i t_i$ for the resolution used
there.
\end{corollary}
\begin{proof}
The Witt restriction step (F10.6) gives every coefficient of the
Swan divisor of $p^e\chi$ at most
$\lfloor T_0/p^e\rfloor<p$. Apply F13g on the resolution of
the complete local normalization underlying the cover. Its
exceptional matrix is negative definite, so that Swan divisor
must be zero. Purity extends the character to the resolution.
The divisible unramified character group from F13c gives the
same twist decomposition with exponent $p^e$.
For F10's abelian inertia kernel, all characters of its
$p$-primary part are killed after multiplying by $p^e$;
character separation gives the last assertion exactly as in
F10.
\end{proof}

\paragraph{Numerical consequence.}
F13f's sufficient hypothesis can therefore be weakened from
$t_i<1$ to $t_i<p$ for every $i$. Under the latter hypothesis,
F13g forces every rank-one inducing character's $p$-primary
part to be unramified; the rest of the proof of F13f is
unchanged. Neither inequality is asserted to hold for an
arbitrary normal surface after the fixed cover.

\section{Canonical points on polarization torsors}
\subsection{Universal sections of polarization torsors}
\paragraph{Conventions.}
Throughout this subsection $2$ is invertible on the base.  We use the
fppf relative Picard sheaf, $B=A^\vee=\operatorname{Pic}^0_{A/S}$,
the usual Poincaré biduality $A=B^\vee$, and
\[
 \lambda_L(a)=[t_a^*L\otimes L^{-1}].
\]
Thus an ample $L$ gives a positive polarization $\lambda_L:A\to B$.
The torsor $T_\delta=\operatorname{Pic}^{\delta}_{A/S}$ is a torsor under
$B$ by tensor product.  A section here means a section of this sheaf
over the moduli stack, and hence a construction invariant under
isomorphisms and commuting with arbitrary base change.  It is not the
assertion that each individual algebraically closed fibre has a point.

\paragraph{G1 (Proved: the complete existence criterion).}
Let $D=(d_1\mid\cdots\mid d_n)$ be a polarization type, and put
$r=\#\{i: d_i\text{ is odd}\}$.  The universal $T_\delta$ over the
moduli stack of type $D$ over $\mathbb Z[1/2]$ has a section if and only
if $r\leq1$.  The construction in the affirmative cases is canonical.
It requires no inversion of the odd divisors of $\prod_i d_i$.
More intrinsically, its hypothesis is that the alternating pairing
on $A[2]$ induced by $\lambda$ has rank at most two.  This formulation
also specifies the assertion on integral models of the polarization
moduli problem at primes dividing its degree: it depends only on the
locally constant type of the étale group $A[2]$, not on terminology for
the non-étale part of the polarization kernel.

We first recall precisely the input about symmetric line bundles.
The symmetric classes in $T_\delta$ form an étale $B[2]$-torsor.  The
normalized symmetry of a line bundle identifies this torsor with the
torsor of quadratic refinements of
\[
 b_\lambda(x,y)=e_2(x,\lambda(y)),\qquad x,y\in A[2].
\]
We write $\mu_2$ additively as $\mathbb F_2$, so the refinement identity
is $q(x+y)=q(x)+q(y)+b_\lambda(x,y)$.  This is the scheme-valued
statement of Poonen--Rains, Proposition 3.6 and diagram (3.7),
pp.~1312--1313 \cite{DGJ_PoonenRains2011}; the exact sequence preceding
it identifies its fibre over $\lambda$ with the stated torsor.
In particular this is an identification of sheaves over arbitrary
bases, rather than a bijection only on geometric points.

Let $R$ be the radical of $b_\lambda$.  It is a finite étale subgroup
of $A[2]$, because the rank of this map of finite étale groups is locally
constant.  If $r=0$, take $q=0$.  If $r=1$, define
\[
 q(x)=\begin{cases}0&x\in R,\\1&x\notin R.\end{cases}
\]
This is a morphism of finite étale schemes: both loci are open and
closed.  It is a quadratic refinement, as can be checked on the
two-dimensional quotient $A[2]/R$.  In that quotient distinct nonzero
vectors have pairing one, and their sum is the third nonzero vector.
The other cases in the refinement identity are immediate.  The
definition uses no choices and therefore gives the asserted section
of the symmetric Picard torsor, and of $T_\delta$.

For necessity we may pull a proposed section back to characteristic
zero.  The automorphism $[-1]$ of every polarized abelian scheme
forces the proposed section to be symmetric.  Over the Siegel space
of marked complex polarized abelian varieties, the resulting section
of a finite étale torsor is constant.  Descent would therefore give
a quadratic refinement fixed by the integral group preserving the
alternating lattice with basis $e_i,f_i$ and
$E(e_i,f_j)=d_i\delta_{ij}$.  This description of the complex moduli
stack as the quotient of Siegel space by this group is recorded, with
the definition of the group, in \cite[\S\S5.3--5.4, pp.~30--33]{DGJ_Laga2025}.

For every integral lattice vector $v$, the integral transvection
$x\mapsto x+E(x,v)v$ preserves $E$ and has integral inverse
$x\mapsto x-E(x,v)v$.  In the subspace $W$ of $A[2]$ spanned by the
pairs with odd $d_i$, its reductions include the transvection along
every $v\in W\setminus\{0\}$.  If $q$ were invariant, choose $x\in W$
with $b_\lambda(x,v)=1$.  Invariance and the quadratic identity give
$q(x)=q(x+v)=q(x)+q(v)+1$, hence $q(v)=1$ for every nonzero $v\in W$.
When $r\geq2$, choose independent orthogonal vectors $v,w\in W$.
Then $q(v+w)=q(v)+q(w)=0$, a contradiction.  Thus no section exists
in these cases, already over $\mathbb C$.

The determinant comparisons below use the following uniqueness fact.  There is
no nonzero natural section of the universal $A[2]$ or $A^\vee[2]$ over
the moduli problem of polarized abelian schemes of given dimension
and polarization degree, or over its locus where the rank of
$b_\lambda$ is fixed.  In characteristic zero, the integral lattice
group contains an $\operatorname{SL}_2(\mathbb Z)$ on each pair
$e_i,f_i$, even when $d_i$ is even.  Its action on the corresponding
two-dimensional summand modulo two has no fixed nonzero vector.
There are therefore no invariant vectors in the whole lattice modulo
two, or in its dual.  This proves the assertion on each complex
component.  It implies the characteristic-zero assertion over any
field by spreading an object to a finitely generated field and
embedding that field in $\mathbb C$.
For a polarized abelian variety in positive characteristic we use
the polarized lifting theorem of Mumford--Norman--Oort, stated as
Theorem (8.3).1 in \cite[p.~29]{DGJ_Oort2014}: it lifts over an integral
mixed-characteristic ring.  The degree and the rank on two-torsion
are constant on this ring.  Two sections of a finite étale separated
scheme that agree generically agree everywhere on this connected
base, since their equalizer is open and closed.  The asserted section
is thus zero also on the special fibre.  Finally equality of sections
of a finite étale scheme can be tested on all geometric points, so
this proves the assertion over arbitrary, including nonreduced,
bases.  Thus two constructions of sections of $T_\delta$ which are
natural on this entire degree locus (or the indicated rank locus)
are equal: their difference is such a two-torsion section because
each is fixed by $[-1]^*$.  This is the form of uniqueness needed
below; no claim about extra vertical components in an unspecified
integral model of a fixed-type moduli stack is used.

The finite computation accompanying this paper enumerates refinements
and transvection invariants in ranks zero, two, four and six, with
radicals of dimensions zero or two.  It confirms the small cases;
the argument above proves every rank and does not rely on extrapolation.

\subsection{The determinant construction and its signs}
\paragraph{G2a (Proved: R9's boxed relation, with typed signs).}
Let $\mathcal P$ be the normalized Poincaré bundle on $A\times_S B$.
For a relatively ample line bundle $L$ of class $\delta$, write
\[
 \Phi(L)=R\operatorname{pr}_{B*}
       (\mathcal P\otimes\operatorname{pr}_A^*L),\qquad
 F_\delta(L)=\det\Phi(L)^{-1},\qquad \chi=\chi(L).
\]
These are relative constructions.  The complex is perfect; in the
ample case it is a vector bundle of rank $\chi$.  Tensoring $L$ by a
bundle from $S$ changes its determinant by a bundle from $S$ only,
so $F_\delta$ defines a morphism of relative Picard sheaves.
The theorem of the square gives, up to bundles from the base,
\[
 (1,\lambda_L)^*\mathcal P
   =m^*L\otimes\operatorname{pr}_1^*L^{-1}
          \otimes\operatorname{pr}_2^*L^{-1}.
\]
Changing variables $(x,a)\mapsto(x+a,a)$ therefore gives
$\lambda_L^*\Phi(L)=\pi_*L\otimes L^{-1}$ up to the same harmless
base factors.  Taking determinants shows
$\lambda_L^*F_\delta(L)=L^\chi$ in the relative Picard sheaf.
Consequently, if $\delta'$ is the class of $F_\delta(L)$, then
\[
 \lambda_\delta\lambda_{\delta'}=[\chi]_B,
 \qquad
 \lambda_{\delta'}\lambda_\delta=[\chi]_A.
\]
Indeed the formula for a pulled-back line bundle first gives
$\lambda_\delta\lambda_{\delta'}\lambda_\delta
=\chi\lambda_\delta$; the isogeny $\lambda_\delta$ is faithfully flat
and can be cancelled on the right.  Duality gives the other equality.
Thus $\delta'$ is positive, whereas the determinant itself has class
$\varphi(\delta)=-\delta'$.  This proves precisely the negative sign
in the box of 59:10.

The resulting torsor morphism and its equivariance are
\[
 F_\delta:T_\delta\longrightarrow T_{\delta'},\qquad
 F_\delta(L\otimes\beta)=F_\delta(L)\otimes\lambda_{\delta'}(\beta),
 \quad \beta\in B.
\]
This follows from
$\Phi(L\otimes\beta)=t_\beta^*\Phi(L)$.
In particular, with our stated translation convention the linear
part is $+\lambda_{\delta'}$.  The page's use of $T_{-\delta}$ can be
typed without an ambiguous suppression of signs: the canonical
$B$-equivariant isomorphism
\[
 T_\delta\longrightarrow T_{-\delta},\qquad
 L\longmapsto([-1]^*L)^{-1}
\]
has linear part $+1$, whereas $L\mapsto L^{-1}$ has linear part $-1$.
Using the first gives the displayed version of (*) on 59:12.

The equalities also give
$\chi(\delta')=\chi^{n-1}$ and
$\delta''=\chi^{n-2}\delta$ for $n\geq2$, as equalities of
Néron--Severi sections.  To see the first, take degrees of the
isogenies and use positivity to select the positive square root.
The second follows after applying the same relation to $B$ and
using the injection of the Néron--Severi sheaf into symmetric
homomorphisms.  In dimension one the double transform is the
canonical degree-one polarization, so the expression with a negative
exponent is interpreted in this sense, not as tensoring a torsor
by a nonintegral integer.

\paragraph{G2b (Proved: the ``à vérifier'' comparisons).}
For every Néron--Severi section $\eta$, the norm
\[
 c_\eta=L\otimes[-1]^*L\in T_{2\eta}
\]
is independent of the local choice of $L\in T_\eta$: a change by
$\beta\in A^\vee$ cancels with $[-1]^*\beta=\beta^{-1}$.
It is a canonical section, compatible with arbitrary base change.
The determinant commutes with $[-1]^*$, as is seen by pulling back
the Poincaré bundle and changing variables on $A$.  It therefore
sends symmetric classes to symmetric classes.  Choose such a class
étale locally, which is possible because $2$ is invertible.  Then
$2L=c_\delta$ and $2F_\delta(L)=c_{\delta'}$.  Thus the isomorphism
of squared torsors induced by $F_\delta$ carries the canonical
section to the canonical section.  This proves the margin question
at 59:12, including its formulation as an isomorphism of the
corresponding two-torsion torsors.

We spell out why the additional determinant trivializations on
59:13--14 have the same normalization, rather than inferring it from
an equality of torsor classes.  Determinants of perfect complexes,
tensor powers, pullbacks, and their compositions are all natural
constructions on the full polarized moduli problem.  The comparison
of $\delta$ with $-\delta$ used there produces a section of
$T_{2\delta'}$; relative Serre duality supplies the comparison also
for the antiample class.  Namely, up to a base line,
\[
 F_{-\delta}(L^{-1})
       =\bigl([-1]^*F_\delta(L)\bigr)^{(-1)^{n-1}}.
\]
Both this comparison section and $c_{\delta'}$ are invariant under
$[-1]$; their difference is consequently a natural section of
$A[2]$.  The uniqueness argument in G1 makes it zero.  This proves
the margin question on the canonical square at 59:13.

For the double transform, put $a=\chi^{n-2}$ with $n\geq2$.
The composite $F_{\delta'}F_\delta$ has target $T_{a\delta}$ and
linear part $[a\chi]_B$.  Comparing it with the tensor-power map
$L\mapsto L^a$ therefore supplies a natural trivialization of
$T_{a(\chi-1)\delta}$, with reversal of the entire sign if the
opposite comparison is chosen.  The preceding square construction
also trivializes $T_{2a\delta}$.  If $\chi$ is even, $\chi-1$ is odd;
combining these trivializations with integral Bézout coefficients
gives a natural section of $T_{a\delta}$.  It is invariant under
$[-1]$, as is every step of the construction.  When $n\geq3$, $a$
is even and its comparison with $c_{(a/2)\delta}$ again differs by
a universal two-torsion section; hence they agree.  When $n=2$,
$a=1$ and $\chi$ even forces at most one odd elementary divisor.
The resulting section therefore agrees with the section of G1,
by the same uniqueness lemma on the full even-degree surface locus.
Changing Bézout coefficients gives the same result by this
uniqueness.  These statements establish the determinant
trivializations as actual sections, not merely vanishing classes
in $H^1(S,B)$.

There are two literal slips to retain in the record.  The word
``ppcm'' on 59:13--14 must be ``pgcd'' for the Bézout deduction:
for $n\geq2$ every $m^n-m^{n-1}$ is even and the value at $m=-1$
is $\pm2$.  Also the assertion on 59:13 that an even $\chi(\delta)$
implies divisibility of $\delta$ by two is false.  On $E\times E$,
the exterior product of degree-one and degree-two line bundles has
type $(1,2)$ and Euler characteristic two, but its restriction to
the first factor has odd degree, so its class is not twice another
class.  G1 supplies exactly the needed correction in dimension two.
The $B$ subscripts for the torsors of class $\delta$ on 59:14 must
be $A$; the class $\delta$ was defined on $A$.

\subsection{The Albanese sign and the canonical theta class}
\paragraph{G3 (Proved).}
At 59:23 the selected convention is
$\delta(x,y)=[N_yN_x^{-1}]$, followed by
``Nous prendrons la première convention''.  For a smooth proper
geometrically connected curve $C/S$, let $J=\operatorname{Pic}^0_{C/S}$
and let $\lambda_\Theta:J\to J^\vee$ be its positive theta
polarization with our translation convention.  Then the map of
59:24 is
\[
 c:J\longrightarrow J^\vee,\qquad c=-\lambda_\Theta.
\]
Thus the correction on that page, ``Non, c'est sans doute la
symétrique de la $\Theta$-polarisation'', has the right sign.
To verify it, choose a section $x$ locally and let
$a:C\to J$, $y\mapsto[y-x]$.  The definition by the Poincaré
bundle identifies $c$ with the inverse of
$a^*:J^\vee\to J$.  Jacobian autoduality gives
$a^*\lambda_\Theta=-1$; see Milne, Lemma 6.9 and
Remark 6.10(c), p.~23 of the revised notes
\cite{DGJ_MilneJac}.  We inspected the rendered page because text
extraction drops the minus sign in this lemma.  The equality over
families follows from its fibrewise equality and rigidity of
homomorphisms of abelian schemes.  This includes nonreduced bases;
the relative Hom functor is unramified.  In particular transporting
$\operatorname{Alb}^1_{C/S}$ from $J^\vee$ to $J$ by
$\lambda_\Theta^{-1}$ gives $\operatorname{Pic}^{-1}_{C/S}$, in
agreement with the sign warning on 59:9.

For a principally polarized abelian scheme $(A,\delta)$ there is a
canonical principally polarized torsor $(P,\Theta)$ as described
on 59:26--27.  One construction is $P=T_\delta$, with its structural
group transported from $A^\vee$ to $A$ by $\lambda_\delta$.
The universal line bundle has a rank-one pushforward, and its
evaluation at the zero section defines the theta divisor on $P$.
For any other such torsor, sending a point $p$ to the line bundle
pulled back from $\Theta$ by $a\mapsto a+p$ gives the required
isomorphism with $T_\delta$; its linear part is
$\lambda_\delta$, hence it is an isomorphism.  Its uniqueness
follows since a translation preserving $\Theta$ lies in the kernel
of this principal polarization.  These arguments are fppf local
and hence descend to arbitrary bases.

The involution $L\mapsto[-1]^*L$ of $T_\delta$ preserves this
theta divisor: it preserves the condition that the unique section
vanish at zero.  It has linear part $-1$ on its structural group.
By the just-proved uniqueness, it is precisely the theta involution
used on 59:28 to define $K\in P^{\otimes2}$.  Consequently
\[
 K=L+[-1]^*L=c_\delta\in T_{2\delta}.
\]
This proves that page's normalization question directly.

For a curve the theta torsor is
$\operatorname{Pic}^{g-1}_{C/S}$ with its effective theta divisor.
Relative Serre duality identifies its theta involution as
\[
 L\longmapsto\omega_{C/S}\otimes L^{-1}.
\]
The equality holds scheme-theoretically, since the theta divisor
is the determinant divisor of cohomology and Serre duality
identifies the corresponding determinants.  Uniqueness of the
theta involution now gives
$K=[\omega_{C/S}]$ in $\operatorname{Pic}^{2g-2}_{C/S}$.
No assertion about generators of the Picard group of the moduli
stack of curves is needed.  The opposite structural-group
identifications mentioned on 59:27 are related by this very
involution; the minus sign of $c$ must not be erased when writing
an equivariant map.

\section{Higher extensions of the pro-generalized Jacobian}\label{sec:H}
\subsection{The exact category and the positive result}
Throughout this section $k$ is algebraically closed of characteristic $p>0$,
and $\mathcal P$ is Serre's abelian category of commutative
pro-quasi-algebraic groups. Algebraic groups are viewed through their
quasi-algebraic perfection. We do not identify this category with the
category of all sheaves on a perfect \emph{big} \'etale site.
For a smooth curve with function field $K$, use its smooth projective
completion $C$. Write $J_{\mathfrak n}$ for the connected degree-zero
generalized Jacobian of a nonzero effective modulus $\mathfrak n$.

\begin{theorem}[H1: higher Ext vanishing; Proved]\label{h:vanishing}
The object $J_\infty=\varprojlim_{\mathfrak n}J_{\mathfrak n}$ has
projective dimension at most one in $\mathcal P$. Thus
\[
 \Ext^i_{\mathcal P}(J_\infty,G)=0\qquad(i\geq2)
\]
for every $G\in\mathcal P$, including all the algebraic targets in Part H
and characteristics two and three.
\end{theorem}
\begin{proof}
For $\mathfrak n=\sum_{x\in S}n_x[x]$, put
\[
 R_{\mathfrak n}=\prod_{x\in S}
  (\OO_{C,x}/\mathfrak m_x^{n_x})^\times,
 \qquad L_{\mathfrak n}=R_{\mathfrak n}/\GG_m,
\]
with diagonal scalar units. The generalized-Jacobian sequence is
\[
 0\longrightarrow L_{\mathfrak n}\longrightarrow J_{\mathfrak n}
 \longrightarrow J=\Pic^0(C)\longrightarrow0.
\]
Its algebraic-group assertion is \cite[Chapter I, \S1, Theorem 3,
p.~12]{SerreGACC}. Equivalently, rigidifications of a degree-zero line
bundle along $\mathfrak n$ give this sequence; rigidifications of the
trivial bundle are $R_{\mathfrak n}$ modulo its scalar automorphisms.
Forgetting rigidifications gives compatible transition morphisms.
Inverse limits are exact in $\mathcal P$
\cite[\S2.5, Proposition 10, p.~19]{Serre1960}, so
\[
 0\longrightarrow L\longrightarrow J_\infty\longrightarrow J
 \longrightarrow0,\qquad L=\varprojlim L_{\mathfrak n}.
\]
Choose $x_0\in C(k)$ and restrict to the cofinal moduli containing it.
Residue at $x_0$ splits the diagonal scalar subgroup, compatibly with
all truncations. Consequently
\[
 L\simeq
 \left(\prod_{x\ne x_0}\GG_m\right)
 \times\prod_{x\in C(k)}U_x,
 \qquad U_x=1+t_x k[[t_x]].
\]
This equality is an equality of proalgebraic groups: each finite set of
points and each finite collection of jet orders occurs in a modulus.
The parameters $t_x$ are choices, irrelevant to the vanishing.
The big-Witt power-series isomorphism and its $p$-typical decomposition
\cite[Propositions 1.14 and 1.10, pp.~11 and 8]{HesselholtWitt} give
\[
 U_x\simeq\prod_{r\geq1,\ p\nmid r}W,
 \qquad W=\varprojlim_n W_n.
\]
Here $W$ is the infinite Witt group, not its universal cover.

Serre proves that $\GG_m$ has projective dimension one (\S7.5,
pp.~50--51), that $W$ has projective dimension one (\S8.6,
Corollary 3, p.~58), and that an abelian variety has projective dimension
at most one (\S9.2, Proposition 4, p.~61) \cite{Serre1960}.
For a product $T=\prod_i T_i$ of copies of $\GG_m$ or $W$, and a
\emph{quasi-algebraic} target $B$, Proposition 7 of \S3.4, p.~27, gives
\[
 \Ext^2_{\mathcal P}(T,B)
 =\varinjlim_{F\text{ finite}}
  \Ext^2_{\mathcal P}\left(\prod_{i\in F}T_i,B\right)=0.
\]
The proposition passes an inverse limit in the first variable to a
direct limit of Ext groups; its target hypothesis is essential and is
satisfied here. Finite products are finite direct sums. In particular,
the displayed vanishing holds for every elementary target, including
the finite groups $\ZZ/\ell$. Proposition 9 on the same page then
implies $\operatorname{pd}_{\mathcal P}T\leq1$ for \emph{all} targets.
Thus $\operatorname{pd}L\leq1$. The long exact Ext sequence for
$0\to L\to J_\infty\to J\to0$ now gives the result in every degree
at least two.
\end{proof}

\begin{proposition}[H2: the additive target; Proved]\label{h:additive}
For every $P\in\mathcal P$ and every connected quasi-algebraic $B$,
$\Ext^i_{\mathcal P}(P,B)=0$ for $i\geq2$.
In particular, $G=\GG_a$ is not the unresolved target suggested in
Part H's wording.
\end{proposition}
\begin{proof}
For an elementary source $\GG_a$, the assertion in degree two is
\cite[\S8.6, Corollary 5, p.~58]{Serre1960}. The other elementary
sources---finite groups, $\GG_m$, and abelian varieties---have
projective dimension at most one; for finite groups see \S4.4,
Corollary 2, pp.~34--35. Every quasi-algebraic group has a finite
composition series with these quotients (\S1.3, Corollary to
Proposition 7). The long exact sequence proves the assertion for a
quasi-algebraic source, and \S3.4, Proposition 7 passes to an arbitrary
proalgebraic source. Degrees at least three vanish by \S10.1,
Theorem 1. No inverse limit in the target has been used.
\end{proof}

\begin{proposition}[H3: augmentation cannot be omitted; Refuted comparison]
\label{h:augmentation}
The equality $\Ext^*(J_\infty,G)=H^*(K,G)$ is false in degree zero
for the connected $J_\infty$ above.
\end{proposition}
\begin{proof}
For $G=\ZZ/\ell$, a morphism $J_\infty\to G$ factors through a
connected $J_{\mathfrak n}$ and hence is zero. But
$H^0(K,\ZZ/\ell)=\ZZ/\ell$. On 43:80 the earlier object
$J'_{X/k}$ is explicitly an extension of $\ZZ$ by a connected group.
The boxed passage to $K$ must retain this augmentation; it cannot be
read as a statement about the connected object alone. The infinite
discrete group $\ZZ$ is not an object of Serre's $\mathcal P$.
Statement~\ref{h:vanishing} therefore settles the higher-Ext question
in the specified Serre category, without proving an equivalence of
Ext theories or the augmented all-degree comparison on another site.
\end{proof}

\subsection{The local-Jacobian lead and its topology}
\begin{proposition}[H4: an \'etale local-cohomology obstruction; Proved]
\label{h:local}
Let $A=k[[s,t]]$ and $\mathfrak m=(s,t)$. Then
\[
 H^2_{\mathfrak m,\mathrm{et}}(A,\ZZ/p)=0,
 \qquad
 H^3_{\mathfrak m,\mathrm{et}}(A,\ZZ/p)
 \simeq\bigoplus_{\substack{i,j>0\\p\nmid\gcd(i,j)}}(k,+)
\]
as $\mathbf F_p$-vector spaces. All other degrees vanish.
Thus the concentration used in SA19-H4 cannot be transferred from
Zariski cohomology to \'etale cohomology.
\end{proposition}
\begin{proof}
Quasi-coherent \'etale comparison and the \v Cech localization model
give $H^i_{\mathfrak m}(A,\GG_a)=0$ for $i\ne2$ and
\[
 E=H^2_{\mathfrak m}(A,\GG_a)
   =A_{st}/(A_s+A_t)
   =\bigoplus_{i,j>0}k e_{i,j},\qquad e_{i,j}=s^{-i}t^{-j}.
\]
All sums have finite support. Frobenius acts by
$F(ae_{i,j})=a^p e_{pi,pj}$. Partition the pairs into orbits
$p^r v$, $r\geq0$, where the two coordinates of $v$ are not both
divisible by $p$. On each orbit $F-1$ is injective: the last nonzero
coefficient of a proposed kernel element produces an uncancelled
coefficient one step later. Its cokernel is $(k,+)$, by the quotient
map
\[
 \sum_r a_r e_{p^r v}\longmapsto\sum_r a_r^{p^{-r}}.
\]
The map kills $F-1$, is surjective, and its kernel is exactly the
finite telescoping relations $F(ae_{p^r v})-ae_{p^r v}$.
The Artin--Schreier exact sequence now gives the asserted result.
The opened inputs are Stacks, Theorem 59.22.4, Section 59.63, and
Lemma 51.2.1 \cite{StacksEtaleQC,StacksAS,StacksLocalCech}.
\end{proof}

The ordinary functor $G\mapsto H^2_{\mathfrak m}(A,G)$ does not even
descend to $\mathcal P$ with its usual action on $\GG_a$:
Frobenius is an isomorphism in $\mathcal P$ and is not surjective on
$E$. If a projective representing object with this value and the
usual Artin--Schreier action were asserted, its Hom functor would make
$F-1:E\to E$ surjective, contradicting $e_{1,1}\notin(F-1)E$.
Even replacing $E$ by $\varinjlim_F E$ leaves a nonzero cokernel of
$F-1$: the orbits become indexed by $\ZZ$, still with finite support,
and the same telescoping map applies. This observation is conditional
on that particular attempted perfection, not a perfect-site comparison.

The reading explicitly describes 43:70--74 as Zariski theory.
Artin--Schreier is not exact on the Zariski site. The calculation
therefore refutes the proposed transfer and projectivity inference,
not an unspecified Zariski representability theorem. Completion from
$k[s,t]_{(s,t)}$ carries any alleged Artin--Schreier preimage to one
over $A$, so the obstruction also applies to that essentially finite
type example. The finite checks in \texttt{finite\_checks.py} retain
the extra target exponent and verify the obstruction for $p=2,3$;
the finite-support proof above establishes it for every $p$.

\begin{corollary}[H5: the perfect \'etale category; Proved]\label{h:perfect}
On the site $S=(\operatorname{Perf}/k)_{\mathrm{et}}$, for the
connected $J_\infty$ above and every quasi-algebraic target $G$,
\[
 \Ext^n_{\operatorname{Ab}(S)}(J_\infty,G)
   \simeq\Ext^n_{\mathcal P}(J_\infty,G)=0\qquad(n\geq2).
\]
Thus the higher vanishing in Part H holds in both stated categories,
with algebraic groups interpreted through perfection in those categories.
\end{corollary}
\begin{proof}
Suzuki's Proposition 3.8.1, p.~47 of the opened author version,
proves comparison for affine pro-quasi-algebraic sources and affine
quasi-algebraic targets. His Remark 3.8.4, pp.~55--56, extends it to
sources that are extensions of a perfected abelian variety by an
affine proalgebraic group, and arbitrary quasi-algebraic targets
\cite{SuzukiExt2022}. The sequence $0\to L\to J_\infty\to J\to0$
constructed above places $J_\infty$ in exactly that category.
Apply the comparison, then the proved projective-dimension bound.
\end{proof}

The comparison is not a formal consequence of a fully faithful exact
embedding: Suzuki explicitly states that proalgebraic realization on
$S$ is only left exact (p.~53). His proof passes through the perfect
pro-fppf site and the inverse-source limit theorem, Lemma 3.8.2,
pp.~48--53. No interchange of an arbitrary inverse limit in the target
is being asserted here.

For any actual exact sequence of sheaves
$0\to J_\infty\to J_{\mathrm{aug}}\to\ZZ\to0$, the same
higher vanishing follows. Indeed,
$\Ext^n_S(\ZZ,G)=H^n(S,G)=0$ for $n>0$: an \'etale cover of
the terminal object $\Spec k$ splits, so global sections are exact.
Apply the long exact sequence in the first variable. This does not
identify an unspecified augmented pro-object with ordinary
$H^*(K,G)$; the geometric argument and coefficient perfection still
have to be specified, especially for rational points in degree zero.

\begin{proposition}[H4b: a local Jacobian is not projective; Refuted claim]
Let $A_0=k[s,t]_{(s,t)}$ and let $x$ be its closed point.
The ordinary local Jacobian $J_x$ in
$\operatorname{Pro}(\operatorname{AbSch}_k)$, corepresenting
$G\mapsto H^2_{x,\mathrm{fppf}}(A_0,G)$ on commutative finite-type
$k$-group schemes, is not projective. In fact
\[
 \Ext^1_{\mathrm{naive}}(J_x,\ZZ/p)
 \simeq\bigoplus_{\substack{i,j>0\\p\nmid\gcd(i,j)}}(k,+).
\]
\end{proposition}
\begin{proof}
Existence and the stated universal property follow from To\"en's
construction and Definition 3.5, pp.~17--18, applied to the
codimension-two point of the smooth surface $\mathbb A^2_k$
\cite{Toen2026}. Let $f:J_x\to\GG_a$ correspond to
$e=s^{-1}t^{-1}\in E_0=H^2_x(A_0,\GG_a)$, and form the
pullback of the Artin--Schreier epimorphism along $f$:
\[
 Q=J_x\mathbin{\times}_{\GG_a,f,F-1}\GG_a,
 \qquad 0\longrightarrow\ZZ/p\longrightarrow Q
              \longrightarrow J_x\longrightarrow0.
\]
A splitting is exactly a morphism $g:J_x\to\GG_a$ with
$(F-1)g=f$, hence an element $h\in E_0$ with $(F-1)h=e$.
Completion carries this equation to the impossible equation in
$H^2_{(s,t)}(k[[s,t]])$ computed above. Thus the extension is
nonsplit, proving nonprojectivity in the stated category.

Local cohomology commutes with flat completion; since $E_0$ is
$\mathfrak m$-torsion, completion leaves it unchanged. Hence it is
the previously computed module $E$ with its Frobenius action.
To\"en's Lemma 3.9, pp.~21--22, gives
$\Ext^1_{\mathrm{naive}}(J_x,\GG_a)=0$.
The Artin--Schreier long exact sequence identifies the displayed
Ext group with $E/(F-1)E$, whose full orbit decomposition was proved
above. This works for every $p$, including $2$ and $3$.
\end{proof}

\paragraph{Scope of the modern geometric comparison.}
To\"en's Theorem 0.3 and Proposition 3.7 compare the full complex
of local Jacobians $J_*(X)$ with fppf cohomology for smooth $X$,
with $G$ affine if the cohomological degree is positive and unipotent
if it exceeds one. Section 3.2, p.~19, explicitly distinguishes
Yoneda Ext in pro-group schemes from Ext of all fppf sheaves;
the realization is not fully faithful in higher degrees.
Remark 3.10(2), p.~22, explains the missing restrictions in the old
formulation by its use of Zariski cohomology \cite{Toen2026}.
This is not a category-independent all-target comparison.

Cais--Otabe's Theorem 1.3, p.~3, and formula (1.8), p.~2,
prove
\[H^1_{\mathrm{fppf}}(U,G)=\varinjlim_{\mathfrak m}
\Ext^1(J_{\mathfrak m},G)\]
for every finite commutative group
scheme $G$, including nonreduced ones \cite{CaisOtabe2026}.
Their p.~3 discussion notes that Deligne's 1974 letter asserted
arbitrary commutative $G$ but left details in a two-variable
effaceability argument. Their finite-target theorem is therefore
not used as an unqualified all-degree comparison.

\begin{proposition}[H3b: augmentation alone does not repair ordinary degree zero; Refuted claim]
On the perfect \'etale site, no source object $P$ can give
isomorphisms
\[
 \Hom(P,G^{\mathrm{perf}})\simeq H^0(k(t),G)
\]
natural in all commutative algebraic $k$-groups $G$, where the
right side means ordinary rational points. This applies also to
any proposed augmented $J_{\mathrm{aug}}$.
\end{proposition}
\begin{proof}
Take $G=\GG_a$ and its homomorphism $F:x\mapsto x^p$.
The induced endomorphism of $G^{\mathrm{perf}}$ is an isomorphism,
so it acts bijectively on $\Hom(P,G^{\mathrm{perf}})$ for every
$P$. On $H^0(k(t),\GG_a)=k(t)$ it is not surjective: if
$h^p=t$, the valuation at $t=0$ gives $p\,v_0(h)=1$, a
contradiction. Naturality would conjugate these two maps, which
is impossible. Thus both the augmentation and the perfection of
the geometric/coefficient argument must be specified before an
all-degree comparison can be asserted. This does not affect
higher Ext vanishing in the precisely stated categories.
\end{proof}


\subsection{SA19 H1--H2: mathematical identification and the chronology gate}

\paragraph{H1: component-group duality.}
The boxed statement on 43:117 identifies the component groups of dual
abelian varieties as dual with values in $\QQ/\ZZ$. This matches
SGA~7~I, Expos\'e~IX, (1.2.1) and Conjecture~1.3, printed p.~323
(PDF p.~328 in the opened original scan): the pairing is constructed,
while its perfectness is stated separately as a conjecture
\cite{SGA7I}. The original page was visually checked.
For a complete discrete valuation field with perfect residue field of
characteristic $p>0$, perfectness is now Suzuki's Theorem~C,
p.~4 of the opened author version; see its setup on p.~2 and
Theorems~A--B on pp.~3--4 \cite{SuzukiPairing2020}.
An algebraically closed residue field satisfies this hypothesis.
The derived duality and descent in those theorems give a substantive
modern counterpart to the programme on 43:116--120. This is a
known theorem, not a new proof here. The residue-field hypothesis
cannot simply be dropped: Bosch--Lorenzini's Lemma~6.1 and
Example~6.2, pp.~393--395, give Jacobians over imperfect residue
fields for which the pairing has nonzero radical
\cite{BoschLorenzini2002}; these examples are outside the manuscript's
algebraically closed setting.

\paragraph{H2: three-variable Ext and biextensions.}
The boxed three-variable Ext on 43:16 is labelled the Lang pairing.
On 43:21 the tensor product of two resolutions, the first cohomology
of a Hom complex, and the equivariant cup product give additional
evidence for a biextension interpretation. Some decorations and
inputs remain uncertain in the transcription; three-variable
notation alone does not specify an exact category or a construction.
The precise published comparison is
\[
 \operatorname{Biext}^{1}(P,Q;G)
   \simeq \operatorname{Ext}^{1}(P\mathbin{\otimes}^{\mathbf L}Q,G),
\]
SGA~7~I, Expos\'e~VII, Corollary~3.6.5, printed p.~205
(PDF p.~210), following the categorical construction in 3.6.4,
p.~204 \cite{SGA7I}. The original formula was visually checked;
the corollary states degree one, not an unrestricted all-degree
comparison.

Mumford's \emph{Bi-extensions of formal groups}, \S\S2--3,
original pp.~310--314, gives the torsor definition with two compatible
extension laws and then the formulation for group functors
\cite{Mumford1969}. The official TIFR retypeset volume retains
original page numbers in its margins; its running article pages
321--337 differ from the original 307--322. Expos\'e~VII's introduction,
printed p.~133, explicitly attributes the introduction of
biextensions in the formal-group setting to Mumford. Any historical
priority claim must address that contemporaneous attribution.

\paragraph{The dating inference is invalid; priority remains Undetermined.}
The dates cited in SA19 H1--H2 do not date the relevant scientific
leaves. On 43:144, June~1962 dates Shatz's cited Harvard thesis.
On 43:87, May~1964 dates Lang's cited Bourbaki talk; that bibliography
also cites N\'eron's 1964 paper. These are bibliographic dates,
which do not give an upper bound for the dates of other leaves.
On 43:118, September~1963 occurs in a cancelled administrative list
mentioning James and Bott. Adjacency to 43:117 does not establish
when the latter's scientific content was written. The supplied
transcriptions of 43:16, 43:21 and 43:117 carry no explicit dates,
and the inspected cover sheets provide none.

The official Montpellier catalogue gives cote~43 the tentative span
\emph{[à partir de 1961-vers 1968]}, marking the notes themselves
undated. Footnote~37 records a verso with Bourbaki exposé~338
(1967--1968); the catalogue also warns that versos are often
unrelated to the mathematical notes \cite{MontpellierCatalogue}.
This does not support a dossier-wide upper bound of 1964 or 1965.
Nor does that late verso date the particular scientific leaves.
The catalogue dates cote~44 separately to 1964; this cannot be
transferred to cote~43.

Thus the mathematical identifications are supported, but the offered
1962--64 chronology does not prove pre-SGA-7 priority. This refutes
the proposed inference, not the possibility of early composition.
The exact remaining gate is independently secure dating or provenance
for the particular scientific leaves, together with an argument
connecting that evidence to those texts. A dated quotation or
independently dated matching draft could close the gate; a repeated
dossier-wide date range cannot.


\section{The endomorphism algebra of the residual complex}\label{sec:I}
The finite topology on $\Hom_A(M,N)$ means pointwise convergence
when the modules are discrete. The universal problem below uses
separated complete graded topological $A$-algebras, continuous
multiplication, and continuous graded homomorphisms, with $A$ discrete
as base. The prescribed degree zero has the product of local adic
topologies. This specifies an interpretation of completion; the
uncertain transcription on 43:65 does not uniquely specify a category.

\begin{proposition}[I1: degree-zero diagonality; Proved]\label{i:zero}
For a regular noetherian local ring, normalize the residual complex by
$K^i=\bigoplus_{\operatorname{ht}\mathfrak p=i}E_A(A/\mathfrak p)$.
Then $\End_A^0(K)=\prod_{\mathfrak p}\widehat{A_{\mathfrak p}}$
as topological algebras, and $\End_A^i(K)=0$ for $i<0$.
\end{proposition}
\begin{proof}
Write $E_{\mathfrak p}=E_A(A/\mathfrak p)$. If
$a\in\mathfrak p\setminus\mathfrak q$, every element of
$E_{\mathfrak p}$ is killed by a power of $a$, while multiplication
by $a$ is invertible on $E_{\mathfrak q}$. Thus the corresponding Hom
vanishes unless $\mathfrak p\subseteq\mathfrak q$. Distinct primes
of the same height are incomparable; projections detect every element
of a direct sum. This proves diagonality and negative-degree vanishing.
Matlis' endomorphism theorem gives the displayed rings
\cite[Lemma 2, proof, p.~123]{Sastry1992}. Annihilators of finite
subsets of $E_{\mathfrak p}$ give exactly the local adic topology:
the submodules killed by powers of the maximal ideal have finite
length over $A_{\mathfrak p}$ and exhaust the hull. The direct sum
gives the product topology.
\end{proof}

\begin{theorem}[I2: complete surface presentation; Proved]\label{i:surface}
Let $A=k[[s,t]]$, $F=\operatorname{Frac}(A)$, and
$\mathfrak m=(s,t)$. Then $\End_A^*(K)$ with differential $\theta$
is universal among the separated complete graded topological
$A$-algebras with degree zero
\[
 B^0=F\times\prod_{\operatorname{ht}\mathfrak p=1}
                 \widehat{A_{\mathfrak p}}\times A,
\]
one specified element of degree one, its square zero, and the
immediate-specialization incidence relations. In particular the
canonical map from the algebra defined by this precise universal
property is an isomorphism, including degree two.
\end{theorem}
\begin{proof}
Let $\mathcal P_1$ denote the height-one primes, put
$R_{\mathfrak p}=\widehat{A_{\mathfrak p}}$, choose a generator
$f_{\mathfrak p}$ of $\mathfrak p$, and put
$F_{\mathfrak p}=\operatorname{Frac}(R_{\mathfrak p})$.
The factorial Gorenstein ring $A$ has residual resolution
\begin{equation}\label{i:res}
 0\to A\to F\xrightarrow{d^0}
 I=\bigoplus_{\mathfrak p\in\mathcal P_1}E_{\mathfrak p}
 \xrightarrow{d^1}E\to0,
\end{equation}
where $E=E_A(k)$, $E_{\mathfrak p}=F/A_{\mathfrak p}$, and
$d^0$ consists of quotient maps. Write $\delta_{\mathfrak p}$ for
the components of $d^1$ and $D(-)=\Hom_A(-,E)$.

\emph{Individual arrows.}
By \cite[Theorem 1.1, p.~1]{Schenzel2013}, $D(E_{\mathfrak p})$
is free of rank one over $R_{\mathfrak p}$: $A/\mathfrak p$ is
complete of dimension one. Its element $\delta_{\mathfrak p}$ is
a basis. Indeed, applying $\Hom_A(A/\mathfrak p,-)$ to
\eqref{i:res} leaves only
$\kappa(\mathfrak p)\to\Hom_A(A/\mathfrak p,E)$.
The free resolution of $A/(f_{\mathfrak p})$ shows this map is
surjective with kernel isomorphic to $A/\mathfrak p$, so it is
nonzero. Dualizing
$0\to\kappa(\mathfrak p)\to E_{\mathfrak p}
\xrightarrow{f_{\mathfrak p}}E_{\mathfrak p}\to0$
shows that the kernel of restriction of $D(E_{\mathfrak p})$ to
$D(\kappa(\mathfrak p))$ is
$f_{\mathfrak p}D(E_{\mathfrak p})$. Hence
$\delta_{\mathfrak p}$ is a unit times a basis. We use it to identify
\begin{equation}\label{i:arrows}
 \Hom_A(E_{\mathfrak p},E)=R_{\mathfrak p},\qquad
 \Hom_A(F,E_{\mathfrak p})=F_{\mathfrak p}.
\end{equation}
For the second identity, apply $\Hom_{A_{\mathfrak p}}(-,E_{\mathfrak p})$
to $0\to A_{\mathfrak p}\to F\to E_{\mathfrak p}\to0$
and localize at $f_{\mathfrak p}$. The middle Hom is already an
$F$-vector space. Explicitly $a\in F_{\mathfrak p}$ sends
$f\in F$ to $fa\bmod R_{\mathfrak p}$, using
$E_{\mathfrak p}=F_{\mathfrak p}/R_{\mathfrak p}$.

\emph{All graded pieces.}
Put
\[
 P=\prod_{\mathfrak p}R_{\mathfrak p},\qquad
 \mathbb A=\left\{(a_{\mathfrak p})\in\prod_{\mathfrak p} F_{\mathfrak p}:
 a_{\mathfrak p}\in R_{\mathfrak p}\text{ for almost all }\mathfrak p\right\}.
\]
For $S=A\setminus\{0\}$, factoriality gives $\mathbb A=S^{-1}P$:
clear its finitely many negative valuations by a product of powers of
$f_{\mathfrak p}$. Conversely every denominator has finitely many
prime divisors. A family $F\to E_{\mathfrak p}$ has values in the
direct sum precisely when its coefficients are integral almost
everywhere: evaluate at $1$ for necessity and use the finitely many
poles of each $f\in F$ for sufficiency. Consequently
\begin{equation}\label{i:pieces}
 \End_A^1(K)=\mathbb A\oplus P,\qquad
 \End_A^2(K)=\mathbb A/F.
\end{equation}
For the second equality dualize \eqref{i:res}, using $D(E)=A$:
\begin{equation}\label{i:dual}
 0\to A\to P\to D(F)\to E\to0.
\end{equation}
The first map is diagonal, since its components are
$a\delta_{\mathfrak p}$. Localizing kills $E$ and leaves $D(F)$
unchanged, proving the equality. For $c\in P$ representing $I\to E$
and $a\in\mathbb A$ representing $F\to I$, composition is
\begin{equation}\label{i:product}
 c\circ a=[(c_{\mathfrak p}a_{\mathfrak p})_{\mathfrak p}]
                         \in\mathbb A/F.
\end{equation}
All other products of positive-degree arrows vanish by endpoints;
degrees above two vanish. The differential is $(1,1)$, with square
$[1]=0$. Thus the degree-two relation is precisely diagonal $F$,
not an independent zero relation for each intermediate prime.

\emph{The topologies in degree one.}
The first summand has the usual restricted-product topology: $P$ with
its product adic topology is open. Indeed, vanishing at $1/u$ means
membership in $uP$, and these form a basis. The second summand has
the product of topologies $\tau_{\mathfrak p}$ with bases
\begin{equation}\label{i:mixed}
 uA+f_{\mathfrak p}^n R_{\mathfrak p}\ (u\notin\mathfrak p),
 \quad\text{or equivalently}\quad
 \mathfrak m^N A+f_{\mathfrak p}^n R_{\mathfrak p}.
\end{equation}
To prove this, exactness and factoriality give
\[
 \ker\delta_{\mathfrak p}
  =(A[1/f_{\mathfrak p}]+A_{\mathfrak p})/A_{\mathfrak p}.
\]
A fraction whose other components vanish has no other prime factors
in its denominator. Any finite set in $E_{\mathfrak p}$ lies in the
cyclic $A$-module generated by $1/(u f_{\mathfrak p}^n)$ for some
$u\notin\mathfrak p$. Its annihilator in the Hom space is exactly
$uA+f_{\mathfrak p}^nR_{\mathfrak p}$: modulo
$f_{\mathfrak p}^n$, the element $c/u$ must come from $A$.
Since $(u,f_{\mathfrak p}^n)$ is $\mathfrak m$-primary, some
$\mathfrak m^N$ is contained in it; conversely choose
$u\in\mathfrak m^N\setminus\mathfrak p$. This proves the
equivalent basis. At $\mathfrak p=(s)$, it is coefficientwise
$t$-adic convergence on $k((t))[[s]]$. Hence
$t^N\delta_{\mathfrak p}\to0$, although $t^N$ is always an
$s$-adic unit. The two occurrences of $P$ require different topologies.

\emph{Degree two.}
Write $P_\tau$ for the product with \eqref{i:mixed}. The subgroup
$H=\{h\in D(F):h(1)=0\}$ is open and topologically $P_\tau/A$.
For the quotient assertion, if $M'\subset M$ and $E$ is injective,
restriction $D(M)\to D(M')$ is an open surjection for the finite
topologies over a noetherian ring. The image of the annihilator of a
finitely generated $N\subset M$ is exactly the annihilator of
$M'\cap N$: extend a map from $M'$ by zero on $N$, and then use
injectivity. The intersection is finitely generated, proving openness.
Apply this to $F/A\subset I$; the kernel is $D(E)=A$.
Furthermore $D(F)=\bigcup_{u\ne0}u^{-1}H$, since every element of
$E$ is killed by a nonzero element of $A$, and multiplication by $u$
is a topological automorphism. This specifies the topology on
$\mathbb A/F$ without an unjustified ordinary-adelic quotient.

\emph{The universal morphism.}
Let $C$ have the prescribed degree zero and $T\in C^1$ satisfying
the relations. Write $e_\eta,e_P,e_m$ for the three block
idempotents. Finite sums of the individual height-one idempotents
converge to $e_P$, so continuity and incidence give
\begin{equation}\label{i:T}
 T=e_PTe_\eta+e_mTe_P.
\end{equation}
For $a=b/u\in\mathbb A$, $b\in P$, $u\ne0$, and $c\in P$, set
\begin{align}\label{i:normal}
 \alpha_C(a)&=bTe_\eta u_\eta^{-1},&
 \beta_C(c)&=e_mTc,\\
 \gamma_C([a])&=e_mTbTe_\eta u_\eta^{-1}.\nonumber
\end{align}
The elements of $P$ occupy the height-one block, and $u_\eta^{-1}$
means $1/u$ in the generic block. Clearing denominators and using
$A$-linearity proves independence of representations. For diagonal
$a=v/u\in F$, the last expression is
$v e_mTe_PTe_\eta u_\eta^{-1}=0$, by \eqref{i:T} and $T^2=0$.
Thus $\gamma_C$ descends to $\mathbb A/F$.
These formulas give an algebra homomorphism, the identity in degree
zero: the only composable positive-degree product is
$\beta_C(c)\alpha_C(a)=\gamma_C([ca])$, agreeing with
\eqref{i:product}. All others vanish by endpoints. The diagonal
actions agree directly, using $A$-linearity for the closed-point
action and the right $F$-action at the generic point.

We check continuity carefully. The map $\alpha_C$ is continuous
on the open subgroup $P$ with product adic topology, and hence on all
$u^{-1}P$. For $\beta_C$, given a zero-neighborhood $U$, choose
$V$ with $V+V\subset U$. Continuity from product adic degree zero
gives a finite set $T_0$ of primes and powers
$f_{\mathfrak p}^{n_{\mathfrak p}}$ such that tuples satisfying
those divisibilities map into $V$; all remaining coordinates are
unrestricted. Each map
$A\to C^1$, $a\mapsto a e_mTe_{\mathfrak p}$, is continuous
for the $\mathfrak m$-adic topology. Choose $N$ so that the sum of
its images of $\mathfrak m^N$ over $\mathfrak p\in T_0$ lies
in $V$. Tuples satisfying
$c_{\mathfrak p}\in\mathfrak m^N A+
f_{\mathfrak p}^{n_{\mathfrak p}}R_{\mathfrak p}$ split into a
tuple controlled by the first condition and finitely many such
$\mathfrak m^N A$ terms. Their images lie in $U$. This proves
continuity from $P_\tau$.
The continuous map $c\mapsto\beta_C(c)Te_\eta$ vanishes on
diagonal $A$, hence descends to the open quotient $H=P_\tau/A$.
Multiplication by $u_\eta^{-1}$ gives continuity of $\gamma_C$
on every $u^{-1}H$ and globally.

The morphism sends $\theta$ to $T$. It is unique because
\eqref{i:normal}, with $T=\theta$, represents every positive-degree
endomorphism. Finally, all Hom spaces in the finite topology are
closed subgroups of products of discrete target modules, defined by
closed $A$-linearity equations, so are separated and complete.
Composition is jointly continuous by controlling finitely many inputs
and their finitely many intermediate images. Statement~\ref{i:zero}
identifies the degree-zero topology. This proves the theorem.
\end{proof}

\begin{theorem}[I3: every regular local surface; Proved]\label{i:noncomplete}
Theorem~\ref{i:surface} holds for every regular noetherian local ring
$A$ of dimension two, with closed-point diagonal algebra $\widehat A$.
The completion required in degree one is specified below.
\end{theorem}
\begin{proof}
Put $B=\widehat A$, $S=A\setminus\{0\}$ and retain
$F,E,E_{\mathfrak p},R_{\mathfrak p},P$ from the preceding proof.
Let $C_{\mathfrak p}$ be the separated completion of
$R_{\mathfrak p}$ in the topology $\tau_{\mathfrak p}$ of
\eqref{i:mixed}, now defined using $A$, and put
$C=\prod_{\mathfrak p}C_{\mathfrak p}$. We will prove
\begin{equation}\label{i:noncompletepieces}
 C_{\mathfrak p}=\Hom_A(E_{\mathfrak p},E),\qquad
 \End_A^1(K)=S^{-1}P\oplus C,\qquad
 \End_A^2(K)=S^{-1}C/S^{-1}B.
\end{equation}

\emph{The completion map has no extra kernel.}
The map $R_{\mathfrak p}\to\Hom_A(E_{\mathfrak p},E)$,
$c\mapsto\delta_{\mathfrak p}c$, is injective: a nonzero $c$ is
a unit times a power of $f_{\mathfrak p}$, and therefore acts
surjectively on $E_{\mathfrak p}$. Its composite with the nonzero
$\delta_{\mathfrak p}$ cannot vanish. Exactness and factoriality
still give the same formula for $\ker\delta_{\mathfrak p}$,
so the induced finite topology on $R_{\mathfrak p}$ is exactly
\eqref{i:mixed}.

For density, take $h:E_{\mathfrak p}\to E$ and finitely many
inputs, contained in the cyclic $A$-module generated by
$e=1/(u f_{\mathfrak p}^n)$, $u\notin\mathfrak p$.
Applying $\Hom_A(A/f_{\mathfrak p}^n,-)$ to the residual resolution
leaves a surjection
\[
 E_{\mathfrak p}[f_{\mathfrak p}^n]
 \xrightarrow{\delta_{\mathfrak p}} E[f_{\mathfrak p}^n],
\]
since $A/f_{\mathfrak p}^n$ has a length-one free resolution.
The source is cyclic over $A_{\mathfrak p}$, generated by $e$.
Choose $c\in A_{\mathfrak p}$ with
$\delta_{\mathfrak p}(ce)=h(e)$. Then
$\delta_{\mathfrak p}c$ agrees with $h$ on all the selected
inputs. The Hom space is separated complete, being closed in a
product of discrete copies of $E$. It is therefore precisely the
separated completion $C_{\mathfrak p}$. Taking products gives
$C=\Hom_A(I,E)$ and dense $P_\tau\subset C$.

The first degree-one summand is unchanged. Dualizing the residual
resolution, now using $D(E)=B$, gives
\[
 0\longrightarrow B\longrightarrow C\longrightarrow D(F)
 \longrightarrow E\longrightarrow0.
\]
The first map sends $b$ to $(b\delta_{\mathfrak p})_{\mathfrak p}$.
Localization proves the last equality of \eqref{i:noncompletepieces};
in particular its entire kernel is $S^{-1}B$. The open-restriction
argument above identifies the open subgroup
$H=\{h\in D(F):h(1)=0\}$ with the topological quotient $C/B$,
and again $D(F)=\bigcup_{u\ne0}u^{-1}H$.

\emph{The universal morphism extends over exactly this completion.}
Let $Q$ be a separated complete target algebra with the prescribed
degree zero $F\times P\times B$ and an element $T$ satisfying the
relations. The map $P\to Q^1$, $c\mapsto e_mTc$, is continuous
for $P_\tau$ by the finite-coordinate and $\mathfrak m^NA$ argument
in the preceding proof. A continuous additive homomorphism is
uniformly continuous, so it extends uniquely to a continuous map
$\beta_Q:C\to Q^1$. This extension is right $P$-linear by
density and continuity. It is also left $B$-linear. Indeed,
$A$-linearity first extends from $P_\tau$ to $C$. For $b\in B$
approximate $b$ by $a_j\in A$ in the $\mathfrak m$-adic topology.
Multiplication by $a_j$ converges pointwise to multiplication by $b$
on $\Hom_A(I,E)$, because every tested value in $E$ is killed
by a power of $\mathfrak m$; multiplication in $Q$ is continuous.
Passing to limits proves the assertion. In particular
\[
 \beta_Q((b\delta_{\mathfrak p})_{\mathfrak p})=b e_mTe_P.
\]
Use the preceding formula for $\alpha_Q:S^{-1}P\to Q^1$, and set
\[
 \gamma_Q([c/u])=\beta_Q(c)Te_\eta u_\eta^{-1}
 \qquad(c\in C,\ u\in S).
\]
Clearing denominators proves independence of representation.
Every element $b/u$ from $S^{-1}B$ maps to
$b e_mTe_PTe_\eta u_\eta^{-1}=0$, by $T^2=0$.
The composition rule is
$\beta(c)\alpha(b/u)=[cb/u]$ for the natural right $P$-action
on $C$; the formulas respect this and all diagonal actions.
Continuity of $\gamma_Q$ on $H=C/B$ follows from the quotient
topology and then extends over each $u^{-1}H$. These maps give the
required continuous algebra homomorphism. Density of $P_\tau$ in
$C$ gives uniqueness: all positive-degree terms are topologically
generated by degree zero and $\theta$. This proves the theorem.
\end{proof}

\paragraph{Remaining gate.}
The presentation is proved in every regular local surface, including
the requested $k[[s,t]]$ case and degree two, under the stated
completion convention. In higher dimension a longer residual
resolution has further intermediate terms. Propositions I4--I7
prepare the proof, and Theorem I8 below identifies all higher relations
in the complete open-left-ideal topology category.
Constructing a residual complex or proving its residue theorem does
not establish that assertion. Yekutieli explicitly specifies
semi-topological rings and a tensor topology
\cite[Definitions 1.2.1 and 1.2.11, pp.~15 and 18]{Yekutieli1992};
his construction is not an opened presentation theorem for the whole
endomorphism algebra. Schenzel's completeness criterion explains why
simply identifying the noncomplete arrow module with
$R_{\mathfrak p}$ would have failed; the dense completion in
Theorem~\ref{i:noncomplete} supplies the missing argument.

\begin{proposition}[I4: adjacent arrows in every dimension; Proved]\label{i:adjacent}
Let $\mathfrak p\subset\mathfrak q$ be an immediate specialization
in a regular noetherian local ring $A$. Put $R=A_{\mathfrak q}$,
$\mathfrak m=\mathfrak qR$, and $T=\widehat{A_{\mathfrak p}}$.
For the residual differential $\delta_{\mathfrak p\mathfrak q}$,
the map
\[
 T\longrightarrow\Hom_A(E_{\mathfrak p},E_{\mathfrak q}),
 \qquad c\longmapsto\delta_{\mathfrak p\mathfrak q}c
\]
is injective and dense. Its induced topology has basis
\[
 \mathfrak m^N R+(\mathfrak pR_{\mathfrak p})^n T,
 \qquad N,n\geq1.
\]
The Hom module is its separated completion for this topology.
\end{proposition}
\begin{proof}
Both hulls are $R$-modules, so $A$-linearity is equivalent to
$R$-linearity. Set $h=\operatorname{ht}(\mathfrak pR)$; then
$\dim R=h+1$. Choose parameters $x_1,\ldots,x_h$ in the regular
local ring $R_{\mathfrak p}$. In its top local cohomology,
which is its injective hull, the cyclic submodules generated by
\[
 e_n=[1/(x_1\cdots x_h)^n],\qquad
 J_n=(x_1^n,\ldots,x_h^n)R_{\mathfrak p},
\]
exhaust $E_{\mathfrak p}$, and $\operatorname{Ann}(e_n)=J_n$.
The ideals $J_n$ are cofinal with powers of $\mathfrak pR_{\mathfrak p}$.
For $h=0$ use $e_n=1$ and $J_n=0$ in the fraction field.
These are the usual Gorenstein and local-cohomology identifications
\cite[Lemmas 47.21.3, 47.21.7, and proof of 47.21.10]{StacksGorenstein}.
Clearing finitely many coefficients over $R_{\mathfrak p}$ shows
that every finite set in $E_{\mathfrak p}$ lies in
$R(e_n/u)$ for suitable $n$ and $u\in R\setminus\mathfrak pR$.

Let $I_n=J_n\cap R$. This ideal is $\mathfrak pR$-primary,
so $R/I_n$ has dimension and depth one. The finite-free-resolution
bound for a regular local ring gives
$\operatorname{pd}_R(R/I_n)\leq h$
\cite[Proposition 10.110.1]{StacksRegularResolution}.
Applying $\Hom_R(R/I_n,-)$ to the residual resolution leaves only
\[
 E_{\mathfrak p}[I_n]\xrightarrow{\delta_{\mathfrak p\mathfrak q}}
 E_{\mathfrak q}[I_n]
\]
in degrees $h,h+1$. The last map is surjective. Its kernel
$W_n=\Ext_R^h(R/I_n,R)$ is finite over $R$.

Given a homomorphism $f$ and inputs in $R(e_n/u)$, lift
$f(e_n/u)$ through this surjection. Since the Artinian complete
intersection $R_{\mathfrak p}/J_n$ is Gorenstein,
$E_{\mathfrak p}[I_n]$ is generated by $e_n/u$ over
$R_{\mathfrak p}$. Thus some $c\in R_{\mathfrak p}$ gives
$\delta_{\mathfrak p\mathfrak q}(ce_n/u)=f(e_n/u)$.
This proves density. Every nonzero $c\in T$ acts surjectively on
$E_{\mathfrak p}$: its Matlis dual is multiplication by $c$ on
the regular complete local domain $T$. Since the displayed
surjection shows $\delta_{\mathfrak p\mathfrak q}\ne0$, the
coefficient map is injective.

Identify $E_{\mathfrak p}[I_n]$ with $(R/I_n)_{\mathfrak p}$
using $e_n$. The finite submodule $W_n$ localizes to this whole
ring, since $E_{\mathfrak q}[I_n]$ vanishes at $\mathfrak p$.
It is therefore commensurable with $R/I_n$: clearing its finite
set of denominators, and then lifting $1$ after localization,
gives $a_n,b_n\in R\setminus\mathfrak pR$ with
\[
 a_n(R/I_n)\subseteq W_n\subseteq b_n^{-1}(R/I_n).
\]
Vanishing on $e_n/u$ means that $c\bmod J_nT$ belongs to
$uW_n$. As $u$ varies, these subgroups are cofinal with
$uR+J_nT$. The ideal $(u,I_n)$ is $\mathfrak m$-primary, and
$\mathfrak m^N$ contains an element outside $\mathfrak pR$.
Thus this family is cofinal with
$\mathfrak m^N R+J_nT$. Cofinality of the $J_n$ with the
$\mathfrak p$-adic powers proves the asserted basis.
Finally the Hom module, being a closed subgroup of a product of
discrete copies of $E_{\mathfrak q}$, is separated complete.
The injective dense map therefore identifies it with the separated
completion just described.
\end{proof}

\begin{proposition}[I5: dense generation in every dimension; Proved]\label{i:dense}
The graded subalgebra of $\End_A^*(K)$ generated by its diagonal
degree-zero algebra and $\theta$ is dense in each degree for the
finite topology.
\end{proposition}
\begin{proof}
First fix $\mathfrak p\subseteq\mathfrak q$. Induct on the height
gap to approximate every map $E_{\mathfrak p}\to E_{\mathfrak q}$
on a finite set by finite sums of saturated paths of residual
arrows with diagonal coefficients. Gap zero follows from
Proposition~\ref{i:zero}; gap one is Proposition~\ref{i:adjacent}.
For larger gap work over $R=A_{\mathfrak q}$, of dimension $d$,
and put the prescribed inputs in a cyclic $R$-submodule $Re$,
where $e=e_n/u$ is as in the preceding proof. Its annihilator $I$
is $\mathfrak pR$-primary. Hence $R/I$ has positive depth, and
\cite[Proposition 10.110.1]{StacksRegularResolution} gives
$\operatorname{pd}_R(R/I)\leq d-1$.
The last differential of $\Hom_R(R/I,K_R)$ is consequently
surjective:
\[
 \bigoplus_{\substack{\operatorname{ht}(\mathfrak rR)=d-1\\
                      \mathfrak p\subseteq\mathfrak r\subset\mathfrak q}}
 E_{\mathfrak r}[I]\longrightarrow E_{\mathfrak q}[I].
\]
Lift $f(e)$; its preimage has finitely many nonzero components
$y_{\mathfrak r}$. The maps $Re\to E_{\mathfrak r}$ sending
$e$ to $y_{\mathfrak r}$ extend to maps
$g_{\mathfrak r}:E_{\mathfrak p}\to E_{\mathfrak r}$ by
injectivity. Thus $\sum_{\mathfrak r}
\delta_{\mathfrak r\mathfrak q}g_{\mathfrak r}$ agrees with
$f$ on $Re$. Approximate each $g_{\mathfrak r}$ on the prescribed
finite inputs by induction, and compose with
$\delta_{\mathfrak r\mathfrak q}$.

For the whole complex, finitely many inputs involve only finitely
many source summands, and their images involve only finitely many
target summands. Apply the preceding construction to these finitely
many component maps, insert endpoint idempotents, and sum. The
result has exactly the prescribed values on the original inputs.
Every path lies in the algebra generated by the diagonal algebra
and $\theta$, proving density.
\end{proof}

\paragraph{Refined higher-dimensional gate.}
Dense generation does not prove a presentation of a completed
algebra. It remains to establish the quotient topology and uniform
continuity needed to extend the universal map from the dense path
subalgebra, and to identify the entire closed ideal of relations
with the ideal generated by $\theta^2$ and incidence.
Yekutieli's Theorem~1.14, p.~9 of the opened author version,
constructs an adelic DG action; Lemma~1.11, pp.~7--8, is an
approximation statement \cite{Yekutieli2003}. An action does not
itself identify its kernel or supply the required quotient topology.
The full presentation in dimension at least three remains
\emph{Undetermined}.

\begin{proposition}[I6: closed-image exactness of residual columns; Proved]\label{i:closedexact}
Let $R$ be a regular noetherian local ring of dimension $d$,
$\mathfrak p$ a nonmaximal prime of height $h$, and $K_R$ its
residual resolution. With finite topologies and differentials by
postcomposition, for $j>h$ one has
\[
 \overline{\operatorname{im}\bigl(
  \Hom(E_{\mathfrak p},K_R^{j-1})\to
  \Hom(E_{\mathfrak p},K_R^j)\bigr)}
 =\ker\bigl(\Hom(E_{\mathfrak p},K_R^j)\to
             \Hom(E_{\mathfrak p},K_R^{j+1})\bigr).
\]
The kernel at $j=h$ is zero.
\end{proposition}
\begin{proof}
Continuity gives the inclusion from left to right. For the other
inclusion, take a cycle $f:E_{\mathfrak p}\to K_R^j$ and put
finitely many prescribed inputs in $Re$ with $e=e_n/u$ as in
Proposition~\ref{i:adjacent}. Its annihilator $I$ is
$\mathfrak p$-primary and
$I R_{\mathfrak p}=(x_1^n,\ldots,x_h^n)$.
The finite module $M=\Ext_R^j(R/I,R)$ vanishes after localization
at $\mathfrak p$ when $j>h$, by the length-$h$ Koszul resolution
there. Thus some $v\notin\mathfrak p$ annihilates $M$.

The element $e/v\in E_{\mathfrak p}$ has the same annihilator
$I$. The restriction of $f$ to $R(e/v)\simeq R/I$ is a cycle.
Its restriction to $Re$ multiplies the cohomology class by $v$,
because $f(e)=v f(e/v)$. This class is zero. Therefore some map
$b:Re\to K_R^{j-1}$ has differential equal to $f$ on $Re$.
Its cyclic source implies that $b$ has finitely many nonzero
components. Extend each component to $E_{\mathfrak p}$ by
injectivity and sum, obtaining
$g:E_{\mathfrak p}\to K_R^{j-1}$. Its differential agrees with
$f$ on the prescribed inputs, proving density in the kernel.

At $j=h$, only the diagonal summand receives a map from
$E_{\mathfrak p}$, so the Hom group is $\widehat{R_{\mathfrak p}}$.
Choose a prime immediately above $\mathfrak p$. The injectivity
in Proposition~\ref{i:adjacent} gives zero kernel for the next
differential.
\end{proof}

For a height-gap-two interval $\mathfrak p\subset\mathfrak q$,
apply this with $R=A_{\mathfrak q}$ and $j=h+1$. It identifies
the kernel of postcomposition from the full adjacent-arrow family
to $\Hom(E_{\mathfrak p},E_{\mathfrak q})$ as the closure of
the family obtained from the preceding residual differential and
$\widehat{A_{\mathfrak p}}$. This is an exact relation statement
in the actual Hom topologies. It does not identify the closure
inside an independently completed free algebra. The quotient-topology
and universal-map continuity gate therefore remains.

\begin{proposition}[I7: strict exactness above the first arrow over a complete ring; Proved]
Let $A$ be a complete regular local ring, let
$E=E_A(A/\mathfrak m)$, and let $\mathfrak p\ne\mathfrak m$.
Write $h=\operatorname{ht}\mathfrak p$ and
$C^j=\Hom_A(E_{\mathfrak p},K_A^j)$ with the finite topology.
The complex $C^\bullet$ is exact as a complex of abstract modules.
Moreover, for every $j>h+1$, its differential induces an open
surjection
\[
 C^{j-1}\longrightarrow Z^j:=\ker(C^j\to C^{j+1}).
\]
No assertion that the first-arrow topology is the local adic
topology is included.
\end{proposition}
\begin{proof}
Matlis' endomorphism theorem identifies
$A=\Hom_A(E,E)$ because $A$ is complete. For a projective
resolution $P_\bullet\to E_{\mathfrak p}$, tensor--Hom adjunction
and exactness of $\Hom_A(-,E)$ give
\[
 \Ext_A^j(E_{\mathfrak p},A)
   =\Hom_A\bigl(\operatorname{Tor}_j^A(E_{\mathfrak p},E),E\bigr).
\]
This uses no finite-generation hypothesis on $E_{\mathfrak p}$:
the adjunction is applied to the whole projective resolution.
Choose $u\in\mathfrak m\setminus\mathfrak p$. The source
$E_{\mathfrak p}$ is an $A_u$-module, whereas $E_u=0$.
Flat localization and derived tensor associativity show
\[
 E_{\mathfrak p}\mathbin{\otimes}_A^{\mathbf L}E
  \simeq
 E_{\mathfrak p}\mathbin{\otimes}_{A_u}^{\mathbf L}E_u=0.
\]
All the displayed Ext groups vanish. Since $K_A^\bullet$ is an
injective resolution of $A$, this proves exactness of $C^\bullet$.

To prove openness, fix a finite set of inputs in $E_{\mathfrak p}$
and put them in a cyclic submodule $N=Ae$ as in
Proposition~\ref{i:adjacent}. Its annihilator $I$ is
$\mathfrak p$-primary and $I A_{\mathfrak p}$ is generated by
powers of a regular system of parameters. If $j>h+1$, the finite
module $\Ext_A^{j-1}(A/I,A)$ vanishes after localization at
$\mathfrak p$. Choose $v\notin\mathfrak p$ annihilating it, and
put $N'=A(e/v)$.

Let $f\in Z^j$ vanish on $N'$. Abstract exactness gives
$g\in C^{j-1}$ with $dg=f$. The restriction of $g$ to $N'$
is a cycle. Restricting further to $N$ multiplies its cohomology
class by $v$, so the class is zero. There is therefore
$b:N\to K_A^{j-2}$ with $db=g|_N$. This map has only finitely
many nonzero injective components; extend each to
$E_{\mathfrak p}$ and sum to obtain $\widetilde b$.
Then $g-d\widetilde b$ vanishes on $N$ and still maps to $f$.
Thus the image of the annihilator of $N$ contains the open
subgroup $Z^j\cap\operatorname{Ann}(N')$. This is exactly
openness of the surjection.
\end{proof}

This removes an abstract exactness and quotient-topology obstruction
for complete-ring columns above their first arrow. It does not
supply the missing uniform construction on arbitrary families of
intermediate primes in an independently defined universal algebra.
In particular it does not promote the full higher-dimensional
presentation to a proved result.


\begin{theorem}[I8: the full presentation in the finite-topology category; Proved]\label{i:linearuniversal}
Let $A$ be a regular noetherian local ring of finite dimension, and
let $K$ be its residual resolution, normalized by height. Put
$B=\prod_{\mathfrak p}\widehat{A_{\mathfrak p}}$.
The finite-topology algebra $\End_A^*(K)$ is universal among
separated complete graded topological $A$-algebras whose topology
has a basis of open left ideals, equipped with the prescribed
continuous degree-zero map from $B$ and an element $T$ of degree one
satisfying
\[
 T^2=0,\qquad e_{\mathfrak q}Te_{\mathfrak p}=0
 \quad\text{unless }\mathfrak p\subset\mathfrak q
 \text{ is an immediate specialization}.
\]
In particular the completed presentation in this category is an
isomorphism in every degree and in every dimension. The universal
property allows a continuous map from $B$; requiring it to identify
degree zero gives the version in the question.
\end{theorem}

\begin{proof}
We first prove a statement about discrete representations that does
not assume the presentation.

\emph{Cousin complexes.}
Call a complex $M$ a Cousin complex here if
\[
 M^i=\bigoplus_{\operatorname{ht}\mathfrak p=i}M_{\mathfrak p},
\]
where $M_{\mathfrak p}$ is an $A_{\mathfrak p}$-module all of whose
elements are killed by powers of $\mathfrak pA_{\mathfrak p}$.
Its differentials are $A$-linear. A map between two such point
modules is zero unless the source prime is contained in the target
prime, by the element-annihilator argument of
Proposition~\ref{i:zero}. Thus the differential has the required
incidence property. The action of
$\widehat{A_{\mathfrak p}}$ on $M_{\mathfrak p}$ is canonical:
on each element it is determined modulo a power of
$\mathfrak pA_{\mathfrak p}$. Every $A$-linear morphism of Cousin
complexes commutes with these completed diagonal actions.

\emph{A universal property of a complete-intersection Cousin
resolution.}
Let $J=(x_1,\ldots,x_h)$ be generated by an $A$-regular sequence,
put $N=A/J$, and put
\[
 C_J=\Hom_A(N,K)=K[J].
\]
The Koszul resolution and its self-duality give
$\Ext_A^i(N,A)=0$ for $i\ne h$ and
$\Ext_A^h(N,A)\simeq N$
\cite[Lemma 48.15.5, tag 0BQZ]{StacksKoszulDuality}.
Choose this last identification, so that
$N$ identifies with $Z^h(C_J)$. Every prime containing $J$ has
height at least $h$, so $C_J^i=0$ for $i<h$. In degree $h$ its
terms identify with
\[
 C_J^h=\bigoplus_{\substack{\mathfrak r\supset J\\
                            \operatorname{ht}\mathfrak r=h}}
 N_{\mathfrak r},
\]
with the augmentation $N\to C_J^h$ given by localization.
Indeed $N$ is Cohen--Macaulay, has no embedded primes, and the
Koszul self-duality identifies the localized top Ext at each
minimal prime with the corresponding Artinian complete intersection.

For every Cousin complex $M$, an $A$-linear map
$N\to Z^h(M)$ extends uniquely to a morphism
\[
                         C_J\longrightarrow M.          \tag{I8.1}
\]
Here is the full extension argument. In degree $h$, localize the
given map at each height-$h$ prime and take their direct sum. Suppose
degrees through $j-1$ have been defined compatibly with the
differentials, where $j>h$. For a prime $\mathfrak q$ of height
$j$, project the composite
$C_J^{j-1}\to M^{j-1}\to M^j$ to $M_{\mathfrak q}$ and
localize at $\mathfrak q$. It kills the kernel of
\[
 (C_J^{j-1})_{\mathfrak q}\longrightarrow
 (C_J^j)_{\mathfrak q}=C_{J,\mathfrak q}^j.
\]
For $j=h+1$ this follows because the prescribed map on $N$ is a
cycle; for larger $j$ it follows from the already established chain
identity and $d_M^2=0$. The displayed arrow is surjective because
the localized complex has no cohomology in degree $j>h$ and no
terms above $j$. Consequently the composite factors uniquely
through $C_{J,\mathfrak q}^j$. Take the direct sum over
$\mathfrak q$. Each source element has finitely many components,
so this gives a well-defined morphism, and comparison of all target
components proves the chain identity. Surjectivity in this
construction also proves uniqueness at every stage. This proves
(I8.1).

\emph{These resolutions generate every Cousin complex.}
Take $v\in M_{\mathfrak p}$ and set
$h=\operatorname{ht}\mathfrak p$. Its differential has finitely
many nonzero components, at primes $\mathfrak q$ strictly containing
$\mathfrak p$. For each such component choose an element of
$\mathfrak q\setminus\mathfrak p$ and a power killing that
component. The product $a$ of these powers is outside
$\mathfrak p$ and satisfies $d_M(av)=0$.
The ideal $\operatorname{Ann}_A(v)$ is $\mathfrak p$-primary:
elements outside $\mathfrak p$ act invertibly, and a power of
$\mathfrak p$ kills $v$. Since $A$ is regular, hence
Cohen--Macaulay, this ideal contains an $A$-regular sequence of
length $h$. Explicitly, after choosing $i<h$ terms, the quotient
is Cohen--Macaulay and its associated primes have height $i$ in
$A$. The height-$h$ annihilator is contained in none of them; prime
avoidance supplies the next nonzerodivisor. This proves the assertion
inductively, using the Cohen--Macaulay quotient property
\cite[Lemma 10.104.2, tag 02JN]{StacksCMRegularSequence}.
Choose such a sequence and its ideal $J$.
The map $N=A/J\to Z^h(M)$ sending $1$ to $av$ extends by
(I8.1) to $C_J\to M$.

The $\mathfrak p$-component $z_{\mathfrak p}$ of the image of
$1\in N$ in $C_J^h$ generates
$C_{J,\mathfrak p}^h\simeq N_{\mathfrak p}$. The extension sends
$z_{\mathfrak p}$ to $av$, and therefore sends
$z_{\mathfrak p}/a$ to $v$. Thus every vector is in the image of
one such morphism. Taking the sum of these maps gives a surjection
of Cousin complexes
\[
                    P:=\bigoplus_\lambda C_{J_\lambda}
                    \longrightarrow M.                  \tag{I8.2}
\]

\emph{All residual endomorphisms act on the quotient.}
Every $C_J=K[J]$ is preserved by all $A$-linear graded
endomorphisms of $K$. Hence $P$ is a discrete continuous
$\End_A^*(K)$-module: a finite set of vectors of $P$ involves only
finitely many inputs in copies of $K$.
The kernel of (I8.2) is preserved by $B$ and by $\theta$ because
the map is a morphism of Cousin complexes. It is consequently
preserved by every word in these operators. It is also preserved by
every $f\in\End_A^*(K)$. To see this without assuming any
relation theorem, take one vector in the kernel, involving
finitely many inputs $z_1,\ldots,z_r$ in copies of $K$.
Proposition~\ref{i:dense} supplies a word polynomial $w$ in $B$
and $\theta$ agreeing with $f$ on all $z_i$. Thus the two operators
give exactly the same value on this vector, and $w$ preserves the
kernel. So $f$ preserves it as well.

It follows that $M$ has a continuous action of the whole
$\End_A^*(K)$ extending its diagonals and differential.
Continuity follows by lifting any finite set of vectors through
(I8.2) and using the corresponding finite set of inputs in $K$.
Uniqueness follows from density. The same density argument shows
that every morphism of Cousin complexes commutes with this action.
This construction proves all relations on every discrete Cousin
representation; it uses density to preserve a kernel in a known
discrete representation, not as a substitute for a presentation.

\emph{Passage to the complete target algebra.}
Let $Q$ be a target in the statement. Replace it by the closed
subalgebra generated by $B$ and $T$. Insert finite sums of the
point idempotents between factors in a word and pass to their
limits. The incidence relations show that a word of degree greater
than $\dim A$ vanishes. Thus this closed subalgebra has finitely
many nonzero graded pieces. Its topology still has a basis of open
left ideals. Replacing an open left ideal $L$ by
$\bigoplus_i(L\cap Q^i)$ gives a smaller open homogeneous left
ideal, since there are finitely many grades. These may therefore
be used as a basis.

For such an ideal $L$, the discrete graded left module $Q/L$
decomposes under $B$ into point-local torsion summands. In fact the
annihilator in $B$ of each vector is open, so it contains the full
factors outside a finite set of primes and suitable powers of the
maximal ideals at the remaining primes. The operator $T$ is its
differential. Decomposing by the constant offset
\emph{graded degree minus prime height} expresses $Q/L$ as a sum
of shifted Cousin complexes. By the result just proved it has a
unique continuous $\End_A^*(K)$-action extending its given
$B,T$-action.

Define
\[
 \phi_L:\End_A^*(K)\longrightarrow Q/L,
 \qquad f\longmapsto f(1+L).
\]
These maps are continuous and compatible as $L$ shrinks, because
the quotient maps are morphisms of the Cousin representations.
Completeness and separatedness give
$Q\simeq\varprojlim_L Q/L$ as topological additive groups.
The $\phi_L$ therefore define a continuous graded map
$\phi:\End_A^*(K)\to Q$, sending $B$ to its prescribed image
and $\theta$ to $T$. It is multiplicative: this holds on word
polynomials, which are dense, and both multiplications and $\phi$
are continuous. Density proves uniqueness as well.

Finally the finite topology on $\End_A^*(K)$ has a basis of open
homogeneous left ideals, namely annihilators of finite sets of
homogeneous inputs. It is separated and complete as proved above.
Thus it belongs to the target category and has exactly the asserted
universal property.
\end{proof}

\paragraph{Completion convention after I8.}
Theorem~\ref{i:linearuniversal} closes every higher-dimensional
relation and quotient-topology issue for the completion whose
neighborhoods are open left ideals, the intrinsic finite topology of
an endomorphism ring. This is stronger than an adelic action or dense
generation statement: it constructs the unique continuous map to
every complete target in that category. It does not identify an
unspecified completion with this one. The transcription of 43:66
says separated complete graded topological algebra but does not
unambiguously state the open-left-ideal convention; 43:64 also
mentions several possibly incompatible linear topologies. If
``topological algebra'' there is instead taken to allow arbitrary
topologies without a basis of open left ideals, that enlarged
universal problem is not settled by I8. The surface proof I2--I3
already works in that enlarged category.


\begin{proposition}[I7b: strict dense column quotients without completeness]
Let $A$ be a regular noetherian local ring and let $K$ be its residual
resolution, normalized by heights. Put
$C_h^j=\operatorname{Hom}_A(K^h,K^j)$ and give it the finite topology.
For $j>h+1$, postcomposition with the differential induces an open
surjection
\[
 C_h^{j-1}\longrightarrow
 \operatorname{im}(d:C_h^{j-1}\to C_h^j),
\]
where the image has its subspace topology. That image is dense in
$Z_h^j=\ker(d:C_h^j\to C_h^{j+1})$.
\end{proposition}
\begin{proof}
First fix a prime $\mathfrak p$ of height $h$ and write
$C_{\mathfrak p}^j=\operatorname{Hom}_A(E_{\mathfrak p},K^j)$.
Finitely many prescribed inputs lie in a cyclic submodule $N=Ae$
whose annihilator $I$ is $\mathfrak p$-primary and whose localization
is generated by powers of a regular system of parameters of
$A_{\mathfrak p}$. Since $j-1>h$, the finite module
$\operatorname{Ext}_A^{j-1}(A/I,A)$ vanishes after localization at
$\mathfrak p$. Choose $v\notin\mathfrak p$ annihilating it and
put $N'=A(e/v)$.

If $f$ belongs to the image of $d$ and vanishes on $N'$, choose an
existing preimage $g$. Then $g|_{N'}$ is a cycle. Restriction to
$N$ multiplies its cohomology class by $v$, so there is
$b:N\to K^{j-2}$ with $db=g|_N$. Its image has finitely many
injective components; extend them to $E_{\mathfrak p}$ to obtain
$\widetilde b$. The map $g-d\widetilde b$ vanishes on $N$ and
still maps to $f$. Therefore
\[
 \operatorname{im}(d)\cap\operatorname{Ann}(N')
       \subseteq d\bigl(\operatorname{Ann}(N)\bigr),
\]
which proves openness. Membership in the image supplied the
preimage, so completeness of $A$ was not used. Density in cycles
is Proposition~I6.

Finally,
\[
 C_h^j=\prod_{\operatorname{ht}\mathfrak p=h}
             C_{\mathfrak p}^j
\]
topologically: a finite set of inputs in $K^h$ meets finitely many
source summands. The product of open surjections is open, and the
product of dense subgroups is dense. This proves the assertion.
\end{proof}

\paragraph{Exact remaining gate.}
If $Q$ is a separated complete topological abelian group and a
continuous additive map $C_h^{j-1}\to Q$ kills $d(C_h^{j-2})$,
then I6, the proposition, and completeness extend it uniquely to
$Z_h^j$. This is a conditional descent result: it requires a
continuous map on the entire preceding-degree family. It neither
constructs that map for arbitrary adjacent-arrow families nor
controls the assembly of endpoint-localized families in a target
without open left ideals. The first-arrow case is excluded, and
the full broad-topology presentation remains undecided.


\section{The programme of 1964}
\subsection{What the seven statements actually say}
\paragraph{J0 (Textual correction and scope).}
We keep the uncertain exponent in the transcription of Theorem~3
at 44:2 uncertain.  The letter at 44:8 unambiguously writes
``$J^i$ et $J^{n+1-i}$ sont canoniquement duales l'une de l'autre'',
so $n+1-i$ is the mathematically intended complementary index.
This does not resolve an illegible handwritten character by fiat.
The same leaf says that the equivalence relation defining $J^i$ is
one ``que je ne sais pas bien identifier pour l'instant''.
Consequently the literal quotient by the manuscript's Albanese
equivalence is not automatically identified with Murre's modern
universal regular quotient.  That identification needs proof.

For reference, the seven assertions separate as follows.  This table
records present proof status, rather than representing unresolved
assertions as completed items.
\begin{description}
\item[Theorem 1, 44:2.]
{\raggedright The Codex transcription reads ``$N(\zeta)\subset\operatorname{Alb}(S)$
n'est autre \emph{ensemblistement} que le noyau'', while Pro reads
\emph{en substance}; this conflict is retained. The common
incidence-kernel interpretation is proved in J1i below, even
scheme-theoretically. The common kernel for all regular tests also
has the stated set-theoretic property. Identifying either with the
particular undefined manuscript $N(\zeta)$ remains a separate
textual and mathematical gate.\par}
\item[Theorem 2, 44:2.]
The quotient dimensions are bounded, ``en fait, par
$\frac12b^{2i-1}(X)$''.  For a modern algebraic representative in
codimensions $1,2,n$ this bound follows by taking a miniversal
correspondence and its injective rational Tate realization; see
\cite[Proposition 8.2, p.~31]{DGJ_ACMVFunctor}.  Specifically choose
a prime different from the characteristic and the degree of the
miniversal correspondence; the injection has source of rank
$2\dim\operatorname{Ab}^i(X)$ and target of rank $b^{2i-1}$.
For incidence equivalence, Saito's Proposition~3.4 gives this bound
in every codimension; J1i proves the exact family-quotient form.
For the equivalence defined by all regular homomorphisms, the
general uniform-boundedness problem remains unresolved.
\item[Theorem 3, 44:2.]
The quotient is naturally an abelian variety, with complementary
duality.  The incidence quotient exists in every codimension and has
complementary isogeny duality by Saito's Theorem~3.6 and
Proposition~4.4. Existence of the modern universal regular
representative is known in codimensions $1,2,n$; its general
existence, identification with the incidence quotient, and the
claimed integral duality are separate issues.  Under the modern
algebraic-representative interpretation explicitly requested in
J.2, the literal integral duality is false by J2a below, already
in characteristic five for a three-dimensional Jacobian.  This
does not identify the manuscript's less precisely specified
incidence quotient.  General regular representability and that identification remain
unresolved; J1i gives the exact established incidence version and
its specified isogeny normalization.
\item[Theorem 4, 44:2.]
The raw middle self-duality is ``associée à une polarisation''.
This is false, even allowing either overall sign: J2b below.
The primitive qualification in the margin of 44:8 is essential.
There is also an unresolved margin-reading conflict at 44:2:
Codex reads a form of ``Annulé/Cancelé tel quel'', whereas Pro
reads ``Connu tel quel''; this is not silently resolved.
\item[Theorem 5, 44:4.]
The Néron--Tate pairing is $(-1)^mZ\cdot Z'$ for cycles of
codimensions $i,j$ with $i+j=n$, whose restrictions to the generic
member of a hyperplane pencil are algebraically trivial, with all
pencil fibres geometrically integral. This literal formula is false:
J5 below gives a split quadric fourfold in characteristic three and
four explicit planes with both generic restrictions literally zero,
while their signed intersection has degree two. Here $i=j=m=2$,
so the retained sign variable causes no ambiguity. The example
also meets every explicit integral-fibre hypothesis. A stronger
height comparison with admissibility or vertical-correction
conditions is a different statement and is not refuted.
\item[Theorem 6, 44:4.]
The Codex transcription reads \emph{isogénie}, whereas the Pro
transcription reads \emph{isomorphisme}, if $2i\leq n-1$;
without the scan this disagreement must remain explicit.  In
either reading, its image
is the fixed factor if $2i=n$; it is onto if $2i\geq n+1$.
These are assertions about algebraic representatives, not merely
cohomological weak Lefschetz. The literal middle clause is false:
J6C gives a smooth cubic fourfold and a pencil with all fibres
geometrically integral whose target has a constant elliptic factor,
although the source representative is zero. The axis is singular,
which the manuscript does not exclude. This works for both the
incidence and universal regular quotients. J6P proves the
codimension-two fixed-image clause for a smooth pencil axis, and
generic stable-range surjectivity in every characteristic; the
finite-kernel step retains an explicit torsion hypothesis. J6Z
proves surjectivity in the highest nonzero target codimension.
J6L below proves
the stable isogeny under the relevant algebraic Lefschetz lowering
operator on $X$, and surjectivity under the corresponding operator
on $Y$, with no claim that those hypotheses are automatic.
\item[Theorem 7, 44:6.]
``L'équivalence numérique $=$ $\tau$-équivalence'' for all cycles
on a smooth projective variety.  It is false in every intermediate codimension in which such
a counterexample can occur: the explicit Fermat-quartic Ceresa
cycle below is numerically trivial but is not $\tau$-equivalent
to zero.
\end{description}

\subsection{The precise incidence quotient and the remaining regular-quotient problem}
\paragraph{J1i (Known construction, with the family-kernel consequence proved here).}
Let $k$ be an algebraically closed field of any characteristic, let
$X/k$ be smooth projective of pure dimension $n$, and write
$A^i(X)=\operatorname{CH}^i(X)_{\rm alg}$. Define
\[
 I^i(X)=\bigcap_{T,w}\ker\left(
 w_*:A^i(X)\longrightarrow\operatorname{Pic}^0(T)(k)\right),
 \qquad w\in\operatorname{CH}^{n-i+1}(X\times T),
\]
where $T$ runs through smooth projective $k$-varieties; on $k$-points
$\operatorname{Pic}^0(T)$ means its reduced identity component.
This is \emph{incidence equivalence}, also called Picard equivalence
in the older sources. It is not a definition by all regular
homomorphisms to abelian varieties.

Kleiman explicitly separates the two relations in
\cite[pp.~446--447]{DGJ_Kleiman1971}. His p.~446 attributes the
boundedness argument for incidence equivalence to an unpublished
1964 observation of Grothendieck, and p.~447, footnote~5, records
the complementary duality with finite kernel. Thus the historical
connection with 44:2--9 has a primary-source basis, without
identifying two different equivalence relations.
The complete construction is Saito's
\cite[Definition 3.1, p.~108; Proposition 3.4, pp.~109--110;
Theorem 3.6, pp.~110--111]{DGJ_Saito1979}: there is an abelian
variety $P^i_{\rm inc}(X)$ and a surjective regular Picard
homomorphism
\[
 \theta_X^i:A^i(X)\longrightarrow P^i_{\rm inc}(X)(k),
 \qquad\ker\theta_X^i=I^i(X),
 \qquad 2\dim P^i_{\rm inc}(X)\leq b^{2i-1}(X).
\]
Here the Betti number is the dimension of $\ell$-adic cohomology
for any prime $\ell\ne\operatorname{char}k$. The map is universal
among Picard homomorphisms, which are the homomorphisms obtained
by a complementary correspondence to the Picard variety of a
smooth projective variety, with the image viewed as an abelian
subvariety. All these statements hold in every codimension,
including characteristics two and three. Saito's bound is written
using $b^{2n-2i+1}$; Poincaré duality gives the displayed index.

Here is the exact family-kernel statement which follows. Let
$(S,s_0)$ be smooth projective and pointed, and let
$\xi\in\operatorname{CH}^i(S\times X)$. Correspondence
functoriality supplies a homomorphism of $k$-abelian varieties
\[
 f_\xi:\operatorname{Alb}(S)\longrightarrow P^i_{\rm inc}(X),
 \qquad
 f_\xi([s]-[s_0])=\theta_X^i(\xi_s-\xi_{s_0}).
\]
For each $w$ in the definition of $I^i(X)$, composition with
$\xi$ is a divisor correspondence on $S\times T$ and defines
$g_{w,\xi}:\operatorname{Alb}(S)\to\operatorname{Pic}^{\rm var}(T)$.
Define $N_{\rm inc}(\xi)$ as the scheme-theoretic intersection of
all their kernels. Then
\[
 N_{\rm inc}(\xi)=\ker f_\xi,
 \qquad
 \operatorname{Alb}(S)/N_{\rm inc}(\xi)
     \simeq\operatorname{Im}(f_\xi),
 \qquad
 \dim\operatorname{Alb}(S)/N_{\rm inc}(\xi)
       \leq\tfrac12 b^{2i-1}(X).                 \tag{J1i.1}
\]
The quotient and the image here are abelian varieties; the kernel
may be nonreduced. To prove the equality, universality makes every
$g_{w,\xi}$ factor through $f_\xi$, giving one inclusion.
Conversely, $\theta_X^i$ is itself a Picard homomorphism
\cite[(3.5), p.~110]{DGJ_Saito1979}. Thus for some $T_0,w_0$,
its composition with a closed immersion
$P^i_{\rm inc}(X)\hookrightarrow\operatorname{Pic}^{\rm var}(T_0)$
is $(w_0)_*$. The corresponding $g_{w_0,\xi}$ has kernel exactly
$\ker f_\xi$, giving the other inclusion as group schemes, not
only on geometric points. The fppf quotient by a homomorphism's
kernel is its abelian-variety image. The dimension assertion now
follows from Saito's bound. This proves (J1i.1).

Saito further constructs, using his explicitly specified
miniversal-cycle integers, an isogeny
\[
 P^i_{\rm inc}(X)\longrightarrow
          P^{n+1-i}_{\rm inc}(X)^\vee;
\]
see \cite[(4.1)--(4.2), p.~112; Proposition 4.4 and
Corollary 4.5, p.~113]{DGJ_Saito1979}. This is a known
\emph{isogeny} statement. Its integer factors must be retained;
it does not supply the unscaled integral intersection isomorphism
asserted in J.2. Therefore the incidence reading of the
representability and boundedness portions of Theorems~1--3 is
settled by the cited construction and (J1i.1). The stronger
integral duality and the identification with a quotient defined
by all regular maps are separate assertions.

\paragraph{The regular interpretation and its exact remaining gate.}
A homomorphism $h:A^i(X)\to B(k)$ to a $k$-abelian variety is
\emph{regular} when every pointed smooth projective family of
cycles induces a morphism to $B$. Put
\[
 I^i_{\rm ab}(X)=\bigcap_h\ker h.
\]
This is Samuel's abelian equivalence. Every Picard homomorphism is
regular, so
$I^i_{\rm ab}(X)\subseteq I^i(X)$. A universal regular
homomorphism, if it exists, therefore has a canonical surjective
homomorphism from its abelian variety to $P^i_{\rm inc}(X)$.
This observation does not prove that the kernel is finite or zero.

For a fixed family $\xi$, let $N_{\rm ab}(\xi)$ be the
scheme-theoretic intersection of the kernels of all homomorphisms
$h_\xi:\operatorname{Alb}(S)\to B$ obtained from regular $h$.
Then the induced group homomorphism
\[
 \operatorname{Alb}(S)(k)\longrightarrow
                    A^i(X)/I^i_{\rm ab}(X)
\]
has kernel $N_{\rm ab}(\xi)(k)$, without assuming a universal
regular representative. Indeed, every composite on degree-zero
zero-cycles of $S$ is regular and factors through its Albanese;
hence the Albanese kernel is killed modulo $I^i_{\rm ab}(X)$.
The displayed map is well defined because the Albanese map on
zero-cycles is surjective over an algebraically closed field, and
its kernel is the stated intersection by definition. The
intersection is a closed subgroup scheme: on a noetherian
scheme the sum of the defining ideal sheaves is generated by a
finite subfamily. Thus each individual family has an abelian
quotient. What is missing is a uniform bound as the family varies.

More precisely, the existence of such a uniform dimension bound
is equivalent to existence of a universal regular homomorphism.
The bounded-target criterion is
\cite[Theorem 2.2(ii), p.~103 and proof pp.~103--106]{DGJ_Saito1979}
for the class of all regular homomorphisms, identified in his
Example (1.3.1), p.~101. To apply it to the family formulation,
any surjective regular $h$ admits a family on an abelian variety
whose induced map to its target is surjective
\cite[Proposition 1.2, pp.~98--100]{DGJ_Saito1979}; this target
is a quotient of the corresponding family quotient, so a uniform
family bound bounds all regular targets. Conversely, a universal
regular representative is itself among the tests defining every
$N_{\rm ab}(\xi)$, so each such quotient is its abelian
subvariety image. This proves the claimed equivalence of the two
boundedness problems. A bound by the particular Betti number is
stronger than mere existence.

The known existence cases are $i=1,2,n$.
For codimension two the complete repair of Murre's proof is
\cite[Theorem 1 and its two proofs, author version
pp.~2--4]{DGJ_Kahn2021}; for all three indices the Betti bound is
\cite[Proposition 8.2, p.~31]{DGJ_ACMVFunctor}.
That proposition explicitly restricts its cohomological injection
to $i=1,2,n$. Applying it in arbitrary codimension would assume
the missing injectivity of the higher Bloch map and would falsely
close the remaining gate.

A concrete unresolved test is Kahn's Conjecture~1, author version
p.~5: take a generic abelian threefold $A$ defined over a field
$k_0$ finitely generated over $\mathbb Q$, a nonisotrivial elliptic
curve $E/k_0(t)$, and
\[
 V=(A_{k_0(t)}\times E)_{\overline{k_0(t)}}.
\]
The conjecture states that a universal regular homomorphism on
$A^3(V)$ does not exist. Its preceding infinite-torsion result
alone does not prove this conjecture. Although every surjective
regular homomorphism is surjective on prime-to-characteristic
primary torsion \cite[Remark 2, p.~4]{DGJ_Kahn2021}, an infinite
torsion kernel may be killed in every finite-dimensional target.
To refute universal existence one must still construct regular
quotients of unbounded dimension, or another actual obstruction
satisfying the regularity condition. Neither the torsion argument
nor the opened motivic constructions provides that step here.
Accordingly the all-regular version in codimension three on
fourfolds, and its requested general form, remain unresolved.

\paragraph{A single self-incidence kernel is not the common kernel.}
The letter at 44:9 constructs $N$ from one symmetric
self-correspondence on a parameter variety. Such a kernel cannot
be substituted for $N_{\rm inc}(\xi)$ without an additional
argument. The distinction is visible in the same explicit
threefold as J2b. Let $X=E\times\mathbb P^1\times\mathbb P^1$,
let $S=E^\vee$, and set $\xi=c_1(\mathcal P)h_1$, where
$\mathcal P$ is the normalized Poincaré bundle. Its
self-incidence divisor is zero because $h_1^2=0$, so its
self-kernel is all of $S$. But pairing with the independent
family $c_1(\mathcal P)h_2$ gives the elliptic duality
$S\to S^\vee$, up to sign, which is an isomorphism. Hence
$N_{\rm inc}(\xi)=0$. Equivalently the family map is the
inclusion of the first elliptic factor in $P^2_{\rm inc}(X)=S^2$.
The projective-bundle decomposition and the two complementary
Poincaré tests prove this identification integrally, just as in
J2b. Thus one self-pairing can miss an entire abelian family that
other incidence tests detect. This refutes that particular
shortcut from 44:9 to Theorem~1; it does not decide an undefined
$N(\xi)$ on 44:2 by assigning it a new definition.

\paragraph{Exact status of the remaining Lefschetz assertions.}
The literal middle clause of Theorem~6 is refuted by J6C below.
The stable- and high-range clauses remain unresolved in their
full stated generality. Even after choosing the incidence quotient,
they concern algebraic correspondences and abelian varieties.
Cohomological weak Lefschetz identifies the ambient
cohomology in its stable range, but does not show that an inverse
on the required algebraic part is induced by correspondences.
For the universal regular quotient, general existence is already
the preceding gate. For the literal manuscript, the two
transcriptions additionally disagree between an isomorphism and
an isogeny in the stable range. Neither a cohomological
isomorphism nor Saito's complementary isogeny proves the stronger
integral restriction claim. These are retained as explicit
unresolved assertions. J6L, J6P, and J6Z establish the precisely
stated conditional and restricted cases below.

\subsection{An explicit counterexample to raw positivity}
\paragraph{J2b (Refuted, with a complete counterexample).}
Let $k=\overline{\mathbb F}_3$ and let
\[
 E:\ y^2z=x^3-xz^2\subset\mathbb P^2_k,
 \qquad X=E\times\mathbb P^1\times\mathbb P^1.
\]
The curve is nonsingular: in characteristic three the derivative
of $y^2z-x^3+xz^2$ with respect to $x$ is $z^2$; if all partials
vanished, $z=0$, the equation would force $x=0$, and the derivative
with respect to $z$ would then force $y=0$, which is not a
projective point.  Its point at infinity is the origin.  Thus $X$
is a smooth projective threefold satisfying exactly the hypotheses
of J.2.  In fact the same construction works with any elliptic
curve in every characteristic.

Let $h_1,h_2$ be the two hyperplane classes.  The projective bundle
formula for Chow groups \cite[Tag 02TX]{DGJ_StacksPB} gives
\[
 \operatorname{CH}^2(X)=
   \operatorname{Pic}(E)h_1\oplus\operatorname{Pic}(E)h_2
          \oplus\mathbb Z h_1h_2.
\]
Consequently its algebraically trivial subgroup is
$\operatorname{Pic}^0(E)h_1\oplus\operatorname{Pic}^0(E)h_2$.
The decomposition and its inverse are given by multiplication by
the hyperplane classes and proper pushforward along the
projections.  They apply to families and preserve regular
homomorphisms.  The universal property of the Picard variety
therefore proves, as an algebraic representative and not just as
an abstract group of points,
\[
 \operatorname{Ab}^2(X)=E^\vee\times E^\vee\simeq E^2.
\]
The last identification uses the standard principal polarization
of the elliptic curve.

Write $T=(E^\vee)^2$, with origin $a=0$, and let
$\mathcal P_1,\mathcal P_2$ be the pullbacks to $T\times E$ of
the normalized Poincaré bundle on $E^\vee\times E$.
The completely specified family
\[
 z=c_1(\mathcal P_1)h_1+c_1(\mathcal P_2)h_2
       \in\operatorname{CH}^2(T\times X)
\]
has $z(0)=0$ and induces the identity on the displayed algebraic
representative.  Form exactly the divisor correspondence in the
letter at 44:9:
\[
 D=r_*(p^*z\cdot q^*z)\in\operatorname{CH}^1(T\times T).
\]
Because $h_1^2=h_2^2=0$ and $\int_{\mathbb P^1\times\mathbb P^1}
h_1h_2=1$, the two diagonal entries of its homomorphism
$\rho:T\to T^\vee$ vanish and the two off-diagonal entries are
the elliptic intersection duality.  Thus, up to one overall sign
depending on the Poincaré convention,
\[
 \rho=\begin{pmatrix}0&\lambda_E\\\lambda_E&0\end{pmatrix}.
\]
There is no hidden scalar two in either entry: the two cross terms
give the two distinct off-diagonal entries, once each.
This homomorphism is an isomorphism, with inverse obtained by
interchanging the two dual factors using $\lambda_E^{-1}$.

Pull it back along the diagonal $u\mapsto(u,u)$ and the
antidiagonal $u\mapsto(u,-u)$ of $E^2$.  The resulting symmetric
homomorphisms are respectively $2\lambda_E$ and $-2\lambda_E$,
with both signs reversed if $-\rho$ is used.  The restriction of
an ample line bundle to either elliptic subgroup must have positive
degree.  One of these two degrees is negative for either overall
sign.  Hence neither $\rho$ nor $-\rho$ is a polarization.
The same argument rules out positivity of a rational
Néron--Severi class inducing this correspondence.  Since $\rho$
has zero kernel, passing to the letter's quotient $T/\ker\rho$
does not change the counterexample.

This refutes the literal unqualified Theorem~4 and the positivity
assertion in J.2, including the proposed single sign from
Theorem~5.  It does not refute the primitive correction which
Grothendieck himself adds at 44:8:
``il faut tenir compte de la décomposition primitive''.
Indeed, relative to the ample class $h_1+h_2$ on the surface
factor, the sum and difference directions have opposite
intersection signs.  Their Gram matrix is
$\operatorname{diag}(2,-2)$; changing the sign separately on the
primitive direction makes it positive.  The sum/difference change
of basis has determinant $-2$, so this observation is an isogeny
decomposition, not an integral direct-sum assertion at two.

The supplied program verifies smoothness of the displayed elliptic
curve, the intersection matrix, its determinant $-1$, and the
sum/difference diagonalization.  Its finite computations certify
exactly these algebraic assertions; the universal-property and
ampleness arguments are given above.

\paragraph{Known scope and a discarded integral-lattice route.}
The preceding example refutes positivity while possessing an
intersection isomorphism.  We next give a separate counterexample
to the universal isomorphism assertion itself.
For comparison, Achter--Casalaina-Martin--Vial prove a canonical
symmetric purely inseparable isogeny for geometrically rationally
chain connected threefolds, and an isomorphism for geometrically
stably rational threefolds; see \cite[Theorem 1, p.~6 and
Theorem 12.12, p.~54]{DGJ_ACMVDiagonal2025}.  Their Question 13.4,
p.~57, concerns positivity in the remaining nonliftable setting.
Those are stronger-hypothesis statements; they do not establish
the assertion of J.2 for arbitrary threefolds.

We tested a possible integral obstruction using
$E\times\operatorname{Km}(F\times G)$ over a finite field,
with distinct ordinary elliptic isogeny factors.  Its
Néron--Severi lattice can have discriminant sixteen, so the raw
cohomological restriction is not unimodular.  The attempted
deduction fails at identification of the integral Tate lattice
of $\operatorname{Ab}^2$: the assertion that the universal regular
homomorphism is an isomorphism on all prime-to-characteristic
torsion is a named \emph{standard assumption} in
\cite[Definition 6, p.~4; Proposition 3.8, p.~17]{DGJ_ACMVBloch},
not an unconditional theorem for all smooth threefolds in
positive characteristic.  The universal integral assertion in
the sweep's discussion of 44:8 must therefore not be used.
This route remains conditional; it is not a counterexample.


\subsection{A generic hyperelliptic Jacobian obstructs integral autoduality}
\paragraph{J2a (Refuted, with a complete counterexample).}
Let $a_0,\ldots,a_6$ be algebraically independent over $\mathbb F_5$,
put $K=\mathbb F_5(a_0,\ldots,a_6)^{\mathrm{perf}}$ and
$k=\overline K$, and let
\[
 C/K:\quad y^2=x^7+a_6x^6+\cdots+a_1x+a_0,
 \qquad A=\operatorname{Jac}(C),\qquad X=A_k.
\]
Here $C$ means the smooth projective model and $\infty$ is its
unique point above infinity.  Its genus is three, and $X$ is a
smooth projective threefold in characteristic five.  We prove
that there is no isomorphism from $\operatorname{Ab}^2(X)$ to its
dual induced by the integral intersection of families of cycles.
We use the precise required compatibility: for every family
$z\in\operatorname{CH}^2(T\times X)$, $z(0)=0$, on an abelian
parameter variety $T$, if $u_z:T\to\operatorname{Ab}^2(X)$ is its
regular homomorphism, the homomorphism defined by
$r_*(p^*z\cdot q^*z)$ must equal $u_z^\vee\rho u_z$.
This is exactly the construction in 44:9; no arbitrary rescaling
of intersection is allowed.

Write $B=\operatorname{Ab}^2(A/K)$.  This abelian variety exists,
and its base change to $k$ is the usual algebraic representative,
by \cite[Theorem 5.10, pp.~25--26; Theorem 6.1, pp.~26--27]{DGJ_ACMVFunctor}.
The compatibility defining $\rho$, if satisfied, determines it
uniquely: applying it to a miniversal family on $B$ gives
$[r]^\vee\rho[r]$ for some $r\ne0$, and $\operatorname{Hom}(B,B^\vee)$
is torsion-free.  Therefore every Galois conjugate of $\rho$ is the
same $\rho$.  It descends to $K$.  The perfectness of $K$ makes
$k/K$ separable, so the stated base-change and descent theorem
applies without an inseparable-field qualification.

\paragraph{The dimension of $B$ is at most three.}
Choose a miniversal family for $B$, and a prime $\ell\ne2,5$
not dividing its degree.  Its rational Tate realization is
injective by \cite[Proposition 8.2, p.~31]{DGJ_ACMVFunctor}.
After dualizing and applying Poincaré duality on the threefold,
this makes $H^1(B,\mathbb Q_\ell)$ a Galois quotient of
$H^3(A,\mathbb Q_\ell)(1)$.
The geometric Zariski monodromy on $V=H^1(A,\mathbb Q_\ell)$ is
$\operatorname{Sp}_6$.  Indeed, this is
\cite[Theorem 3.4 and Corollary 3.5, pp.~13--14]{DGJ_AchterPries2006}
for the hyperelliptic moduli space.  Here is the precise transfer to our coefficient space
$S=\operatorname{Spec}\mathbb F_5[a_0,\ldots,a_6,\Delta^{-1}]$.
Let $\mathcal H_3^w$ be the stack of smooth hyperelliptic
curves with one marked Weierstrass point.  The morphism
$\mathcal H_3^w\to\mathcal H_3$ forgetting the point is finite
étale of degree eight: its pullback to any family is that
family's reduced branch divisor.  Thus its geometric fundamental
group has finite-index image.  Finite-index subgroups of a
Zariski-dense subgroup of the connected group
$\operatorname{Sp}_6$ are still Zariski dense.

There is, more precisely, a presentation
$\mathcal H_3^w=[S/G]$ with connected smooth group
\[
 G=\{(b,a,c):b\in\mathbb G_a,\ a,c\in\mathbb G_m,
                    \ c^2=a^7\}
       \simeq\mathbb G_a\rtimes\mathbb G_m.
\]
The coordinate changes are $x\mapsto ax+b$, $y\mapsto cy$;
the monic leading coefficient imposes $c^2=a^7$.
The isomorphism on the multiplicative factor is
$s\mapsto(a,c)=(s^2,s^7)$, with inverse $s=c/a^3$.
Every pointed hyperelliptic curve has this form fppf locally:
the marked branch point trivializes the quotient conic locally,
we choose it as infinity, and we choose the anti-invariant
coordinate $y$ and normalize its leading coefficient.  Any
isomorphism of two such presentations has exactly the displayed
form, since functions with a pole of order at most two at the
marked point are spanned by $1,x$, and the anti-invariant
functions with pole order at most seven are spanned by $y$.
This proves the claimed presentation, also for families by the
corresponding locally free sheaves and base change.
The atlas $S\to[S/G]$ has geometrically connected fibres $G$,
so it surjects on geometric fundamental groups: a connected
finite étale cover of the stack stays connected after this
pullback, which is the finite-cover criterion for surjectivity.
Consequently the geometric Zariski monodromy on $S$ is
$\operatorname{Sp}_6$.  Removing a proper closed subset of
normal $S$ gives a surjection on fundamental groups, and
purely inseparable extension to $K$ leaves the étale Galois
representations unchanged.

The symplectic exterior-cube decomposition is
\[
 \bigwedge^3V(1)=V\oplus P,
 \qquad \dim V=6,\quad \dim P=14,
\]
where the displayed $V$ denotes its Lefschetz copy and $P$ is
primitive.  Both summands are absolutely irreducible for the
geometric symplectic group and are preserved by the arithmetic
Galois action.  A direct verification of the only less familiar
assertion is as follows.  With symplectic basis $e_i,f_i$, $P$
is the kernel of symplectic contraction.  Its fourteen distinct
torus weights consist of the eight $\pm\epsilon_1\pm\epsilon_2
\pm\epsilon_3$ and the six $\pm\epsilon_i$.  The first weight
vectors are wedges choosing one of $e_i,f_i$ for each $i$; the
second are $v_i\wedge(e_j\wedge f_j-e_h\wedge f_h)$ for
$\{i,j,h\}=\{1,2,3\}$ and $v_i=e_i$ or $f_i$.
Symplectic permutations and interchanges generate all eight
vectors of the first kind from any one.  For an explicit root operation, let $N(e_2)=e_3$,
$N(f_3)=-f_2$, and let $N$ vanish on the remaining basis
vectors.  It belongs to the symplectic Lie algebra and sends
$e_1\wedge e_2\wedge f_3$ to
$-e_1\wedge(e_2\wedge f_2-e_3\wedge f_3)$.
The opposite root $N'(e_3)=e_2$, $N'(f_2)=-f_3$ sends the
latter second-kind vector to $-2e_1\wedge e_2\wedge f_3$.
Permuting and interchanging symplectic pairs gives all six
second-kind vectors and the converse transitions.  Thus a nonzero
subrepresentation, which contains a torus weight vector, contains
all fourteen vectors.  This proves irreducibility over
$\mathbb Q_\ell$, where $2\ne0$.
Consequently the quotient $H^1(B)$ can have only the constituents
$V$ and $P$, each at most once.

The constituent $P$ cannot occur.  We give the spreading
step explicitly.  Write
$K_n=\mathbb F_5(a_0^{1/5^n},\ldots,a_6^{1/5^n})$, so
$K=\bigcup_nK_n$.  The abelian variety $B$, its group laws,
and the finitely many components and coefficients of a
miniversal cycle $\Gamma\in\operatorname{CH}^2(B\times A)$
are of finite presentation; hence they descend to some $K_n$.
The model of $A$ is the Frobenius pullback of its model over
$S$.  After removing a proper closed subset of this Frobenius
pullback $S_n$, the descended $B$ extends to an abelian scheme
$\mathcal B\to U$, and the closures of the cycle components
give a cycle on $\mathcal B\times_U\mathcal A$.  To obtain
the abelian scheme, extend projective equations and the group
laws and their inverse, shrink so that the group identities
hold, and shrink again to the smooth locus with geometrically
connected fibres.  Properness is retained by the projective
model.  No extension of the universal regular property to every
fibre is being assumed.

The generic Tate injection of Proposition~8.2 is equivariant
for $\operatorname{Gal}(K_n)$: purely inseparable base change
identifies this group and its étale cohomology with those over
$K$.  Both its source and target are the generic fibres of
lisse $\mathbb Q_\ell$-sheaves on $U$, and therefore factor
through $\pi_1(U)$.  Since
$\operatorname{Gal}(K_n)\to\pi_1(U)$ is surjective, their
equivariant map is a map of those lisse sheaves.  Its rank is
locally constant.  Dualizing gives a surjection on every fibre
of $U$, just as at the generic point.  Its kernel is generically
one of the sums of the two intrinsic Lefschetz summands;
equality of morphisms of lisse sheaves at a geometric generic
point is equality everywhere, so its fibrewise quotient has the
same summands.  The Frobenius pullback and shrinking preserve
connected geometric monodromy by the preceding fundamental-group
argument.  The ordinary locus of the Jacobian is
nonempty.  For an explicit check, the smooth curve
$y^2=x^7-x$ over $\mathbb F_5$ has Cartier--Manin matrix
\[
 \begin{pmatrix}0&0&1\\0&3&0\\1&0&0\end{pmatrix},
 \qquad\det=2\pmod5.
\]
This follows directly from applying Cartier to
$x^{j-1}dx/y=y^{-5}x^{j-1}(x^7-x)^2dx$:
the $(i,j)$ entry is the coefficient of $x^{5i-j}$ in
$(x^7-x)^2$.  Thus Cartier is invertible, so the Jacobian is
ordinary.  Nonvanishing of the same determinant defines a
nonempty open ordinary locus in the coefficient space.  It
meets the open set where the preceding objects have spread;
choose a closed point there, over a finite field $\mathbb F_q$.

For clarity, the passage from the Cartier calculation to
Frobenius slopes can be made entirely in terms of characteristic
polynomials.  The Manin congruence in
\cite[\S2.5, p.~7; Cartier formula in \S3.1, pp.~7--8]{DGJ_AchterHowe2018}
says that the Frobenius polynomial modulo five is $T^3$ times
the characteristic polynomial of an invertible iterate of the
Cartier operator (equivalently its adjoint Frobenius operator).
That second polynomial has nonzero constant coefficient.
Hensel factorization therefore gives exactly three unit roots.
The principal polarization pairs each root $\alpha$ with
$q/\alpha$; see also \cite[Lemma 19.2, p.~41]{DGJ_MilneAV}.
Hence the remaining three roots have normalized valuation one.
Thus the six Frobenius eigenvalues on $H^1(A)$ have three
$q$-adic slopes zero and three slopes one.  Their triple
products on $H^3(A)(1)$ include one slope $-1$ and one slope $2$.
The Lefschetz summand has only slopes zero and one, so both
extreme slopes belong to $P$.  The surjection on $\ell$-adic fibres gives equality of
Frobenius characteristic polynomials for $H^1(B)$ and the
indicated summand or sum of summands.  These polynomials have
rational coefficients: the primitive polynomial is the quotient
of the exterior-cube polynomial by the Lefschetz polynomial.
Equality in $\mathbb Q_\ell[T]$ is therefore equality in
$\mathbb Q[T]$.  Embed their algebraic roots in
$\overline{\mathbb Q}_5$ and normalize $v(q)=1$.
If $P$ occurred in the quotient $H^1(B)$, its negative-valuation
root would then be a Frobenius eigenvalue on $H^1(B)$.  This
contradicts the algebraic integrality of Frobenius eigenvalues
on abelian varieties, as stated in
\cite[Theorem 19.1 and its proof, pp.~40--41]{DGJ_MilneAV}.
This argument compares characteristic polynomials; it makes no
assertion that the correspondence admits a particular crystalline
realization.  Hence only the
six-dimensional constituent can occur, and $\dim B\leq3$.
The finite computation records the full primitive slope
multiplicities $1,6,6,1$ in slopes $-1,0,1,2$; the exclusion uses
only the first of these.

\paragraph{The translation family has intersection $\pm2\lambda_A$.}
Embed $C$ by $Q\mapsto[Q-\infty]$ and let
\[
 z_t=[C+t]-[C],\qquad t\in A.
\]
Equivalently $z=m_-^*[C]-\operatorname{pr}_X^*[C]$ on
$A\times A$, where $m_-(t,x)=x-t$.  This is a specified family
of algebraically trivial codimension-two cycles with $z_0=0$.
It induces a $K$-homomorphism $u:A\to B$.
Let $\theta$ denote the cohomology class of the principal theta
polarization.  The Poincaré formula is $[C]=\theta^2/2$ in
cohomology.  It follows as well by characterizing $[C]$ by
$\int_A[C]\alpha\beta=\int_Ca^*\alpha\,a^*\beta$ and the
symplectic pairing defining the principal Jacobian polarization.
With symplectic basis $e_i,f_i$ and
$\theta=\sum e_i\wedge f_i$, the part of the class of $z$ in
$H^1(A_{\rm par})\otimes H^3(A_{\rm space})$ is obtained by
expanding
\[
 \frac12\left(\sum_{i=1}^3
 (e_{i,x}-e_{i,t})\wedge(f_{i,x}-f_{i,t})\right)^2
       -\frac12\left(\sum e_{i,x}\wedge f_{i,x}\right)^2.
\]
All terms of parameter degree zero cancel.  Therefore in
$r_*(p^*z\cdot q^*z)$ only the product of the two parameter-degree-one
parts contributes.  Wedge multiplication and integration of
$e_{1,x}f_{1,x}e_{2,x}f_{2,x}e_{3,x}f_{3,x}$ give
\[
 -2\sum_{i=1}^3
 (e_{i,t}\wedge f_{i,s}-f_{i,t}\wedge e_{i,s}).
\]
This is $-2$ times the normalized Poincaré pairing of the
principal polarization, with the sign reversed if the opposite
Poincaré convention is chosen.  Faithfulness of the Tate
realization on homomorphisms of abelian varieties
\cite[Theorem 12.5, pp.~21--22]{DGJ_MilneAV} gives the
actual homomorphism $\pm2\lambda_A$, not merely its isogeny class.
The wedge calculation is also saved as an exact integer
calculation in \texttt{check\_generic\_jacobian.py}.

\paragraph{The impossible degree.}
If an intersection-induced isomorphism $\rho:B_k\to B_k^\vee$
existed, its defining compatibility and the preceding calculation
would give
\[
 u^\vee\rho u=\pm2\lambda_A.
\]
The right side is an isogeny.  Thus $u$ has finite kernel;
since $\dim B\leq3=\dim A$, it is an isogeny onto $B$.
Taking degrees yields
\[
 (\deg u)^2=\deg(2\lambda_A)=2^6,
 \qquad\deg u=8.
\]
We show that the $K$-abelian variety $A$ admits no isogeny of
degree eight.  The splitting field of its generic septic has
Galois group $S_7$: equivalently, start with seven independent
roots and take the invariant field of their permutations.
Passing to the perfect closure is purely inseparable and leaves
this group unchanged.  If $r_1,\ldots,r_7$ are the roots and $W_1,\ldots,W_7$ are the finite
Weierstrass points, the elements $[W_i-\infty]$ generate $A[2](k)$,
with their sum as the only relation.  To check the last assertion,
a putative proper-subset relation would make
$\prod_{i\in I}(x-r_i)$ a square in $k(C)$.  In the quadratic
extension $k(C)/k(x)$, a rational function whose square lies in
$k(x)$ is in $k(x)$ or in $y\,k(x)$.  The odd valuations at the
chosen roots exclude the first case unless $I$ is empty and
exclude the second unless $I$ is the full set.  There are
therefore six independent classes, as required by the rank of
$A[2]$.

It follows that $A[2](k)$ is the six-dimensional $S_7$-module
$\mathbb F_2^7/\langle(1,\ldots,1)\rangle$, equivalently the
subspace of vectors with even coordinate sum.  This module is
irreducible.  Indeed any nonzero even-sum vector has both a zero
and a one coordinate; subtracting its image under the
transposition of those coordinates yields a vector supported
on exactly two positions.  The $S_7$-orbit of that vector spans
the whole even-sum subspace.  Every nonzero invariant subspace
is consequently the whole module.

The kernel of a degree-eight isogeny in characteristic five is
an étale group of order eight.  Its subgroup killed by two is a
nonzero Galois-stable subgroup of $A[2](k)$, and has at most
eight elements.  Irreducibility would force it to be all of
$A[2](k)$, which has sixty-four elements.  This contradiction
proves the claimed impossibility of an integral
intersection-induced isomorphism.

This argument uses only the rational miniversal Tate embedding,
not the unproved integral torsion standard assumption.  It
rules out the raw integral autoduality demanded in J.2.  It does
not rule out a rational pairing, an isogeny, a polarization after
an appropriate Lefschetz normalization, or an unrelated abstract
isomorphism of $B$ with its dual.


\subsection{An explicit counterexample to numerical equals tau equivalence}
\paragraph{J1-Th7 (Refuted).}
The manuscript statement at 44:6 is
``L'équivalence numérique $=$ $\tau$-équivalence''.
Its valid general direction is that $\tau$-equivalence implies
numerical equivalence: intersection numbers are constant in a
connected algebraic family, and dividing by a nonzero integer
does not change vanishing of integer intersection numbers.
The reverse implication fails already for codimension two on
an abelian threefold over $\overline{\mathbb Q}$.

Here are all data of one counterexample.  Put
\[
 k=\overline{\mathbb Q},\qquad
 C=V(x^4+y^4-z^4)\subset\mathbb P^2_k,\qquad
 P=(1:0:1),\qquad X=\operatorname{Jac}(C).
\]
The three partial derivatives have no simultaneous projective
zero, so $C$ is a smooth plane quartic of genus three.  In
particular $X$ is a smooth projective threefold.  Let
$a_P:C\hookrightarrow X$ be $Q\mapsto[Q-P]$, and define
\[
 Z=(a_P)_*[C]-[-1]_* (a_P)_*[C]\in\operatorname{CH}^2(X).
\]
The action of $[-1]$ on $H^4(X,\mathbb Q_\ell(2))$ is $+1$,
since the cohomology algebra of an abelian variety is the
exterior algebra of $H^1$ and $[-1]$ acts by $-1$ on $H^1$.
Thus the class of $Z$ is zero.  Compatibility of cycle classes
with products and degree gives $\deg(Z\cdot D)=0$ for every
divisor $D$ on $X$.  This is precisely numerical triviality
of this codimension-two cycle.

Bloch's \emph{Algebraic cycles and values of L-functions},
introduction p.~95 and Theorem~(4.1), p.~104
\cite{DGJ_Bloch1984}, proves that the class of this very
Fermat-quartic cycle has infinite order modulo algebraic
equivalence.  His definition of that quotient is on p.~94,
and p.~95 states that its class is independent of the chosen
Abel embedding.  The proof on p.~104 works over an algebraic
closure of the number field, so it applies to the field specified
above; no passage to a generic curve or unspecified very general
parameter is involved.  Consequently there is no nonzero integer
$N$ for which $NZ$ is algebraically trivial.  Equivalently,
$Z$ is not $\tau$-equivalent to zero.  The use of the published
theorem is exactly the non-algebraicity step; the data,
smoothness, homological triviality, and numerical consequence
have been supplied explicitly.

For completeness the same example proves strictness in every
intermediate codimension $2\leq i\leq n-1$ in every dimension
$n\geq3$.  Put $m=n-3$, $r=i-2$, $X_n=X\times\mathbb P^m$,
and $Z_i=\operatorname{pr}_X^*Z\cdot h^r$.  It is numerically
trivial by the projective bundle formula.  If a nonzero multiple
of $Z_i$ were algebraically trivial, intersecting it with
$h^{m-r}$ and pushing forward to $X$ would show that the same
multiple of $Z$ were algebraically trivial, a contradiction.
The operations used preserve algebraic equivalence.  Hence the
correct unrestricted comparison is a strict implication in
these intermediate degrees, not equality.

The saved calculation \texttt{check\_fermat.py} verifies the
singular ideal is the unit ideal in each of the three projective
charts, verifies the specified point, and records genus three.
It does not purport to compute Bloch's infinite-order theorem.
The original scanned source pages were rendered and visually
read, including pp.~94--95 and 104; the complete volume PDF
has an offset of seven pages from its printed numbering.


\subsection{SA19 H6: what the planning leaves establish}
\paragraph{H6a (Textual and variance correction, proved).}
The morphism from an open curve is explicitly present:
43:27 boxes $\varphi_S:X-S\to\mathscr P^1_{S,X/Y}$.
At 43:28 its universal property is called the theorem of
Rosenlicht.  Corollary~3 at 43:30 states an ``hom. canonique
\emph{injectif}''
\[
 \operatorname{Ext}^1(\mathscr P_{S,X/Y},G)
        \longrightarrow H^1(X-S,G).
\]
The next sentence says only ``Il \emph{semble} que c'est
\emph{[bijectif|surjectif]}'' when $S\to Y$ is surjective.
The supplied transcription leaves the bracketed word uncertain;
either reading describes the conjectured stronger direction.
43:31 proposes reductions, and 43:34 writes ``et sans doute''
for the higher-degree comparison.  These qualifications must be
retained in a documentary assessment.

The sweep's stated inference to injectivity on abelian
fundamental groups has its variance reversed.  When the two
sides are interpreted as finite-constant torsors, the map is
pullback of characters:
\[
 \operatorname{Hom}_{\rm cont}(\pi_1(J)^{\rm ab},M)
 \longrightarrow
 \operatorname{Hom}_{\rm cont}(\pi_1(U)^{\rm ab},M).
\]
Its injectivity for every finite abelian $M$ detects
\emph{surjectivity} of $\pi_1(U)^{\rm ab}\to\pi_1(J)^{\rm ab}$.
Indeed a proper closed image in a profinite abelian group is
detected by a nonzero character of a finite quotient, and
conversely a surjection gives injective pullback.  It says
nothing by itself about the kernel.  The quotient
$\mathbb Z/4\to\mathbb Z/2$ has injective pullback of characters
for every finite coefficient group and a nonzero kernel.
If the stronger bijectivity were separately proved, it would
yield an isomorphism of the represented profinite groups.
Neither conclusion is a point-to-section injectivity theorem.

\paragraph{H6b (Complete scoped documentary audit).}
The exact lead extends its negative finding to the full range
43:1--208. The audit now covers all 195 supplied nonblank
transcription files: DGJ read 184 complete files, and FHI
explicitly certified a complete rereading of the eleven files
43:56--66. No readable rational-point-to-section injectivity
argument was found. The saved per-file inventory records source
paths, SHA-256 hashes, topic groups, readers, and uncertainty
counts in \texttt{work/dgj/outputs/c43\_audit.tsv}; the complete
scope discussion is \texttt{work/dgj/audit\_SA19H6.md}.
The earlier 71 keyword matches served only to plan the reading.

The near matches have different subjects. At 43:77 a rational
section of a torsor extends; at 43:103 the reused printed text
concerns cross-sections of Stiefel manifolds. At 43:145 the
injectivity is a localization statement in $H^1$. At
43:167--170 a normalized torsor section proves injectivity from
group extensions to primitive torsor classes. At 43:204--205,
sections of a dual abelian scheme define the Tate pairing;
43:207 uses the Albanese map to define a zero-cycle kernel and
a pairing. None is the sought rational-point-to-section map.

This finishes the \emph{readable supplied-transcription} audit.
It does not prove absence in original handwriting, unresolved
alternate readings, illegible strings, unsupplied material, or
other archival items. The corpus retains 2,620 uncertainty and
48 illegibility macros. Thirteen scan labels have no separate
file: 14, 39, 44, 83, 88, 99, 111, 127, 157, 166, 175, 179,
and 194; their contents are not inferred. Thus the scoped
negative finding is complete, while a stronger historical
nonexistence assertion remains unestablished.

% BEGIN J5 completed quadric counterexample
\subsection*{J5: the height formula fails for a split quadric fourfold}
\textbf{Status: Refuted.} Root independently audited the complete proof and both supplied transcriptions of 44:4.
The failure below concerns the literal statement of Theorem5, 44:4.
Its insertion reads ``section hyperplane générique d'un pinceau dont
[les] fibres sont géom. intègres'', and it concludes
$\langle\zeta,\zeta'\rangle=(-1)^mZ\mathbin{.}Z'$.
The example satisfies this statement,
including the inserted requirement that all pencil fibres be
geometrically integral. It does not depend on identifying the
manuscript's equivalence relation or constructing any nonzero height:
both restricted cycles are literally zero. The original middle
case $n=2m=4$, $i=j=m=2$ also makes its displayed sign unambiguous.

Let $k=\overline{\mathbb F}_3$, and let
\[
 X=V(x_0x_3+x_1x_4+x_2x_5)\subset\mathbb P^5_k.
\]
The six partial derivatives are $x_3,x_4,x_5,x_0,x_1,x_2$,
so $X$ is a smooth projective fourfold. Consider the pencil
\[
 H_t=V(x_2-tx_5),\quad t\in\mathbb A^1,
 \qquad H_\infty=V(x_5).
\]
Its base locus in $X$ is the smooth quadric surface
$B=V(x_2,x_5,x_0x_3+x_1x_4)\subset\mathbb P^5$.
Blowing up $B$ resolves the pencil and gives a smooth projective
total space over $\mathbb P^1$; the fibres remain $X\cap H_t$,
because $B$ is a Cartier divisor in each member. For finite $t$,
the fibre has equation
\[
                 x_0x_3+x_1x_4+t x_5^2=0
                 \quad\text{in }\mathbb P^4.
\]
It is smooth for $t\ne0$ and is a rank-four quadric cone for $t=0$.
The fibre at infinity is another rank-four quadric cone, with
equation $x_0x_3+x_1x_4=0$ and extra coordinate $x_2$.
All fibres are geometrically integral: a quadratic form which is
a product of two linear forms has rank at most two, whereas these
have rank at least four, over every algebraic extension.
In particular this is a Lefschetz pencil with smooth base and two
ordinary quadratic singular fibres, rather than a pencil whose
total family is degenerate.

Take the four planes
\[
 \alpha=V(x_3,x_4,x_5)\subset X,
 \qquad
 \beta=V(x_2,x_3,x_4)\subset X,
\]
\[
 \alpha'=V(x_0,x_1,x_5)\subset X,
 \qquad
 \beta'=V(x_0,x_1,x_2)\subset X.
\]
Put $Z=[\alpha]-[\beta]$, $Z'=[\alpha']-[\beta']$.
If $K=k(t)$ and $Y=X\cap H_t$ is the generic
member, then the proper scheme-theoretic intersections are
\[
 \alpha\cap Y=\beta\cap Y
   =L:=V(x_2,x_3,x_4,x_5)\simeq\mathbb P^1_K,
\]
and
\[
 \alpha'\cap Y=\beta'\cap Y
   =L':=V(x_0,x_1,x_2,x_5)\simeq\mathbb P^1_K,
\]
all with multiplicity one. Indeed $H_t$ restricts on $\alpha$
to $x_2=0$, and on $\beta$ to $-t x_5=0$, where $t$ is a unit
in $K$; the same statements hold for $\alpha',\beta'$.
Therefore
\[
                  \zeta=Z\cdot Y=0,
                  \qquad \zeta'=Z'\cdot Y=0             \tag{J5.1}
\]
already in the group of cycles, not only modulo an equivalence.
Any Neron--Tate bilinear pairing on their associated points is zero.

Here is a direct calculation using only proper transverse
intersections. The crossed pairs $\alpha\cap\beta'$ and
$\beta\cap\alpha'$ are empty. The other two intersections are
\[
 \alpha\cap\alpha'=\{[e_2]\},\qquad
 \beta\cap\beta'=\{[e_5]\}.
\]
At $[e_2]$, the chart $x_2=1$ writes the smooth fourfold with local
coordinates $x_0,x_1,x_3,x_4$ and
$x_5=-x_0x_3-x_1x_4$. The ideals of the two planes are
$(x_3,x_4)$ and $(x_0,x_1)$. They therefore intersect transversely
with local multiplicity one. At $[e_5]$, the chart $x_5=1$
gives exactly the same coordinate calculation, with
$x_2=-x_0x_3-x_1x_4$. Consequently
\[
                  Z\cdot Z'=1+1-0-0=2.                 \tag{J5.2}
\]
The integer two in (J5.2) is the degree of a zero-cycle; it is not
to be reduced modulo the characteristic.

Equations (J5.1)--(J5.2) contradict exactly the proposed formula
\[
             \langle\zeta,\zeta'\rangle
                       =(-1)^m Z\cdot Z',
                  \qquad m=2,
\]
by giving $0=2$. Changing one overall sign cannot fix it.
Both cycles have algebraically trivial generic restrictions, and
every member of the pencil is geometrically integral, as required
by the revised ink. Thus the literal higher-codimension Theorem5
is false even in this smallest balanced case beyond surfaces.

\paragraph{What the integral-fibre condition misses.}
The obstruction is exact: generic restriction does not determine
intersection numbers of higher-codimension extensions. The class
$Z$ is in the kernel of generic restriction but pairs nontrivially with $Z'$,
so this intersection pairing cannot descend to that quotient.
In this displayed example each plane meets the pencil base in a line,
so those intersections have excess dimension. The literal statement
does not forbid this. An additional condition on extensions or their
incidence with the base would be a stronger assertion requiring a
separate analysis. After blowing up the pencil base, the strict transform of $\alpha$
lies in the fibre at infinity and that of $\beta$ in the fibre at
zero; their common exceptional correction cancels in the difference.
Geometric integrality eliminates independent vertical divisor
components, but it does not eliminate these vertical codimension-two
classes through the singularities of the quadric fibres.

The following precise elementary replacement records the necessary
correction. For any rational vector space $V$ with a bilinear form
$b$ and a restriction map $r:V\to W$, let $K=\ker r$.
The form factors through $r(V)$ if and only if
$b(K,V)=b(V,K)=0$. Necessity follows by pairing a zero image with
every image; sufficiency follows by changing either lift by an
element of $K$. For a symmetric form, its restriction always
descends to
\[
                K^\perp/(K\cap K^\perp).
\]
Apply this to rational cycle classes on a smooth projective model,
their intersection form, and restriction to the generic fibre.
A higher-codimension height comparison needs an admissibility or
vertical-correction condition that supplies this orthogonality;
geometric integrality of the fibres alone does not supply it.
This counterexample does not refute a comparison with such
additional conditions, nor the classical divisor comparison for
surfaces with integral fibres.

\paragraph{Finite check.}
The saved script \texttt{check\_quadric\_height.py} checks the
gradient, the pencil ranks, all four literal restriction ideals,
the two transverse point intersections and the two empty intersections,
and the integer intersection value. An additional normal-bundle and
moving-plane calculation independently checks $Z^2=2$ as well.
The written proof supplies the geometric and height conclusions.

% END J5 completed quadric counterexample

% BEGIN J6L independently audited conditional theorem
\subsection{A precise conditional reduction for the Lefschetz assertions}
\paragraph{J6L (Conditional theorem, proved).}
Let $k$ be algebraically closed, of any characteristic. Let
$j:Y\hookrightarrow X$ be a smooth hyperplane section of a smooth
projective $n$-fold, and use the incidence abelian varieties of J1i.
Fix $\ell\ne\operatorname{char}k$ and put $r=2i-1$. For a smooth
projective $d$-fold $W$, write $L_W$ for cup product by the chosen
hyperplane class. The convention for the Lefschetz lowering
operator is
\[
 \Lambda_W(L_W^a v)=
 \begin{cases}L_W^{a-1}v,&a>0,\\0,&a=0,\end{cases}
 \qquad v\in H^b_{\rm prim}(W,\mathbb Q_\ell).
\]
Thus $\Lambda_W$ lowers degree by two and Tate twist by one; it
has no extra $\mathfrak{sl}_2$ coefficient. Let $B(W)$ mean that
this operator is induced by a rational algebraic correspondence
in $\operatorname{CH}^{d-1}(W\times W)_{\mathbb Q}$.
With this explicit convention the following implications hold:
\begin{enumerate}
\item If $2i\leq n-1$ and $B(X)$ holds, then
$j^*:P^i_{\rm inc}(X)\to P^i_{\rm inc}(Y)$ is an isogeny.
\item If $2i\geq n+1$ and $B(Y)$ holds, then the same homomorphism
is surjective.
\end{enumerate}
Only the relevant degree of $\Lambda_X$, respectively
$\Lambda_Y$, is needed. The cohomological inputs are Deligne's hard Lefschetz theorem
\cite[Theorem 4.1.1, p.~217]{DGJ_Deligne1980} and weak
Lefschetz \cite[(4.1.6), pp.~218--219]{DGJ_Deligne1980}.
The coefficient-free lowering convention is also spelled out in
\cite[Section 4.1, pp.~6--7]{DGJ_Peters2026}.
No integral inverse, equality of
incidence and universal regular equivalence, or middle-degree
fixed-part assertion is part of this conclusion.

\emph{Cohomological detection lemma.}
If $U,V$ are smooth projective, $d=\dim U$, and
$\Gamma\in\operatorname{CH}^{d+b-a}(U\times V)_{\mathbb Q}$
acts as zero on
\[
 H^{2a-1}(U,\mathbb Q_\ell(a))
       \longrightarrow H^{2b-1}(V,\mathbb Q_\ell(b)),
\]
then the induced rational homomorphism
\[
 P^a_{\rm inc}(U)\longrightarrow P^b_{\rm inc}(V)
\]
is zero. Here a rational homomorphism is an element of
$\operatorname{Hom}_{k\text{-av}}(-,-)\otimes\mathbb Q$.
To prove the lemma, set $A=P^a_{\rm inc}(U)$, $e=\dim A$.
Choose a miniversal cycle
$u\in\operatorname{CH}^a(A\times U)$ inducing $[m]$ on $A$,
where $m>0$; see \cite[Lemma 2.3, pp.~104--105]{DGJ_Saito1979}.
Choose a Picard realization of the universal map for $V$, namely
a correspondence $w$ and a closed immersion
$P^b_{\rm inc}(V)\hookrightarrow\operatorname{Pic}^{\rm var}(T)$,
as in \cite[(3.5), p.~110]{DGJ_Saito1979}. The composite
\[
 D=w\circ\Gamma\circ u\in
               \operatorname{CH}^1(A\times T)_{\mathbb Q}
\]
induces the composite of $[m]$, the proposed homomorphism, and
this immersion. Its action on
$H^{2e-1}(A,\mathbb Q_\ell(e))\to
H^1(T,\mathbb Q_\ell(1))$ is zero by the hypothesis.
For a divisor correspondence on $A\times T$, this action is the
Poincaré-dual form of its homomorphism on rational Tate modules.
Consequently the displayed composite homomorphism is zero:
a nonzero homomorphism of abelian varieties has a positive
dimensional abelian-variety image and hence has nonzero rational
Tate-module image. The closed immersion and the invertibility
of $[m]$ in the isogeny category give the lemma. Equivalently one
may clear all denominators first; torsion-free Hom groups give
the same conclusion. The case $A=0$ is immediate.

We now prove the two assertions. The cohomology conventions give
\[
 j_*j^*=L_X,\qquad j^*j_*=L_Y,
 \qquad
 j_*:H^q(Y,\mathbb Q_\ell)
         \longrightarrow H^{q+2}(X,\mathbb Q_\ell(1)).
\]
In the first range, $r\leq n-2$, weak Lefschetz makes $j^*$ an
isomorphism on $H^r$. Every Lefschetz string in $H^r(X)$ has
another step above it, so $\Lambda_X L_X=1$ there. The rational
correspondence
\[
 \beta=\Lambda_X\circ j_*\in
           \operatorname{CH}^{n-1}(Y\times X)_{\mathbb Q}
\]
therefore induces the inverse of $j^*$ on $H^r$:
$\beta j^*=\Lambda_XL_X=1$, and $j^*\beta=1$ follows since
$j^*$ is an isomorphism. The detection lemma applied to the two
composites minus the diagonals proves that $\beta$ is a two-sided
inverse on incidence abelian varieties in the isogeny category.
After clearing denominators, each composite is multiplication
by the same positive integer. Thus the actual restriction
homomorphism is an isogeny.

In the second range, $r\geq n>\dim Y$, every Lefschetz string in
$H^r(Y)$ has a step below it. Hence $L_Y\Lambda_Y=1$ there.
The rational correspondence
\[
 \beta=j_*\circ\Lambda_Y\in
           \operatorname{CH}^{n-1}(Y\times X)_{\mathbb Q}
\]
satisfies $j^*\beta=L_Y\Lambda_Y=1$ on $H^r(Y)$. Detection
therefore gives a right inverse in the isogeny category.
Clearing denominators makes the composite multiplication by a
positive integer on $P^i_{\rm inc}(Y)$, so the actual
restriction is surjective. If $i>\dim Y$, the target is zero
and surjectivity is automatic. This proves J6L.

The proof specifies exactly what cohomological weak Lefschetz
alone lacks: algebraic correspondences for the relevant lowering
operators. It also explains why the stable-range conclusion here
is an isogeny and does not settle the alternate transcription
``isomorphisme''. The middle fixed-factor clause and removal of
the stated algebraicity hypotheses remain active gates.

For comparison, in characteristic zero there is a further precise
reduction relating incidence equivalence to the algebraic part of
the Abel--Jacobi map. Peters' Corollary~7.8, p.~16, states that
if all primitive Lefschetz components of the relevant algebraic
odd cohomology remain algebraic, then these two quotients agree
up to isogeny \cite{DGJ_Peters2026}. Standard conjecture $B(X)$
implies that primitive-component condition; the codimension-two
condition is known, as summarized in his Section~7, pp.~15--16.
This result concerns complex Abel--Jacobi and incidence quotients.
It does not identify the universal regular quotient in all
codimensions and does not remove the positive-characteristic
algebraicity gate in J6L.

% END J6L

% BEGIN J6P: smooth-axis trace theorem and top-target surjectivity
\subsection{Codimension-two pencils and the trace image}
\paragraph{J6P (Proved with the stated smooth-axis hypothesis).}
The middle assertion at 44:4 reads ``le facteur fixe de $P^i(Y)$
est identique à l'image de $P^i(X)$''. Here $\operatorname{Ab}^2$ means the universal regular algebraic
representative, not an unidentified manuscript quotient.
Let $k$ be algebraically closed, let $X/k$ be smooth projective of
dimension $n\geq4$, and let a hyperplane pencil have a smooth
codimension-two base locus $B$. Assume its generic member $Y/K$,
$K=k(t)$, is smooth. Write $j:Y\hookrightarrow X_K$ and
$A=\operatorname{Ab}^2_{Y/K}$. Then
\begin{equation}\label{eq:J6P-trace}
 \operatorname{im}\bigl(j^*:\operatorname{Ab}^2_{X/k,K}\longrightarrow A\bigr)
 =\operatorname{im}\bigl(\tau:T_K\longrightarrow A\bigr),
 \qquad T=\operatorname{Tr}_{K/k}(A),
\end{equation}
as abelian subvarieties of $A$. If $n\geq5$, both images are $A$,
so $j^*$ is surjective. No standard conjecture or torsion
injectivity assumption is used for these two conclusions.
The smoothness of the pencil axis is an additional hypothesis:
it is not inferred from the manuscript's integral-fibre condition.

The representative and its separable base-change compatibility
are supplied by \cite[Theorems 6.1 and 5.10, pp.~25--27]{DGJ_ACMVFunctor}.
The trace exists by \cite[Definition 6.1 and Theorem 6.2, p.~20]{DGJ_Conrad2006}.
We use its image, since its kernel can be infinitesimal in positive
characteristic; see \cite[Theorem 6.12, pp.~22--23]{DGJ_Conrad2006}.

\emph{Step 1: spreading a trace family.}
Set $\widetilde X=\operatorname{Bl}_B X$, with blowdown $p$,
exceptional divisor $e:E\hookrightarrow\widetilde X$, and
$q:E\to B$. The pencil gives $f:\widetilde X\to\mathbb P^1$.
Its generic fibre is $Y$: blowing up a member along its Cartier
divisor $B$ changes nothing. Both $\widetilde X$ and $B$ are
smooth projective over $k$.

Choose a miniversal cycle
$\Gamma\in\operatorname{CH}^2(A\times_KY)$ whose induced map
under the universal regular homomorphism $\Phi_Y$ is $[r]_A$,
where $r>0$. This exists over $K$ itself, including when $K$ is
imperfect, by \cite[Lemma 4.7, pp.~18--19]{DGJ_ACMVFunctor}.
Pull it back along $\tau\times1$ to
$u\in\operatorname{CH}^2(T_K\times_KY)$. Normalize at the zero
of $T$ if necessary. Then
\[
 \Phi_Y(u_t-u_0)=r\tau(t).
\]
Take closures of cycle representatives to obtain
$Z\in\operatorname{CH}^2(T\times_k\widetilde X)$ restricting
to $u$ over the generic point of $\mathbb P^1$. Refined fibres
$Z_t-Z_0$ are algebraically trivial and give a regular
homomorphism $v:T\to\operatorname{Ab}^2_{\widetilde X/k}$.
The generic restriction $\delta$ therefore satisfies
\begin{equation}\label{eq:J6P-spread}
 \delta\circ v_K=r\tau,
 \qquad
 \delta:\operatorname{Ab}^2_{\widetilde X/k,K}\longrightarrow A.
\end{equation}
This equality holds as morphisms: it is the equality of regular
maps induced by the generic restriction of the same Chow class.

\emph{Step 2: the exceptional summand.}
The codimension-two blowup formula is
\[
 \operatorname{CH}^2(\widetilde X)
 =p^*\operatorname{CH}^2(X)\oplus e_*q^*\operatorname{CH}^1(B).
\]
Its inverse is $(p_*,-q_*e^*)$. Indeed,
$p_*p^*=1$, $p_*e_*q^*=0$, $q_*e^*p^*=0$, and
$q_*e^*e_*q^*=-1$, the last identity following from
$N_{E/\widetilde X}=\mathcal O_E(-1)$ and
$q_*c_1(\mathcal O_E(1))=1$. The general formula is opened in
\cite[Section 2.3, Theorem 2.8(3)]{DGJ_Jiang2023}, which states
Fulton's blowup formula. All these maps preserve algebraic
equivalence and commute with the refined fibres in a smooth
parameter space. Consequently the decomposition is compatible
with regular homomorphisms and gives an actual isomorphism
\begin{equation}\label{eq:J6P-blowup}
 \operatorname{Ab}^2_{\widetilde X/k}
 \simeq\operatorname{Ab}^2_{X/k}\times\operatorname{Pic}^{\rm var}_{B/k}.
\end{equation}
One can also obtain this directly from the universal properties:
the displayed inverse Chow operations give inverse homomorphisms
of the three universal regular representatives.

Let $b:B_K\hookrightarrow Y$. Under
\eqref{eq:J6P-blowup}, $\delta$ is the sum of $j^*$ and the
homomorphism induced by $b_*$ on algebraically trivial divisors.
There is no excess term here: $E$ and the smooth generic fibre
meet in the Cartier divisor $B_K$, so refined base change gives
$\delta e_*q^*=b_*$.

Choose a smooth member $Y_0/k$ of the pencil; one exists since
the smooth locus of the proper pencil contains its generic point.
Applying weak Lefschetz to $X\supset Y_0\supset B$ shows that
restriction
\[
 a:\operatorname{Pic}^{\rm var}_{X/k}
       \longrightarrow\operatorname{Pic}^{\rm var}_{B/k}
\]
is an isogeny. Indeed, $n\geq4$ makes both restrictions
isomorphisms on $H^1(-,\mathbb Q_\ell(1))$ and hence on the
rational Tate modules of these Picard varieties. Write $H$ for
the chosen hyperplane class. The projection formula gives,
for every algebraically trivial line bundle $L$ on $X$,
\[
 b_*(L|_{B_K})=H|_Y\cdot L|_Y=j^*(H\cdot L).
\]
Thus, as homomorphisms of abelian varieties,
$b_*a_K=j^*\beta_K$, where
$\beta:\operatorname{Pic}^{\rm var}_{X/k}\to
\operatorname{Ab}^2_{X/k}$ is induced by $L\mapsto H\cdot L$.
Since $a$ is surjective, the exceptional summand has image in
$\operatorname{im}(j^*)$. Therefore
$\operatorname{im}(\delta)=\operatorname{im}(j^*)$.

Equation \eqref{eq:J6P-spread} and the surjectivity of $[r]_T$
now imply $\operatorname{im}(\tau)\subseteq\operatorname{im}(j^*)$.
Conversely, the map $j^*$ has constant source; the universal
property of $T$ factors it through $\tau$. This proves
\eqref{eq:J6P-trace}.

\emph{Step 3: why the trace fills the stable-range target.}
Assume $n\geq5$. Choose a prime $\ell$ different from the
characteristic and coprime to the integer of a miniversal cycle
for $A$. By \cite[Proposition 8.2, p.~31]{DGJ_ACMVFunctor}, there
is a Galois-equivariant injection
\[
 T_\ell A\hookrightarrow H^3(Y_{K^{\rm sep}},\mathbb Z_\ell(2))_{
 \rm torsion\,free}.
\]
Weak Lefschetz identifies its rational target with
$H^3(X_{K^{\rm sep}},\mathbb Q_\ell(2))$. This Galois
representation is constant because $X$ is defined over the
algebraically closed field $k$. Thus $T_\ell A$ has trivial
Galois action. All $\ell$-power torsion points of $A$ are
$K$-rational, and their union is a divisible group.

By Lang--N\'eron,
$A(K)/\tau(T(k))$ is finitely generated
\cite[Theorem 7.1, p.~23]{DGJ_Conrad2006}. A divisible group
has zero image in a finitely generated abelian group, so
$A(K)[\ell^\infty]\subseteq\tau(T(k))$. Prime-to-characteristic
power torsion is Zariski dense in an abelian variety. For
completeness, if a proper abelian subvariety contained all this
torsion, its positive-dimensional abelian quotient would have
no $\ell$-power torsion; this is impossible since torsion lifts
through a surjective homomorphism of abelian varieties over an
algebraic closure. Hence $\operatorname{im}(\tau)=A$, as claimed.

\emph{The finite-kernel boundary.}
For $n\geq4$, the same $j^*$ has finite kernel if the universal
map for $Y$ is injective on $\ell$-primary cycle torsion for a
prime $\ell$ avoiding the characteristic and the integer of a
chosen miniversal cycle for $X$. Here is the full implication.
If the reduced identity component $C$ of its kernel were
positive-dimensional, restrict a miniversal cycle $\Gamma_X$
to $C$. The map
\[
 T_\ell C\longrightarrow T_\ell A^2(X_{K^{\rm sep}})
 \longrightarrow H^3(X_{K^{\rm sep}},\mathbb Z_\ell(2))_{
 \rm torsion\,free}
\]
is injective by Proposition 8.2. Its pullback cycle family on
$Y$ is killed by $\Phi_Y$ since the representative map is
$[r]j^*|_C=0$. The stated torsion injectivity then makes its
Tate module of cycles zero. Functoriality of the Bloch map
\cite[Section 8.1, pp.~30--31]{DGJ_ACMVFunctor} makes the
displayed cohomology classes restrict to zero on $Y$. But weak
Lefschetz is injective on $H^3(X)\to H^3(Y)$ for every $n\geq4$,
a contradiction after tensoring with $\mathbb Q_\ell$.

This torsion hypothesis is known in characteristic zero by
Murre's theorem, as recorded in
\cite[Remark 8.6, p.~33]{DGJ_ACMVFunctor}. It therefore gives,
in characteristic zero and with the smooth-axis hypothesis,
an isogeny in the stable range and an isogeny onto the trace
image when $n=4$. Proposition 8.2 alone does not remove the
hypothesis: its injection is constructed from a miniversal
cycle and does not assert that $\Phi_Y$ is injective on torsion.
The same remark explicitly leaves that injectivity unsettled
in general positive characteristic.

\paragraph{J6Z (The highest nontrivial target codimension).}
Let $X/k$ be smooth projective of dimension $n\geq3$ and let
$j:Y\hookrightarrow X$ be any smooth hyperplane section. Then
\[
 j^*:P^{n-1}_{\rm inc}(X)\longrightarrow
       P^{n-1}_{\rm inc}(Y)=\operatorname{Alb}(Y)
\]
is surjective in every characteristic. The analogous assertion
holds for universal regular representatives whenever the source
exists; in particular it holds for
$\operatorname{Ab}^2(X)\to\operatorname{Alb}(Y)$ on a threefold.

To prove it, restriction
$\operatorname{Pic}^{\rm var}(X)\to\operatorname{Pic}^{\rm var}(Y)$
is an isogeny by weak Lefschetz in degree one. Intersection by
$H_X^{n-2}$ defines a homomorphism from the former Picard
variety to $P^{n-1}_{\rm inc}(X)$, and its restriction to $Y$
is intersection by $H_Y^{n-2}$ followed by the Albanese map.
The induced map
\[
 \operatorname{Pic}^{\rm var}(Y)\longrightarrow\operatorname{Alb}(Y)
\]
is an isogeny: on rational Tate modules, the codimension-one
and zero-cycle Bloch isomorphisms identify it with
\[
 L_Y^{n-2}:H^1(Y,\mathbb Q_\ell(1))
       \xrightarrow{\ \sim\ }H^{2n-3}(Y,\mathbb Q_\ell(n-1)),
\]
which is hard Lefschetz. See
\cite[Section 8.1, pp.~30--31]{DGJ_ACMVFunctor} for those
Bloch isomorphisms and
\cite[Theorem 4.1.1 and (4.1.6), pp.~217--219]{DGJ_Deligne1980}
for hard and weak Lefschetz. A surjective homomorphism thus
factors through $j^*$, proving surjectivity. In this top target
codimension, the common incidence quotient is the Albanese:
the Poincar\'e family of divisors on $Y$ detects exactly its
kernel. This argument uses neither an inverse Lefschetz
correspondence nor the torsion standard assumption.

\emph{Exact scope.}
The fixed-part result proves the modern codimension-two
interpretation of the middle clause for pencils with smooth
axis. It does not identify that interpretation with the
manuscript's undefined equivalence relation, cover singular
pencil axes, prove positive-characteristic finite kernel, or
settle the higher-codimension stable and middle assertions.

% END J6P

% BEGIN J6C
\subsection{A cubic pencil refuting the literal middle fixed-part clause}
\paragraph{J6C (Refuted by an explicit cubic pencil).}
At 44:4, Theorem~6 says, under the conditions of Theorem~5,
``si $2i=n$ le facteur fixe de $P^i(Y)$ est identique à
l'image de $P^i(X)$''. The inserted condition on $Y$ in
Theorem~5 is ``section hyperplane générique d'un pinceau''
whose fibres are geometrically integral. No smoothness of the
pencil axis is imposed in either supplied transcription.
The word ``générique'' modifies the hyperplane section of the
pencil; the pencil itself is not required to be general.
The following example meets those literal hypotheses and has
zero source but a positive-dimensional fixed target factor.
The cycle hypotheses in Theorem~5 can, if needed, be met by
taking $Z=Z'=0$ with $i=j=2$.

\emph{Existing-source check.}
The underlying elliptic-cone geometry is classical:
\cite[Lemma 8.1, pp.~314--315]{DGJ_ClemensGriffiths1972}
describes the cone of lines through an Eckardt point, and
\cite[Proposition 10, p.~384; Lemma 29, p.~398]{DGJ_Roulleau2009}
studies the associated elliptic curve and homomorphisms on the
complex Fano surface. The explicit pencil and positive-characteristic
calculation below are supplied with a complete proof; no priority
claim is made for this construction.

Let $k=\overline{\mathbb F}_5$, and let $X\subset\mathbb P^5_k$
be the cubic fourfold
\begin{equation}\label{eq:J6C-fourfold}
 x_0^3+x_1^3+x_2^3+x_3^2x_4+x_4^3+x_5^3=0.
\end{equation}
Consider the hyperplane pencil $x_5=t x_4$, with parameter
$t\in\mathbb A^1_k$, and the member $x_4=0$ at infinity.
Its axis is the integral surface
\[
 B:\quad x_4=x_5=0,\qquad x_0^3+x_1^3+x_2^3=0.
\]
It is the cone over the smooth plane cubic
\[
 E:\quad u_0^3+u_1^3+u_2^3=0\quad\text{in }\mathbb P^2_k,
\]
with vertex $p=[0:0:0:1:0:0]$. Choose
$e_0=[1:-1:0]$ on $E$ as origin. Let $K=k(t)$, and let
$Y/K$ be the generic pencil member. The bundle
$\mathcal O_X(1)$ is the very ample bundle of this displayed
embedding. The morphisms below are all defined over $K$.

\emph{Smooth total variety and integral pencil members.}
The gradient of \eqref{eq:J6C-fourfold} is
\[
 (3x_0^2,3x_1^2,3x_2^2,2x_3x_4,
   x_3^2+3x_4^2,3x_5^2).
\]
A common zero forces $x_0=x_1=x_2=x_5=0$, then
$x_3x_4=0$ and $x_3^2+3x_4^2=0$, and finally
$x_3=x_4=0$. Thus $X$ is smooth. For a finite parameter,
the member is the cubic in $\mathbb P^4$ with equation
\begin{equation}\label{eq:J6C-member}
 x_0^3+x_1^3+x_2^3+x_3^2x_4+a x_4^3=0,
 \qquad a=1+t^3.
\end{equation}
If $a\ne0$, the same derivative calculation shows that it
is smooth. If $a=0$, its singular locus has the single
geometric point $[0:0:0:0:1]$. The member at infinity is
the cone, with vertex $p$, over the smooth Fermat cubic
surface $x_0^3+x_1^3+x_2^3+x_5^3=0$; its singular locus
also has just one point.

Every member is geometrically integral. Indeed, over an
algebraically closed field a reducible hypersurface in
$\mathbb P^4$ has two components meeting in dimension at
least two, and their intersection is singular for the
product equation. A repeated component is singular along
a three-dimensional subset. Both possibilities contradict
the computed empty or zero-dimensional singular locus.
This also proves geometric integrality of the generic
member after extending $K$ to an algebraic closure. The
three finite singular parameters are distinct, since
$1+t^3$ is relatively prime to $3t^2$ in characteristic five.

The pencil is the incidence family in $X\times\mathbb P^1$.
Because $x_4,x_5$ form a regular sequence along the
codimension-two axis, the two-generated ideal is of linear
type \cite[Proposition IV-25, pp.~172--173]{DGJ_EisenbudHarris2000};
its blowup is that incidence family. Thus its scheme
fibres are exactly the displayed hyperplane sections,
with no additional fibre components. Equivalently, each
local incidence equation restricts to a nonzerodivisor on
each fibre of $X\times\mathbb P^1$, so this is a flat
hyperplane family. The singular axis is the sole missing
hypothesis relative to the positive theorem J6P.

\emph{The source representative vanishes.}
Weak Lefschetz for the smooth ample hypersurface
$X\subset\mathbb P^5$ gives
\[
 H^3(X,\mathbb Q_\ell)=H^3(\mathbb P^5,\mathbb Q_\ell)=0
 \qquad(\ell\ne5).
\]
Consequently
\begin{equation}\label{eq:J6C-zero-source}
 \operatorname{Ab}^2(X)=0,\qquad P^2_{\rm inc}(X)=0.
\end{equation}
For the first conclusion use the miniversal-cycle injection
\cite[Proposition 8.2, p.~31]{DGJ_ACMVFunctor}; for the
second use Saito's incidence bound
\cite[Proposition 3.4, pp.~109--110]{DGJ_Saito1979}.
Weak Lefschetz is cited in
\cite[(4.1.6), pp.~218--219]{DGJ_Deligne1980}.

\emph{A constant elliptic factor in the target.}
Blow up the $K$-point $p$ of the smooth threefold $Y$:
\[
 \pi:\widehat Y=\operatorname{Bl}_pY\longrightarrow Y,
 \qquad D\simeq\mathbb P^2_K.
\]
Let $i:S\hookrightarrow\widehat Y$ be the strict transform
of $B_K$. This is the smooth ruled resolution of the
elliptic cone, with projection $q:S\to E_K$. Concretely,
its fibre over $e=[u_0:u_1:u_2]$ is the strict transform
of the ruling line
\[
 [\lambda u_0:\lambda u_1:\lambda u_2:\mu:0:0]
 \quad\text{in }B_K.
\]
Thus $S$ is a $\mathbb P^1$-bundle over $E_K$, and every
ruling fibre $\ell$ satisfies
$\pi^*H\cdot\ell=1$ and $D\cdot\ell=1$.

The multiplicity of $B_K$ as a divisor of $Y$ at $p$ is
three. To check it, use $x_3=1$ and put $v=x_4$; equation
\eqref{eq:J6C-member} becomes
\[
 v+a v^3+u_0^3+u_1^3+u_2^3=0.
\]
The three $u_j$ are regular parameters at $p$, and
$v(1+a v^2)=-(u_0^3+u_1^3+u_2^3)$, with the first factor
in parentheses a unit. Hence $v$ has order three. As a
Cartier divisor on $\widehat Y$, therefore,
\[
 S=\pi^*H-3D,
 \qquad
 \deg(N_{S/\widehat Y}|_\ell)=1-3=-2.
\]

Put $J=\operatorname{Pic}^0(E)$, a constant elliptic
curve over $k$. The algebraic-correspondence operations
\[
 i_*q^*:A^1(E_K)\longrightarrow A^2(\widehat Y),\qquad
 q_*i^*:A^2(\widehat Y)\longrightarrow A^1(E_K)
\]
preserve algebraic equivalence and regularity. Here $i^*$
is the refined pullback for the Cartier embedding $i$;
$q_*$ sends zero-cycles on the surface $S$ to divisor
classes on $E_K$. They induce homomorphisms of abelian
varieties
\[
 \widehat u:J_K\longrightarrow\operatorname{Ab}^2(\widehat Y),
 \qquad
 \widehat v:\operatorname{Ab}^2(\widehat Y)\longrightarrow J_K.
\]
The Cartier self-intersection formula
\cite[Definition 42.29.1(1)(b)]{DGJ_StacksGysin}
and the projection formula
\cite[Lemma 42.26.4]{DGJ_StacksProjection} give,
for every $L\in\operatorname{Pic}^0(E_K)$,
\begin{align*}
 q_*i^*i_*q^*L
 &=q_*\bigl(c_1(N_{S/\widehat Y})\cdot q^*L\bigr)\\
 &=\bigl(q_*c_1(N_{S/\widehat Y})\bigr)\cdot L=-2L.
\end{align*}
Thus
\begin{equation}\label{eq:J6C-minus-two}
 \widehat v\widehat u=[-2]_{J_K}.
\end{equation}
This sign is an equality of divisor-class homomorphisms
on $J$, not an unstated choice of sign for an elliptic
principal polarization.

The point-blowup formula gives
\[
 \operatorname{CH}^2(\widehat Y)
    =\pi^*\operatorname{CH}^2(Y)\oplus\mathbb Z[\text{line in }D].
\]
The exceptional integer is constant under algebraic
equivalence, so $\pi^*$ and $\pi_*$ are inverse on $A^2$.
They also preserve regular families and therefore induce
inverse isomorphisms of the universal regular
representatives. This is the $r=3$ case of the blowup
formula \cite[Theorem 2.8(3)]{DGJ_Jiang2023}.
Define
\[
 u=\pi_*\widehat u:J_K\longrightarrow\operatorname{Ab}^2(Y),
 \qquad
 v=\widehat v\pi^*:\operatorname{Ab}^2(Y)\longrightarrow J_K.
\]
Equations above give $vu=[-2]_{J_K}$. In characteristic
five this map is an isogeny of degree four. Hence $u$ has
finite kernel and a one-dimensional abelian image. Since
its source is constant and $u$ is defined over $K$, the
universal property of the $K/k$-trace puts this image in
the fixed part of $\operatorname{Ab}^2(Y)$. That fixed
part is therefore positive-dimensional, contradicting
\eqref{eq:J6C-zero-source}.

The factor is detected by incidence equivalence as well:
$q_*i^*\pi^*$ is an actual correspondence from $Y$ to the
curve $E_K$, taking codimension-two classes to divisor
classes. It is one of the Picard tests defining the
incidence quotient. Thus the constant family $i_*q^*$
cannot vanish in that quotient, since composing with
this test is $[-2]$. In particular the counterexample
does not depend on the disputed integral autoduality
of a universal representative.

\emph{Conclusion and remaining scope.}
For $n=4,i=2$, a generic hyperplane section of this
integral-fibre pencil has a nonzero fixed elliptic
factor, although the ambient fourfold has zero
codimension-two representative. The literal middle
clause of Theorem~6 is therefore false. This does not
refute J6P, whose smooth-axis hypothesis is essential,
or settle the separate stable- and high-range assertions
in all codimensions. The exact gradient ideals,
singular supports, origin, parameter roots, and normal
degree are checked by \texttt{check\_cubic\_trace.py};
the constant-factor conclusion is the correspondence
proof above, not an output of the finite computation.

% END J6C

% BEGIN JWalker
\subsection{The remaining regular-representability gate}
The August 5, 2026 revision of
\cite[Introduction, p.~1]{DGJ_ACMVTwist2026} still distinguishes
universal algebraic representatives in codimensions $1,2,\dim X$
from merely regular algebraic intermediate-Jacobian maps in
the other codimensions. This is a precise current source check,
not an inference of open status from an unsuccessful search.
The final published version of Kahn's concrete fourfold
expectation is now also opened: Conjecture~3.1, p.~847
\cite{DGJ_Kahn2021}. It corresponds to Conjecture~1, p.~5 in
the six-page author version cited above.

For complex varieties one should additionally distinguish
universality among all regular homomorphisms from universality
among isogeny lifts. Write
\[
 \psi:A^i(X)\longrightarrow J_a^{2i-1}(X)(\mathbb C),\qquad
 \omega:A^i(X)\longrightarrow J_W^{2i-1}(X)(\mathbb C)
\]
for the usual and Walker Abel--Jacobi maps, and
$\alpha:J_W^{2i-1}(X)\to J_a^{2i-1}(X)$ for their isogeny.
By \cite[Theorem 2.3, p.~12]{DGJ_ACMVWalker}, $\omega$ is
initial among regular lifts of $\psi$ along isogenies. This
theorem does not say that $\omega$ is universal among all
regular homomorphisms.

Here is an exact elementary test for that stronger assertion.
For every regular homomorphism
$\phi:A^i(X)\to B(\mathbb C)$, let
\[
 C_\phi=\operatorname{im}(\psi,\phi)
       \subset J_a^{2i-1}(X)\times B.
\]
The image is the group of complex points of an abelian
subvariety, by the image theorem for regular homomorphisms
\cite[Lemma 4.7, pp.~18--19]{DGJ_ACMVFunctor}.
Then $\omega$ is universal among all regular homomorphisms
if and only if
\begin{equation}\label{eq:JWalker-gate}
 \dim C_\phi=\dim J_a^{2i-1}(X)
 \quad\hbox{for every regular }\phi.
\end{equation}
Indeed, if \eqref{eq:JWalker-gate} holds, the surjective
projection $C_\phi\to J_a^{2i-1}(X)$ is an isogeny. Initiality
among isogeny lifts gives a homomorphism
$J_W^{2i-1}(X)\to C_\phi$ inducing $(\psi,\phi)$; its second
projection factors $\phi$ through $\omega$. This factorization
is unique because $\omega$ is surjective. Conversely, if
$\omega$ is universal, $(\psi,\phi)$ factors through $J_W$,
so its image has dimension at most $\dim J_W=\dim J_a$;
the surjective projection to $J_a$ gives the reverse inequality.

The unresolved issue in this route is therefore a possible
positive-dimensional kernel of $C_\phi\to J_a$ for a regular
test $\phi$. The divisible-kernel statement in Theorem~2.3
does not rule it out: divisible groups can have nonzero maps
to groups of points of abelian varieties. Likewise, infinite
torsion in $A^3$ does not by itself construct regular targets
of unbounded dimension. The test above is not asserted for
Kahn's fourfold, and neither its proof nor its failure has
been obtained here. This leaves the general statement in
Theorems~1--3 explicitly unresolved.

% END JWalker

\section{Numbered statement register}\label{sec:K}
This register lists every current result row, including unresolved parent tasks.
The identifiers are stable across the report, calculations and check-event table.
The full typed statements, diplomatic French quotations and proofs occur at the
locations given below. A proved special case does not change an unresolved parent
task to proved. Historical attempts are retained with their dated verification state.
\begingroup\small\raggedright
\par\medskip\noindent\textbf{K.1. A1 --- Undecided.}\par
Picard isomorphism for every smooth hyperplane in dimension at least four.
\par\emph{Statement/proof location:} Part A, A1r and A4.
\emph{Manuscript reference:} 57:97.
\par Reduced to the dimension-four case; finite multiplicative obstruction remains.
\par\medskip\noindent\textbf{K.2. A2 --- Undecided.}\par
Injectivity and the least universal dimension.
\par\emph{Statement/proof location:} Part A, A2b and A4b.
\emph{Manuscript reference:} 57:97.
\par Universal point injectivity and its least dimension remain undecided. The kernel in dimension at least three is a finite p-group; for each nonzero order-p class its bad smooth hyperplanes form the finite determinantal locus A2D. A2V forces a Cartier-exact leading term and Frobenius kernel on negative H1 of the hyperplane, with vanishing order divisible by p and a nonnormal mu\_p torsor. A2W forces rank(d:H0Omega1->H0Omega2)>=n-1. A2Y additionally requires noninjective Frobenius on H1(O\_Y) and no smooth W2 lifting of Y. In dimension3 the target is of general type by applying A3s to L\textasciicircum{}a. No global realization or universal exclusion of this remaining gate is proved.
\par\medskip\noindent\textbf{K.3. A3 --- Proved.}\par
There exists a smooth projective surface with a very ample L and nonzero H1(L\textasciicircum{}-1); the least requested dimension is two.
\par\emph{Statement/proof location:} sections/A.tex, a:very-ample-example.
\emph{Manuscript reference:} 57:99-100.
\par Explicit characteristic-two Raynaud surface followed by 556 simultaneous ninth-root covers; h0(L)=562 and h1(L\textasciicircum{}-1)=1. Complete GF(2) basis and finite transition certificate included; three targeted independent audits passed.
\par\medskip\noindent\textbf{K.4. A4 --- Proved.}\par
Exact finite multiplicative obstruction description of kernel and cokernel.
\par\emph{Statement/proof location:} Part A, Statement a:finite.
\emph{Manuscript reference:} 57:97.
\par Uses the actual multiplicative tower; not just dimensions or a direct sum of additive layers.
\par\medskip\noindent\textbf{K.5. B1 --- Refuted.}\par
Seshadri without Pic as formulated in prompt.
\par\emph{Statement/proof location:} work/bce/section\_B.tex; section.tex.
\emph{Manuscript reference:} 57:110.
\par Split projective quadric xy=z(z-tw), D=(x,z), all characteristics; original manuscript retains essential properness of Pic functor
\par\medskip\noindent\textbf{K.6. B2 --- Refuted as a fibre-type-only criterion; fixed-family iff proved.}\par
Exact conditions in Question 2 and necessity of R2.
\par\emph{Statement/proof location:} sections/B.tex, B2 exact obstruction and paired projective families.
\par Identical regular base and abstract geometric fibre types give opposite extension properties; no criterion from those data alone exists. The exact deformation-sensitive Cl/Pic kernel criterion is proved for regular base and normal integral geometric fibres. R2 is not necessary and requires the original lci hypothesis for sufficiency.
\par\medskip\noindent\textbf{K.7. C1 --- Proved.}\par
Projectivity descends from completion.
\par\emph{Statement/proof location:} work/bce/section\_C.tex; Artin1969 Thm3.5 pp31-32.
\emph{Manuscript reference:} 57:227-231;272.
\par Artin approximation of an ample bundle on the closed fibre; relative ampleness over a local base
\par\medskip\noindent\textbf{K.8. C2 --- Proved.}\par
Picard comparison when R1 has finite length.
\par\emph{Statement/proof location:} sections/C.tex, Theorem C2.G.
\par Full excellent henselian comparison proved by principal finite determination, blow-up and Stein factorization, conductor descent, and approximation; independent proof audit passed.
\par\medskip\noindent\textbf{K.9. C3 --- Proved.}\par
Punctured henselian Pic equals completed Pic.
\par\emph{Statement/proof location:} work/bce/section\_C.tex; Elkik1973 III.1 Thm3 p569 and Pic cor p573.
\emph{Manuscript reference:} 57:194;227-231.
\par Any noetherian henselian pair; no isolated-singularity or characteristic restriction
\par\medskip\noindent\textbf{K.10. C4 --- Proved in original cone scope; arbitrary-local extrapolation refuted.}\par
Local Neron-Severi comparison modulo p groups.
\par\emph{Statement/proof location:} sections/C.tex, cone comparison and general-local counterexample.
\par The source fixes the local cone at its vertex and then compares punctured Picard groups. That statement is proved modulo bounded p-groups. A characteristic-three factorial local hypersurface with henselian/completed punctured Pic Z refutes the arbitrary-local extrapolation; different intrinsic etale-local NS definitions are stated separately.
\par\medskip\noindent\textbf{K.11. D1-positive --- Proved.}\par
Positive quotient Hodge asymmetry.
\par\emph{Statement/proof location:} work/dgj/section\_D.tex.
\emph{Manuscript reference:} 57:287.
\par Characteristic two shearing involutions with F and V; dimensions (2,3),(3,2), specializing to (3,3).
\par\medskip\noindent\textbf{K.12. D1-negative --- Proved.}\par
Negative quotient Hodge asymmetry.
\par\emph{Statement/proof location:} work/dgj/section\_D.tex.
\emph{Manuscript reference:} 57:287.
\par Characteristic two shearing involutions with F and V; dimensions (2,3),(3,2), specializing to (3,3).
\par\medskip\noindent\textbf{K.13. D1-family --- Proved.}\par
Variation of quotient Hodge asymmetry.
\par\emph{Statement/proof location:} work/dgj/section\_D.tex.
\emph{Manuscript reference:} 57:287.
\par Characteristic two shearing involutions with F and V; dimensions (2,3),(3,2), specializing to (3,3).
\par\medskip\noindent\textbf{K.14. D2-flatness --- Proved.}\par
Nonflat fixed groups and Picard torsion.
\par\emph{Statement/proof location:} work/dgj/section\_D.tex.
\emph{Manuscript reference:} 57:288.
\par Fixed lengths 8 and 16; algebra R[z]/(z\textasciicircum{}4,t z\textasciicircum{}2).
\par\medskip\noindent\textbf{K.15. D2-specialization --- Proved.}\par
Picard variety and specialization.
\par\emph{Statement/proof location:} work/dgj/section\_D.tex.
\emph{Manuscript reference:} 57:288.
\par Fibrewise image of (F\textasciicircum{}2,[2]) has kernel degree 2 then 4.
\par\medskip\noindent\textbf{K.16. D2-Igusa --- Proved.}\par
Modular family of the Igusa surface.
\par\emph{Statement/proof location:} work/dgj/section\_D.tex D4o and D4s; section\_D\_supersingular.tex.
\emph{Manuscript reference:} 57:311.
\par Both algebraized miniversal hulls computed: ordinary W[[t,s,u,v]]/(s\ensuremath{{}^2}+2s); supersingular W[[t,s,u,v]]/(s\ensuremath{{}^4}−t*s\ensuremath{{}^2}−2s). Canonical fixed elliptic subgroup and all-Artinian Pic-tau flatness proved; root/fhi independent audits pass.
\par\medskip\noindent\textbf{K.17. E1 --- Refuted as stated; intended valuative version proved.}\par
Properness of Pic0 over a trait.
\par\emph{Statement/proof location:} sections/E.tex, E1 literal counterexample and intended valuative theorem.
\par A normal total-space elliptic degeneration refutes the literal assertion. The version with geometrically normal fibres is proved over a trait. Arbitrary-base representability questions are distinct E4 rows.
\par\medskip\noindent\textbf{K.18. E2 --- Proved.}\par
Corr is separated under proper flat universal H0=O.
\par\emph{Statement/proof location:} work/bce/section\_E.tex; Raynaud1970 3.2(c),5.2;Murre1966 Thm4.
\emph{Manuscript reference:} 57:22;32;313.
\par Trait normalization into étale closure E of identity; connected fibres force map zero; Murre yields unramified separated scheme
\par\medskip\noindent\textbf{K.19. E3 --- Proved.}\par
Problem A universal flattening monomorphism of finite presentation.
\par\emph{Statement/proof location:} work/bce/section\_E.tex; Stacks05PS Lem38.27.6(3);0CWQ Lem77.11.8(3).
\emph{Manuscript reference:} 57:73-79.
\par Proper finitely presented X and finitely presented F; all base changes
\par\medskip\noindent\textbf{K.20. E4 --- Undecided.}\par
Remaining structure questions on 57:308-315.
\par\emph{Statement/proof location:} sections/E.tex, individual questions 57:308-315.
\par E4.T proves Pic-tau closedness without projectivity for proper flat finitely presented geometrically integral families over arbitrary scheme bases. E4.S settles the fibrewise smoothness/defect claims; E4.Curve proves full Pic schematicity for smooth proper geometrically connected curves. Remaining mathematical gate: full Pic schematicity for geometrically integral nonnormal families and normal higher-dimensional families over arbitrary bases. The scan310 finite-component iff retains an explicitly recorded reading/scope uncertainty; multiplication nonopenness is already refuted by E4.H.
\par\medskip\noindent\textbf{K.21. F --- Undecided.}\par
Finite generation of the full local etale fundamental group.
\par\emph{Statement/proof location:} work/fhi/F\_proof.md section F4.
\emph{Manuscript reference:} 58:19.
\par For complete normal local surfaces with algebraically closed residue field, the unrestricted conjecture remains. In equal characteristic F8/F9b settle every geometric quotient with normal pro-p kernel over the full tame quotient. F10 proves finite abelian inertia using Krishna–Majumder2026v2; F10c gives all-prime rank-linear H1 bounds for abelian monodromy, and F11 proves uniform finite generation of metabelian quotients. F12 gives the characteristic-p bound for metabelian coefficients; F12b settles pro-p closed second derived subgroup. F13 proves a fixed log-conductor bound after one fixed cover. F13c/d remove unramified character order and reduce the remaining coefficient-prime-not-p estimate to a degree-linear tame-generator bound for cyclic covers of uniformly bounded p-power degree (F13.6). F13e bounds refined-Swan zero length, but not cleanup complexity. F13.5/F13.6, further solvable and nonsolvable wild quotients, and mixed characteristic remain unresolved. The actual wild kernel is not assumed pro-p.
\par\medskip\noindent\textbf{K.22. G1 --- Proved.}\par
Universal Pic-delta section exists iff at most one elementary divisor is odd.
\par\emph{Statement/proof location:} work/dgj/section\_G.tex.
\emph{Manuscript reference:} 59:10-14.
\par Over Z[1/2], construction uses rank of induced pairing on A[2].
\par\medskip\noindent\textbf{K.23. G2 --- Proved.}\par
Determinant constructions agree with canonical sections.
\par\emph{Statement/proof location:} work/dgj/section\_G.tex.
\emph{Manuscript reference:} 59:12-14.
\par Includes square and even-Euler double transform comparisons; corrects ppcm and even-chi divisibility claim.
\par\medskip\noindent\textbf{K.24. G3-Albanese --- Proved.}\par
Sign of the Albanese torsor.
\par\emph{Statement/proof location:} work/dgj/section\_G.tex.
\emph{Manuscript reference:} 59:23-28.
\par Minus sign visually verified; canonical norm on PPAV theta torsor.
\par\medskip\noindent\textbf{K.25. G3-principal --- Proved.}\par
Canonical class of principally polarized abelian schemes.
\par\emph{Statement/proof location:} work/dgj/section\_G.tex.
\emph{Manuscript reference:} 59:23-28.
\par Minus sign visually verified; canonical norm on PPAV theta torsor.
\par\medskip\noindent\textbf{K.26. H-Ga --- Proved.}\par
Higher Ext vanishing for the additive group.
\par\emph{Statement/proof location:} work/fhi/H.tex; H\_proof.md.
\emph{Manuscript reference:} 43:80.
\par Corrects assertion that Ga remains difficult
\par\medskip\noindent\textbf{K.27. H-all --- Proved.}\par
Higher Ext vanishing for every algebraic commutative group.
\par\emph{Statement/proof location:} work/fhi/H.tex; H\_proof.md.
\emph{Manuscript reference:} 43:24-35;43:80.
\par All proalgebraic targets; all p
\par\medskip\noindent\textbf{K.28. I-2d --- Proved.}\par
Residual complex endomorphisms in dimension two, degree two.
\par\emph{Statement/proof location:} work/fhi/I.tex; I\_proof.md.
\emph{Manuscript reference:} 43:65-68;SA19-H3.
\par Precise finite-topology universal interpretation; mixed arrow topology
\par\medskip\noindent\textbf{K.29. I-all --- Undecided.}\par
Completed algebra presentation in every dimension.
\par\emph{Statement/proof location:} work/fhi/I\_universal\_linear.tex; I\_proof.md.
\emph{Manuscript reference:} 43:65-68;SA19-H3.
\par All dimensions are proved for separated complete graded algebras with a basis of open left ideals; all surfaces are proved for arbitrary jointly continuous topology. The original43:66 specifies the broader category without the open-left-ideal hypothesis. I7b now proves open dense differential-image quotients for whole source columns over arbitrary regular local rings. Remaining: construction and continuity of the entire adjacent-arrow family and its endpoint assembly in arbitrary complete targets, in dimension at least3.
\par\medskip\noindent\textbf{K.30. J1-Th1 --- Proved for common incidence or regular-test kernels; manuscript identification unresolved.}\par
Set-theoretic kernel equality for manuscript Albanese equivalence.
\par\emph{Statement/proof location:} work/dgj/section\_J.tex J1i; section\_J\_incidence.tex.
\emph{Manuscript reference:} 44:2.
\par Common incidence kernels agree scheme-theoretically with the represented family map. Common all-regular kernels agree on points. Undefined manuscript N and en substance/ensemblistement variants remain explicit.
\par\medskip\noindent\textbf{K.31. J1-Th2 --- Known and proved for incidence; all-regular generality unresolved.}\par
Bound for all manuscript quotients.
\par\emph{Statement/proof location:} work/dgj/section\_J.tex J1i; section\_J\_incidence.tex.
\emph{Manuscript reference:} 44:2.
\par Every incidence family quotient has dimension at most b(2i-1)/2 in all codimensions and characteristics. For all regular maps, general uniform boundedness is equivalent to universal existence, with codimension3/fourfold gate.
\par\medskip\noindent\textbf{K.32. J1-Th3 --- Incidence existence and isogeny duality known; modern integral duality refuted; general regular existence unresolved.}\par
General representability and complementary integral duality.
\par\emph{Statement/proof location:} work/dgj/section\_J.tex J1i and J2a.
\emph{Manuscript reference:} 44:2;44:8.
\par Saito constructs the incidence quotient in all codimensions with specified isogeny normalization. The modern Ab integral assertion is false. Kahn2021 Conjecture1 gives a concrete remaining regular-representability test.
\par\medskip\noindent\textbf{K.33. J1-Th4 --- Refuted.}\par
Raw middle intersection autoduality is positive up to one sign.
\par\emph{Statement/proof location:} work/dgj/section\_J.tex.
\emph{Manuscript reference:} 44:2;44:8-9.
\par Primitive-qualified question is not refuted.
\par\medskip\noindent\textbf{K.34. J1-Th5 --- Refuted.}\par
Higher codimension height/intersection comparison as revised.
\par\emph{Statement/proof location:} work/dgj/section\_J.tex J5; section\_J5\_quadric.tex.
\emph{Manuscript reference:} 44:4.
\par Over Fbar3: Q4, pencil x2=t*x5, all fibres geometrically integral. Both restricted cycles literally zero; Z.Zprime=2 from two transverse points. i=j=m=2 fixes the sign. Excess intersections with the pencil base are permitted by the literal statement; stronger admissibility conditions remain distinct.
\par\medskip\noindent\textbf{K.35. J1-Th6 --- Refuted as stated in the middle clause; full stable and high clauses unresolved.}\par
Lefschetz for all algebraic representatives.
\par\emph{Statement/proof location:} work/dgj/section\_J.tex J6C/J6L/J6P/J6Z; section\_J6\_cubic.tex.
\emph{Manuscript reference:} 44:4.
\par J6C: over Fbar5 a pencil with every fibre geometrically integral has zero ambient Ab2 and Pinc2 but a constant elliptic factor in the generic member, detected by composition[-2]. No smooth-axis hypothesis occurs in44:4. J6P proves the smooth-axis codimension-two correction; J6L and J6Z give conditional/general and unconditional top-target cases. Full stable/high, positive-characteristic finite-kernel and integral-isomorphism readings remain gates.
\par\medskip\noindent\textbf{K.36. J1-Th7 --- Refuted.}\par
Numerical equals tau equivalence for all cycles.
\par\emph{Statement/proof location:} work/dgj/section\_J.tex.
\emph{Manuscript reference:} 44:6.
\par Explicit Fermat quartic Ceresa cycle over Qbar; Bloch1984 pp94-95 and Theorem4.1 p104 visually verified. Extends to all intermediate codimensions by products.
\par\medskip\noindent\textbf{K.37. J2-isomorphism --- Refuted.}\par
Threefold algebraic representative canonical self-duality.
\par\emph{Statement/proof location:} work/dgj/section\_J.tex.
\emph{Manuscript reference:} 44:2;44:8-9.
\par Generic monic septic Jacobian over the algebraic closure of F5(a0,...,a6)\textasciicircum{}perf; intersection forces a degree8 K-isogeny, ruled out by irreducible S7 action on A[2]. No integral torsion standard assumption used.
\par\medskip\noindent\textbf{K.38. J2-positivity --- Refuted.}\par
Threefold algebraic representative polarization.
\par\emph{Statement/proof location:} work/dgj/section\_J.tex.
\emph{Manuscript reference:} 44:2;44:8-9.
\par Explicit E/F3 times P1 times P1; matrix [[0,lambda],[lambda,0]].
\par\medskip\noindent\textbf{K.39. K --- Current register complete; final audit pending.}\par
Complete numbered statement of results.
\par\emph{Statement/proof location:} sections/K.tex, complete current numbered register.
\par Every current A-J and lead row is numbered with status and proof location. Final completeness requires resolution or the user-defined stopping review for every remaining research gate.
\par\medskip\noindent\textbf{K.40. SA14-H1 --- Proved.}\par
Generic Pic-tau restriction kernel and existence of its nonzero infinitesimal part.
\par\emph{Statement/proof location:} sections/A.tex, a:generic-unipotent and a:very-ample-example.
\emph{Manuscript reference:} 57:99-100.
\par For normal projective X the geometric-generic Pic-tau restriction kernel is finite connected unipotent; its tangent space is H1(L\textasciicircum{}-1). A3 gives an explicit nonzero example in dimension two.
\par\medskip\noindent\textbf{K.41. SA14-H2 --- Refuted as prompt formulated.}\par
Seshadri without Pic.
\par\emph{Statement/proof location:} work/bce/additional\_leads.tex; section.tex.
\emph{Manuscript reference:} 57:110.
\par Same as B1; omitted essential properness in source must be restored
\par\medskip\noindent\textbf{K.42. SA14-H3 --- Refuted as recalled; target undecided.}\par
Universal surjectivity without cohomological vanishing and recalled local analogue.
\par\emph{Statement/proof location:} sections/A.tex, A1r and tested approaches.
\emph{Manuscript reference:} 57:97.
\par Bhatt-de Jong local theorem is injectivity up to p-power torsion, not the recalled isomorphism; no global solution follows.
\par\medskip\noindent\textbf{K.43. SA14-H4 --- Undecided.}\par
Boxed universal injectivity in dimension at least four.
\par\emph{Statement/proof location:} sections/A.tex, A2b and A4b.
\emph{Manuscript reference:} 57:97.
\par Universal point injectivity and its least dimension remain undecided. The kernel in dimension at least three is a finite p-group; for each nonzero order-p class its bad smooth hyperplanes form the finite determinantal locus A2D. A2V forces a Cartier-exact leading term and Frobenius kernel on negative H1 of the hyperplane, with vanishing order divisible by p and a nonnormal mu\_p torsor. A2W forces rank(d:H0Omega1->H0Omega2)>=n-1. A2Y additionally requires noninjective Frobenius on H1(O\_Y) and no smooth W2 lifting of Y. In dimension3 the target is of general type by applying A3s to L\textasciicircum{}a. No global realization or universal exclusion of this remaining gate is proved.
\par\medskip\noindent\textbf{K.44. SA14-H5 --- Proved.}\par
All henselian claims.
\par\emph{Statement/proof location:} work/bce/additional\_leads.tex; section\_C.tex;Artin-Popescu;Elkik1973.
\emph{Manuscript reference:} 57:227-231;272.
\par C1 projectivity, full C2.G Picard completion for arbitrary excellent henselian local base and proper X with finite-length H1, and C3 punctured Pic by Elkik.
\par\medskip\noindent\textbf{K.45. SA14-H6 --- Proved with Pic algebraic-space hypothesis.}\par
Pic/Pic\ensuremath{\tau} exists as an algebraic space if tau flat fp; stronger closed version is schematic.
\par\emph{Statement/proof location:} work/bce/additional\_leads.tex; Stacks06PH;0418.
\emph{Manuscript reference:} 57:277.
\par Flat fp translation quotient is an algebraic space without closedness (prompt summary). HUNCHES additionally assumes tau closed: since tau is open, the quotient is separated unramified and a scheme. Both formulations preserved.
\par\medskip\noindent\textbf{K.46. SA14-H7 --- Unsettled as composite.}\par
Normal trait properness and scheme representability.
\par\emph{Statement/proof location:} work/bce/additional\_leads.tex; section.tex.
\emph{Manuscript reference:} 57:314-315.
\par Literal total-normal assertion false; intended connected relative-normal trait valuation proved; arbitrary-base geomnormal Pic schematicity not yet settled
\par\medskip\noindent\textbf{K.47. SA14-H8 --- Proved with complete dimension correction.}\par
Exceptional divisor lattice injects into NS exceptional fibre.
\par\emph{Statement/proof location:} work/bce/additional\_leads.tex; KollarFamilies Lem11.60 p436.
\emph{Manuscript reference:} 57:8.
\par Kernel exactly free lattice on one-dimensional branches; zero if every component has dimension>=2; use ±D nef and normalization
\par\medskip\noindent\textbf{K.48. SA14-H9 --- Proved.}\par
Corr always separated.
\par\emph{Statement/proof location:} work/bce/section\_E.tex; Raynaud1970 3.2(c),5.2;Murre1966 Thm4.
\emph{Manuscript reference:} 57:22;32;313.
\par Trait normalization into étale closure E of identity; connected fibres force map zero; Murre yields unramified separated scheme
\par\medskip\noindent\textbf{K.49. SA14-H10 --- Proved.}\par
SA14 lead H10.
\par\emph{Statement/proof location:} work/bce/additional\_leads.tex; Deligne1968 Thm5.5(i) p123;SGA3 VI\_B Prop1.6 p327.
\emph{Manuscript reference:} 57:64.
\par Deligne Thm5.5 gives conormal=(R1f*O)\textasciicircum{}dual; SGA3 criterion yields smooth Pic0; proper neutral locus yields abelian scheme
\par\medskip\noindent\textbf{K.50. SA14-H11 --- Proved.}\par
SA14 lead H11.
\par\emph{Statement/proof location:} work/bce/additional\_leads.tex; SGA3 XVI4.2 pp521-522;VI\_B1.4.2 pp325-326;SGA6 XII1.1-1.2.
\emph{Manuscript reference:} 57:74;78.
\par Kernel quotient and fibrewise-finite monomorphism argument, full proof Theorem H11.G.
\par\medskip\noindent\textbf{K.51. SA14-H12 --- Proved in original cone scope; refuted as arbitrary-local extrapolation.}\par
Local NS invariance modulo p-groups.
\par\emph{Statement/proof location:} work/bce/additional\_leads.tex; section.tex.
\emph{Manuscript reference:} 57:193-194.
\par 57:184–192 explicitly fixes cone vertex S;193 compares actual punctured Pic groups modulo identity components,194 first introduces henselization. Cone comparison is modulo bounded p-power. Arbitrary punctured-NS comparison fails 0\ensuremath{\to}Z\ensuremath{\to}Z in char3; a strict-henselian definition names a different invariant.
\par\medskip\noindent\textbf{K.52. SA14-H13 --- Proved/refuted components.}\par
R2 refinement and necessity.
\par\emph{Statement/proof location:} work/bce/additional\_leads.tex; section.tex.
\emph{Manuscript reference:} 57:111.
\par Sufficient with lci regular base; unnecessary; without lci not sufficient
\par\medskip\noindent\textbf{K.53. SA19-H1 --- Known; Undetermined chronology.}\par
Component-group pairing identified with SGA7 IX1.3 and Suzuki perfect-residue theorem; priority date not established.
\par\emph{Statement/proof location:} work/fhi/priority.md; priority.tex.
\emph{Manuscript reference:} 43:117; SA19-H1.
\par Bibliographic dates on 87 and144 and administrative list118 do not date scientific leaf117
\par\medskip\noindent\textbf{K.54. SA19-H2 --- Known identification; Undetermined chronology.}\par
Three-variable Ext matches biextension formalism; claimed pre-SGA7 priority not established.
\par\emph{Statement/proof location:} work/fhi/priority.md; priority.tex.
\emph{Manuscript reference:} 43:16;43:21; SA19-H2.
\par SGA7VII p133 attributes introduction to Mumford; precise archival dating remains gate
\par\medskip\noindent\textbf{K.55. SA19-H3 --- Undecided.}\par
SA19 lead H3.
\par\emph{Statement/proof location:} See the item discussion; no completed proof is claimed..
\par All dimensions are proved for separated complete graded algebras with a basis of open left ideals; all surfaces are proved for arbitrary jointly continuous topology. The original43:66 specifies the broader category without the open-left-ideal hypothesis. I7b now proves open dense differential-image quotients for whole source columns over arbitrary regular local rings. Remaining: construction and continuity of the entire adjacent-arrow family and its endpoint assembly in arbitrary complete targets, in dimension at least3.
\par\medskip\noindent\textbf{K.56. SA19-H4 --- Refuted.}\par
SA19 lead H4.
\par\emph{Statement/proof location:} See the item discussion; no completed proof is claimed..
\par The actual ordinary local Jacobian of k[s,t]\_(s,t) has an explicit nonsplit Z/p extension and is not projective; higher Ext vanishing for the curve pro-Jacobian is separately proved in Serre category and on the perfect etale site.
\par\medskip\noindent\textbf{K.57. SA19-H5 --- Refuted as stated; corrected programme unresolved.}\par
The sweeping open-status classification of the seven theorems.
\par\emph{Statement/proof location:} work/dgj/section\_J.tex.
\emph{Manuscript reference:} 44:2-9.
\par Universal integral intersection autoduality, raw positivity, literal higher-codimension height formula, middle fixed-part clause without smooth pencil axis, and numerical=tau are refuted. Primitive normalization, admissible height corrections, and manuscript-equivalence identification remain separate. Incidence existence/bound are known in every codimension. General regular representability and full stable/high Lefschetz assertions remain unresolved.
\par\medskip\noindent\textbf{K.58. SA19-H6 --- Complete scoped audit; inferred pi1 injectivity refuted.}\par
No readable rational-point-to-section injectivity proof in the supplied cote43 corpus; character-variance correction.
\par\emph{Statement/proof location:} work/dgj/section\_J.tex H6a-H6b; audit\_SA19H6.md; outputs/c43\_audit.tsv.
\emph{Manuscript reference:} 43:1-208 supplied readable transcriptions; exact positive locators43:27-35.
\par 184 complete files read by DGJ and11 by FHI. No sought proof found;2620 uncertainty macros,48 illegibility macros and13 absent scan labels limit the conclusion. Ext1 injectivity at43:30 detects surjectivity of the fundamental-group map, not injectivity. Original scans and historical nonexistence not established.
\par\medskip\noindent\textbf{K.59. A4b --- Proved.}\par
Finiteness and explicit p-power exponent for the restriction kernel and cokernel.
\par\emph{Statement/proof location:} Part A, Corollary a:pgroups.
\emph{Manuscript reference:} 57:97.
\par Both groups for n>=4; kernel for n=3. Exponent p\textasciicircum{}e for p\textasciicircum{}e>=N+1.
\par\medskip\noindent\textbf{K.60. A4c --- Known.}\par
Picard restriction under the stated negative-cohomology vanishings.
\par\emph{Statement/proof location:} Part A, Corollary a:vanishing.
\emph{Manuscript reference:} 57:97.
\par Hypotheses verified for smooth X when the stated vanishings hold; not supplied in A1.
\par\medskip\noindent\textbf{K.61. A3r --- Proved.}\par
Any very ample H1-negative counterexample yields a smooth surface counterexample.
\par\emph{Statement/proof location:} Part A, Proposition a:surface.
\emph{Manuscript reference:} 57:100.
\par Strict reduction of the existence gate; not an existence theorem.
\par\medskip\noindent\textbf{K.62. A1r --- Proved.}\par
The universal dimension>=4 Picard isomorphism reduces to dimension four.
\par\emph{Statement/proof location:} Part A, Proposition a:fourfold.
\emph{Manuscript reference:} 57:97.
\par Auxiliary high-degree cut preserves the original hyperplane embedding on the cut.
\par\medskip\noindent\textbf{K.63. A2b --- Proved.}\par
A smooth Segre surface hyperplane has restriction kernel Z(1,-1).
\par\emph{Statement/proof location:} Part A, Proposition a:segre.
\emph{Manuscript reference:} 57:97.
\par Valid in every characteristic; proves only lower bound three for a possible universal threshold.
\par\medskip\noindent\textbf{K.64. R3 --- Proved.}\par
Tangent of the Picard restriction kernel equals h H1(X,L\textasciicircum{}-1).
\par\emph{Statement/proof location:} Part A, Proposition a:tangent.
\emph{Manuscript reference:} 57:99-100.
\par The equivalence K=0 iff tangent=0 requires K finite connected; generic connectedness audited separately.
\par\medskip\noindent\textbf{K.65. R4 --- Proved.}\par
Restriction kernel/cokernel are bounded finite p-groups in the dimensions stated.
\par\emph{Statement/proof location:} Part A, Corollary a:pgroups.
\emph{Manuscript reference:} 57:97.
\par The phrase built from layers requires obstruction maps and extension data, not direct sums.
\par\medskip\noindent\textbf{K.66. B2-necessity --- Refuted.}\par
Is R2 necessary?.
\par\emph{Statement/proof location:} work/bce/section\_B.tex; section.tex.
\emph{Manuscript reference:} 57:111.
\par Projective regular smoothing xy-z\ensuremath{{}^2}-tw\ensuremath{{}^2}=0 has non-R2 special quadric cone but locally factorial total space; compare xy-z\ensuremath{{}^2}-t\ensuremath{{}^2}w\ensuremath{{}^2}=0 over same trait
\par\medskip\noindent\textbf{K.67. B2-sufficiency --- Proved with original lci hypothesis.}\par
R2 fibres over regular base suffice.
\par\emph{Statement/proof location:} work/bce/section\_B.tex; SGA2 XI 3.13(ii).
\emph{Manuscript reference:} 57:111.
\par Parafactoriality in local dimension >=4; without lci explicit Segre-cone projective family refutes it
\par\medskip\noindent\textbf{K.68. B2-exact --- Proved impossibility of a fibre-type-only criterion; fixed-family iff proved.}\par
Determine exact condition on fibres and S.
\par\emph{Statement/proof location:} work/bce/section\_B.tex; section.tex.
\emph{Manuscript reference:} 57:111.
\par Proper lci families X1 and X2 have identical base and geometric fibre types but opposite extension properties. Thus these data cannot determine the answer. For regular S and normal integral geometric fibres the exact deformation-sensitive obstruction is kernel Cl/Pic\ensuremath{\to}Cl\_generic/Pic\_generic.
\par\medskip\noindent\textbf{K.69. C2-char0 --- Proved in characteristic zero.}\par
Pic(X)=Pic(completion) when R1 has finite length.
\par\emph{Statement/proof location:} work/bce/section\_C.tex; Artin1969 Thm3.5;PerlingSchroer Prop1.2 proof.
\emph{Manuscript reference:} 57:227-231;272.
\par Eventual H1(m\textasciicircum{}n O)=0; formal log; finite determination plus Artin density
\par\medskip\noindent\textbf{K.70. C2-mixed --- Proved in all characteristics for every excellent henselian local base.}\par
Same over excellent DVR bases.
\par\emph{Statement/proof location:} work/bce/section\_C.tex; section\_C.tex Theorem C2.G;Stacks02OB(3),02OS,03H0;PerlingSchroer Prop1.2 E=O;Artin-Popescu.
\emph{Manuscript reference:} 57:227-231;272.
\par Full Theorem C2.G: blow up maximal ideal, Stein factor gives Cartier ideal with R1 supported exceptionally; principal-ideal finite determination, auxiliary Stein H1 and H2 finite length, conductor-unit descent via uniform formal functions, Artin approximation. No flatness assumption on X.
\par\medskip\noindent\textbf{K.71. C4-cones --- Proved in original cone context.}\par
Local cone comparison modulo p-groups.
\par\emph{Statement/proof location:} work/bce/section\_C.tex; section.tex.
\emph{Manuscript reference:} 57:193-194.
\par Split formal neighbourhood; finite filtration by additive H1(O\_X(n)) killed by p; Serre vanishing bounds exponent
\par\medskip\noindent\textbf{K.72. C4-general --- Scope-dependent refutation.}\par
General local NS comparison modulo p-groups.
\par\emph{Statement/proof location:} work/bce/section\_C.tex; section.tex.
\emph{Manuscript reference:} 57:193-194.
\par Char3 factorial local hypersurface xy=z\ensuremath{{}^2}+t\ensuremath{{}^2}(1+t) has completed/henselian punctured Pic Z; intrinsic étale-local NS conventions must be distinguished
\par\medskip\noindent\textbf{K.73. E1-literal --- Refuted.}\par
Normal total space implies extension of algebraically trivial generic bundle.
\par\emph{Statement/proof location:} work/bce/section\_E.tex; section.tex.
\emph{Manuscript reference:} 57:314-315.
\par Normal projective cubic Y\ensuremath{{}^2}Z=X\ensuremath{{}^3}+X\ensuremath{{}^2}Z+t\ensuremath{{}^2}Z\ensuremath{{}^3}; P-O generic degree zero does not extend
\par\medskip\noindent\textbf{K.74. E1-relative --- Proved in intended valuative sense.}\par
Proper flat geometrically normal connected family over trait.
\par\emph{Statement/proof location:} work/bce/section\_E.tex; Raynaud1970 Thm3.3.1;FGA VI Cor2.3.
\emph{Manuscript reference:} 57:314-315.
\par Separated Pic algebraic space is scheme over trait; proper neutral subset yields extension with special class in Pic0
\par\medskip\noindent\textbf{K.75. E1-functor --- Refuted in general.}\par
Pic0 subfunctor is a proper algebraic space.
\par\emph{Statement/proof location:} work/bce/section\_E.tex; LaurentSchroer2024 Prop2.2.
\emph{Manuscript reference:} 57:314-315.
\par Fibrewise Pic0 may not be representable even for smooth projective Enriques families
\par\medskip\noindent\textbf{K.76. SA14-H10-main --- Proved.}\par
Pic0 abelian for smooth projective char0 arbitrary base.
\par\emph{Statement/proof location:} work/bce/additional\_leads.tex; Deligne1968 Thm5.5(i) p123;SGA3 VI\_B Prop1.6 p327.
\emph{Manuscript reference:} 57:64.
\par Deligne Thm5.5 gives conormal=(R1f*O)\textasciicircum{}dual; SGA3 criterion yields smooth Pic0; proper neutral locus yields abelian scheme
\par\medskip\noindent\textbf{K.77. SA14-H10-conormal --- Proved.}\par
Locally free conormal at unit in char0 implies smoothness without flatness.
\par\emph{Statement/proof location:} work/bce/additional\_leads.tex; SGA3 VI\_B Prop1.6 original p327.
\emph{Manuscript reference:} 57:64.
\par For group lfp along identity over Q-base, omega locally free iff smooth there; no flatness assumed
\par\medskip\noindent\textbf{K.78. SA14-H11-flattening --- Proved.}\par
Flattening problem A.
\par\emph{Statement/proof location:} work/bce/additional\_leads.tex; Stacks05PS;0CWQ.
\emph{Manuscript reference:} 57:73-79.
\par Same as E3
\par\medskip\noindent\textbf{K.79. SA14-H11-affinity --- Proved over every integral one-dimensional noetherian base after shrinking.}\par
Replace quasi-affine Pic pullback by affine after shrinking dimension1 base.
\par\emph{Statement/proof location:} work/bce/additional\_leads.tex; SGA3 XVI4.2;VI\_B1.4.2;SGA6 XII1.1-1.2.
\emph{Manuscript reference:} 57:74;78.
\par Full Theorem H11.G via affine flat kernel quotient and fibrewise-finite monomorphism; no normality/projectivity assumption.
\par\medskip\noindent\textbf{K.80. E4-308 --- No mathematical question.}\par
Title leaf.
\par\emph{Statement/proof location:} work/bce/section\_E.tex; transcription.
\emph{Manuscript reference:} 57:308.
\par Title only
\par\medskip\noindent\textbf{K.81. E4-309a --- Refuted.}\par
Pic disjoint sum of finite-type schemes?.
\par\emph{Statement/proof location:} work/bce/section\_E.tex; Kleiman2005 Ex4.14 and deduction.
\emph{Manuscript reference:} 57:309.
\par Semistable split conic has connected component with infinitely many origins
\par\medskip\noindent\textbf{K.82. E4-309aprime --- Refuted.}\par
Closure of quasi-compact subset quasi-compact?.
\par\emph{Statement/proof location:} work/bce/section\_E.tex; Kleiman2005 Ex4.14 and deduction.
\emph{Manuscript reference:} 57:309.
\par Closure of generic point in conic Pic contains infinitely many origins
\par\medskip\noindent\textbf{K.83. E4-309adoubleprime --- Refuted generally.}\par
Pic\ensuremath{\tau} closed in Pic?.
\par\emph{Statement/proof location:} work/bce/section\_E.tex; Kleiman2005 Ex4.14;FGA VI Cor2.3.
\emph{Manuscript reference:} 57:309.
\par Conic Pic tau identity copy is open not closed; normal morphism positive subset result retained separately
\par\medskip\noindent\textbf{K.84. E4-309b --- Refuted in written generality.}\par
Pic/Pic\ensuremath{\tau} quotient scheme?.
\par\emph{Statement/proof location:} work/bce/section\_E.tex; Theorem E4.Q; Kleiman2005 Example4.14 pp31–33.
\emph{Manuscript reference:} 57:309.
\par Nonsplit Mumford conic over R[[t]] has Pic\ensuremath{\tau}=trivial subgroup S but full Pic is a nonschematic algebraic space; quotient equals Pic. Even flat tau is insufficient without closedness. Stronger variants with full Pic assumed schematic or closed nonflat tau are separate.
\par\medskip\noindent\textbf{K.85. E4-309c --- Refuted, including smooth normal families.}\par
Pic\ensuremath{\tau} universally open?.
\par\emph{Statement/proof location:} work/bce/section\_E.tex; section.tex.
\emph{Manuscript reference:} 57:309.
\par E4.H has a clopen character component supported on the closed point. E4.R also refutes openness with reduced fibres by dimension jump.
\par\medskip\noindent\textbf{K.86. E4-309cflat --- Refuted.}\par
Margin: Pic\ensuremath{\tau} flat?.
\par\emph{Statement/proof location:} work/bce/section\_E.tex; PartD.
\emph{Manuscript reference:} 57:309.
\par E4.H gives smooth projective X/W(k), char(k)=2, with Pic\ensuremath{\tau}=\ensuremath{\mu}2\ensuremath{{}^2}×(SpecR disjoint Speck), coordinate module R\textasciicircum{}4\ensuremath{\oplus}k\textasciicircum{}4. Explicit degree768 construction and root audit completed.
\par\medskip\noindent\textbf{K.87. E4-309d --- Proved under standard Pic representability hypotheses.}\par
Pic\ensuremath{\tau} quasi-compact hence finite type?.
\par\emph{Statement/proof location:} work/bce/section\_E.tex; LaurentSchroer2024 Thm2.1.
\emph{Manuscript reference:} 57:309.
\par Open subspace of finite presentation for proper flat fp cohflat0
\par\medskip\noindent\textbf{K.88. E4-309e --- Refuted even when represented.}\par
Pic\ensuremath{\tau} separated when representable?.
\par\emph{Statement/proof location:} work/bce/section\_E.tex; LaurentSchroer2024 Prop2.3(i).
\emph{Manuscript reference:} 57:309.
\par Free K3 involution combined with t=u\ensuremath{{}^2} produces regular projective X with double special fibre; O(Z) has order2 and generic restriction trivial but nontrivial special restriction; Raynaud5.3 ensures full Pic is a scheme
\par\medskip\noindent\textbf{K.89. E4-309bprime --- Mathematical readings settled; pale abbreviation retained.}\par
Semicontinuity of unsure\{Min.|Dim.\} Lissité.
\par\emph{Statement/proof location:} work/bce/section\_E.tex; Official MPL57 scan309; FGA VI Cor2.7 p234; Theorem E4.S; DGJ targeted audit.
\emph{Manuscript reference:} 57:309.
\par Original scan309 prefers Dim. Lissite but remains pale. Fibre dimension is upper semicontinuous; relative smooth locus is open. E4.S gives geometrically integral reduced surface family with nonsmooth generic Pic-tau and smooth special Pic-tau, refuting openness of fibrewise smoothness and upper semicontinuity of tangent-dimension defect. No missing direction is invented.
\par\medskip\noindent\textbf{K.90. E4-309f --- Refuted generally; known positive under geometric normality.}\par
Continuity of ell-rank of Pic fibres.
\par\emph{Statement/proof location:} work/bce/section\_E.tex; FGA VI Cor2.7.
\emph{Manuscript reference:} 57:309.
\par Projective cubic y\ensuremath{{}^2}z=x\ensuremath{{}^3}+x\ensuremath{{}^2}z+tz\ensuremath{{}^3} over C[[t]] has smooth elliptic generic fibre and integral nodal special fibre. Pic[\ensuremath{\ell}] has geometric F\ensuremath{\ell} dimension2 generically and1 specially; normal-family constancy remains FGA VI Cor2.7.
\par\medskip\noindent\textbf{K.91. E4-309g --- Proved under universal H0=O.}\par
Pic\ensuremath{\tau} product formula.
\par\emph{Statement/proof location:} work/bce/section\_E.tex; section.tex forthcoming.
\emph{Manuscript reference:} 57:309.
\par E2 represents Corr by separated unramified scheme; based maps to affine n-torsion prove its geometric fibres torsion-free; fibrewise vanishing descends. Connectedness alone false via dual-number thickening
\par\medskip\noindent\textbf{K.92. E4-309h --- Delegated to root Part A.}\par
Hyperplane Pic\ensuremath{\tau} bijective/isomorphism and surjectivity.
\par\emph{Statement/proof location:} work/bce/section\_E.tex; PartA.
\emph{Manuscript reference:} 57:309.
\par Conditions explicitly left unspecified in manuscript; preserve unsure reading
\par\medskip\noindent\textbf{K.93. E4-3101 --- Proved in represented setting.}\par
Pic\ensuremath{\tau} closed if f normal?.
\par\emph{Statement/proof location:} work/bce/section\_E.tex; FGA VI Cor2.3 p232.
\emph{Manuscript reference:} 57:310.
\par FGA VI Cor2.3; does not imply representability of Pic0
\par\medskip\noindent\textbf{K.94. E4-3102 --- Refuted.}\par
Pic\ensuremath{\tau} flat?.
\par\emph{Statement/proof location:} work/bce/section\_E.tex; PartD.
\emph{Manuscript reference:} 57:310.
\par Same fully specified Heisenberg degeneration as E4.H; independent of the rejected ordinary-Igusa tangent argument.
\par\medskip\noindent\textbf{K.95. E4-3103 --- Refuted in general and for normal families in mixed characteristic.}\par
Pic\ensuremath{\tau} universally open incl normal and char0 cases.
\par\emph{Statement/proof location:} work/bce/section\_E.tex; FGA VI Rem2.9.
\emph{Manuscript reference:} 57:310.
\par E4.H settles the residue-characteristic2 normal-family case negatively. The normal characteristic-zero positive statement remains valid.
\par\medskip\noindent\textbf{K.96. E4-3104 --- Refuted even for a smooth projective family.}\par
Multiplication by every n universally open?.
\par\emph{Statement/proof location:} work/bce/section\_E.tex; section.tex;FGA VI Thm2.5.
\emph{Manuscript reference:} 57:310.
\par E4.H is killed by2, and its identity section is not open: the four geometric generic branches of \ensuremath{\mu}2\ensuremath{{}^2} meet at the special identity. Thus [2] is not open. The cuspidal Ga example gives a simpler field counterexample; invertible n is étale.
\par\medskip\noindent\textbf{K.97. E4-3104ft --- Known with exact theorem.}\par
Margin: multiplication by n finite type?.
\par\emph{Statement/proof location:} work/bce/section\_E.tex; SGA6 XIII Thm3.6 p635;Kleiman Rem6.28.
\emph{Manuscript reference:} 57:310.
\par n nonzero; proper X and represented Pic
\par\medskip\noindent\textbf{K.98. E4-3105ft --- Proved.}\par
Pic\ensuremath{\tau} finite type?.
\par\emph{Statement/proof location:} work/bce/section\_E.tex; LaurentSchroer2024 Thm2.1.
\emph{Manuscript reference:} 57:310.
\par Same as309d
\par\medskip\noindent\textbf{K.99. E4-3105proper --- Refuted.}\par
Pic\ensuremath{\tau} proper for normal morphism?.
\par\emph{Statement/proof location:} work/bce/section\_E.tex; LaurentSchroer2024 Prop2.3(ii).
\emph{Manuscript reference:} 57:310.
\par Explicit family of diagonal K3 three-quadric intersections modulo sign involution: generic Enriques double-cover class of order2 cannot extend across four fixed nodes; all geometric fibres normal integral
\par\medskip\noindent\textbf{K.100. E4-3105projective --- Refuted normal branch; smooth branch positive.}\par
Pic\ensuremath{\tau} projective for smooth or normal projective morphism?.
\par\emph{Statement/proof location:} work/bce/section\_E.tex; Kleiman Ex6.18.
\emph{Manuscript reference:} 57:310.
\par Same normal projective quotient family has nonproper tau, hence nonprojective; smooth projective branch remains the classical positive theorem
\par\medskip\noindent\textbf{K.101. E4-3105prime --- Refuted.}\par
Universal openness implies finite type/proper under normality?.
\par\emph{Statement/proof location:} work/bce/section\_E.tex; transcription.
\emph{Manuscript reference:} 57:310.
\par Normal K3-quotient family in characteristic0 has universally open tau but nonproper tau
\par\medskip\noindent\textbf{K.102. E4-3105diagram --- Refuted for multiplication openness; reading qualifier uncertain.}\par
Margin diagram and multiple-p surjective/open assertion.
\par\emph{Statement/proof location:} work/bce/section\_E.tex; Official MPL57 scan310; enlarged crop MPL57\_310\_diagram.png; Theorem E4.H.
\emph{Manuscript reference:} 57:310.
\par Theorem E4.H explicitly makes multiplication by2 nonopen for a smooth projective family. The mathematical openness assertion is refuted. Original scan310 confirms Non exemple and the branch diagram, but the surrounding surjectivity/component qualifier remains illegible; no additional assertion is invented.
\par\medskip\noindent\textbf{K.103. E4-3106 --- Refuted.}\par
Only finitely many S-points killed by n?.
\par\emph{Statement/proof location:} work/bce/section\_E.tex; section.tex.
\emph{Manuscript reference:} 57:310.
\par Cuspidal cubic Ga over infinite algclosed charp has infinitely many p-torsion points
\par\medskip\noindent\textbf{K.104. E4-3106criterion --- Picard-scope conditions proved; exact iff reading uncertain.}\par
Bracketed criterion for finite type via components over constructible dense opens.
\par\emph{Statement/proof location:} work/bce/section\_E.tex; Official MPL57 scan310; new\_E4\_scope\_310.tex; Laurent-Schroer2.1.
\emph{Manuscript reference:} 57:310.
\par For proper flat fp cohomologically-flat-in-degree-zero families over noetherian S, Pic-tau is fp; both defensible above-open and outside-open finite-component conditions therefore hold for closed irreducible base changes. The exact broader iff quantifiers remain unreadable. An explicit infinite-branch etale group refutes the generic-open component-count criterion for arbitrary lft groups; it is not a Picard realization.
\par\medskip\noindent\textbf{K.105. E4-3107 --- Known with exact theorem.}\par
NS ranks uniformly bounded?.
\par\emph{Statement/proof location:} work/bce/section\_E.tex; SGA6 XIII Thm5.1 p650;Kleiman Rem6.19.
\emph{Manuscript reference:} 57:310.
\par Proper family over noetherian base; rank bounded though not constructible
\par\medskip\noindent\textbf{K.106. E4-3108 --- Refuted.}\par
Connected components of Pic finite type?.
\par\emph{Statement/proof location:} work/bce/section\_E.tex; Kleiman Ex4.14;Cor6.25.
\emph{Manuscript reference:} 57:310.
\par Split conic infinite-origin component; positive for projective geomint family
\par\medskip\noindent\textbf{K.107. E4-3108copies --- Refuted.}\par
Pic disjoint union of copies of Pic\ensuremath{\tau}?.
\par\emph{Statement/proof location:} work/bce/section\_E.tex; Kleiman Ex4.14 and deduction.
\emph{Manuscript reference:} 57:310.
\par Split conic: one connected component has infinitely many origins; tau is one trait
\par\medskip\noindent\textbf{K.108. E4-3109 --- Proved standard quotient.}\par
Definition of Pic/Pic\ensuremath{\tau} over field.
\par\emph{Statement/proof location:} work/bce/section\_E.tex; H6 over field;SGA6 XIII.
\emph{Manuscript reference:} 57:310.
\par Etale group scheme with geometric points NS/torsion; over nonclosed field retain Galois action
\par\medskip\noindent\textbf{K.109. E4-310Howard --- Multiplicity invariance refuted; existence/descent words not a typed statement.}\par
Multiplicity compatibility with specialization; existence/descent.
\par\emph{Statement/proof location:} work/bce/section\_E.tex; PartD;transcription.
\emph{Manuscript reference:} 57:310.
\par E4.H: generic four reduced points versus two special components of length4. Total length4\ensuremath{\to}8, component count4\ensuremath{\to}2, neutral geometric length1\ensuremath{\to}4. Further object for existence/descent is unspecified.
\par\medskip\noindent\textbf{K.110. E4-311open --- Refuted, including smooth normal families.}\par
Pic\ensuremath{\tau} universally open?.
\par\emph{Statement/proof location:} work/bce/section\_E.tex; section.tex.
\emph{Manuscript reference:} 57:311.
\par E4.H directly disproves openness by the special-only character component; E4.R supplies a reduced-fibre dimension-jump example.
\par\medskip\noindent\textbf{K.111. E4-311flat0 --- Refuted in full generality.}\par
Pic0 flat?.
\par\emph{Statement/proof location:} work/bce/section\_E.tex; LaurentSchroer Prop2.2;PartD.
\emph{Manuscript reference:} 57:311.
\par May not even be represented; Part D simple Igusa gives fibre rank jump
\par\medskip\noindent\textbf{K.112. E4-311flatimp --- Refuted.}\par
Pic0 flat implies Pic\ensuremath{\tau} flat?.
\par\emph{Statement/proof location:} work/bce/section\_E.tex; Theorem E4.H;Raynaud1979;explicit character scheme calculation.
\emph{Manuscript reference:} 57:311.
\par Basechange Theorem E4.H to A=W(k)/4. Pic0=\ensuremath{\mu}2\ensuremath{{}^2}\_A flat rank4; Pic\ensuremath{\tau} has additional open component \ensuremath{\mu}2\ensuremath{{}^2}\_k, whose coordinate module k4 is not A-flat. Parent audit passed.; explicit degree768 norm system and first-jet proof, root structural audit accepted
\par\medskip\noindent\textbf{K.113. E4-311canonicalArtin --- Refuted by an explicit dual-number Igusa deformation.}\par
Canonical abelian Pic subgroup on Artin local base?.
\par\emph{Statement/proof location:} work/bce/section\_E.tex; Part D D4a; Katz1981 Theorem2.1(3),(4) pp148–149.
\emph{Manuscript reference:} 57:311.
\par D4a gives k=Fbar2, R=k[epsilon]/epsilon\textasciicircum{}2 and ordinary Serre-Tate parameters q12=1+epsilon, all other qij=1. A compatible abelian Picard-variety subgroup would force 1=q12 by Katz morphism lifting. Pic-tau remains flat. Root, BCE and FHI targeted audits passed; the uncertain manuscript notation is retained.
\par\medskip\noindent\textbf{K.114. E4-311canonicalReduced --- Refuted.}\par
Canonical abelian Pic subgroup on reduced base?.
\par\emph{Statement/proof location:} work/bce/section\_E.tex; Part D D3.
\emph{Manuscript reference:} 57:311.
\par Part D3 cyclic shear gives an explicit reduced trait base with no embedded abelian subgroup having the fibrewise Picard varieties; kernel lengths of (F\textasciicircum{}2,2) change from 2 to 4.
\par\medskip\noindent\textbf{K.115. E4-311NS --- Refuted.}\par
Continuity of torsion of NS.
\par\emph{Statement/proof location:} work/bce/section\_E.tex; Part D D2 and product Igusa computation.
\emph{Manuscript reference:} 57:311.
\par Part D4 product Igusa family has ordinary generic NS torsion Z/2 and supersingular special NS torsion zero; the component-group assertion is false.
\par\medskip\noindent\textbf{K.116. E4-311Igusa --- Proved.}\par
Study modular family of Igusa surface.
\par\emph{Statement/proof location:} work/dgj/section\_D.tex D4o and D4s; section\_D\_supersingular.tex.
\emph{Manuscript reference:} 57:311.
\par Both algebraized miniversal hulls computed: ordinary W[[t,s,u,v]]/(s\ensuremath{{}^2}+2s); supersingular W[[t,s,u,v]]/(s\ensuremath{{}^4}−t*s\ensuremath{{}^2}−2s). Canonical fixed elliptic subgroup and all-Artinian Pic-tau flatness proved; root/fhi independent audits pass.
\par\medskip\noindent\textbf{K.117. E4-312 --- No Picard question.}\par
Tensor-product insertion.
\par\emph{Statement/proof location:} work/bce/section\_E.tex; transcription.
\emph{Manuscript reference:} 57:312.
\par Unrelated Bourbaki tensor-product leaf retained in inventory
\par\medskip\noindent\textbf{K.118. E4-313i --- Proved.}\par
Pic separated if fibres geometrically locally integral?.
\par\emph{Statement/proof location:} work/bce/section\_E.tex; LaurentSchroer Prop2.3(i);Raynaud1970 6.1.9.
\emph{Manuscript reference:} 57:313.
\par Finite-etale Stein reduction makes locally integral connected fibres integral; Pic separation descends
\par\medskip\noindent\textbf{K.119. E4-313ii --- Proved.}\par
Pic\ensuremath{\tau} separated if fibres geometrically reduced?.
\par\emph{Statement/proof location:} work/bce/section\_E.tex; Raynaud1970 Ex6.1.6 p49;Thm6.4.1(3) p56;Cor6.4.5 p57.
\emph{Manuscript reference:} 57:313.
\par Geometrically reduced fibres imply propertyN and dprime=1 on alltraits; Raynaud criterion yields separated tau
\par\medskip\noindent\textbf{K.120. E4-313ntorsion --- Proved under geometrically reduced fibres.}\par
n-torsion Pic separated under same fibre hypothesis?.
\par\emph{Statement/proof location:} work/bce/section\_E.tex; Raynaud1970 Cor6.4.5 and deduction.
\emph{Manuscript reference:} 57:313.
\par n-torsion is a subfunctor of separated tau
\par\medskip\noindent\textbf{K.121. E4-313dim --- Refuted even for geometrically reduced fibres.}\par
Constancy of dim Pic0.
\par\emph{Statement/proof location:} work/bce/section\_E.tex; section\_E.tex Theorem E4.R;SGA3 VI\_B4.2;FGA VI2.7.
\emph{Manuscript reference:} 57:313.
\par Two ruled surfaces glued along a split Poincare ribbon give a projective flat geometrically reduced connected family with universal H0=O and Pic0=E times Ga\textasciicircum{}\{h0(P\_s)\}, dimensions1to2. Normal/no-additive case remains constant.
\par\medskip\noindent\textbf{K.122. E4-313alpha --- Proved in represented proper-flat Picard setting.}\par
Abelian-rank inequality alpha<=alpha prime.
\par\emph{Statement/proof location:} work/bce/section\_E.tex; Artin1969 Thm3.1 p30;GilleGuralnick2023 Thms1.1,4.1;Raynaud1970 3.3.1,5.2.
\emph{Manuscript reference:} 57:313.
\par Unprimed=special. Closure G of generic tau is flat; prime-to-p Kummer torsor lifting makes Pic0\_special/G0\_special unipotent; constant dimension and Gille-Guralnick upper semicontinuity of maximal affine commutative subgroup dimension give alpha\_special<=alpha\_generic.
\par\medskip\noindent\textbf{K.123. E4-313lambda --- Proved in represented proper-flat Picard setting.}\par
Additive-rank inequality lambda>=lambda prime.
\par\emph{Statement/proof location:} work/bce/section\_E.tex; GilleGuralnick2023 Thms1.1,4.1 pp1,7-8;Raynaud1970 3.3.1,4.1,5.2.
\emph{Manuscript reference:} 57:313.
\par Unprimed=special. Gille-Guralnick maximal-unipotent-dimension upper semicontinuity applies directly to separated tau. For represented nonseparated Pic use Raynaud separated quotient and reduced neutral groups.
\par\medskip\noindent\textbf{K.124. E4-313ellrank --- Proved in the stated geometrically reduced setting.}\par
Known inequality 2alpha+mu<=2alpha prime+mu prime.
\par\emph{Statement/proof location:} work/bce/section\_E.tex; FGA VI Remark2.8 original pp234–235, visually verified; Theorem2.5 p233; E4.rank.2.
\emph{Manuscript reference:} 57:313.
\par Unprimed is special and primed generic. Choose ell away from residue characteristic and both NS torsion orders; separated etale Pic[ell] lifts distinct special torsion points to distinct generic ones, giving 2alpha\_special+mu\_special<=2alpha\_generic+mu\_generic. Also valid whenever the relevant prime-to-characteristic torsion subspaces are separated.
\par\medskip\noindent\textbf{K.125. E4-313Corr --- Proved.}\par
Corr exists for geomint proper flat families.
\par\emph{Statement/proof location:} work/bce/section\_E.tex; Murre Thm4 plus E2 proof.
\emph{Manuscript reference:} 57:313.
\par E2 stronger hypotheses on H0 and flatness give separated unramified scheme
\par\medskip\noindent\textbf{K.126. E4-314open --- Proved under standard cohomological representability hypotheses.}\par
Pic\ensuremath{\tau} open of finite presentation.
\par\emph{Statement/proof location:} work/bce/section\_E.tex; LaurentSchroer Thm2.1.
\emph{Manuscript reference:} 57:314.
\par Exact modern theorem; arbitrary Pic functor openness also classical
\par\medskip\noindent\textbf{K.127. E4-314locus --- Proved.}\par
Tau-triviality locus open retrocompact.
\par\emph{Statement/proof location:} work/bce/section\_E.tex; LaurentSchroer Thm2.1.
\emph{Manuscript reference:} 57:314.
\par Pullback of open fp Pic\ensuremath{\tau} by the section representing L
\par\medskip\noindent\textbf{K.128. E4-314closedProjective --- Proved over arbitrary scheme base.}\par
Tau locus open closed for projective geomint family.
\par\emph{Statement/proof location:} work/bce/section\_E.tex; Zheng2009 Lemma3.8 pp313–314; Kleiman2005 Th6.3(h) pp53–54 and B.27 p79; Kollar2022 Lemma11.60 p436; Stacks054F,05F4,055B,04AI.
\emph{Manuscript reference:} 57:314.
\par Theorem E4.T: proper flat fp geometrically integral fibres suffice, without projectivity. Regular Galois alteration, Hodge-index vertical correction, valuation transitivity and negativity prove trait specialization; noetherian approximation gives arbitrary scheme base. Separate root and DGJ AI audits passed.
\par\medskip\noindent\textbf{K.129. E4-314removeProjective --- Proved over arbitrary scheme base.}\par
Remove projectivity from closed tau locus under geomint fibres.
\par\emph{Statement/proof location:} work/bce/section\_E.tex; Zheng2009 Lemma3.8 pp313–314; Kleiman2005 Th6.3(h) pp53–54 and B.27 p79; Kollar2022 Lemma11.60 p436; Stacks054F,05F4,055B,04AI.
\emph{Manuscript reference:} 57:314.
\par Theorem E4.T: proper flat fp geometrically integral fibres suffice, without projectivity. Regular Galois alteration, Hodge-index vertical correction, valuation transitivity and negativity prove trait specialization; noetherian approximation gives arbitrary scheme base. Separate root and DGJ AI audits passed.
\par\medskip\noindent\textbf{K.130. E4-314normal0 --- Proved as subset/valuative statement.}\par
Neutral part closed/proper for normal morphism.
\par\emph{Statement/proof location:} work/bce/section\_E.tex; FGA VI Thm1.1 and Cor2.3.
\emph{Manuscript reference:} 57:314.
\par Not representability of fibrewise neutral subfunctor
\par\medskip\noindent\textbf{K.131. E4-315extension --- Proved intended / refuted literal.}\par
Algebraically trivial generic bundle extends.
\par\emph{Statement/proof location:} work/bce/section\_E.tex; section.tex.
\emph{Manuscript reference:} 57:315.
\par E1 distinction between relative normality and normal total space
\par\medskip\noindent\textbf{K.132. E4-315schemeNormal --- Unsettled arbitrary base.}\par
Pic scheme for proper flat geomnormal fibres.
\par\emph{Statement/proof location:} work/bce/section\_E.tex; Raynaud1970 Thm3.3.1; Theorem E4.Curve; BLR6.6(2c); Cesnavicius2015 provenance.
\emph{Manuscript reference:} 57:315.
\par Projective families, normal noetherian dimension-one bases, and all geometrically normal curve families over arbitrary bases are schematic. The remaining gate is arbitrary-base higher-dimensional geometrically normal families; no genuine Picard counterexample or general proof is obtained. See E4\_schematicity\_attempts.md.
\par\medskip\noindent\textbf{K.133. E4-315schemeIntegral --- Unsettled arbitrary base.}\par
Pic scheme for proper flat geomint fibres.
\par\emph{Statement/proof location:} work/bce/section\_E.tex; Laurent-Schroer2.1; Raynaud1970 Thm3.3.1; Theorem E4.Curve; E4\_schematicity\_attempts.md.
\emph{Manuscript reference:} 57:315.
\par Artin gives a separated algebraic space. Projective families and smooth proper curves over arbitrary bases are schematic. Integral singular curves and higher-dimensional integral families remain. The standard nonsplit conic has reducible geometric special fibre and does not refute this question; arbitrary nonschematic group spaces also do not provide Picard realizations.
\par\medskip\noindent\textbf{K.134. D0-R6 --- Refuted.}\par
The unrestricted full Picard product of R6.
\par\emph{Statement/proof location:} work/dgj/section\_D.tex.
\emph{Manuscript reference:} 57:286.
\par Missing manuscript hypotheses A\_Gamma=0; Poincare line bundle on E x E disproves unrestricted product.
\par\medskip\noindent\textbf{K.135. D0-tau --- Proved under stated hypotheses.}\par
Pic-tau quotient formula for abelian A,B and a constant translation subgroup on B.
\par\emph{Statement/proof location:} work/dgj/section\_D.tex.
\emph{Manuscript reference:} 57:286.
\par All examples satisfy the stated hypotheses; relative fppf sheaf proof includes nonreduced bases.
\par\medskip\noindent\textbf{K.136. D1 --- Proved.}\par
Both signs and variation of quotient Hodge asymmetry.
\par\emph{Statement/proof location:} work/dgj/section\_D.tex.
\emph{Manuscript reference:} 57:287.
\par Characteristic two shearing involutions with F and V; dimensions (2,3),(3,2), specializing to (3,3).
\par\medskip\noindent\textbf{K.137. D2 --- Proved.}\par
A nonflat fixed subgroup and nonflat Pic-tau in a smooth projective quotient family.
\par\emph{Statement/proof location:} work/dgj/section\_D.tex.
\emph{Manuscript reference:} 57:288.
\par Fixed lengths 8 and 16; algebra R[z]/(z\textasciicircum{}4,t z\textasciicircum{}2).
\par\medskip\noindent\textbf{K.138. D3 --- Proved.}\par
No abelian subscheme with fibrewise reduced Picard identity components.
\par\emph{Statement/proof location:} work/dgj/section\_D.tex.
\emph{Manuscript reference:} 57:288.
\par Fibrewise image of (F\textasciicircum{}2,[2]) has kernel degree 2 then 4.
\par\medskip\noindent\textbf{K.139. D4-product --- Proved under stated hypotheses.}\par
Picard scheme and torsion NS in the product Igusa family.
\par\emph{Statement/proof location:} work/dgj/section\_D.tex.
\emph{Manuscript reference:} 57:311.
\par Pic-tau flat; Picvar compatible; NS torsion Z/2 generically and zero at supersingular point.
\par\medskip\noindent\textbf{K.140. D4-versal --- Proved.}\par
Full modular or versal Igusa deformation family.
\par\emph{Statement/proof location:} work/dgj/section\_D.tex D4o and D4s; section\_D\_supersingular.tex.
\emph{Manuscript reference:} 57:311.
\par Both algebraized miniversal hulls computed: ordinary W[[t,s,u,v]]/(s\ensuremath{{}^2}+2s); supersingular W[[t,s,u,v]]/(s\ensuremath{{}^4}−t*s\ensuremath{{}^2}−2s). Canonical fixed elliptic subgroup and all-Artinian Pic-tau flatness proved; root/fhi independent audits pass.
\par\medskip\noindent\textbf{K.141. G2a --- Proved under stated hypotheses.}\par
Determinant boxed relation with explicit sign and typed equivariance.
\par\emph{Statement/proof location:} work/dgj/section\_G.tex.
\emph{Manuscript reference:} 59:10-12.
\par Proved for polarizations, the hypothesis used in Part G.
\par\medskip\noindent\textbf{K.142. G2b --- Proved.}\par
Determinant trivializations agree with canonical sections.
\par\emph{Statement/proof location:} work/dgj/section\_G.tex.
\emph{Manuscript reference:} 59:12-14.
\par Includes square and even-Euler double transform comparisons; corrects ppcm and even-chi divisibility claim.
\par\medskip\noindent\textbf{K.143. G3 --- Proved.}\par
Albanese sign c=-lambda\_Theta and canonical class K=omega\_C.
\par\emph{Statement/proof location:} work/dgj/section\_G.tex.
\emph{Manuscript reference:} 59:23-28.
\par Minus sign visually verified; canonical norm on PPAV theta torsor.
\par\medskip\noindent\textbf{K.144. J0 --- Proved textual scope correction.}\par
Complementary index and distinction of equivalence relations.
\par\emph{Statement/proof location:} work/dgj/section\_J.tex.
\emph{Manuscript reference:} 44:2;44:8.
\par Exponent remains palaeographically uncertain; n+1-i explicit in letter; modern Ab is not silently identified with manuscript quotient. Th1 en substance/ensemblistement and Th4 margin Connu/Annulé or Cancelé transcription conflicts retained.
\par\medskip\noindent\textbf{K.145. J2a --- Refuted.}\par
Universal canonical intersection isomorphism Ab2 to its dual.
\par\emph{Statement/proof location:} work/dgj/section\_J.tex.
\emph{Manuscript reference:} 44:2;44:8-9.
\par Generic monic septic Jacobian over the algebraic closure of F5(a0,...,a6)\textasciicircum{}perf; intersection forces a degree8 K-isogeny, ruled out by irreducible S7 action on A[2]. No integral torsion standard assumption used.
\par\medskip\noindent\textbf{K.146. J2b --- Refuted.}\par
Intersection isomorphism is a polarization up to one sign.
\par\emph{Statement/proof location:} work/dgj/section\_J.tex.
\emph{Manuscript reference:} 44:2;44:8-9.
\par Explicit E/F3 times P1 times P1; matrix [[0,lambda],[lambda,0]].
\par\medskip\noindent\textbf{K.147. F1 --- Known.}\par
Checked prime-to-p, pronilpotent and strongly F-regular finiteness theorems.
\par\emph{Statement/proof location:} work/fhi/F.tex; F\_proof.md.
\emph{Manuscript reference:} 58:19.
\par Full surface group remains undecided
\par\medskip\noindent\textbf{K.148. F2 --- Refuted.}\par
Dimension-free consequence is false for k[[t]] in characteristic p.
\par\emph{Statement/proof location:} work/fhi/F.tex; F\_proof.md.
\emph{Manuscript reference:} 58:19.
\par Correct consequence requires components dimension at least two
\par\medskip\noindent\textbf{K.149. F3 --- Proved.}\par
Metabelian group with pronilpotent quotient Z3 and fg open abelianizations need not be fg.
\par\emph{Statement/proof location:} work/fhi/F.tex; F\_proof.md.
\emph{Manuscript reference:} 58:19.
\par Abstract obstruction; no geometric realization claimed
\par\medskip\noindent\textbf{K.150. F-full --- Undecided.}\par
Full finite generation for all normal complete local surfaces.
\par\emph{Statement/proof location:} work/fhi/F\_proof.md section F4.
\emph{Manuscript reference:} 58:19.
\par For complete normal local surfaces with algebraically closed residue field, the unrestricted conjecture remains. In equal characteristic F8/F9b settle every geometric quotient with normal pro-p kernel over the full tame quotient. F10 proves finite abelian inertia using Krishna–Majumder2026v2; F10c gives all-prime rank-linear H1 bounds for abelian monodromy, and F11 proves uniform finite generation of metabelian quotients. F12 gives the characteristic-p bound for metabelian coefficients; F12b settles pro-p closed second derived subgroup. F13 proves a fixed log-conductor bound after one fixed cover. F13c/d remove unramified character order and reduce the remaining coefficient-prime-not-p estimate to a degree-linear tame-generator bound for cyclic covers of uniformly bounded p-power degree (F13.6). F13e bounds refined-Swan zero length, but not cleanup complexity. F13.5/F13.6, further solvable and nonsolvable wild quotients, and mixed characteristic remain unresolved. The actual wild kernel is not assumed pro-p.
\par\medskip\noindent\textbf{K.151. H1 --- Proved.}\par
Connected pro-generalized Jacobian has projective dimension at most one in Serre category.
\par\emph{Statement/proof location:} work/fhi/H.tex; H\_proof.md.
\emph{Manuscript reference:} 43:24-35;43:80.
\par All proalgebraic targets; all p
\par\medskip\noindent\textbf{K.152. H2 --- Proved.}\par
Ext>=2(P,Ga)=0 for every source P in Serre category.
\par\emph{Statement/proof location:} work/fhi/H.tex; H\_proof.md.
\emph{Manuscript reference:} 43:80.
\par Corrects assertion that Ga remains difficult
\par\medskip\noindent\textbf{K.153. H3 --- Refuted.}\par
Unaugmented all-degree geometric comparison fails in degree zero.
\par\emph{Statement/proof location:} work/fhi/H.tex; H\_proof.md.
\emph{Manuscript reference:} 43:80.
\par Retain augmentation by Z in a specified larger category
\par\medskip\noindent\textbf{K.154. H4 --- Proved.}\par
Etale local cohomology has H3\_m(k[[s,t]],Z/p) nonzero.
\par\emph{Statement/proof location:} work/fhi/H.tex; H\_proof.md.
\emph{Manuscript reference:} 43:70-74;SA19-H4.
\par Refutes topology/category transfer, not original Zariski representability
\par\medskip\noindent\textbf{K.155. I1 --- Proved.}\par
Degree-zero endomorphisms are diagonal in every regular local residual complex.
\par\emph{Statement/proof location:} work/fhi/I.tex; I\_proof.md.
\emph{Manuscript reference:} 43:67.
\par Finite topology gives product adic topology
\par\medskip\noindent\textbf{K.156. I2 --- Proved.}\par
Complete surface presentation including End2=restricted adeles/F.
\par\emph{Statement/proof location:} work/fhi/I.tex; I\_proof.md.
\emph{Manuscript reference:} 43:65-68;SA19-H3.
\par Precise finite-topology universal interpretation; mixed arrow topology
\par\medskip\noindent\textbf{K.157. I-full --- Undecided.}\par
Full presentation under the manuscript completion without a category specification.
\par\emph{Statement/proof location:} work/fhi/I\_universal\_linear.tex; I\_proof.md.
\emph{Manuscript reference:} 43:65-68;SA19-H3.
\par All dimensions are proved for separated complete graded algebras with a basis of open left ideals; all surfaces are proved for arbitrary jointly continuous topology. The original43:66 specifies the broader category without the open-left-ideal hypothesis. I7b now proves open dense differential-image quotients for whole source columns over arbitrary regular local rings. Remaining: construction and continuity of the entire adjacent-arrow family and its endpoint assembly in arbitrary complete targets, in dimension at least3.
\par\medskip\noindent\textbf{K.158. R3g --- Proved.}\par
The generic smooth hyperplane restriction kernel is finite connected.
\par\emph{Statement/proof location:} sections/A.tex, a:generic.
\emph{Manuscript reference:} 57:99-100.
\par Proof uses Kummer Lefschetz, the p-torsion dlog injection, and the incidence variety; targeted independent check by fhi.
\par\medskip\noindent\textbf{K.159. A3f --- Proved.}\par
Direct Raynaud fibre polarizations in the stated degree range are not very ample.
\par\emph{Statement/proof location:} sections/A.tex, a:fibre.
\emph{Manuscript reference:} 57:99-100.
\par Semigroup pole orders force the fibre map to factor through Frobenius; finite check covers 518 pairs, while the proof treats all coprime pairs.
\par\medskip\noindent\textbf{K.160. E4-3103-normal-p --- Refuted.}\par
Pic tau universally open for normal morphism in residue characteristic p.
\par\emph{Statement/proof location:} work/bce/section\_E.tex; Theorem E4.H;Raynaud1979 4.2.1-4.2.6 pp135-139;CotnerZavyalov2024 Thm1.1.8(iii).
\emph{Manuscript reference:} 57:310.
\par Explicit rank8 Heisenberg degeneration over W(k), p2, realized by Raynaud complete intersections; Pic\ensuremath{\tau}=\ensuremath{\mu}2\ensuremath{{}^2} times (SpecR disjoint Speck), so an open component maps only to closed point. Parent audit passed.; explicit degree768 norm system and first-jet proof, root structural audit accepted
\par\medskip\noindent\textbf{K.161. E4-313dim-reduced --- Refuted.}\par
Constancy of Pic dimension with geometrically reduced fibres.
\par\emph{Statement/proof location:} work/bce/section\_E.tex; section\_E.tex Theorem E4.R.
\emph{Manuscript reference:} 57:313.
\par Finite flat double X over a ruled surface, locally z(z-d)=0, glues two reduced components along a doubled section. Universal base-change exactness and units Mayer-Vietoris yield Pic0 dimensions1to2.
\par\medskip\noindent\textbf{K.162. SA14-H11-affinity-normal --- Proved.}\par
Affine Pic pullback after shrinking one-dimensional base, for normal generic fibres.
\par\emph{Statement/proof location:} work/bce/additional\_leads.tex; SGA6 XII1.1;Kleiman4.18.5;Stacks09N9;H6.
\emph{Manuscript reference:} 57:74;78.
\par Full quotient proof in additional\_leads.tex; includes finite/projective/relative-normal candidate and smooth birational candidate
\par\medskip\noindent\textbf{K.163. E4-3105projective-smooth --- Known.}\par
Projectivity of tau for smooth projective morphism.
\par\emph{Statement/proof location:} work/bce/section\_E.tex; LaurentSchroer2.3(ii);Kleiman6.18.
\emph{Manuscript reference:} 57:310.
\par Kleiman2005 Exercise6.18 p59 and Answer6.18 p72 explicitly prove projectivity of tau for smooth projective geometrically integral families over noetherian S; opened primary text. Disconnected smooth fibres reduce via finite étale Stein factor.
\par\medskip\noindent\textbf{K.164. D4-first-order --- Refuted.}\par
Universal first-order Igusa deformation has nonflat Pic-tau.
\par\emph{Statement/proof location:} work/dgj/section\_D.tex D4o; attempt\_D4\_ordinary.tex.
\emph{Manuscript reference:} 57:311.
\par Every Artinian deformation of an Igusa surface, ordinary or supersingular first factor, has flat Pic-tau; full proofs independently audited by root and fhi. Literal Chai2008v2 Theorem2.4(iii) and Chai-Oort2017 first-order assertion are false.
\par\medskip\noindent\textbf{K.165. D4-gap --- Refuted inference; actual Igusa assertion separately refuted by complete proof.}\par
Different absolute tangent dimensions of nonsmooth torsors imply nonflatness.
\par\emph{Statement/proof location:} work/dgj/section\_D.tex D4o; attempt\_D4\_ordinary.tex; finite torsor certificate.
\emph{Manuscript reference:} 57:311.
\par The finite free torsor refutes the tangent inference. Complete ordinary and supersingular deformation calculations additionally refute actual first-order nonflatness; historical unresolved statuses superseded.
\par\medskip\noindent\textbf{K.166. H5 --- Proved.}\par
Higher Ext vanishing also on the perfect etale site with exact category matching.
\par\emph{Statement/proof location:} work/fhi/H.tex; H\_proof.md H5.
\emph{Manuscript reference:} 43:80.
\par Source is extension of abelian variety by affine proalgebraic group; no arbitrary pro-target comparison
\par\medskip\noindent\textbf{K.167. H4b --- Refuted.}\par
Ordinary local Jacobian of k[s,t]\_(s,t) is not projective.
\par\emph{Statement/proof location:} work/fhi/H.tex; H\_proof.md H4b.
\emph{Manuscript reference:} 43:70-74;SA19-H4.
\par Actual nonsplit Z/p extension in Pro(AbSch); all p
\par\medskip\noindent\textbf{K.168. I3 --- Proved.}\par
Universal presentation for every regular local surface via explicit arrow completion.
\par\emph{Statement/proof location:} work/fhi/I.tex; I\_proof.md section I3.
\emph{Manuscript reference:} 43:65-68;SA19-H3.
\par No extra degree-two kernel; arrow completion target identified rigorously
\par\medskip\noindent\textbf{K.169. I4 --- Proved.}\par
Every adjacent residual arrow identifies its full Hom module as a separated coefficient completion.
\par\emph{Statement/proof location:} work/fhi/I.tex; I\_proof.md I4; I\_higher.md.
\emph{Manuscript reference:} 43:65-67;SA19-H3.
\par Topology q\textasciicircum{}N A\_q + p\textasciicircum{}n ÂA\_p; injectivity and density proved; every dimension
\par\medskip\noindent\textbf{K.170. I5 --- Proved.}\par
The diagonal algebra and differential generate all graded endomorphisms densely in every dimension.
\par\emph{Statement/proof location:} work/fhi/I.tex; I\_proof.md I5; I\_higher.md.
\emph{Manuscript reference:} 43:66-67;SA19-H3.
\par Dense generation does not settle quotient topology or all closed relations
\par\medskip\noindent\textbf{K.171. H3b --- Refuted.}\par
Even augmented perfect-site Hom cannot compare naturally to ordinary H0(K,G) for all algebraic groups.
\par\emph{Statement/proof location:} work/fhi/H.tex; H\_proof.md H3b.
\emph{Manuscript reference:} 43:80.
\par Perfection mismatch is independent of missing degree augmentation
\par\medskip\noindent\textbf{K.172. F1b --- Known.}\par
Full punctured etale fundamental group of every complete log terminal surface is finite in all equal and mixed characteristics.
\par\emph{Statement/proof location:} work/fhi/F.tex; F\_proof.md F1b.
\emph{Manuscript reference:} 58:19.
\par Includes p2,p3; log terminality additional; no prime-to-p-order assertion
\par\medskip\noindent\textbf{K.173. A3n --- Proved.}\par
A very ample negative-H1 example requires q\_nil>=h0(L)>=5 on a surface.
\par\emph{Statement/proof location:} sections/A.tex, a:nilpotence.
\emph{Manuscript reference:} 57:99-100.
\par Opened Enokizono Th4.20 with E=0; low projective embeddings have H1(O)=0.
\par\medskip\noindent\textbf{K.174. A3s --- Proved.}\par
An A3 surface is of general type, satisfies explicit numerical bounds, and has no smooth W2 lifting.
\par\emph{Statement/proof location:} sections/A.tex, a:generaltype.
\emph{Manuscript reference:} 57:99-100.
\par Opened Terakawa Th1.6, corrected Zhang Th3.1, and Langer Cor11.2; very ampleness forces degree at least3 on a genus-one fibre.
\par\medskip\noindent\textbf{K.175. A3R --- Proved.}\par
Every positive two-parameter Raynaud polarization with negative H1 fails very ampleness.
\par\emph{Statement/proof location:} sections/A.tex, a:raynaud-all.
\emph{Manuscript reference:} 57:99-100.
\par Connection lemma gives vanishing for a>=ell; fibre semigroup excludes a<ell. This family exclusion is compatible with the separate A3 cover construction.
\par\medskip\noindent\textbf{K.176. C2-DVR --- Proved in all characteristics without flatness of X.}\par
Picard completion isomorphism with finite-length H1 over an actual excellent henselian DVR.
\par\emph{Statement/proof location:} work/bce/section\_C.tex; section\_C.tex Theorem C2.D;Artin1969 Thm3.5.
\emph{Manuscript reference:} 57:272.
\par Complete elementary proof in Theorem C2.D; no H2 finite-length assumption. Flat closure plus Artin-Rees fibre-product patching handles nonflat X.
\par\medskip\noindent\textbf{K.177. SA14-H11-affinity-general --- Proved.}\par
Affine Pic pullback for proper surjections after shrinking a one-dimensional integral base.
\par\emph{Statement/proof location:} work/bce/additional\_leads.tex; SGA3 XVI4.2 pp521-522;VI\_B1.4.2 pp325-326;SGA6 XII1.1-1.2.
\emph{Manuscript reference:} 57:74;78.
\par Kernel quotient and fibrewise-finite monomorphism argument, full proof Theorem H11.G.
\par\medskip\noindent\textbf{K.178. D4-modular-Picvar --- Proved.}\par
A modular ordinary-Igusa degeneration has flat Pic-tau but incompatible fibrewise Picard varieties.
\par\emph{Statement/proof location:} work/dgj/section\_D.tex; section\_D\_modular.tex.
\emph{Manuscript reference:} 57:311; 57:288.
\par Over W(k), E=F:Y2Z+XYZ=X3+XZ2, P=(0:0:1), Q=(-2:1:w) with 2w2-w+8=0. Klein-four free quotient; Pic-tau finite flat degree8 over F-dual; Albanese degrees4 generic and2 special exclude an abelian Picvar subscheme.
\par\medskip\noindent\textbf{K.179. F3b --- Proved.}\par
Explicit metabelian example has H1(G,M\_n)=L\_n\textasciicircum{}n and unbounded normalized twisted cohomology.
\par\emph{Statement/proof location:} work/fhi/F.tex; F\_proof.md F3b.
\emph{Manuscript reference:} 58:19.
\par Necessary d-generator bound dimH1<=d dimM fails; no geometric realization claimed
\par\medskip\noindent\textbf{K.180. I6 --- Proved.}\par
Residual Hom columns are exact after closing images above the source height.
\par\emph{Statement/proof location:} work/fhi/I.tex; I\_proof.md I6; I\_relations.md.
\emph{Manuscript reference:} 43:66-68;SA19-H3.
\par Quadratic kernel identified in actual Hom topology for gap-two intervals; free-algebra topology still separate
\par\medskip\noindent\textbf{K.181. R3u --- Proved.}\par
The generic hyperplane restriction kernel is infinitesimal unipotent for normal projective X.
\par\emph{Statement/proof location:} sections/A.tex, a:generic-unipotent.
\emph{Manuscript reference:} 57:99-100.
\par Rational dlog handles normal X; Raynaud Proposition6.2.3 excludes mu\_p; Cartier duality proves unipotence. Root independently checked the source and proof.
\par\medskip\noindent\textbf{K.182. A-CZ346ii --- Proved.}\par
The additional smooth one-form isomorphism hypothesis in Cotner-Zavyalov Example3.4.6(ii) is incompatible with its alpha\_p kernel.
\par\emph{Statement/proof location:} sections/A.tex, a:cz-vacuity.
\emph{Manuscript reference:} 57:99-100 source route.
\par The conditional auxiliary example is vacuous; this is not a refutation of the paper main theorems and does not settle A1/A2.
\par\medskip\noindent\textbf{K.183. R1 --- Corrected and rechecked.}\par
Supplied divisor counterexample and refined relative-complete-intersection criterion.
\par\emph{Statement/proof location:} sections/B.tex.
\emph{Manuscript reference:} 57:110-111.
\par The displayed affine factorization requires char not2 and i\ensuremath{{}^2}=-1; proper examples in every characteristic supplied. R2 sufficiency retains the manuscript relative lci hypothesis.
\par\medskip\noindent\textbf{K.184. R2 --- Carried as background; not used.}\par
Transgression identity on57:144.
\par\emph{Statement/proof location:} Input readings only.
\emph{Manuscript reference:} 57:144.
\par No result in this run relies on this identity; no new independent verification is claimed.
\par\medskip\noindent\textbf{K.185. R5 --- Carried as background; not used.}\par
Factor2 in the recalled Igusa b2 formula.
\par\emph{Statement/proof location:} Input readings only.
\emph{Manuscript reference:} 57:70.
\par The quotient constructions use directly computed Hodge and fixed-group data, not the supplied b2 formula.
\par\medskip\noindent\textbf{K.186. R6 --- Refuted as stated; corrected in scope.}\par
Unrestricted product formula for the Picard group of a quotient.
\par\emph{Statement/proof location:} sections/D.tex, D0.
\emph{Manuscript reference:} 57:286.
\par Cross-correspondence terms already refute the full Pic product formula for trivial Gamma. The properly scoped Pic-tau formula is proved and used.
\par\medskip\noindent\textbf{K.187. R7 --- Carried as background; not used.}\par
Functional equation determinant identity and correction terms.
\par\emph{Statement/proof location:} Input readings only.
\emph{Manuscript reference:} 60:53.
\par This background claim is not used by PartsA-J proofs in this run; no independent adoption or new verification is claimed.
\par\medskip\noindent\textbf{K.188. R8 --- Corrected and rechecked in stated cases.}\par
Local Lefschetz reductions and finite H1 of open subgroups.
\par\emph{Statement/proof location:} sections/F.tex.
\emph{Manuscript reference:} 58:2-19.
\par Full group finite generation remains open in the run; prime-to-p and pronilpotent results are proved from opened sources. Dimension-one extrapolation is refuted.
\par\medskip\noindent\textbf{K.189. R9 --- Proved with typed signs.}\par
Boxed determinant relation and canonical torsor point criterion.
\par\emph{Statement/proof location:} sections/G.tex, G1 and G2a.
\emph{Manuscript reference:} 59:10-14.
\par Complete criterion and determinant comparison are proved, with quadratic-refinement monodromy obstruction.
\par\medskip\noindent\textbf{K.190. R10 --- Known theorem with scope and date corrections.}\par
Neron component-group pairing and historical comparison.
\par\emph{Statement/proof location:} sections/H.tex, priority appendix.
\emph{Manuscript reference:} 43:117.
\par SGA7IX pairing and its perfectness conjecture are distinguished; current perfect-residue-field theorem cited. Leaf priority is not inferred from dossier dates.
\par\medskip\noindent\textbf{K.191. I7 --- Proved.}\par
Complete regular-ring residual columns are exact and their maps onto cycles are open above the first arrow.
\par\emph{Statement/proof location:} work/fhi/I.tex; I\_complete\_columns.tex; I\_complete\_columns.md.
\emph{Manuscript reference:} 43:66-68;SA19-H3.
\par All dimensions over complete A; uniform universal-family topology remains unresolved
\par\medskip\noindent\textbf{K.192. D4-supersingular --- Proved.}\par
Full supersingular Igusa hull and all-Artinian Pic-tau flatness.
\par\emph{Statement/proof location:} work/dgj/section\_D.tex D4s; section\_D\_supersingular.tex.
\emph{Manuscript reference:} 57:311.
\par Exact ring W[[t,s,u,v]]/(s\ensuremath{{}^4}−t*s\ensuremath{{}^2}−2s), finite free rank4 over W[[t,u,v]]. Pic-tau finite locally free degree8 over F-dual. Complete independent root/fhi audits.
\par\medskip\noindent\textbf{K.193. J1i --- Known, with derived kernel statement Proved.}\par
All-codimension incidence quotient, family kernels, bound and complementary isogeny.
\par\emph{Statement/proof location:} work/dgj/section\_J\_incidence.tex; section\_J.tex.
\emph{Manuscript reference:} 44:2;44:8.
\par Primary1971 record credits Grothendieck1964; all-regular representability not identified. Single-self-pairing shortcut refuted by E×P1\ensuremath{{}^2}.
\par\medskip\noindent\textbf{K.194. I8 --- Proved.}\par
Full residual endomorphism presentation in every dimension for complete algebras with open-left-ideal basis.
\par\emph{Statement/proof location:} work/fhi/I\_universal\_linear.tex.
\emph{Manuscript reference:} 43:64-66.
\par Full proof independently passed by root and DGJ; BCE passed the generator lemma; open-left-ideal category retained; no inference for arbitrary non-linear topology
\par\medskip\noindent\textbf{K.195. E4-309bprime-fibreSmooth --- Refuted.}\par
Is fibrewise Picard smoothness an open condition?.
\par\emph{Statement/proof location:} work/bce/section\_E.tex; new\_E4\_smoothness.tex; code/check\_cuspidal\_igusa.output.json; work/dgj/audit\_E4\_cusp.md.
\emph{Manuscript reference:} 57:309.
\par E4.S: char2 cuspidal Igusa degeneration has Pic-tau=Ydual times G[2], generically Ydual times mu2 times Z2 and specially Ydual times Ga. Original X fibres remain geometrically integral reduced.
\par\medskip\noindent\textbf{K.196. E4-309bprime-defect --- Refuted.}\par
Upper semicontinuity of tangent dimension minus fibre dimension?.
\par\emph{Statement/proof location:} work/bce/section\_E.tex; Theorem E4.S.
\emph{Manuscript reference:} 57:309.
\par Same E4.S has tangent dimension2 on both fibres, geometric dimensions1 generic and2 special; the defect drops1to0.
\par\medskip\noindent\textbf{K.197. E4-309bprime-dimension --- Proved.}\par
Upper semicontinuity of Picard fibre dimension?.
\par\emph{Statement/proof location:} work/bce/section\_E.tex; FGA VI Cor2.7 p234; Stacks fibre dimension theorem.
\emph{Manuscript reference:} 57:309.
\par The relative local dimension function is upper semicontinuous on a locally finitely presented group space; translate to the identity and compose with the identity section.
\par\medskip\noindent\textbf{K.198. E4-309bprime-relativeSmooth --- Proved.}\par
Openness of the relative smooth locus?.
\par\emph{Statement/proof location:} work/bce/section\_E.tex; Standard smooth-locus theorem; E4.S scope distinction.
\emph{Manuscript reference:} 57:309.
\par Smoothness is open for locally finitely presented morphisms. It includes flatness and differs from smoothness of each individual geometric fibre.
\par\medskip\noindent\textbf{K.199. D4-artin-Picvar --- Refuted.}\par
Every Artinian Igusa deformation admits a compatible abelian Picard subscheme.
\par\emph{Statement/proof location:} work/dgj/section\_D.tex D4a; section\_D\_artin\_picvar.tex.
\emph{Manuscript reference:} 57:311.
\par R=Fbar2[epsilon]/epsilon\ensuremath{{}^2}, E0=F0:y\ensuremath{{}^2}+xy=x\ensuremath{{}^3}+1, P=(0,1), q12=1+epsilon, otherq=1. Any abelian Picvar lift would yield f:B\ensuremath{\to}Adual and force1=1+epsilon. Root and bce independently audited; matches FGA VI Remark3.4 p238.
\par\medskip\noindent\textbf{K.200. J6L --- Proved under explicit algebraicity hypotheses.}\par
Conditional Lefschetz theorem for incidence abelian varieties.
\par\emph{Statement/proof location:} work/dgj/section\_J.tex J6L; section\_J\_lefschetz.tex.
\emph{Manuscript reference:} 44:4 conditional stable and high ranges.
\par Cohomologically zero correspondences act trivially on incidence varieties. Lambda\_X j\_* is stable inverse; j\_* Lambda\_Y is high-range right inverse. No extra sl2 coefficients; root independently audited.
\par\medskip\noindent\textbf{K.201. J6P --- Proved under explicit smooth-axis and generic-member hypotheses.}\par
Codimension-two pencil restriction equals the trace image, with stable-range surjectivity.
\par\emph{Statement/proof location:} work/dgj/section\_J\_pencil.tex; section\_J.tex J6P.
\emph{Manuscript reference:} 44:4 middle clause for i=2, n=4; stable surjectivity for i=2, n>=5.
\par All characteristics. Actual Ab2 blowup decomposition absorbs exceptional Picvar B via Picvar X. Finite kernel additionally follows from stated torsion-injectivity assumption, known char0. Independent root audit PASS.
\par\medskip\noindent\textbf{K.202. J6Z --- Proved.}\par
Highest nonzero target codimension restriction Pinc\textasciicircum{}(n-1)X to AlbY is surjective.
\par\emph{Statement/proof location:} work/dgj/section\_J\_pencil.tex; section\_J.tex J6Z.
\emph{Manuscript reference:} 44:4 high range i=n-1 for n>=3.
\par Any smooth hyperplane section in every characteristic. Also universal regular representatives whenever source exists, in particular Ab2 of any smooth projective threefold. Independent root audit PASS.
\par\medskip\noindent\textbf{K.203. F4 --- Proved.}\par
Uniform normalized simple-module H1 bounds imply quantitative finite generation of every prosolvable group.
\par\emph{Statement/proof location:} work/fhi/F\_prosolvable\_criterion.tex; F.tex.
\emph{Manuscript reference:} 58:19.
\par At most ceil(C)+1 generators; geometric uniform bound still unproved; nonsolvable quotients separate
\par\medskip\noindent\textbf{K.204. F1c --- Known.}\par
Mixed-characteristic BCM-regular pairs have finite prime-to-p local fundamental group bounded by inverse perfectoid signature.
\par\emph{Statement/proof location:} work/fhi/F.tex.
\emph{Manuscript reference:} 58:19.
\par Complete normal local ring; algebraically closed residue; effective boundary and Q-Cartier/BCM-regular pair are additional assumptions
\par\medskip\noindent\textbf{K.205. F5 --- Proved.}\par
Uniform rank-linear H1 bound for all abelian prime-to-p monodromy F\_p local systems on complete normal equal-characteristic local surfaces.
\par\emph{Statement/proof location:} work/fhi/F\_tame\_abelian\_bound.tex; F.tex.
\emph{Manuscript reference:} 58:19.
\par Excludes F3 as a geometric quotient in char2; no full prosolvable or full pi1 finite-generation claim
\par\medskip\noindent\textbf{K.206. F6 --- Proved.}\par
Every pro-p-by-abelian pro-prime-to-p quotient of a complete normal equal-characteristic local surface group is uniformly finitely generated.
\par\emph{Statement/proof location:} work/fhi/F\_pro\_p\_by\_abelian.tex; F.tex.
\emph{Manuscript reference:} 58:19.
\par Also virtually abelian tame quotient gives finite generation; arbitrary solvable towers and full group remain undecided
\par\medskip\noindent\textbf{K.207. A2D --- Proved.}\par
For fixed nontrivial p-torsion M and very ample L, smooth hyperplanes trivializing M form a finite determinantal locus, and a nonempty locus requires h1(M)-h1(M(-L))>=h0(L)-2.
\par\emph{Statement/proof location:} sections/A.tex, a:determinantal; work/root/A2\_determinantal.tex.
\emph{Manuscript reference:} 57:97.
\par Complete conormal and maximal-minor proof; BCE independent targeted audit passed. Gives a finite fixed-pair test, not a universal exclusion or a counterexample.
\par\medskip\noindent\textbf{K.208. E4-315schemeCurve --- Proved.}\par
Full Picard scheme for smooth proper geometrically connected curves over arbitrary base.
\par\emph{Statement/proof location:} work/bce/section\_E.tex; Theorem E4.Curve; BLR1990 6.6 Th2(c) pp164–166, 9.4 Prop4 p260; Stacks07S6; prior derivation Cesnavicius2015.
\emph{Manuscript reference:} 57:315.
\par Genus0 or >=2 is projective by the relative canonical bundle. In genus1 C=Pic1 is an E-torsor; Picn=C/E[n] and BLR6.6 Th2(c) gives affine neighborhoods of every finite orbit, so finite-flat quotient is schematic. Noetherian approximation gives arbitrary base.
\par\medskip\noindent\textbf{K.209. F7 --- Proved.}\par
Uniform rank-linear H1 bound for all finite F\_p systems extending over a fixed resolution.
\par\emph{Statement/proof location:} work/fhi/F\_unramified\_coefficients.tex; F.tex.
\emph{Manuscript reference:} 58:19.
\par Includes nonabelian and wild coefficient monodromy; extension over fixed resolution required
\par\medskip\noindent\textbf{K.210. F7a --- Proved.}\par
Relative generation criterion for a normal pro-p kernel over any finitely generated quotient.
\par\emph{Statement/proof location:} work/fhi/F\_unramified\_coefficients.tex; F.tex.
\emph{Manuscript reference:} 58:19.
\par No solvability of quotient required; uniform simple F\_p-module H1 bound required
\par\medskip\noindent\textbf{K.211. F7b --- Proved.}\par
Pro-p extensions of local fundamental-group quotients factoring through a fixed resolution are uniformly finitely generated.
\par\emph{Statement/proof location:} work/fhi/F\_unramified\_coefficients.tex; F.tex.
\emph{Manuscript reference:} 58:19.
\par Quotient may be nonabelian and wild; does not assert actual kernel of pi1(U) to pi1(X) is pro-p
\par\medskip\noindent\textbf{K.212. J6C --- Refuted.}\par
The middle fixed factor equals the image for every integral-fibre hyperplane pencil.
\par\emph{Statement/proof location:} work/dgj/section\_J6\_cubic.tex; section\_J.tex J6C.
\emph{Manuscript reference:} 44:4 middle clause with n=4 and i=2.
\par Characteristic5: X=sum\_\{j=0\}\textasciicircum{}2 x\_j\textasciicircum{}3+x3\textasciicircum{}2*x4+x4\textasciicircum{}3+x5\textasciicircum{}3, pencil x5=t*x4. All fibres geometrically integral. Ab2X=Pinc2X=0. Ruling correspondence from constant Pic0(Fermat cubic) has transpose composite[-2], so target trace is positive. Root and BCE full independent audits PASS; exact finite certificate saved. No priority claim; Eckardt-cone geometry has classical sources.
\par\medskip\noindent\textbf{K.213. D4b --- Refuted.}\par
Every deformation of a Jacobian ordinary Igusa surface lifts its Albanese elliptic fibration.
\par\emph{Statement/proof location:} work/dgj/section\_D\_partsch.tex; section\_D.tex D4b.
\emph{Manuscript reference:} 57:311; associated modern-source verification.
\par Any lift to a smooth proper genus-one curve gives pullback from its abelian Pic0 lifting the special Picvar inclusion, contradicted by q12=1+epsilon. Refutes literal Partsch2013 Proposition6.6 pp102-103, Theorem6.7 p104 and Proposition6.9 p106 for this actual ordinary Jacobian Igusa deformation. Root full audit PASS; finite Z2-to-mu2 map certifies the etale-source inference problem separately.
\par\medskip\noindent\textbf{K.214. JWalker --- Unresolved; exact equivalent dimension test proved.}\par
Walker isogeny-lift universality implies universality for all regular homomorphisms.
\par\emph{Statement/proof location:} work/dgj/section\_J\_walker.tex; section\_J.tex JWalker.
\emph{Manuscript reference:} 44:2 all-regular interpretation of Th1-3 remains a gate.
\par Stronger universality is equivalent to dim image(usual AJ,phi)=dim Ja for every regular test phi. Positive-dimensional kernel of the first projection is not ruled out by divisible AJ kernel. Kahn final Conjecture3.1 p847 and August2026 ACMV version checked; neither conjecture nor general existence settled.
\par\medskip\noindent\textbf{K.215. F8 --- Proved.}\par
Uniform rank-linear H1 bound for all finite F\_p systems tame along a fixed SNC resolution.
\par\emph{Statement/proof location:} work/fhi/F\_tame\_rootstack.tex; F.tex.
\emph{Manuscript reference:} 58:19.
\par Independent root full-source/proof audit and BCE targeted audit passed; full monodromy may have a p-part; no historical priority claim
\par\medskip\noindent\textbf{K.216. F8b --- Proved.}\par
All local quotients with normal pro-p kernel and arbitrary pro-prime-to-p quotient are uniformly finitely generated.
\par\emph{Statement/proof location:} work/fhi/F\_tame\_rootstack.tex; F.tex.
\emph{Manuscript reference:} 58:19.
\par No abelian or solvable hypothesis on prime-to-p quotient; alternating tame/wild normal factors remain outside the conclusion
\par\medskip\noindent\textbf{K.217. F9 --- Known method; proved here.}\par
Uniform generator bound for the full local tame quotient of a fixed SNC resolution.
\par\emph{Statement/proof location:} work/fhi/F\_tame\_quotient.tex; F.tex.
\emph{Manuscript reference:} 58:19.
\par Full root audit passed; d\_X depends only on component genera and exceptional dual graph; no priority claim
\par\medskip\noindent\textbf{K.218. F9b --- Proved.}\par
Local quotients with normal pro-p kernel and quotient factoring through the full tame local group are uniformly finitely generated.
\par\emph{Statement/proof location:} work/fhi/F\_tame\_quotient.tex; F.tex.
\emph{Manuscript reference:} 58:19.
\par Allows tame quotients with p-part and no solvability assumption; actual wild kernel is not assumed pro-p
\par\medskip\noindent\textbf{K.219. A2V --- Proved necessary condition.}\par
A Cartier-fixed regular one-form vanishing along a smooth ample divisor has order pa and a nonzero exact tangential leading term, yielding a Frobenius kernel on H1(Y,O\_Y(-aY)).
\par\emph{Statement/proof location:} sections/A.tex; work/dgj/A2\_cartier\_leading.tex.
\emph{Manuscript reference:} 57:97.
\par Also forces h0(Omega\_X)>=h0(O\_X(pY)) and a nonnormal associated mu\_p torsor. The affine model dlog(1+t\textasciicircum{}p x) disproves naive p-squared order iteration. Root independent audit and finite coefficient checks passed. Full A2 global extension/exclusion remains unresolved.
\par\medskip\noindent\textbf{K.220. A2W --- Proved under explicit closedness hypothesis.}\par
All global one-forms closed implies Picard point injectivity for every smooth very ample divisor in dimension at least3.
\par\emph{Statement/proof location:} work/dgj/A2\_closed\_forms.tex; root integration in Part A.
\emph{Manuscript reference:} 57:97-100 partial exclusion.
\par H0(Omega1 tensor L\textasciicircum{}-1)=0 in dimension at least2. Any point-kernel example in dimension n>=3 must have dim image d:H0Omega1 to H0Omega2 >=n-1; it necessarily has nonclosed global one-forms. The hypothesis is not automatic. Root audited the all-closed proof; quantitative refinement follows by the same wedge rank.
\par\medskip\noindent\textbf{K.221. F10 --- Proved.}\par
Finite kernel of the abelian local fundamental group map to a fixed resolution and uniform rank-one Swan bound.
\par\emph{Statement/proof location:} work/fhi/F\_abelian\_wild\_inertia.tex; F.tex.
\emph{Manuscript reference:} 58:19.
\par Depends on Nov8 2026 preprint in exact regular F-finite category; root full and BCE targeted audits passed; no priority claim
\par\medskip\noindent\textbf{K.222. F10b --- Proved.}\par
One fixed finite abelian cover removes inertia for all original abelian coefficient monodromy systems.
\par\emph{Statement/proof location:} work/fhi/F\_abelian\_wild\_inertia.tex; F.tex.
\emph{Manuscript reference:} 58:19.
\par All primes and ranks; same cover and resolution; no prime-to-coefficient-degree assumption
\par\medskip\noindent\textbf{K.223. F10c --- Proved.}\par
Uniform rank-linear H1 bound over every coefficient prime for abelian monodromy.
\par\emph{Statement/proof location:} work/fhi/F\_abelian\_wild\_inertia.tex; F.tex.
\emph{Manuscript reference:} 58:19.
\par Root full audit passed; fixed-cover restriction need not be injective; finite kernel contribution bounded
\par\medskip\noindent\textbf{K.224. F11 --- Proved.}\par
Every metabelian quotient of the local fundamental group is uniformly finitely generated.
\par\emph{Statement/proof location:} work/fhi/F\_abelian\_wild\_inertia.tex; F.tex.
\emph{Manuscript reference:} 58:19.
\par Derived length two including alternating tame and wild factors; derived length three and arbitrary nonsolvable quotients remain open here
\par\medskip\noindent\textbf{K.225. A2Y --- Proved under either explicit target hypothesis; independent root audit PASS.}\par
Target Frobenius injectivity or a smooth W2 lifting of the target excludes a Picard point kernel.
\par\emph{Statement/proof location:} work/dgj/A2\_target\_frobenius.tex; root integration in Part A.
\emph{Manuscript reference:} 57:97-100 partial exclusion.
\par n>=3, Y smooth very ample. H1(L\textasciicircum{}-a) injects H1(OY) by a smooth connected E in|L\textasciicircum{}a| and multiplication commutes with Frobenius using s\textasciicircum{}p. A W2 lifting of Y alone makes all required negative H1 vanish in total degree1<min(p,dimY); no X/polarization lifting or dimY<p needed.
\par\medskip\noindent\textbf{K.226. F12 --- Proved.}\par
Uniform H1/rank bound for simple coefficient-p systems with metabelian monodromy.
\par\emph{Statement/proof location:} work/fhi/F\_metabelian\_p\_coefficients.tex; F.tex.
\emph{Manuscript reference:} 58:19.
\par Equal characteristic p; root full proof audit passed; restricted system need not be simple and cover degree need not be prime to p
\par\medskip\noindent\textbf{K.227. F12b --- Proved.}\par
All local quotients whose closed second derived subgroup is pro-p are uniformly finitely generated.
\par\emph{Statement/proof location:} work/fhi/F\_metabelian\_p\_coefficients.tex; F.tex.
\emph{Manuscript reference:} 58:19.
\par Second derived subgroup may be nonabelian; full derived-length3 coefficient-ell-not-p gate remains open here
\par\medskip\noindent\textbf{K.228. F13a --- Proved.}\par
Every simple modular representation of an abelian-by-abelian finite group is induced from a character of a subgroup containing the specified abelian normal subgroup.
\par\emph{Statement/proof location:} work/fhi/F\_metabelian\_log\_conductor.tex; F.tex.
\emph{Manuscript reference:} 58:19.
\par DGJ independent proof; root full audit passed; inducing subgroup normal with abelian quotient
\par\medskip\noindent\textbf{K.229. F13b --- Proved.}\par
Rank-one Swan cutoff is invariant under arbitrary connected finite etale covers of the fixed resolution.
\par\emph{Statement/proof location:} work/fhi/F\_metabelian\_log\_conductor.tex; F.tex.
\emph{Manuscript reference:} 58:19.
\par Root and BCE geometry audits passed; complete local Stein normalization checked; no bound on total character order inferred
\par\medskip\noindent\textbf{K.230. F13 --- Proved.}\par
All metabelian coefficient-ell-not-p systems have a fixed highest logarithmic break bound after one fixed cover.
\par\emph{Statement/proof location:} work/fhi/F\_metabelian\_log\_conductor.tex; F.tex.
\emph{Manuscript reference:} 58:19.
\par Root full proof and source audit passed; rank-linear complete-local Betti transfer F13.5 remains open; Hu-Teyssier finite-type theorem not applied beyond its scope
\par\medskip\noindent\textbf{K.231. A2Z --- Proved construction exclusion; independently audited; canonicalized by root integrator.}\par
A finite separable degree-p cover cannot kill a p-torsion class on a smooth ample full inverse image.
\par\emph{Statement/proof location:} work/dgj/A2\_fierce\_cover.tex.
\emph{Manuscript reference:} A2 failed-construction route only.
\par For finite flat generically separable degree p X to B smooth projective, C smooth ample and Y=pi*C smooth integral, Pic(C)[p](k) injects Pic(Y)[p](k). Nonnormal mu\_p torsor case included: effective conductor adds to positive ramification intersection. No closure of general A2.
\par\medskip\noindent\textbf{K.232. A2S-cross --- Known deduction from existing A3s; initial weaker quasi-elliptic formulation superseded.}\par
A threefold point-kernel target is general type and its leading order obeys the A3s numerical bounds.
\par\emph{Statement/proof location:} work/dgj/A2\_target\_surfaces.tex.
\emph{Manuscript reference:} A2 partial target exclusion.
\par For A2V integer a, a\ensuremath{{}^2}L\ensuremath{{}^2}<=B\_p(Y) from existing A3s since L\textasciicircum{}a is very ample. No remaining quasi-elliptic alternative and no new example.
\par\medskip\noindent\textbf{K.233. A2Gs-cross --- Undecided candidate exclusion; independent proof and source audit pending.}\par
No smooth G\_s torsor over a surface with smooth ample full inverse image.
\par\emph{Statement/proof location:} work/dgj/attempt\_A2\_Gs\_smooth\_surface.tex.
\emph{Manuscript reference:} A2 failed-construction route only.
\par No smooth surface Z with smooth ample D and line subbundle O(-D) of TZ tangent to D. Includes G\_s torsors with both total space and full inverse image smooth. Singular normalization and higher dimensions remain outside this argument.
\par\medskip\noindent\textbf{K.234. F13c --- Proved.}\par
Uniform bounded p-power ramified order modulo an unramified twist on every finite etale resolution cover.
\par\emph{Statement/proof location:} work/fhi/F\_bounded\_cyclic\_reduction.tex; F.tex.
\emph{Manuscript reference:} 58:19.
\par Root full proof audit passed; same exponent e on every cover; Swan conductors unchanged by the unramified twist
\par\medskip\noindent\textbf{K.235. F13d --- Proved.}\par
A degree-linear tame-generator bound for cyclic covers of fixed bounded p-power degree suffices for the remaining metabelian coefficient estimate.
\par\emph{Statement/proof location:} work/fhi/F\_bounded\_cyclic\_reduction.tex; F.tex.
\emph{Manuscript reference:} 58:19.
\par Root full proof audit passed; only sufficient reduction proved; geometric bound F13.6 and complete-local Betti bound F13.5 remain Undecided
\par\medskip\noindent\textbf{K.236. F13e --- Proved.}\par
Refined-Swan zero length is bounded linearly in the degree of every finite etale resolution cover.
\par\emph{Statement/proof location:} work/fhi/F\_refined\_swan\_zeros.tex; F.tex.
\emph{Manuscript reference:} 58:19.
\par Root full proof audit passed; component bound d\_a(-B\_ii)(t\_i-s\_a); total at most d sum\_i(-B\_ii)t\_i; no quantitative cleanup or local Betti conclusion inferred
\par\medskip\noindent\textbf{K.237. F13f --- Proved.}\par
A post-cover conductor cutoff below p gives uniform finite generation at derived length three.
\par\emph{Statement/proof location:} work/fhi/F\_small\_conductor\_corollary.tex; F.tex.
\emph{Manuscript reference:} 58:19.
\par Root audit passed; extra hypothesis t\_i<p on the fixed post-F10b resolution; earlier t\_i<1 criterion sharpened; no assertion such a cover always exists
\par\medskip\noindent\textbf{K.238. I7b --- Proved.}\par
Residual Hom-column differentials have open dense image without assuming the regular local ring complete.
\par\emph{Statement/proof location:} work/bce/I\_strict\_dense\_columns.tex; I.tex.
\emph{Manuscript reference:} 43:64-66.
\par BCE proof and root independent audit passed; arbitrary-target extension is conditional on preceding-degree continuity; whole-family continuity and endpoint assembly remain open
\par\medskip\noindent\textbf{K.239. A1Q --- Known; derivation proved.}\par
Picard restriction for quasi-F-split smooth ample targets, and for every smooth Fano threefold target.
\par\emph{Statement/proof location:} sections/A.tex, a:known-targets; work/root/A1\_known\_targets.tex.
\emph{Manuscript reference:} 57:97.
\par Quasi-F-split targets give injectivity for ambient dimension at least3 and isomorphism for dimension at least4. A smooth Fano threefold target gives an isomorphism in every p, including2. Exact target hypotheses and duality indices independently audited by DGJ; no unrestricted A1 or A2 conclusion.
\par\medskip\noindent\textbf{K.240. F13g --- Proved.}\par
A negative Swan intersection forces its positive coefficient to be divisible by p; nonzero Swan divisors have maximum coefficient at least p.
\par\emph{Statement/proof location:} work/fhi/F\_swan\_residue\_gap.tex; F.tex.
\emph{Manuscript reference:} 58:19.
\par Root full proof and exact primary-source audit passed; applies on every finite etale resolution cover
\par\medskip\noindent\textbf{K.241. F13h --- Proved.}\par
The uniform unramified exponent improves to any e with p\textasciicircum{}(e+1) greater than the maximal conductor cutoff.
\par\emph{Statement/proof location:} work/fhi/F\_swan\_residue\_gap.tex; F.tex.
\emph{Manuscript reference:} 58:19.
\par Root audit passed; sharpens the stated sufficient exponent by one p-step where applicable; F13.5 and F13.6 still open
\par\medskip\noindent\textbf{K.242. A2P --- Proved construction exclusion.}\par
A smooth ample relative Frobenius p-multisection of a projective-line bundle cannot kill p-torsion from its base.
\par\emph{Statement/proof location:} sections/A.tex; work/root/A2\_multisection\_exclusion.tex.
\emph{Manuscript reference:} 57:97.
\par For a quotient F*E to N the normal line is Q\textasciicircum{}p and the relative canonical line Q\textasciicircum{}(p-1). Normalization of a mu-p torsor would make the latter anti-effective, contradicting ample normal bundle. Positive-dimensional smooth projective base, quotient convention and full relative-Frobenius inverse section explicit. DGJ independent equation/conductor audit passed. Does not settle A2.
\endgroup

\bibliographystyle{alpha}
\bibliography{RESEARCH_SA14_SA19_references}
\end{document}
